\quellseite{573}
\chapter[C. Die zehn Gebote dem Volk zu Wittenberg gepredigt]{C. Die zehn Gebote dem Volk zu Wittenberg gepredigt durch den Pater Martin Luther, Augustiner.\footnote{Luther predigte unter großem Zulauf (St. Louiser Ausgabe, Bd. XXII, 1716, No. 150) des Volks zu Wittenberg vom Ende Juni 1516 bis zum 24. Februar 1517 über die zehn Gebote. Zu Anfang des vorigen Jahrhunderts war von diesen Predigten noch eine vollständige lateinische Handschrift Luthers (Weim. Ausg., Bd. I, S. 19. Köstlin, Mart. Luther, 3. Aufl., Bd. I, S. 123) vorhanden, die sich im Besitz von Valentin Ernst Löscher befand. Aus diesem Manuscripte hat er in seinen „Reformations-Acta“, Bd. I, S. 277–299 und S. 729–795, eine Anzahl Predigten abgedruckt, von denen er anmerkt: „Diese kleinen Sermone sind nur Exordia gewesen zu den Predigten Luther's über die zehn Gebote, welche er an den Fest- und Sonntagen gehalten hat.“ In der St. Louiser Ausgabe finden sich diese Exordia Band XII, 1266–1271; 1724–1829; Bd. X, 198–209; 976–981; 1212–1219; 1276–1297; Bd. XIX, 736–761. Diese Exordien sind es, welche uns einen festen Anhalt geben für die genaue Zeitbestimmung. Luthers Sermon am Tage St. Matthäi (den 24. Februar) 1517 war, wie in unserer Ausgabe, Bd. XIX, 754, bereits bemerkt ist, die letzte Einleitungspredigt zu Luthers Vorträgen über die zehn Gebote. Ueber den Verbleib der lateinischen Handschrift ist nichts bekannt. Aber auch von der deutschen Bearbeitung, welche Luther in seinem Briefe an Johann Lang vom 4. September 1517 erwähnt: „Ich habe Euch die Auslegung der zehn Gebote darum in beiden Sprachen geschickt, auf daß Ihr, wenn Ihr etwa wolltet, darüber dem Volk predigen könntet“ (Walch, St. Louiser Ausg., Bd. XVIII, 28), haben wir keine Spur mehr. Die spätere deutsche Uebersetzung ist, wie aus der Vorrede Münsters ersehen werden kann, nicht von Luther, sondern von einem Anderen. Luther nahm dann den Predigten die Predigtform, und gab seine Bearbeitung im Jahre 1518 in Druck. Die erste lateinische Ausgabe erschien am 20. Juli 1518 bei Johannes Grünenberg (Viridimontanum) in Wittenberg unter dem Titel: Decem praecepta Wittenbergensi praedicata populo per P. Martinum Luther Augustinianum. Noch in demselben Jahre kamen zwei Nachdrucke bei Valentin Schumann in Leipzig heraus, und bei demselben ein Druck im Jahre 1519. In letzterem Jahre ist auch noch von Melchior Lotther in Leipzig ein Druck veranstaltet. In den Sammlungen findet sich unsere Schrift lateinisch: In der Baseler Sammlung vom October 1518, Bl. Qa ff., und in den späteren Auflagen derselben; in M. Lutherii Augustiniani theologi synceri lucubrationum pars una. Basileae apud Adam Petri 1520. mense Julio. S. 157; in der Wittenberger (1550), Tom. I, fol. 1; in der Jenaer (1579), Tom. I, fol. 116b; in Löscher's „Reformations-Acta“ (verbessert und ergänzt nach der vorhererwähnten Handschrift; doch ist auch vieles daraus weggelassen), Bd. I, S. 578; in der Erlanger, exegetica opera latina, Tom. XII, p. 1, und in der Weimarischen, Bd. I, S. 398. Eine deutsche Uebersetzung unserer Schrift erschien zuerst 1520 in Basel bei Adam Petri unter dem Titel: „Der zehn Gebote eine nützliche Erklärung durch den hochgelehrten D. Martinum Luther Augustinerordens beschrieben und gepredigt, Geistlichen und Weltlichen dienende. Item eine schöne Predigt von den sieben Todsünden, auch durch ihn beschrieben.“ Noch in demselben Jahre kam bei ebendemselben eine neue Bearbeitung dieser Uebersetzung heraus. Die erstere wurde abgedruckt durch Silvanus Ottmar in Augsburg (in demselben Jahre), die andere zu Basel im Jahre 1523, vielleicht von A. Cratander. In den Sammlungen zuerst in dem von Zeidler besorgten Hallischen Theil, S. 2, „aufs neu mit Fleiß zugerichtet, mit dem Latein conferirt und nach Luthers Art zu reden gebessert“. Daraus ist der Text in die Leipziger Ausgabe hinübergenommen, Bd. IV, S. 1, und aus dieser in die alte Ausgabe Walchs. Der Anfang der deutschen Uebersetzung findet sich auch bei Löscher l. c. S. 580–586, „die Fortsetzung der deutschen Edition obgedachter Erklärung der zehen Gebote und vieles aus dem Manuscripte hat weggelassen müssen, weil der Platz zu enge wurde“. (Ibid. S. 728.) Der Schluß der Auslegung der zehen Gebote ist auch in einer Einzelausgabe erschienen (ohne Jahr, Ort und Drucker) unter dem Titel: „Eine hübsche Predigt, wie die zehen Gebote gezogen werden auf die sieben Todsünden.“ — Diese Erklärung der zehen Gebote ist auch ins Böhmische übersetzt. Der Druck dieser Uebersetzung wurde am 3. November 1520 zu Prag vollendet. Luther thut dessen Erwähnung in einem Briefe an Spalatin vom 3. Februar 1521 (Walch, alte Ausgabe, Bd. XXI, 746, § 5). In der Baseler Sammlung vom October 1518 und darnach in den lucubrationes des Adam Petri, in der Wittenberger und in der Jenaer Gesamtausgabe ist der lateinischen Erklärung eine kurze Zuschrift an den Leser vorangestellt. Sie ist auch abgedruckt in Löschers „Reformations-Acta“. Aus derselben hat Walch durch Mißverständnis entnommen, „daß Luther weder selbst das lateinische Exemplar verfertigt, noch den ersten Abdruck besorgt habe, ob er wohl, da er zu Wittenberg herausgekommen, mit dessen Bewilligung geschehen sei“. Diesen Ausspruch Walchs hat sich auch die Erlanger Ausgabe angeeignet, und reproducirt ihn in dem 12. Bande der exeg. opera, p. III. Mit Bezug darauf sagt die Weimarische Ausgabe mit Recht: „Vielmehr kann man aus dem Druckfehlerverzeichniß in der ersten Ausgabe sicher auf Luther als Herausgeber schließen, da so wie hier allein der Verfasser selbst ändern durfte.“ Die Weimarische Ausgabe legt ihrem Texte die erste Ausgabe zu Grunde, berücksichtigt aber dabei auch Löscher um seiner handschriftlichen Quelle willen, wiewohl seine Arbeit wenig kritisch ist. Wir geben eine neue Uebersetzung nach der Weimarischen Ausgabe und fügen nach Löscher die Zeitbestimmungen für die einzelnen Predigten ein, was in der Erlanger Ausgabe nicht geschehen ist.}}

Gepredigt von Ende Juni 1516 bis zum 24. Februar 1517. Gedruckt 1518.

Aus dem Lateinischen neu übersetzt.

\section[{Vorrede [des] B[ruder] S[chastian] M[ünster]}]{Vorrede [des] B[ruder] S[chastian] M[ünster].\footnote{Diese Vorrede befindet sich auf der Titelrückseite der beiden bei Adam Petri in Basel im Jahre 1520 erschienenen deutschen Uebersetzungen. Münster, der dem Orden der Franciscaner (Minorum de observantia) angehörte, war damals Corrector bei Petri. Ueber ihn ist im Erlanger Briefwechsel, Bd. II, S. 360, Folgendes mitgetheilt: „Er ist geboren 1489 zu Ingelheim in der Pfalz, studirte zu Heidelberg und Tübingen, daher Melanchthon bekannt, gehörte dem Franciscanerorden an, den er 1529 verließ, und starb 23. Mai 1552 in Basel, wo er lange Zeit an der Universität als Lehrer der hebräischen Sprache gewirkt hatte.“ — Münster sowohl wie Conrad Pellicanus, Guardian des Klosters zu Pforzheim, waren im Jahre 1514 Schüler des gelehrten Rabbiners Elias Levita zu Neustadt an der Aisch im Baireuthischen in der hebräischen Sprache. (Niedere, nützliche und angenehme Abhandlungen aus der Kirchen-, Bücher- und Gelehrtengeschichte. Altdorf 1768, S. 308.)}}

Es kommen täglich hervor durch den Druck viel Bücher, lateinisch und deutsch, deren viel fast heilsam und nütze sind dem gemeinen Manne; wie man denn Menschen findet, die also hungrig und begierig sind nach dem Worte GOttes, daß sie auch alle übrige Zeit, so sie haben mögen, in heilsamen Büchern seliglich anlegen, Trost und Unterweisung daraus zu schöpfen. Denen zu Liebe ist gedeutschet dieses Buch durch den frommen, geistlichen und hochgelehrten Mann, D. Martinum Luther beschrieben, der also geistlich, christlich und evangelisch die zehn Gebote erklärt, daß man desgleichen nicht findet; wiewohl viel Lehrer darüber geschrieben haben. Da finden Geistliche und Weltliche, Gute und Böse, Betrübte und Angefochtene, Junge und Alte, Oberer und Unterthan, kurz, jedermann, wie er den Weg gehen soll (das denn die zehn Gebote sind), der zu dem ewigen Leben führt. Da findet man, wie ein jeglich Gebot nach dem Buchstaben und nach dem Geiste gehalten wird. Dazu, wie ein jeglich Gebot etliche Grade und Staffeln hat etc. Was ist es noth, hier viel davon zu schreiben? Durchlies du das Buch, so wirst du finden einen rechten christlichen und evangelischen Lehrer. Und laß dich nicht hindern, daß es nicht mit verblümten und hochträchtigen Worten verdeutscht ist; denn es mehr den Einfältigen, weder den Höflichen zu Liebe transferirt ist, und also transferirt, daß etwan die Subtilitäten, die der obgemeldete Doctor darinnen hat laufen lassen, dem Latein gelassen sind. GOtt sei mit dir.

\quellseite{574}
\section[Das erste Gebot]{Das erste Gebot.\footnote{Wahrscheinlich begann Luther am Peter-Pauls-Tag (29. Juni) 1516. (Weim. Ausg.)}}

\subsection[Du wirst nicht andere Götter haben]{Du wirst nicht andere Götter haben.\footnote{Vulgata: Non habebis deos alienos. Wir haben den Text der Vulgata beibehalten müssen, weil sich die Auslegung auf denselben bezieht.}}

Zuerst entsteht die Frage, warum gebietet er nicht in bejahender Weise (affirmative), nämlich so: Du sollst den rechten (proprium) oder den Einen GOtt haben, oder: Bete mich, den einigen GOtt, an? Zweitens, warum spricht er nicht vielmehr in Befehlsform (imperative): Du sollst nicht andere Götter haben, als in anzeigender Weise (indicative)?

Auf beides antworte ich zugleich, daß jedes Gebot GOttes mehr um deswillen gegeben ist, daß es die schon vergangene und die gegenwärtige Sünde anzeige, als daß es der zukünftigen wehre, denn, wie der Apostel sagt [Röm. 3, 20.]: „Durch das Gesetz kommt nur Erkenntniß der Sünde“, und wiederum [Röm. 11, 32.]: „GOtt hat alles unter die Sünde beschlossen, auf daß er sich aller erbarme.“

Warum findet das Gebot GOttes, wenn es kommt, Sünder, und macht ihrer mehr, „auf daß die Sünde mächtiger werde“, Röm. 5, 20. Aber die Gesetze der Menschen werden um der künftigen Sünden willen gegeben. Darum redet der [Heilige] Geist, wie er denn ein überaus gütiger Lehrer ist, mehr in anzeigender Weise (indicative), als ob er sagen wollte: O du armer Mensch, siehe, ich zeige dir deine Verderbtheit an. So solltest du beschaffen sein, daß du keine Götter hättest, den Namen deines GOttes nicht unnützlich führest, den Feiertag heiligtest, nicht tödtetest, dich nicht gelüsten ließest etc. Jetzt bist du aber ganz und gar anders beschaffen und verkehrst. Darum gebietet er endlich auch in verneinender Weise (negative), denn eine Verneinung ist stärker als eine Bejahung. Denn auch die Samariter verehrten vor Zeiten den Einen GOtt, aber auch zugleich ihre Götter, und so thun jetzt die Juden und Heiden, die Ketzer und Bösen. Ja, jeder
\quellseite{575}
Mensch in der ganzen Welt verehrt den Einen GOtt, „der allen offenbar ist“, wie der Apostel Röm. 1, 19. sagt. Aber darin sündigen sie, daß sie ihn so ehren, daß sie zugleich auch ihre Abgötter ehren. So ist auch kein Mensch in der Welt, der nicht in einem gewissen Grade ein jegliches Gebot halte (facia[t]), denn es ist niemand, der nicht bisweilen den Namen Gottes nicht\footnote{Die Erlanger Ausgabe hat hier non ausgelassen, wiewohl Luther selbst es in den Berichtigungen der ersten Ausgabe eingefügt hat.} unnützlich führe, der nicht bisweilen nicht tödte, nicht ehebreche, nicht stehle, aber dennoch zugleich auch in allem das Widerspiel thue, wie wir hernach weiter sehen werden.

\subsection{Folgerungssatz (Corollarium)}

Alle Kinder Adams sind Abgöttische und dieses ersten Gebotes schuldig.\footnote{In der Handschrift ist noch hinzugefügt: „Es ist genugsam offenbar, daß sie von GOtt abgewendet sind auf sich selbst.“} Aber man muß wissen, daß es zwei Arten von Abgötterei gibt, eine äußerlich, die andere innerlich. Die äußerliche ist, da der Mensch Holz, Stein, Thiere, Sterne anbetet, wie aus dem Alten Testamente und den Büchern der Heiden bekannt ist. Diese aber ist aus der innerlichen hervorgegangen. Die innerliche ist, da der Mensch, bewogen entweder durch die Furcht vor Strafe oder durch Geiz (amore) seines Nutzens, zwar auswendig das Anbeten der Creatur ansehen läßt, aber inwendig die Liebe zur Creatur und das Vertrauen auf dieselbe bleibt. Denn was ist das für ein Gottesdienst, wenn man vor Reichthümern und Ehrenbezeugungen etc. die Kniee nicht beugt, und ihnen dennoch das Edelste opfert, was man hat, nämlich das Herz und Gemüth? Das heißt GOtt nach dem Leibe und dem Fleische anbeten, aber inwendig im Geiste die Creatur anbeten. Diese Abgötterei regiert in jedem Menschen, bis daß er gesund wird durch die Gnade in dem Glauben an Jesum Christum, wie der 81. Psalm, V. 9. 10., sagt: „Höre, mein Volk, ich will unter dir zeugen; Israel, du sollst mich hören, daß unter dir kein anderer Gott sei, und du keinen fremden Gott anbetest“, was soviel gesagt ist: Durch deine Bemühungen und Kräfte wirst du niemals dahin gelangen, daß du nicht einen fremden Gott anbetest. Denn wenngleich du keine Bilder verehrst, so wirst du doch im Herzen die Creatur mir vorziehen. Aber dann wirst du keinen fremden Gott anbeten, wenn du mich hörst, das heißt, der Glaube an mein Wort wird dich von fremden Göttern frei und zu einem wahren Anbeter GOttes machen. Denn dieser wird dich von dem Begehren der Dinge und von dem Vertrauen auf dieselben abwenden und dich zu dem Schöpfer ziehen.\footnote{In der Handschrift ist hier noch hinzugefügt: „So heißt es im 45. Psalm, V. 11. ff.: ‚Höre, Tochter, schaue darauf‘ etc. Nicht durch deine Werke wirst du schön sein, sondern durch den Glauben an Christum, jedoch wirst du nicht ohne Sünde sein, sondern auch geschmückt durch Christum, dessen Schmuck dein ist durch den Glauben an ihn.“}

\subsection{Wie kann dies geschehen?}

Es geschieht so: Der Glaube an Christum nimmt alle Zuversicht auf eigene Weisheit, Gerechtigkeit und Tugend hinweg, indem er lehrt, daß, wenn er nicht für dich gestorben wäre und dich erhielte, weder du noch alle Creatur dir helfen könnte, und so entsteht die Verachtung aller Dinge. Aber da du hörst, daß er für dich gelitten hat, und glaubst, entsteht alsbald die Zuversicht zu ihm und eine süße Liebe, und auf solche Weise ist alle Neigung zu allen Dingen, als zu etwas Unnützem, vergangen, und es entsteht die Werthschätzung Christi allein, als eines überaus nothwendigen Dinges, und es ist dir nichts Anderes geblieben als Jesus allein, an dem allein du genug und volles Genüge hast, so daß du an allen Dingen verzweifelst und ihn allein hast, auf ihm alle deine Hoffnung steht, und ihn um dess willen über alle Dinge liebst. Aber Jesus ist der wahre, Eine, alleinige GOtt. Wenn du den hast, so hast du nicht einen fremden Gott. Die Juden aber, welche fürchten, sie möchten einen fremden Gott haben, wenn sie den Menschen Christum anbeten, beten in um so schlimmerer Weise einen fremden Gott an, nämlich die Abgötter ihres Herzens, welche sie sich von GOtt erdichten. Mit Recht spricht GOtt also [Ps. 81, 9.]: „Du sollst mich hören“, als ob er sagen wollte: Es ist unmöglich, daß du nicht einen fremden Gott anbeten solltest, wenn du mich nicht hörst; weil du nicht gedemüthigt wirst, wird dir auch das Vertrauen auf die Creatur nicht entrissen. Bist du aber nicht gedemüthigt, so wirst du mich nicht anbeten, sondern dich und das Deine. Es liegt also klar zu Tage, daß niemand dies Gebot erfüllt
\quellseite{576}
als der, welcher glaubt an Christum, auf ihn hofft, ihn liebt, und frei ist von der Neigung zu allen Dingen, was aber ohne die Gnade GOttes unmöglich ist. Aber es gibt viele Leute, die auf die Frage, ob sie einen fremden Gott haben und Abgöttische seien, die bestimmte Antwort geben: Nein, durchaus nicht. Um diese nun auf ihrer offenbaren Lüge zu ertappen, so habe Acht darauf: ob sie den Dingen so abgestorben und in Christo so getrost seien, daß sie weder durch Reichthum aufgeblasen noch durch Armuth kleinmüthig werden, weder durch Ehre sich aufblähen noch durch Schande den Muth sinken lassen, weder am Leben sich ergötzen noch durch den Tod sich erschrecken lassen, weder an Wohl-lust ihre Freude haben noch durch Leiden betrübt werden, und ob sie ganz und gar nach beiden Seiten hin so unbeweglich und ruhig seien, daß sie, wie ihnen auch das Loos fallen möge, ein Genüge daran haben, daß sie Jesum Christum haben; ja, daß sie Armuth, Schande, Leiden, Tod lieben und wünschen, Reichthum, Ehre, Wohllust, das Leben verschmähen aus Verlangen nach ihrem Gotte Jesum Christo, wie es heißt im 42. Psalm, V. 2.: „Wie der Hirsch schreiet nach frischem Wasser“ etc., und die Braut im Hohenliede [Cap. 5, 8.]: „Ihr Töchter Jerusalems, saget meinem Freunde, daß ich vor Liebe krank liege“, und Paulus [Phil. 1, 23.]: „Ich habe Lust abzuscheiden und bei Christo zu sein.“ Wer aber diese Gesinnung noch nicht bei sich befindet, der bekenne, daß er noch nicht ein reiner Verehrer GOttes sei, sondern daß ihm noch viel Abgötterei anhänge, denn es thut ihm leid, er wünscht, er liebt, er haßt, nicht wie er soll, noch das, was er soll. Aber hier sagt man: Dies kommt den Vollkommenen zu, nicht allen; eine solche Vollkommenheit ist nicht vonnöthen. Ich antworte: Auch wir wissen, daß dies die Vollkommenen angeht (das heißt, nicht die Juden, sondern die Christen),\footnote{Löscher: Hier schiebt die Handschrift mehrere und zwar sehr nützliche Dinge ein.} nicht als ob alle die verdammt wären, welche nicht so vollkommen sind, sondern dieses Ziel und Ende ist uns vorgesteckt. Niemand, der es nicht erreicht, wird entschuldigt werden, es sei denn, daß er mit Seufzen erkennt und bekennt, er sei nicht so beschaffen, und sich täglich bemüht, ein solcher Mensch zu werden, und demüthig bittet, daß ihm das vergeben werde, was er zu wenig thut, indem er spricht [Matth. 6, 12.]: „Vergib uns unsere Schulden“, und [Ps. 51, 12.]: „Schaffe in mir, GOtt, ein reines Herz.“ Dieselben, sage ich, die in Furcht stehen und bekennen, die da suchen und bitten, wird diese ihnen noch anhängende Abgötterei (idolatriae suae mixtura) nicht zugerechnet um Christi willen, an den sie glauben. Aber denen, welche ohne Furcht, ohne es sich angelegen sein zu lassen, daß sie besser werden, in Sicherheit schnarchen, wird sie völlig zugerechnet, und sie sind in Wahrheit Abgöttische. Es wird ihnen auch die Entschuldigung nicht helfen, daß es nicht vonnöthen sei, vollkommen zu sein, als ob dieses Gebot Steinen oder Holz und nicht vielmehr den Menschen gegeben wäre, und so völlig und vollkommen erfüllt werden müßte, daß auch nicht der kleinste Buchstabe noch ein Tüttel vergehen wird [Matth. 5, 18.].

Um klareren Verständnisses willen, so daß auch Schwächere fassen können, wer, und wie man wider dies Gebot sündige, will ich in sehr ausführlicher Weise Beispiele anzeigen, jedoch nicht alle (dann müßte ich Zahlloses aufzählen), sondern viele, damit durch den Vergleich mit diesen auch die übrigen leichter verstanden werden können.

Erstens. Es sind etliche so unverständig, daß es scheint, daß der Teufel sie nicht im Ernst, sondern gleichsam als seine Narren im Scherze verspottet. Zu diesen gehören die Wahrsager, die Zeichendeuter, die Zauberer, die Seegensprecher, die Abergläubigen,\footnote{In der Weimarischen Ausgabe ist (wie wir meinen, nicht gut) das Komma vor superstitiosi ausgelassen.} und von diesen gibt es vielerlei Arten; aber ich will sie zusammenfassen und sie nach den verschiedenen Lebensaltern unterscheiden, damit man es leichter behalten möge.

\section[Das erste Lebensalter ist die Jugend, die an sich]{Das erste Lebensalter ist die Jugend, die an sich schon dazu geneigt ist, sich vom Teufel verführen zu lassen}

Erstlich gibt es solche Leute, welche mit bestimmten Formeln (conceptis verbis) und gewissen Zeichen Schwerter, Geschoße, Büchsen,\footnote{Weimarische und andere Ausgaben: bardas, Büttenberger, Zenaer und Löscher: bombardas.} und alle eisernen Waffen zu segnen, das heißt, zu beschwören verstehen, so daß sie von denselben nicht verwundet werden können, sei es nun im
\quellseite{577}
Kriege oder sonst. Ich selbst habe einen Jüngling gesehen, der ein gezogenes Schwert auf seinen bloßen Leib (ventrem) setzte und es mit solcher Gewalt gegen sich drückte, daß das Schwert sich bog und auf seinem Leibe das Heft an die Spitze stieß (donec capulus reflexus copularetur acumini etc.), und er ließ das Schwert wieder zurückgehen [in seine ursprüngliche Lage], war aber nicht verwundet.

Zweitens [gibt es Leute], welche dieses [Beschwören] entkräften können durch Beschreibung eines Zirkels und Zeichen von Figuren in den Sand, und indem so einer des andern Beschwörungen unkräftig macht, tödten sie einander.

Drittens gibt es Leute, welche Buchstaben oder Zeichen auf Glas schreiben und aus Feuer stellen, und durch dies Vornehmen (machina) pflegen sie die Mädchen dazu zu zwingen, daß sie sie lieben müssen. Aber diese machen es auch so. Etliche können ihre Liebsten (Adonides) über viele Meilen weit rufen, indem sie einen Topf ans Feuer stellen. Andere nehmen den Kopf eines Esels und legen ihn mitten in das Feuer\footnote{oder: mitten auf den Feuerherd (in medio foco).} und zwingen [dadurch] von denen, die ihrer Kunst (mysterii) anhängen oder derselben theilhaftig sind, so viele als sie wollen, zu ihnen zu kommen. Und wunderbar ist es zu sagen, daß diejenigen, welche [so] gezwungen werden zu kommen, sich [davon] nicht helfen noch retten können, bis daß es ihnen vergönnt ist, sich des Feuers zu bemächtigen, an dem der Kopf des Esels dem Feuer ausgesetzt wird (coctore).

Viertens hat man Briefe, die voll sind von vielen heiligen Worten und Zeichen, von denen man sagt, der Papst Leo habe sie dem [Kaiser] Carl in den Krieg geschickt. Außer dem, daß dies ein wichtiges Vorgeben ist, ist es auch eine Lüge, wie aus den Chroniken erhellt, da Leo und Carl nicht Zeitgenossen waren. Dennoch schreiben sie ihnen soviel zu, daß derjenige, welcher sie trage, weder getödtet noch verwundet werden noch eine Beschädigung erleiden könne, so daß sie sich dessen vermessen, auch in ungerechten Sachen vor Gericht zu gewinnen.

Fünftens ermahnen sie gewisse Monate für das Heirathen. Einige thun dies, damit ihnen keine Kinder geboren werden, andere dagegen, um viele Kinder zu bekommen. So schreiben sie in einer unseligen Narrheit nicht Gotte die Schöpfung des Menschen zu, sondern den Monaten, und dennoch gibt es Bücher, die voll sind dieser Possen.

Sechstens. Das geweihte Wachs der Osterkerze wird von vielen verstohlener Weise abgerissen, ich weiß nicht zu wie vielen und wie großen Zaubereien und abergläubischen Handlungen.

Siebentens. Es gibt Leute, die Weihwasser nehmen, und hoffen mit ihren Gebeten in einem Ameisenhaufen ein Ei zu finden, welches sie unsichtbar machen soll, wenn sie es bei sich tragen.

Achtens. Sie suchen einen Dieb [zu entdecken], indem sie die Namen von verdächtigen Leuten auf Zettel schreiben und in Weihwasser werfen, so daß der, dessen Name zu Boden fällt, der gesuchte Dieb sein soll.

Neuntens. Bei der Hochzeit wollen sie durch Beobachtung thörichter Dinge in Erfahrung bringen, wer von beiden den andern überleben werde, als, wer in der ersten Nacht zuerst einschläft, der müsse zuerst sterben. Desgleichen soll die Frau, wenn sie zum ersten Male in das Haus des Mannes geht, mit der Hand die obere Thürschwelle anrühren und dabei Worte\footnote{Der Spruch lautet nach einer handschriftlichen Notiz: „Ich greiff auff die überthür, mein krieg der gee albeg für.“ (Weim. Ausg.)} sprechen, durch welche sie weissagt, daß sie die Herrin des Mannes sein werde.

Hiernach können ähnliche Dinge beurtheilt werden. Wiewohl diese Dinge mehr thöricht als böse sind, und es scheint, wie ich gesagt habe, daß der Teufel damit einen Scherz treibe, so wird doch endlich ein gar großer Ernst daraus, weil sich nämlich die Leute dadurch gewöhnen, in kleinen Dingen auf die Creatur ihr Vertrauen zu setzen, und daher auch in großen Dingen nicht auf GOtt vertrauen. Sie entschuldigen sich freilich damit, daß es heilige Zeichen und heilige Worte seien, als ob sie um deswillen nicht um so mehr anzuklagen wären, daß sie die heiligen Dinge, Zeichen und Worte GOttes gebrauchen zu unheiligen und teuflischen Sachen.\footnote{Löscher: „Hier wird in der Handschrift viel Nutzbares hinzugefügt.“}

\footnote{Löscher: „Das Folgende ist zu der Predigt hinzugefügt, die im Jahre 1516 am Tage der Heimsuchung Mariä [2. Juli] gehalten worden ist, wie die erwähnte Handschrift bezeugt, aus der das Folgende geschöpft ist.“ — Diese Einleitungspredigt findet sich Walch, St. Louiser Ausgabe, Bd. XII, 1722.} Das zweite ist das männliche Alter und derer, die schon in der Ehe sind, wo die all-

\quellseite{578}
zuzärtliche Liebe zu den Kindern und das Hangen am irdischen Gut bewirkt, daß die Weiblein wunderbar leicht zu diesem Werke des Teufels verführt werden können.

Erstens. Sie unternehmen es, die Kinder durch gewisse abergläubische Gebräuche gesund zu machen. Denn es gibt Leute, welche die bezauberten Kinder (denn dies Uebel ist bei uns sehr häufig, so sind auch gar viele böswillige [pestilentes] alte Weiberlein und Teufelsdienerinnen vorhanden, die solches vornehmen) heilen von einer unbekannten Krankheit, die man auf deutsch „die Elbe“\footnote{Die Elbe oder „das Herzgespann“. Diese Krankheit ist genauer beschrieben im dritten Capitel der kürzeren Auslegung des Briefes an die Galater, Walch, St. Louvier Ausg., Bd. VIII, 1462, § 3. — Hier werden „Elbe“ und „Herzgespann“ als zwei Krankheiten unterschieden.} nennt. Für diese schreiben etliche Dorfpfarrer und Küster, welche dies von den Zauberhuren gelernt haben, Briefe, die an den Hals gehängt werden, aber nicht mit gewöhnlicher Tinte, auch nicht auf Papier, noch an jedem Tage; ich weiß nicht, vielleicht auch nicht mit einer Feder, und an einem Schreibpult und an einem gewöhnlichen Orte. Denn es hat diese Sache ihre verborgenen Geheimnisse, und doch geschieht, was sie thun, weil sie glauben. Andere [vertreiben] eine andere Krankheit, welche man auf deutsch „das Herzgespann“ nennt, das heißt eine Ausdehnung der Brust. Das Vorhandensein dieser Krankheit beweisen sie durch diesen Grund, nämlich wenn bei dem Kinde das Maß von einem Ellbogen zum andern nicht gleich ist dem Maße vom Knie bis zu dem Halse. Ferner stillen sie Kopfschmerzen, ich weiß nicht durch wie viele Arten von Gewurmel und Messen mit dem Gürtel; so auch andere Schmerzen an allen Gliedern. Hier möchte ich, daß die Kinder Verstand hätten und reden könnten, damit sie ihre Mütter wegen ihrer überaus großen Narrheit strafen könnten. Denn es ist kein Zweifel, daß sie sich dem widersetzen würden und an den Tag geben, daß sie mehr Einsicht hätten als ihre Mütter.

Sie sagen aber: Wer sollte sich seines kleinen Kindes und der Frucht seines Leibes nicht erbarmen? Ich antworte: Ja, man soll sich derselben erbarmen, aber nicht soweit, daß man sich darüber in des Teufels Dienst begebe; entweder suche natürliche Arznei, oder bitte GOtt in einfältigem Glauben. Warum opferst du das Kind, welches GOtt geschaffen hat, dem Teufel, daß er es erhalte und gesund mache? Aber GOtt läßt zu, daß durch Zauberer dergleichen an den Kindlein geschehe, um die Ungläubigen zu versuchen, damit die den Teufel suchen, welche GOtt verlassen haben. Denn es ist nicht zu bezweifeln, daß dergleichen Plagen der Kinder bisweilen nicht natürlich sind, sondern die Teufel, beschworen durch die Zauberer, peinigen in solcher Weise die Eltern derselben aus einem verborgenen aber ganz gerechten Gerichte GOttes, um ihren oder den Glauben anderer zu prüfen, oder um beide zu strafen.

Zweitens blasen die Väter den kleinen Kindern gegen den Gaumen oder in den Hals, die an sich weiß nicht was für einer Krankheit leiden; sie sprechen aber siebenmal die Worte, welche zur Beschwörung dienen.

Drittens [gibt es Leute], welche ihr Vieh zeichnen und es so bewahren vor Wölfen, Pest, Wasser und Feuer, ja, sie können in solcher Weise eine ganze Stadt oder ein Haus vor Feuersbrunst bewahren. Der unglückliche Wolf kann sich mit Recht beklagen, daß ihm sein von GOtt beschiedenes Theil durch des Teufels Verbündete entrissen sei. Wenn Hiob so gehandelt hätte, würde er sein Vieh behalten haben, aber GOtt würde ihn nicht gelobt haben.

Viertens. Aus derselben Andacht versehen sie die Früchte auf den Feldern und in den Gärten mit Zeichen, und sind so thöricht, daß sie, die doch von GOtt Regen und Gedeihen empfangen haben, den Teufel um seine Hut bitten, indem sie gleichsam wider GOtt streiten, damit er ihnen das nicht wegnehme, was er gegeben hat.

Fünftens. Feuer und Wasser, auch Schlangen verstehen sie zu beschwören, damit sie dem Vieh auf keine Weise Schaden thun, das heißt, daß ihre Sache Gotte nicht unterworfen sei.

Sechstens. Wenn Kranke den Heiligen ein Gelübde thun, so halten sie es darnach für Sünde, Arzneien zu gebrauchen, natürlich, damit sie GOtt versuchen mit ihren thörichten Gelübden. Und wenn ein Weib das Sacrament der letzten Oelung empfangen hat, so darf sie während des Jahres nur schwarze Kleider tragen, nicht im Reigen tanzen, nicht bei ihrem Manne schlafen. Ein so vertrautes Werkzeug hat der Teufel an dem weiblichen Geschlechte, daß er dasselbe annimmt zu Ausrichtung seines Dienstes (in sua sacra), durch dasselbe Gesetze vorschreibt
\quellseite{579}
, seinen Aberglauben aussäet; in allen Stücken GOtt zuwider, der seinen Gottesdienst (sacra) und sein Priesterthum und sein Wort den Männern auferlegt. Aber dieses Priesterthum der Weiber geht viel stärker im Schwange, und hat diese Lande mit zahllosem Aberglauben, [Zauber-]Segen und nichtsnutziger Lehre erfüllt, vor denen sie sich weit mehr fürchten als vor den Gesetzen und Gottesdiensten des Priesterthums der Männer und GOttes.

Siebentes. Hieher gehören diejenigen, welche achten auf Vorbedeutungen (ominationes), Vogelgeschrei, als das Krächzen der Raben, und die Ueberbleibsel oder Exempel von anderen heidnischen Wahrsagereien (portentorum).

Achtens. Die da besondere Tage (dies Aegyptios) beobachten für ihre Reisen, für ihre Rückkehr, beim Bauen, bei Geschäften, beim Anziehen neuer Kleider. Wenn diese den Glauben hätten, daß das Feuer im Winter angenehmer ist, als im Sommer, desgleichen, daß im Sommer die rechte Zeit ist Heu zu machen, die Früchte einzusammeln etc., dann hätten sie einen rechten Glauben; oder, daß es für jeden Menschen gut ist, früh auszugehen, wenn er gesund ist und es für ihn vonnöthen ist, aber nicht gut, wenn er krank ist. Denn was verdient dieser Spötter, der Teufel, anders, als daß wir ihn auch wiederum verlassen? Diesen stehen diejenigen nahe, welche uns lehren, daß das Aderlassen, das Entwöhnen der Kinder, die Heilsamkeit der Bäder, die rechte Zeit zum Säen nicht zu bestimmen sei nach Gelegenheit der Feuchtigkeit und Wärme oder dem Bedürfniß der Leiber, sondern nach dem Einfluß der Planeten und der Sterne, wiewohl die besten Aerzte lehren, daß hier die Noth die einzige Regel sei. Doch vielleicht haben sie vor, uns durch Hunger zu tödten, da sie uns nur an wenigen Tagen gestatten zu säen. Aber es ist gut, daß diese Possen von niemand oder doch nur von sehr wenigen beachtet werden (da GOtt besser für uns sorgt).

Neuntens. Es folgt die herrliche Astrologie oder Mathematik, die gerne eine Wissenschaft sein wollte, aber ihre angeborne Thorheit nicht ausziehen kann. Sie ist es, die uns lehrt, wie beschaffen oder wie groß ein jeglicher sein werde, der geboren ist bei gewisser Stellung der Himmelszeichen (in horoscopis signorum), nämlich sie sind Rathgeber bei dem göttlichen Geheimnisse, welches auch den Engeln nicht bekannt ist. Nur über Eines wundere ich mich, was ihnen zugestoßen sein mag, daß sie den Stern nicht gefunden haben, der da anzeige, wer als ein Gerechter, und wer als ein Sünder werde geboren werden. Denn sie wollen ja, daß die von ihnen beobachtete Stellung der Himmelszeichen (suos horoscopos) auf Menschen überaus großen Einfluß übe; sodann ist Gerechtigkeit, Sünde, Wahrheit, Lüge nicht eine Sache von so geringer Bedeutung, aber auch nicht so selten als ein Bader, ein Sänger, ein Goldwechsler, ein Fischer, ein Redner, ein Buhler, welche ihre eigenen Nativitäten (horoscopos) haben. Warum ist denn da kein Zeichen für die Gerechtigkeit und Wahrheit? Oder wenn es da ist, warum läßt es niemals seine Wirkung spüren? Denn es wird ja ein jeglicher Mensch geboren als ein Sünder, ein Lügner, ein Thor, wiewohl kein Stern dazu gesetzt ist, dieses Geschick herbeizuführen (ad hoc fatum); er wird auch nicht\footnote{In fünf Ausgaben, die uns außer der Weimarschen vorliegen, ist noch unquam hinzugefügt.} geändert, wenn er nicht von der über dem Himmel befindlichen Gnade heimgesucht wird. Oder ist der Himmel gegen die Gerechtigkeit und Wahrheit so feindselig und gleichgültig, daß er seinen Einfluß übt auf ganz geringfügige Dinge, Bäder und Spiele und Viehbehändel, aber ganz und gar nichts weiß von der Gerechtigkeit? Oder ist der Schöpfer so neidisch, daß er kein Zeichen für das Gute, sondern nur für das Böse verordnet hat? Denn es wird ja kein Mensch gut geboren und wird seiner Natur nach mit den Einflüssen böse bleiben. Diese würde ich unter die seinen (subtiles) Narren gerechnet haben, wenn sie nicht gröber wären als die groben.

Aber sehr schön lösen sie die Einwürfe auf, indem sie sagen: Die Einflüsse zwingen (necessitant) nicht zur Sünde, sondern machen nur geneigt dazu etc., gleich als ob nicht gerade dies eine überaus gottlose Meinung wäre, daß GOtt die Creatur gemacht habe, um zur Sünde geneigt zu machen, und nicht vielmehr, damit er die Gerechtigkeit aufrichte, so daß alles für die Menschen mitwirkte zum Guten, nicht zum Bösen; oder gleich als ob irgend ein Mensch sündige, gedrungen durch Nothwendigkeit, und nicht vielmehr immer aus Neigung. Wer möchte sagen, daß ein Mensch wider seinen Willen sündige? Alle böse Neigung ist nicht außer uns,
\quellseite{580}
sondern in uns, wie Christus sagt [Matth. 15, 19. 11.]: „Aus dem Herzen kommen arge Gedanken“, „was zum Munde eingehet, verunreiniget den Menschen nicht“ etc. Und St. Jacobus [Cap. 1, 14.]: „Ein jeglicher wird versucht, wenn er von seiner eigenen Lust gereizet und gelocket wird“, was nicht von dem Schicksal, sondern von der Erbsünde (origine peccati) herkommt. Denn alles, was GOtt gemacht hat, ist gut [1 Mos. 1, 31.], darum können alle Dinge ihrer Natur nach nur eine Neigung zum Guten veranlassen. Wie ein jegliches Ding beschaffen ist, so wirkt es auch. Daß aber die Dinge zum Bösen dienen, geschieht nicht aus ihrer Natur, sondern durch ihren Mißbrauch (injuria), wie Paulus sagt [Röm. 8, 20.]: „Alle Creatur ist der Eitelkeit unterworfen ohne ihren Willen.“ Sie aber machen die Natur derselben zur Eitelkeit, indem sie wollen, daß diese Dinge durch GOttes Ordnung (institutione) das an sich haben, daß sie eine Neigung zum Sündigen geben. Warum haben sie dem Adam und der Eva, ehe die Schlange zu ihnen kam, diese Neigung nicht eingeflößt? Warum nicht Christo? Warum nicht der heiligen Jungfrau? Fort mit dieser lästerlichen Gottlosigkeit! Jenes Wort [welches die Astrologen zu ihrer Entschuldigung vorbringen] ist das Wort unseres Vaters Adam, der auch seine Neigung [zur Sünde] auf das Weib schob, das heißt, auf die Creatur GOttes, und sagte [1 Mos. 3, 12.]: „Das Weib, das du mir zugesellet hast“ etc. Aber wie gewaltig würden sie sich wider mich legen, wenn sie nachweisen könnten, daß irgend ein Heiliger oder Märtyrer sich dieser Dinge bedient habe, oder darüber geschrieben, oder sie gutgeheißen! Nun aber haben sie dieselben nicht allein nicht gebilligt, sondern auch verworfen, besonders St. Augustinus an vielen Stellen, sodann auch St. Gregorius. Und dennoch findet diese faule Narrheit noch fadere Leute, welche sie glauben.

Abraham, sagen sie, lehrte die Egypter die Astrologie, wie Josephus bezeugt,\footnote{Josephus, Antiqu., lib. I, cap. 8, § 2. — Weim.: Josephus.} als ob Josephus nirgend von der Wahrheit abgewichen wäre, zumal er so eifrig darauf aus ist, Ruhm für die Juden zu erjagen. Abraham hat ohne Zweifel die Egypter gelehrt, GOtt zu verehren, und die rechte Weisheit GOttes, wie auch von Joseph gesagt wird, Ps. 105, 22.: „Daß er seine Fürsten unterwiese nach seiner Weise, und seine Aeltesten Weisheit lehrte.“ Es ist auch nicht zu glauben, daß jene heiligen Männer sich auf die Astronomie gelegt haben, viel weniger aber auf die Astrologie, denn das sind überflüssige (otiosa) Studien junger Leute. Aber da Josephus sah, daß die Griechen eine hohe Meinung von dieser Wissenschaft hatten, und dieselbe in Ehren stand, scheint es, daß er hat erdichten wollen, wie die Juden auch in dieser Sache nicht niedriger, sondern höher ständen als die Griechen, was er auch in allen anderen Dingen zu thun sich bemüht hat, die zur eitlen Ehre dienen.

Sodann ist es wunderbar, daß die bösen Geister nicht endlich einmal geändert worden sind, welche den Einflüssen der Gestirne so viele Jahrhunderte lang unterworfen gewesen sind, dann aber auch ihnen näher sind als wir, da sie ja in der Luft wohnen. Daß aber der Einfluß der Sterne nicht über sie gekommen sei, oder daß die Sterne nicht auf dieselbe Weise ihnen geschienen haben sollen, das scheint ungereimt. Denn sie bleiben hartnäckig in ihrer Beschaffenheit, werden auch durch kein Zeichen verändert, während man doch sagt, daß unsere Seelen durch die kürzeste Einwirkung der Einflüsse verändert werden.

Zum letzten. Was sollen wir zu Mose sagen, der im fünften Buche, Cap. 4, 19.\footnote{Weimarische Ausgabe ebenso wie Ausg.: 5 Mos. 4, 10.} spricht: „Daß du nicht deine Augen aufhebest gen Himmel, und sehest die Sonne und den Mond, und die Sterne, das ganze Heer des Himmels, welche der HErr, dein GOtt, verordnet hat zum Dienst allen Völkern“ etc. Wenn „zum Dienst“, wie kann es denn sein zur Herrschaft? Aber auf ganz seine Weise machen sie eine Ausflucht und sagen, gestützt auf das Ansehen ihres Meisters Ptolemäus: Der Weise herrscht über die Sterne, daher kann er dem Einfluß der Sterne zuvorkommen und ihn verhindern. Also sind sie nicht den Weisen zur Herrschaft, sondern nur den Unweisen.\footnote{Diesen Satz hat die Weimarische Ausgabe, Löscher und Petri als Fragesatz.} Was sagt aber Moses? Er sagt: „Allen Völkern zum Dienst.“ Wenn „allen“, so sind entweder alle Völker weise, und so herrschen alle über die Gestirne, oder es sind nur etliche weise, und dann wird es so kommen, daß
\quellseite{581}
sie nicht allen Völkern zum Dienst geschaffen sind. Aber Moses ist wahrhaftig und du bist ein Lügner,\footnote{Vel Moses verax et tu mendax. Dieser Satz ist schon den früheren Uebersetzern eine crux gewesen. Zeidler gibt ihn so wieder: „Entweder Moses leugt oder du leugst.“} nun davon nichts zu sagen, daß, wenn auch der Wahrheit Weise ein Herr wäre über die Gestirne, Moses nichtsdestoweniger etwas Falsches gesagt hätte, da ja die Gestirne auch den Weisen nicht dienen, sondern vielmehr sie verhindern, so daß sie der Schicksalsbestimmung derselben nicht entgehen könnten, wenn sie nicht den Vorzug der Weisheit hätten. So hätte also gesagt werden müssen, die Sterne seien geschaffen, nicht zum Dienst für die Weisen, sondern zum Kriege [gegen sie], und zur Herrschaft und Tyrannei über die Unweisen. Das heißt aber Moses der Lüge beschuldigen und GOtt der Grausamkeit anklagen und geradezu lästern. Aber dies haben andere weitläufiger behandelt. Es möge genug sein, daß wir angezeigt haben, daß diese nichtige Sache nicht erlaubt sei. Beim Propheten Jeremias, Cap. 10, 2., [spricht der HErr]: „Ihr sollt nicht der Heiden Weise lernen,\footnote{Weimarische: dicere statt discere.} und sollt euch nicht fürchten vor den Zeichen des Himmels.“ Denn allein GOtt soll man in allen Dingen fürchten. Alle anderen Dinge sollen wir so ansehen, daß sie zum Dienst verordnet sind und den Auserwählten zum Besten dienen [Röm. 8, 28.].

\footnote{Löscher: „Das folgende ist an eine Predigt angefügt gewesen, welche am 7. Sonntage nach Trinitatis [den 6. Juli 1516] gehalten worden ist.“ Diese Einleitungs-Predigt findet sich Walch, St. Louiser Ausg., Bd. XI, 1724. Dort sollte die Ueberschrift heißen: „Am siebenten Sonntage nach Trinitatis“, denn Luther folgt in diesen Predigten der consuetudo ecclesiae Romanae, nach welcher Matth. 7, 15. ff. das Evangelium für diesen Sonntag ist. (Weim. Ausg., Bd. I, S. 61).} Das dritte Lebensalter ist das der alten Hexen (vetularum), oder derer, die sich mit ähnlichen Dingen wie diese abgeben, die mit dem Teufel ein Bündniß machen, von denen man überall Kunde hat.

Erstens. Durch Zaubereien können sie die Augen beschädigen und blind machen, die Leiber krank machen, die Beine lähmen,\footnote{sagittare = mit Pfeilen schießen.} durch Erscheinungen (imaginibus) bezaubern, und, wie es ihnen beliebt, entweder tödten oder durch eine langsame und unheilbare Krankheit endlich abzehren, wie ich mehrere gesehen habe, die solches erlitten.

Zweitens. Sie können Ungewitter und Donner erregen, die Früchte verderben, Vieh tödten, desgleichen anderen Butter, Milch und Käse stehlen, das heißt, aus einem Pfosten oder einer Axt oder einem Handtuch melken.

Drittens. Hieher gehören diejenigen, welche ein Wissen [von heimlichen Dingen] durch die Notariatskunst\footnote{Wir haben nach Löscher die Lesart Notariam angenommen; in den andern Ausgaben: Notarium. Daß die erste Lesart vorzuziehen sei, scheint uns der Zusatz zu beweisen: de qua in jure satis dicitur. Ars notaria ist die Kunst, schnell und mit Abbreviationen zu schreiben. Löscher merkt an: „Es wird darunter eine Art der Cabala verstanden, welche Notaricon genannt wird.“} zu erlangen trachten, über welche in den Rechtsbüchern (in jure) gar viel gesagt wird. Desgleichen, welche verborgene Dinge suchen durch ein Crystall, einen Fingernagel oder durch ein elfenbeinernes Heft,\footnote{capulus = Degenriff.} welches jedoch zuvor geweiht sein muß. Doch können sie dies nicht, wenn nicht Mädchen oder Knaben, die noch jungfräulich sind, in einen Kreis eingeschlossen sind, damit diesem überaus bösen Werke ein guter Schein nicht mangele. Dennoch sagt man, daß diese Kunst des Erforschens (inspectionis) sehr oft fehlschlage. Hieher gehören diejenigen, welche mit einer Wünschelruthe verborgene Schätze suchen, besonders aber diejenigen, welche man gewöhnlich kluge Männer und Frauen nennt (das ist, Wahrsager und Wahrsagerinnen), deren Aussprüche (oracula) die Leute sich von verschiedenen Orten und weit über Land holen, wenn jemand etwas verloren oder irgend etwas erlitten hat, und nicht weiß, wer diesen Schaden angerichtet hat. Es ist zu verwundern, daß die Bischöfe es ungestraft hingehen lassen, daß man sich öffentlich solcher Dinge unterstehe.

Viertens. Viele glauben, daß sie auf Besen\footnote{Wittenberger: scopa, was die richtige Lesart ist, denn Löscher merkt an: „im Manuscript heißt es: die auf Besen reiten.“ Baseler und Jenaer: scobam; Löscher, Erlanger und Weimarische: scobem.} reiten, andere auf einem Bocke, andere auf anderem Gaukelwerk (nugas), zu einem gewissen Orte, an welchem diejenigen zusammenkommen, welche dieser verborgenen Kunst (mysterii) theilhaftig sind, um eine Gasterei zu halten. Es ist nicht allein unerlaubt, dies zu thun, sondern auch zu glauben, daß dies wahr sei, wie nachher weitläufiger gesagt werden soll, gleichwie auch daß,
\quellseite{582}
daß sich die alten Hexen in Katzen und Kater verwandeln und bei Nacht umherschweifen, nicht geglaubt werden darf, daß es wahr sei.

Fünftens wird gesagt, daß das ganze Jahr hindurch (per annum) eine gewisse Frau, welche einige die Herodias, andere die Frau Hulde, noch andere die Venus nennen, bei etlichen herumfahre und gleichsam ein regelmäßig wiederkehrendes Reinigungsfest (lustralem februum) halte, die bei ihren Wirthen entweder Kohlen oder Kleinigkeiten von ihrem Wagen zurücklasse, welche nachher als Gold und Silber erfunden werden, wie auch in der lateinischen Sprache ein Sprichwort vorhanden ist von Kohlen und Golde,\footnote{Bei Phaedrus, lib. V, fab. VI, findet sich das Sprichwort: Carbonem pro thesauro invenire.} daß vielleicht hieraus entstanden ist.

Sechstens. Etliche haben eine Art Hausteufel, wie man vor Zeiten die Hausgötter hatte, die bisweilen bei Tage erscheinen. Diese heißen im Munde des Volks (vulgo) bei einigen „Wichtlein“, bei anderen „Hehlkäpplein“,\footnote{In den alten Ausgaben geschrieben: „Wichtlein“ und „Heletäpplein“. Hehlkäpplein heißen sie, weil sie sich durch ein Käpplein unsichtbar machen. Vgl. Walch, St. Louiser Ausgabe, Bd. XX, 2112, § 7.} und man glaubt, daß in einem Hause großes Glück sei, welches von solchen Teufelsgeistern bewohnt wird, und man fürchtet sich mehr, diese Teufel zu beleidigen, als GOtt und alle Welt.

Siebentens. Einige gehen, um den Teufel zu finden, rückwärts um eine Kirche herum, und gerade entgegengesetzt, wie man sonst zu thun pflegt, und wenn sie ihn finden, daß er ihnen entgegen geht, so ergeben sie sich ihm, und machen einen Bund mit ihm. Hieher sind auch die zu rechnen, welche es für ein übel Zeichen halten, wenn ihnen früh Morgens ein Priester begegnet, denn sie bezeichnen sich mit dem Kreuze, damit ihnen an dem Tage nichts Uebeles widerfahre, dessen sie sich versehen, wie der Teufel sie lehrt, daß es durch das erste Begegnen des Priesters eintreten werde.

Achtens. Es gibt Leute, die mit den Teufeln in der Gestalt von Männern (incubi) und Weibern (succubi) zuhalten,\footnote{Vgl. Tischreden, Cap. 24, § 94 f. St. Louiser Ausg., Bd. XXII, 754 ff.} wovon ich sogleich weitläufiger handeln werde.

Neuntens. Diese Bosheit des Satans geht sogar so weit, daß die Sacramente geschändet werden. Denn es gibt Leute, die den Kindern, die getauft werden sollen, Eisenkraut, Käse und ich weiß nicht was für andere Dinge anbinden. Und zwar ist dieses Eisenkraut das, was am meisten zu abergläubischen Handlungen gebraucht wird. Mit wunderlichen Ceremonien graben sie es aus, lassen es dann weihen, und mit erschrecklicher Frevelhaftigkeit rufen sie gar viele (tot) Namen GOttes und der Heiligen darüber an, natürlich verführt durch irgend einen Juden, der mit Blendwerk umgeht (praestigioso).

Ich will hier nichts sagen von der Reinigung (menstrua) der Weiber, aus welcher etliche [Herren] Tränke herstellen, die eine unsinnige Liebe hervorrufen, öfter aber noch den Tod verursachen. Einige mischen sie ins Schweinefutter, damit sie leichter fett werden. Ich lasse anstehen, daß sie den Kindern, die am Schminden (tabe = Auszehrung) oder Magerkeit leiden, so zu rathen pflegen, daß sie die Kinder in einen Kessel stellen, Feuer darunter legen, und in warmem Wasser mäßig kochen, indem ein Weib dabei steht, eine andere aber dreimal um das Haus läuft und durch eine [Thür- oder Fenster-]Oeffnung fragt: Was kochst du? und jene darauf antwortet: Ich koche altes Fleisch, damit es neu werde. Desgleichen, daß sie sich fürchten, die entwöhnten Kinder wieder an die Brust zu legen, wegen der Gefahr, daß sie davon Flucher werden möchten oder schändliche Reden führen. Ich sage nichts von den abscheulichen Dingen, die sie in Geburtsnöthen vornehmen. Aufs andächtigste lesen sie die Legende der heiligen Margaretha und ziehen sie sogar dem Leiden Christi vor. Andere hängen der Gebärenden die Hosen des Mannes an den Hals, und erweisen ihren Wahnwitz noch mit anderen Possen.

Außerdem sind etliche Weiber so thöricht, daß sie das, was etwa im Scherz gesagt wird, für Ernst nehmen, als: Wenn du etwas verloren hast und es wiederfinden willst, so gib Almosen, nämlich eine warme Semmel, und thue geweihtes Salz darauf, so wirst du es alsbald finden. Dies hat jener im Scherze so verstanden wissen wollen: Thue Salz darauf, nämlich auf die verlorene Sache, dann wirst du sie finden. Sie aber verstehen es: „auf die Semmel“. Nicht als ob dies etwas Großes wäre, sondern es ist etwas Ungeziemendes, daß der Teufel ein solch Gespött aus dem Almosen und dem geweihten Salze macht. Eben solche Weisheit haben die
\quellseite{583}
Weiber, welche durch einen ähnlichen Scherz betrogen, die Zweige, die am Tage des heiligen Marcus [25. April] auf dem Wege der Procession gestreut waren, nehmen und über die gesäeten Erbsen und Bohnen ziehen; nämlich weil gesagt wird: Wenn die Saaten mit solchen Zweigen gesteckt werden, so werden sie sicher davor sein, daß die Hühner oder die Vögel des Himmels sie nicht wegnehmen, verstehen sie nicht: „Wenn (das heißt, zu der Zeit, wenn) sie kehren, werden die Vögel dies nicht thun.“ Und wer könnte alle die närrischen, lächerlichen, falschen, nichtigen und abergläubischen Dinge aufzählen, mit welchen dieses leicht verführbare Geschlecht umgeht? Von ihrer Mutter (ex prima) Eva ist es ihnen angeboren, daß sie sich betrügen und zum Narren haben lassen. Hieraus mögen ähnliche Dinge beurtheilt werden. Aber weil es Leute gibt, welche nicht glauben, daß dergleichen möglich sei, andere aber allzuviel glauben, darum müssen wir darüber noch ein wenig reden.

Erstens. Einige glauben nicht, daß diese Hexen (Magas) und Wahrsagerinnen so viel vermögen, daß sie Ungewitter machen, dem Leibe Schaden thun, dem Vieh, den Gütern, den Kindern Uebel zufügen etc., indem sie diesen Grund für ihre Meinung haben, daß sie nicht glauben, daß die Hexen so große Gewalt über die Creaturen GOttes haben, und [glauben,] daß den Teufeln so viel von GOtt nicht zugelassen werde. Und dennoch, wiewohl sie so glauben, daß alles von GOtt auferlegt werde, so nehmen sie doch nicht zu GOtt ihre Zuflucht, wenn sie etwas der Art erleiden. Deshalb soll man sie dessen überreden, daß sie müssen sollen, daß die Teufel unter GOttes Zulassung solches sehr wohl vermögen, was aus der Schrift genugsam bewiesen werden kann.

Erstens. Hiob 1, 16., da Feuer vom Himmel fiel und das Vieh sammt den Knechten tödtete, und ein Wind sich erhob, und das Haus umstieß und seine Söhne erschlug, desgleichen endlich ihn an seinem Leibe mit den schlimmsten Schwären schlug; zuletzt ihn auch an der Seele mit der schwersten Anfechtung zur Verzweiflung versuchte; und dennoch drückt der Text daselbst klar beides aus, nämlich, daß der Teufel dieses gewirkt habe, nachdem er die Erlaubniß von GOtt erhalten hatte, und daß Gott ebendaselbe gethan habe, wie es ebendaselbst heißt: „Das Feuer GOttes fiel vom Himmel“ etc. Und er selbst sagt nachher [Hiob 1, 21.]: „Der HErr hat’s gegeben, der HErr hat’s genommen“ etc., wiewohl doch der HErr daselbst zum Teufel sagt [Hiob 1, 12.]: „Siehe, alles was er hat, sei in deiner Hand; ohne allein an ihn selbst lege deine Hand nicht.“ Wenn also der Satan Feuer und Blitze vom Himmel werfen und so großes Uebel thun kann, wo GOtt es ihm zuläßt, warum sollte er denn nicht auch jetzt dasselbe vermögen, wenn er beschworen wird, daß er seinen Bund halten soll?

Zweitens. Daß die Plagen der Egypter durch böse Engel ausgerichtet sind, bezeugt der 78. Psalm, V. 49.: „Er sandte böse Engel unter sie, und ließ sie toben und wüthen und Leide thun.“ Aber jene Plagen waren gar größer als die, welche jetzt irgend welche Hexen thun können.

Drittens. Von wie viel Besessenen wird im Evangelio erzählt! Wie viele Seelen besitzt er noch täglich mit Irrthum und Sünde! Oder ist es geringer, wenn er eine unsterbliche Seele verderbt, als wenn er ein sterbliches Glied oder die Luft beschädigt?

Viertens. Hat er nicht alle heiligen Märtyrer selbst verfolgt und getödtet? Ja, er hat den HErrn selbst genommen und von der Wüste auf einen Berg und von dem Berge auf die Zinne des Tempels geführt [Matth. 4, 5. 8.]. Desgleichen streitet er [Offenb. 12, 7.] mit Michael und seinen Engeln.

Fünftens. Das Leiden Christi beweist dies alles überaus stark. Denn wenn der Sohn GOttes so viel erlitten hat von den Teufeln und ihren Gliedern, was ist es Wunder, daß er uns unsere Scheuern zerstören kann und die Glieder beschädigen. durch dieselbe Zulassung GOttes?

Sechstens. Siehe doch das an, wie Gott so viele bösen Menschen Macht gibt, der Güter, des Reichthums, der Gewalt, der Herrschaft zu mißbrauchen, anderen zum Verderben und ihnen selbst zur Sünde. Viel mehr läßt er den Teufeln zu, daß sie wider uns der Dinge mißbrauchen, sei es nun, um uns zu strafen, oder um uns zu unterweisen. Es ist aber wahr, daß sie ohne GOttes Erlaubniß auch nicht ein Blatt am Baume bewegen können, denn, wie Christus [Matth. 10, 29. f.] sagt, es fällt auch nicht ein Blatt vom Baume auf die Erde ohne den Willen des Vaters. Darum müssen wir Christen wissen,
\quellseite{584}
daß diese Uebel zwar durch die Teufel und ihre Herren angethan werden, daß es aber dennoch von GOtt so verordnet sei. Das Gute thut GOtt durch sich selbst, das Böse durch Böse. Daher hat Hiob [Cap. 1, 21.] (wie St. Augustinus bemerkt) nicht gesagt: Der HErr hat’s gegeben, der Teufel hat’s genommen, sondern: „Der HErr, der es gegeben hat, der hat es auch genommen.“ Aber davon hernach mehr.

Zweitens. Allzuviel glauben diejenigen, welche glauben, daß sich die alten Vetteln in Kater verwandeln können oder irgendwelche andere Gestalten annehmen und bei Nacht zu ihren Gelagen fahren. Und hierüber ist [im geistlichen Rechte] ein ausdrücklicher Text Frage 15, Cap. 6. Und es ist nicht zu verwundern, daß es niemand glaubt. Denn es sind Täuschungen des Teufels, aber nicht eine wirkliche Sache, wie schon erhellt aus einem Exempel, welches Johannes von Kaisersberg\footnote{Johannes Geiler von Kaisersberg, so genannt nach dem Wohnsitze seines Großvaters, der ihn erzog, wurde am 16. März 1445 in Schaffhausen geboren. Er starb als Prediger am Straßburger Münster den 10. März 1510. (König, Deutsche Litteraturgeschichte, S. 194.)} erzählt: daß eine alte Here, um einen Prediger der Lüge zu überführen, da er gelehrt hatte, solche Dinge seien nicht wahr (falsa), ihn spät Abends rufen ließ, vor seinen Augen sich mit Salbe einschmierte, in einem Backtrog saß und nun davon fahren wollte: siehe, da schlief sie alsobald ein und machte allerlei heftige Bewegungen, bis daß sie von der Bank fiel und eine Wunde am Kopf erhielt. Endlich kam sie wieder zu sich und wurde durch die Wunde und den Fall überführt, daß sie sich getäuscht hatte, so sehr sie sich auch ihrer wunderbaren Gesicht gerühmt hatte.

Dasselbe erhellt aus den Lebensbeschreibungen der Väter, da von einer Jungfrau erzählt wird, welche ihre Eltern zu St. Macarius brachten.\footnote{Vergl. Walch, St. Louiser Ausg., Bd. VIII, 346, Anm. 2.} Vor seinen Augen war sie ein Mädchen, während sie den anderen eine Kuh zu sein schien. Und solche Blendwerke geschehen vielfach auch noch heute, wie jenem geschah, der einem Juden schuldig war und im Schlafe das erlitt, daß ihm das ganze Bein vom Leibe abgerissen wurde; damit er den Juden zu Schaden machen sollte. Es wird auch von einem Mönche erzählt, der fast einen Wagen voll Heu verzehrt hat etc.

Ueberall wird viel von solchen Dingen berichtet. Denn es kann sich der Teufel von einer Person in eine andere verwandeln und auf wunderbare Weise die Sinne täuschen, wovon in der Legende des heiligen Martin gesagt wird, wo einer, der vom Teufel mit einem gewissen Kleide unter dem Namen Christi beschenkt worden war, die Sinne aller täuschte. Aehnliches wird erzählt von einer gewissen Nonne, welche von anderen gesehen wurde, wie sie halbnackt in einem zerfetzten Mantel in ihrer Zelle saß und einen Haufen Kuhmist auf ihrem Kopfe hatte, aber dennoch glaubte, daß sie in dem glänzendsten Kleide und von Christo mit einer goldenen Krone geziert da säße. Desgleichen auch das in der Legende des heiligen Andreas, wie der Teufel in der Gestalt eines schönen Weibes einen so großen Mann betrog. Desgleichen von jenem Ritter, welcher spät Abends in eine Herberge kam, und meinte, daß er in angenehmer Gesellschaft unter Mädchen und Neigen und bei köstlicher Mahlzeit wäre, früh Morgens aber in tiefem Kothe lag; und sein Pferd war an einem Baum angebunden, ungefüttert. Das Aller stärkste aber ist das in der Legende des heiligen Germanus, des Bischofs von Altisiodorum. Siehe im Verzeichniß (catalogo) der Heiligen Buch 6, Cap. 9. Da er einmal eingekehrt war, und man noch [spät], da es schon Schlafenszeit war, einen Tisch zurichtete für die Nachbarn (wie ihm auf sein Befragen gesagt wurde), die in der Nacht da sein würden: so beschwor er die, welche kamen, weckte den Hausvater mit den Seinen auf, und erkundigte sich, ob sie sie kennten. Als sie dieses bejahten, so befahl er ihnen zu den Nachbarn zu gehen, und zuzusehen, ob sie zu Hause wären, und sie fanden sie alle in ihren Betten schlafend, und so wurde es klärlich offenbar, daß es ein Spuk der Teufel war, und er verschwand auf seinen Befehl.\footnote{Dieselbe Erzählung ist ausführlicher und verständlicher in den Tischreden Cap. 24, § 132. Walch, St. Louiser Ausg., Bd. XXII, 776.}

Doch ich will auch denen ein Genüge thun, die Lust haben, darüber zu zanken. Denn sie sagen, daß Weiber in Kater verwandelt worden wären, sei um deswillen wahr, daß ein gewisser Wagehals, der in einem leeren Hause geschlafen habe, mehrere solche Kater verwundet habe; da es aber Morgen geworden war, habe man viele alte Weiber so verwundet gefunden. [Dies hat
\quellseite{585}
keinen Grund,\footnote{Von uns eingefügt.} denn entweder ist es erdichtet, oder die Teufel haben die rechten alten Weiber selbst verwundet, um einer lügenhaften Sache (Glauben zu verschaffen, als wenn jener die rechten alten Weiber verwundet hätte, der da schien die Kater (das ist, die Teufel) zu verwunden. Es ist möglich, sage ich, daß die alten Weiber im Schlafe oder in einer Ohnmacht (exstasi) gemeint haben, daß sie umherlaufen und verwundet werden, während sie doch ruhig im Bette liegen und von den Teufeln in Wahrheit verwundet werden, so daß sie nachher glauben, sie seien wirklich umhergelaufen und von einem anderen verwundet worden. Ebenso ist es auch mit dem Reiten auf dem Besen. Von der Hulde (Herodiade) und den Kobolden (laribus) ist genügend offenbar, daß sie Teufel sind.

Drittens schreiben treffliche Männer von den Teufeln, die als Mann (incubis) oder Weib (succubis) handeln, daß der Satan in Gestalt eines Weibes mit einem Manne zu schaffen haben kann (succumbere), und wenn er den Samen empfangen hat, wiederum bei einem Weibe liegen (incumbere), und so zeugen.\footnote{Die Wittenberger und die Jenaer fügen hinzu: „aber Mißgeburten (monstra), nicht einen rechten Menschen“.} Denn in solcher Weise hat er einen gewissen Einsiedler zu verführen gesucht, daß er ihn zur Unzucht brächte, und ist dann verschwunden, wie St. Johannes, der erste Einsiedler, bei Hieronymus erzählt. Ja, es wird gesagt, daß der Teufel einmal in Gestalt eines Kindes gewesen sei, und fünf Säugammen hätten ihm kaum genug geben können.

Viertens. An denen, die auf einem Mantel fahren, zweifle ich fast. Es ist möglich, daß sie eine kurze Strecke fahren, aber ob auch fernhin, das weiß ich nicht, wenigstens nicht in so kurzer Zeit, wie man sagt. Das weiß ich, wenn sie auch nicht wirklich fahren, kann der Teufel alle Sinne so äffen, daß derjenige, welcher fährt, ruhig ist, und sich doch dünken läßt, er fahre zu seiner Geliebten (dominam) und vollbringe mit ihr alles, was aber nichts als bloße Vorstellungen (phantasmata) sein mögen.

Ein gleiches Urtheil fälle ich über die, welche von fernher Speise und Trank bringen; und dem steht das nicht entgegen, daß sie nachher die Oerter wirklich wiederkennen, an denen sie, wie sie sagen, gewesen sind, weil der Teufel ihnen ähnliche Gestaltungen vor die Sinne bilden konnte. Ebenso [urtheile ich] auch von den Dienern, welche er herbeibringt, wofür mir das ein Anzeichen zu sein scheint, daß niemandem gestattet wird, sie anzureden, damit sie nämlich nicht verschwinden wie der Rauch.

Zum letzten von dem, was in den Crystallen oder Elfenbein oder in einem Fingernagel gesehen wird, ist es gewiß, daß es ein Blendwerk des Teufels ist, welcher Figuren in den Crystallen (ingentis), so viel ihm zugegeben wird. Denn er kann dies nicht immer, noch vor allen, wie man genugsam erfahren hat.

Aber hieraus mag anderes beurtheilt werden. Denn die Arten dieser Gaukeleien (portentorum) sind unzählig und mehren sich noch täglich, weil die Arbeiter im Weinberge des HErrn nicht ihre Schuldigkeit thun.

Dies alles aber ist ganz unbedeutend und gering bei dem Teufel, nämlich es sind Dinge, durch welche er die leiblichen Sinne und Vorstellungen bethört, die niedrigsten Stücke am Menschen. Aber mit viel größerer Feinheit, wie auch größerer Gefahr, bethört er die Gemüther und den Verstand. Und wiewohl sehr viele sind, deren Sinne seinen Blendwerken offen stehen, so bethört er doch noch viel mehr am Gemüthe, nämlich die Gelehrten in der Schrift, die Heiligen in den Werken, so daß er von dem Apostel [2 Cor. 11, 3.] mit seinem eigentlichen Titel bezeichnet wird als ein Verrücker der Gemüther, gleichwie er Eva mit seiner Schlauheit verführte, so sehr, daß er sich auch zum Engel des Lichts verstellt [V. 14.]. Daher sagt auch St. Ambrosius in der Predigt von den Bischöfen: Feindselige Wegelagerer\footnote{supplantatores, eigentlich: die jemandem ein Bein stellen, um ihn zu Falle zu bringen.} stellen an allen Wegen Stricke, die sehr zu fürchten sind wegen der zahllosen Todesgefahr (mortibus), kleine Thiere mit den großen und unzählige kriechende Thiere, so daß es unmöglich ist, daß irgend ein Mensch sicher sei in irgend einem Werke, in seinem Sinne und Verstande, es sei denn, daß er sich beständig mißtraue und sich in allen Dingen fürchte, und mit Hiob sich scheue vor allen seinen Werken. Aber ach, daß es GOtt erbarme, wie sicher sind wir heutzutage, die Niedrigsten ebensowohl als die Höchsten, Gelehrte ebenso wie die Ungelehrten, als ob der Teufel gestorben wäre, in solchem
\quellseite{586}
Grabe, daß wir für unsern Wahn auch die blutigen Kriege führen oder ohne Ende hadern und streiten! So hat es unsere Gottlosigkeit verdient, daß wir in einen verkehrten Sinn dahingegeben würden.

Nun, von den gröberen und sinnlichen Dingen sei dies genug. Es folgt eine andere Art der Uebertretung, nämlich derer, die durch einen schönen Schein betrogen werden. Derer sind wiederum zwei Arten. Die eine besteht aus denen, die hinsichtlich der Reliquien und der Verehrung der Heiligen von der rechten Bahn abweichen, die andere aus denen, welche in Weisheit und eigener Gerechtigkeit wider Gott hoffärtig sind. Zuerst wollen wir von dem Ersten handeln.

Auf zweierlei Weise werden von uns die Heiligen verehrt und angerufen.\footnote{Löscher bemerkt: Luther redet gar nüchtern von der Verehrung der Heiligen, aber doch noch befangen im römischen Aberglauben, den er noch nicht abgelegt hatte.} Erstlich nur wegen zeitlicher und leiblicher Dinge. Und diese verehren die Heiligen in falscher Weise, ja, sie ehren sich selbst mehr in ihnen, weil sie das Ihre suchen, nicht das, was GOttes ist, und deshalb machen sie sich die Heiligen fast zu Abgöttern. Diese verwirft St. Augustinus nicht ganz und gar, sondern lobt sie nur in solcher Weise, daß er aussagt, sie seien besser als diejenigen, welche durch Bündniß mit dem Teufel das Zeitliche suchen. Denn es ist besser, daß man auch das Zeitliche von GOtt begehre als vom Teufel. Aber dadurch sind sie nicht gelobt, ja, nicht einmal Christen. Es ist ein gar geringes Lob, ja, ein großer Tadel, sie seien nicht gut, wenn sie nicht mit den Allerhöchsten verglichen werden. Und doch erhört sie GOtt, der auch den jungen Raben Speise gibt, auch in diesem Dienste (cultu), und gibt ihnen das, was sie bitten, wie aus der Schrift bewiesen wird.

Erstens. Durch Naeman, den Syrer, gab GOtt dem ganzen Syrien Heil, wiewohl er ein Abgöttischer war, doch hochgehalten bei dem Könige von Syrien. [2 Kön. 5, 1. ff.].

Zweitens. Die Leute in Samarien, da sie anfingen GOtt zu dienen, wiewohl sie zugleich ihren Göttern anhingen (was doch unmöglich Gotte gefallen kann), sind von den Löwen befreit worden, von denen sie erwürgt wurden [2 Kön. 17, 24. ff.].

Drittens. Dem Ahab und vielen abgöttischen und bösen Königen gab er viele Siege und Wohlthaten wegen etlicher Werke, die äußerlich vor den Menschen einen Schein haben.

Viertens. St. Augustinus spricht es in seiner Schrift De civitate Dei, Buch 1, Cap. 8, aus, daß den Römern von GOtt ein so überaus großes Reich gegeben worden sei wegen etlicher trefflichen Tugenden, die sie hatten, während sie doch Gotte nicht gefallen konnten, sondern sich dadurch nur bei den Menschen verdient machten.

Fünftens. Der Juden, welche nur aus Furcht vor der Strafe die Sünde mieden, hat er oft verschont, wiewohl sie ein Volk waren, das im Herzen immer irre ging und Gotte nicht gefiel.

Sechstens. Und auch jetzt wendet GOtt öfters Pestilenz, Krieg und Hungersnoth ab um einer Besserung willen, die allein aus Furcht eingetreten ist, wiewohl Gotte doch eine gezwungene Reue nicht gefallen kann, die weniger ist als eine halbe Reue. Dies ist ganz offenbar, weil sie, wenn die Strafe nachläßt, alsbald wieder zu ihrer vorigen Gesinnung und Lebensweise zurückkehren, ja, noch ärger werden. Alle diese, wie sie ihn als einen GOtt verehren und fromm leben um zeitlicher Dinge willen, so empfangen sie auch ihren Lohn. So verfahren wir auch mit den Heiligen GOttes, die wir nicht weiter ehren, als wenn uns ein Fuß oder der Kopf wehtut, oder der Geldbeutel leer ist. Und wiewohl uns dies unnütz und thöricht ist vor GOtt, auch Gotte und den Heiligen nicht gefällt, so dient es doch zur Ehre GOttes, der aus allem Dingen etwas Gutes hervorzubringen kann, ja, dessen Wesen es ist, daß er auch den Thörichten, ja, seinen Feinden und Lästerern wohlthut, regnen läßt über Dankbare und Undankbare, um uns gerade dadurch zu lehren, daß dies nicht die Dinge sind, um derer willen er oder ein Heiliger verehrt werden solle, da wir sehen, daß er diese Dinge auch allen bösen Buben reichlich gibt, ohne daß er darum gebeten oder verehrt wird.

Und damit dies klarer verstanden werde, wollen wir etliche der Heiligen aufzählen, deren abergläubische Verehrung allen bekannt ist.

Der erste ist St. Antonius, welcher von vielen wegen des „heiligen Feuers“\footnote{ignis sacer = die Rose, der Rothlauf, St. Antonsfeuer.} (denn sonst würde man ihn nicht ansehen) verehrt wird und in solcher
\quellseite{587}
Weise verehrt wird, daß man in der Gottlosigkeit fast so weit gegangen ist, daß man glaubt, der heilige Mann könne selbst, aus eigener Kraft oder durch die Gabe, die GOTT ihm gegeben habe, diese Heilungen verrichten. So gar erbitten sie nicht von GOTT durch die Fürbitte des heiligen Antonius, sondern gleichsam allein von Antonius selbst, was sie begehren, daß sie gar nicht an seine Fürbitte denken. Sodann suchen sie diese Gnade der Heilung in solcher Weise allein bei ihm, daß sie entweder nicht zu wissen oder nicht zu glauben scheinen, daß GOTT dasselbe durch irgend jemand anders verleihen könnte. Unterdessen hat man aber noch von niemandem gehört, der den heiligen Antonius darum angerufen hätte, daß er auch zu dem Verdienst kommen möchte, daß er dem Antonius in Freudigkeit des Geistes und in den anderen Tugenden nachfolgte, davon doch in der Legende desselben Erwähnung gethan wird; von dem Feuer aber sagt sie gar nichts.

Der zweite ist der heilige Sebastian, der vor allen wegen der Pestilenz verehrt wird; erst neulich hat man ihm den heiligen Martin als einen Gefährten beigegeben,\footnote{Der folgende Abschnitt, den wir in Klammern gesetzt haben, fehlt in den Ausgaben Schumanns und Luthers. Die Weimarische hat ihn nur am Rande mit der Bemerkung, daß sich dieser Abschnitt in der ersten Ausgabe finde, aber Luther habe diese Stelle, wie sich aus dem Druckfehlerverzeichnis ergebe, am liebsten gestrichen gesehen. In den Ausgaben sind darin viele verschiedene Lesarten, je nachdem sie Luthers Verbesserungen mehr oder weniger benutzt haben.} [und den heiligen Rochius, die ein gar feines Leben geführt und ein hohes Ansehen genossen haben. Das heilige Leben und den Glauben dieser drei Männer sucht sich durchaus niemand anzueignen. Und um auch meinen Verdacht herauszuplaudern, der freilich ungereimt erscheinen mag, so hat es den Schein, daß einige Heilige vor anderen nur wegen des Klanges (allusione) ihrer Namen zu diesem Dienste angerufen worden sind. Denn es hat der Name S. Tonii im Italienischen fast den Klang (Italica allusione propesonat), als wenn man sagte: sanctus ignis, das ist, das heilige Feuer, als ob er wegen seines Namens auch der Nothhelfer (patronus) wäre, der das heilige Feuer heilen könnte, wie auch die heidnischen Römer das Fieber\footnote{Febris wurde als Göttin verehrt, wie Plinius berichtet.} und andere Dinge, die in ihrem Namen einen Anklang hatten (allusiones rerum), den Göttern zuschrieben.] So auch Sebastian „Sanct Pastian“, als ob er allein gegen die „Pest“ helfen könnte, während doch in ihren Lebensbeschreibungen nichts von der Pest gelesen wird. Daher werden auch in Italien diese beiden vor anderen nicht anders geehrt, als wenn du sähest, daß unter den Heiden die übelwollenden Gottheiten begütigt werden. Nun hat der heilige „Rochius“ einen Namen, der im Deutschen lautet „Roche“ und Born, als ob er deshalb dazu nützlich sei, den Born GOttes abzuwenden, weil er entweder in Wahrheit Rochius genannt wird, oder nur erdichtet wird. Doch die Legende desselben, wenn sie anders glaubwürdig ist (si recipitur), erzählt etwas Derartiges).\footnote{Dieser letzte Satz ist von Luther in dem Druckfehlerverzeichnis hinzugefügt.}

Der dritte ist Valentin, der Patron (praefectus) der fallenden Krankheit. Da wir nichts lesen, daß er mit dieser Krankheit zu thun gehabt habe, so möchte ich fast schwören, daß er durch den Gleichklang in der deutschen Sprache zu dem Amte gekommen ist, daß er hier Hülfe leisten muß. Denn cadere heißt „fallen“, was ganz ähnlich klingt wie „Valentin“. Und man darf sich darüber nicht wundern, da die abergläubischen Weiblein auch den heiligen Vincentius dazu angestellt haben, verlorene Sachen zu finden, geleitet durch ebendenselben Gleichklang; denn im Deutschen heißt invenire „vinden“, und darum haben sie den Vincentius zu einem Diener und Wächter über die verlorenen Dinge gesetzt.

Es ist nicht etwas Großes noch Wunderbares, daß GOTT zuläßt, daß dies geschehe, und diese Dinge noch mit seinen beständigen Wohlthaten begleitet, und die Kirche es gleicher Weise billigt, besonders, was den heiligen Antonius betrifft, denn, wie ich gesagt habe, es liegt GOTT wenig daran, diese Dinge auszutheilen, die er auch seinen Feinden gibt. Und er läßt unterdessen sowohl die Thörichten als auch die Schwachen im Glauben sich an diesen Gütern erfreuen, so lange sie nicht besserer Dinge werth sind, oder bis daß sie besserer Güter fähig werden.

Der vierte\footnote{Löscher: „Das Folgende hat Luther am Festtage des heiligen Jacobus [25. Juli] 1516 als Predigt vorgetragen, nachdem er eine kurze Vorrede vorausgeschickt hatte über die Worte [Matth. 20, 22.]: ‚Ihr wisset nicht, was ihr bittet.‘“ — Diese Einleitung ist uns nicht erhalten. (Weim. Ausg.)} ist Christophorus, nach einer
\quellseite{588}
Legende, die nicht sowohl apokryphisch als verdächtig ist, der jedoch so große Ehre hat, daß kein Apostel mit ihm verglichen werden kann, wenn auch seine Thaten mitten in der Bibel geschrieben stehen. Doch wird er nicht um desswillen geehrt, weil er den Menschen Vergebung der Sünden und den Glauben oder die Gnade GOttes erlangt, sondern, wie die Verslein besagen, die sich zu diesem Gottesdienst überaus wohl schicken:

\begin{quote}
Heiliger Christophorus, so große Kraft wohnt dir inne: Wer früh Morgens dich sieht, deß Mund ist zur Nachtzeit voll Lachens. Satan kann ihn nicht schlagen, durch jähen Tod nicht verderben.
\end{quote}

Erstlich. Siehe, wie hier der Gottlose solche große Kraft nicht GOtte, sondern einem hölzernen und gemalten Bilde zuschreibt; er schreibt es nicht einmal dem Glauben zu, welcher allein macht, daß wir es empfangen, wenn die Heiligen etwas erlangen.

Zweitens. Ich möchte gern belehrt werden, durch welche Schriftstelle und durch welche Wunderwerke es bewährt worden sei, daß das Bildniß des Christophorus dies ausrichten könne. So ist unter dem Schein der Gottseligkeit ein ungeheures Meer von Aberglauben eingerissen, da es niemand hindert, indem die Priesterschaft schläft und begierig nach Gewinn trachtet.

Drittens. Warum lacht derjenige gerade am Abend, welcher ihn früh Morgens gesehen hat? Wie? wenn er ihn des Abends ansähe? Hat er denn seine Kraft so gar an das Ansehen am Morgen gebunden, daß derjenige, welcher es früh Morgens vergessen hat oder zu spät gekommen ist, ihn vergeblich angesehen hätte, wenn er ihn darnach zu irgend einer andern Tageszeit angesehen hätte?

Viertens. Warum heißt es nicht: Wer deinem Glauben nachfolgt, der wird immer lachen? Wer mit dir leiden wird, der wird sich mit dir freuen? Aber das sei ferne, daß die ihm nachfolgen wollten, welche jene Dinge suchen. Auch das sei ferne, daß diejenigen, welche ihm nachfolgen wollen, jene Dinge suchen sollten, oder in solcher Weise suchen sollten.

Fünftens. Warum macht nicht das Kreuz Christi die lachen, welche es ansehen, da es das Zeugniß der Schrift hat? Ist denn etwa das Bild des Christophorus besser als das Christi, da doch das Kreuz Christi selbst allein in Wahrheit eine Art Christophorus [Kreuzträger] ist?

Sechstens. Das Allerärgste ist, daß auf diese Weise die Menschen ihn mehr lieben, ihn mehr verehren, fleißiger ansehen und mehr auf ihn vertrauen als auf Christum selbst, damit sie zeigen, daß sie GOtt nicht loben in seinen Heiligen, sondern sich selbst, weil sie nicht die Ehre GOttes, sondern ihren eigenen Nutzen suchen; und dies öffentlich mit frecher Stirn klar bekennen, sie thäten es um der Krankheiten willen. Es wäre besser, man folgte den Heiligen nach, würde durch einen plötzlichen Tod dahingerafft, und würde selig, als daß man ihnen nicht nachfolgte, stürbe eines natürlichen Todes, und würde verdammt. Auch dieser Heilige scheint von Soldaten und Edelleuten und Vornehmen (generosis) zuerst erfunden zu sein, da ihnen ja in den Kriegen ein plötzlicher Tod, zu deutsch „der gehende Tod“ zu befürchten steht, besonders nachdem die Schießgewehre erfunden sind. Aber sie denken nicht daran, daß die Leute daselbst viel häufiger stante morte, zu deutsch „des stehenden Todes“ sterben, und nicht allein „des gehenden Todes“.

Siebentens. Daß die Menschen suchen, hierdurch sicher zu sein, und ohne die Furcht GOttes zu leben, während doch alle Verehrung der Heiligen eingesetzt ist für die Furcht GOttes, daß sie die Menschen zur Buße reize, sie dazu bringe, daß sie des Todes gedenken, und sie hinweise auf das künftige Leben. Aber sie gebrauchen derselben nur, um diese Dinge zu fliehen, jene aber festzuhalten. Deshalb sollte man Christo sagen, daß er sein Wort zurücknehmen müsse, da er sagt [Matth. 25, 13.]: „Wachet, denn ihr wisset weder Tag noch Stunde“, denn wir haben den Christophorus gefunden, der für uns wachen muß, während wir schnarchen, und wir können den ganzen Tag sicher sein, geschweige, daß wir die Stunde nicht wissen sollten. Wie? wenn nun auch Christophorus einen anderen Christophorus gefunden hätte, den er angesehen hätte, und selbst sicher geworden wäre, daß GOtt ihn nicht schlagen würde: wie wäre er denn zu seinem Märtyrertode gekommen? da es ja nun, nicht mit der Schwachheit, sondern mit der Bosheit schließlich so weit gekommen ist, daß GOtt seiner Herrschaft beraubt wird, daß er diejenigen, welche den Christophorus ansehen
\quellseite{589}
, an dem Tage weder durch Wasser noch durch Feuer noch durch irgend welche Creatur tödten kann.

Der fünfte ist St. Laurentius, den sie erstlich mit Fasten ehren, damit er ihr Haus vor Feuer bewahre, darnach aber an seinem Feste den ganzen Tag nicht ein Fünklein Feuers im Hause leiden. Jedoch in anderen Häusern gebrauchen sie Feuer, um damit die Speisen zu kochen. Aber so fürchten und verehren sie nicht den Laurentius, sondern vielmehr das Feuer. Aber vielleicht sind sie so thöricht, daß sie glauben, das Feuer sei dem heiligen Laurentius verhaßt, weil er in demselben verbrannt und geröstet worden ist, so daß, wenn er ein Feuer sieht, er des erlittenen Unrechts gedenke, und wider die zürne, welche dasselbe haben. Wenn sich das so verhielte, dürfte kein Christ an irgend einem Tage Feuer haben. Oder ist er um deswillen ein Herr des Feuers geworden, weil er das Feuer überwunden hat? Das ist freilich wahr [er hat es überwunden], ja, er ist auch ein Herr des Teufels, der Sünde, der Welt und des Fleisches und alles Uebels geworden. Warum wird er also allein wegen des Feuers verehrt? Oder warum verehrst du ihn nicht, damit er dich behüte vor Gottlosigkeit und Grimm, vor Zorn und Hoffahrt, durch welche er von Decius und seinen Dienern, die dies ausführten, auf dem Feuer geröstet ist? Diese hat er endlich überwunden, und kann sie jetzt noch in dir überwinden durch seine Fürbitte. Aber dies möchte man auch anderswo suchen als bei St. Laurentius; warum muß er nur der Beschirmer wider das Feuer sein. Ihm hat man den heiligen Florian zum Genossen gegeben, dessen Leben und Name auch unbekannt ist, und bei ihnen selbst auch nicht auf irgend eine andere Weise bekannt, als daß sie ihn auf dem Papier gemalt sehen, wie er Wasser auf ein brennendes Haus ausgießt. Denn damit sind sie zufrieden, weil sie auch dies allein bei ihm suchen.

Die übrigen überlaufen wir in der Kürze. Der heilige Veit hat seine ihm zugewiesene Stellung bei dem jammervollen Veitstanze. St. Erasmus ist der Beschirmer der Geizigen, doch in solcher Weise, daß er mit einer gewissen Anzahl von Gebeten und mit Wachskerzen verehrt werde: dann wird er sie sicherlich mit Reichthum überhäufen. Denn was sollte er sonst vornehmen, da er nichts zu thun hat? Aber auch der heilige Ludwig vom Orden der Minoriten,\footnote{Barfüßer, Franciscaner.} weil er einmal in seinem Leben [schlechtes] Bier gut gemacht hat, darf nun, da er todt und selig ist, durchaus nichts Anderes thun [als Bier brauen]. Der heilige Wendelin, ein Hirte, ist noch nützlicher wider alle Wölfe als alle Hunde. Und da er lebte, hütete er nur Eines Menschen Vieh, nämlich sein eigenes, jetzt, da er todt ist, wird er gezwungen, das Vieh aller Leute zu hüten. Noch schämen wir Christen uns nicht, die Geschäfte in zeitlichen Dingen so unter die Heiligen zu vertheilen, als wenn sie jetzt Knechte und Handwerksgesellen wären, so daß beinahe der Pfuhl abergläubischer Gebräuche wiedergekehrt ist, daß wir wiederum das Gewirr (Cahos) der römischen Götter und eine Art Pantheon von neuem erbaut haben, und zwar aus keiner andern Ursache, als damit wir nur hier ein gutes Leben haben möchten.

Nun wollen wir noch einige weibliche Heilige anfügen. Unter denselben ist St. Anna, wie sie zugleich eine neuere und eine ältere ist, so auch mit Recht die erste. Wiewohl ihre Legende überall zweifelhaft ist, so ist sie dennoch in dieser Hinsicht wohl die angenehmste gewesen und ist es noch, daß sie als eine neue Frau, weil wir anfangen der alten Dinge überdrüssig zu werden und derer, welche eine ganz gewisse Verehrung haben. Aber sie war noch viel angenehmer, weil sie nicht leer kam, sondern voller Reichthümer, da sie sonst durchaus nichts hätte ausrichten können, wenn sie Armuth brächte. Zweitens, daß sie fast über die heilige Jungfrau erhoben wird. So thun neue Feste immer den alten Abbruch, und wir heben die neuen hoch, mehr geleitet durch den Zulauf des großen Haufens als durch rechte Andacht. Drittens, daß es klar und offenbar ist, daß man angefangen habe, sie zu verehren wegen des Reichthums, nach der apokryphischen und sehr verdächtigen Legende, in welcher jenem verarmten Spieler (auch vom Himmel) wieder geholfen worden ist, daß er reich wurde. Denn allein dies hat das gottlose Auge des Geizes in dieser ganzen Legende beachtet. Und es ist nothwendig, daß wegen des Festes dieser heiligen Mutter jetzt die Feste aller Heiligen etwas verdunkelt werden, besonders die der Apostel, vielleicht weil sie die Armuth gelehrt haben. Und wollte doch GOtt, daß sie
\quellseite{590}
die heilige Anna so des Reichthums halber verehren, daß sie auch anderes suchen, was des Geistes ist, wie die Kirche thut; wie die sogenannten Collecten dieses Tages zeigen. Aber sie glauben, daß sie dazu Feste genug haben, die schon lange auch bis zum Ueberdruß gefeiert worden sind.

Die zweite ist die heilige Barbara, nach deren Legende viele andere erdichtet sind, als die der Catharina, Dorothea und Margaretha. Aber diejenigen, welche sie aufs andächtigste ehren, suchen dadurch, daß sie nicht ohne Sacrament sterben, was nicht so gar verdammlich sein möchte; andere begehren, daß eine schwangere Frau bei der Geburt nicht in Gefahr geraten möge. Denn man liest, daß diese heilige Märtyrin dies und andere Dinge bei ihrem Abscheiden (transitu) erbeten habe, daß derjenige, welcher ihr Gedächtniß halte, dies erlangen sollte. Erstlich ist nach dem Zeugniß Aller diese Legende sowohl apocryphisch als auch verdächtig, zumal da geschrieben wird, daß sie eine Jüngerin des Origenes gewesen sei, und Eusebius einer so großen Jungfrau nicht gedacht hat. Ferner, daß in die Legende der heiligen Catharina gar vieles eingemengt sei, das merkt jeder, der nur eine Nase hat. Ich aber nehme es mir heraus, mich darüber zu verwundern, daß die heiligen Märtyrinnen, die um Christi willen sterben, so anmaßend gewesen sein sollen oder sich so gestellt haben, daß sie gewünscht haben sollten, von den Menschen verehrt zu werden, oder daß sie gewiß gewesen sein sollten, daß sie künftig verehrt werden würden. Aber siehe auch das an: da sie sterben sollen, erbitten sie nicht für die, welche sie verehren, daß sie gerecht und selig werden, sondern daß sie Reichthum, Gesundheit, Sicherheit (was sie alles durch ihren Tod verachtet haben) erlangen mögen, indem sie vielleicht nicht dieselben Güter für andere zu erlangen trachteten, zu denen sie selbst nun eingingen. Daher ist es ohne Zweifel gekommen, daß das Gedächtniß des Leidens der Barbara und ähnlicher Heiligen fleißiger gehalten wird als das Christi selbst, welches doch die Kraft ist aller Leiden der Heiligen. Doch wohl, wir wollen geflissentlich blind sein, sonst könnten wir es mit Händen greifen und erkennen, daß heutzutage in der Kirche die Heiligen um so mehr Ehre haben, je mehr Fabeln in ihren Legenden sind oder Gewinn oder Nutzen. Also müssen wir auch die Schwachheit dieser Leute tragen, bis daß sie unterrichtet werden, und eine bessere Gesinnung eine Gestalt in ihnen gewinne.

Drittens. Die heilige Juliana und Ottilia sind Augenärzte, denn sie werden von niemandem verehrt als von denen, welche an den Augen leiden, und durchaus wegen keiner anderen Sache als wegen der Augen, indem sie aus einem einmal geschehenen Wunder ein natürliches Recht und eine unfehlbare Folgerung machen, was, wie es mir scheint, nicht weit entfernt ist von einem Versuchens Gottes.

Die heilige Apollonia wird wegen des Zahnwehs ohne Unterlaß verehrt und sonst für nichts Anderes. Ihrer langjährigen Jungfrauschaft und ihres überaus brünstigen Glaubens gedenkt niemand.

Die heilige Scholastica ist zur Herrin gemacht über den Donner, und ich weiß, daß dieser Glaube völlig unter den Weibern herrscht, daß, wenn sie hören, irgend ein Heiliger habe mit einer beliebigen Sache zu schaffen gehabt, sie glauben, er vermöge in dieser Sache sehr viel, und sei gleichsam zum Herrn darüber gesetzt. Und etliche Priester unterdrücken diesen Aberglauben nicht, besonders wenn sich spüren läßt, daß er auch etwas einbringen werde.

Nun, wie ich gesagt habe, GOtt läßt dies zu und schenkt das, was man bittet, nicht weil es etwas Großes ist, sondern weil er gut und gnädig ist, daß er sowohl die Schwachen in solcher Weise trägt, als auch die Unwürdigen so beschenkt, deren träges Herz nicht weiß, um bessere und ewige Dinge zu bitten, um welche es sich geziemt, den wahren GOtt zu bitten, und daran er ein Wohlgefallen hat, wenn man ihn darum bittet, und er sie schenken kann. Denn diejenigen, welche GOtt und seine Heiligen recht verehren, suchen an den Heiligen zuerst die Ehre GOttes, und wie sie ihnen nachfolgen mögen, wie die Kirche thut in ihren Collecten, erbitten daher solche Dinge nicht, ja, wenn sie dieselben haben, werden sie ihnen genommen, damit sie ihren Sinn allein auf das Geistliche richten; es können auch an keinem heisseren Zeichen die rechten Verehrer erkannt und von den falschen unterschieden werden, als wenn jenen die zeitlichen Güter genommen werden, diesen aber gegeben in reichem Maaße.

Denn gleichwie GOtt der äußerlichen Gerechtigkeit äußerliche Güter verleiht und die innerlichen
\quellseite{591}
wegnimmt, so verleiht er der innerlichen Gerechtigkeit innerliche Güter und nimmt die äußerlichen weg. Dies ist offenbar an allen Heiligen; wenn GOtt ihnen die zeitlichen Güter auch nicht in der That (effectu) entzogen hat, so hat er doch ihr Herz (affectu) davon abgezogen. Wiewohl deshalb solche Leute nichts im Himmel verdienen, so haben sie doch einen Lohn auf Erden, und sie thun darin recht, daß sie dies nicht von dem Teufel begehren. Denn in solcher Weise gibt er auch vielen Priestern und Mönchen fette und liebliche Weide, wiewohl sie doch in ihrem ganzen Leben nichts Anderes thun, als daß sie mit kaltem und unwilligem Herzen ihre vorgeschriebenen Gebete (horas canonicas) immer lesen und nie beten. Wie nun sie sich dünken lassen, daß sie GOtte damit ein Genüge geleistet haben, so läßt GOtt sich auch dünken, daß er ihnen damit genug gethan habe, wenn er sie in diesem Leben also Mastvieh mästet.

Aber einige Zänker befürchten, daß ich darin frevelhaft handele, daß ich verbiete, die Heiligen wegen leiblicher Noth anzurufen, da auch die Kirche bete, daß sie durch der Heiligen Fürbitte von Fährlichkeit der Seele und des Leibes möge befreit werden, als in der Collecte: „und daß wir befreit werden mögen von Widermärtigkeiten am Leibe, und gereiniget werden von bösen Gedanken im Herzen“ und dergleichen. Zweitens schmücken sie ihre gute Meinung (wie sie es nennen), daß sie das Zeitliche nicht weiter suchen, als sofern sie dadurch bequemer das Geistliche suchen könnten, und wollen um deswillen länger leben und gesund sein, damit sie Buße thun und besser werden könnten.

Drittens. Daß etliche die Ueberzeugung haben (wie Johannes Gerson dafürhält), den Heiligen im Himmel seien derartige Gaben in solcher Weise zugetheilt, wie sie auf der Erde verschiedene Gaben des Geistes hatten, nach dem Worte 1 Cor. 12, 9.: „Einem andern die Gabe gesund zu machen“ etc.

\subsection{Hierauf antworte ich:}

Erstens sehe ich dieses Wort Christi dagegen [Matth. 6, 33.]: „Trachtet am ersten nach dem Reich GOttes und nach seiner Gerechtigkeit, so wird euch solches alles zufallen.“ Und wiederum [V. 32.]: „Nach solchem allem trachten die Heiden.“ Wenn Christus diejenigen Heiden nennt, welche nach solchen Dingen trachten, so hoffe ich, daß es mir erlaubt sei, sie schwache Christen zu nennen oder wenigstens träge Menschen, die ein geringes Verlangen haben [nach geistlichen Gütern].

Zweitens. Ich wehre auch nicht, daß man die Heiligen um zeitliche Dinge anrufe, aber das verwerfe ich, daß sie dies allein um der zeitlichen Dinge willen thun, und ihrer Tugenden und Exempel nicht achten. Zeige mir nur Einen aus der Zahl aller Heiligen, der um seiner Demuth willen oder wegen einer anderen Tugend einen so großen Zulauf, Verehrung und Namen hätte, als St. Valentin hat um der fallenden Krankheit willen.

Drittens. Ich gestehe zu, daß die Kirche in ihren Litan[e]ien [um Behütung] vor Blitz und Ungewitter und anderen leiblichen Nöthen bittet, aber sie thut dies in rechter Ordnung, indem sie zuerst um Barmherzigkeit und Vergebung der Sünden bittet. Denn auch Paulus hat geboten, daß wir bitten sollen, damit wir in diesem Leben Frieden und Ruhe haben möchten [1 Tim. 2, 2.]. So thue du auch, und du wirst leben, und du wirst nicht hören, daß ich etwas wider dich rede.

Viertens. Ich gebe zu, daß die Schwachen im Glauben um Gesundheit und Leben bitten dürfen, damit sie darnach ein desto besseres Leben führen. Aber hier muß ich zweierlei sagen. Wer weiß, oder wer wagt es von sich selbst zu behaupten, daß er ein so reines Herz habe, daß er dies in so einfältiger Meinung begehre? Der Erfolg beweist, daß der größere Theil sich vielmehr nur so stellt, als wolle er sein Leben bessern. Denn fast immer werden diejenigen ärger, welche durch ihr ungestümes Bitten von der Strafe oder Züchtigung GOttes befreit werden, so daß es scheint, daß GOtt uns gerade dadurch unserer Thorheit überführen wolle, nämlich daß wir nicht verstehen, daß er gesund macht, wenn er züchtigt, und uns vielmehr recht krank werden läßt, wenn er nach unserem Wunsche seine [strafende] Hand abzieht. Zweitens. Wenn du dein Leben bessern willst, so kannst du das nie besser thun, als wenn du unter der angefangenen Züchtigung (das heißt, in der recht eigentlichen Besserung und Arznei des Lebens) ausharrst bis an den Tod. Denn diese Besserung des Lebens ist um so gründlicher und beständiger, da sie nicht aus dir ihren Anfang genommen hat, sondern von
\quellseite{592}
GOtt, der dir weit zuvorkommt, ehe du bittest. Oder wie kann ich jetzt alle Herrlichkeit und Würde des Lebens beschreiben, welches im Leiden dahergeht (passivae)? Wir alle sind thöricht, die wir uns dessen vermessen, unser Leben zu bessern, wenn wir in guter Gesundheit stehen, da wir nicht erkennen, daß es in der Krankheit angefangen hat, besser zu werden, und daß das Leben im Leiden viel vortrefflicher ist als das Leben, welches im Thun steht (activa). Daher wäre den Christen vielmehr das zu rathen, daß sie nicht allein ihre Krankheiten gern tragen sollten, sondern auch wünschen, daß der Tod je eher je lieber komme. Dem Christen ist nichts mehr nütze, als daß er so bald wie möglich sterbe, wie St. Cyprian lehrt. Aber wir hören lieber den Juvenal: Man muß beten, daß ein gesunder Geist in einem gesunden Leibe sei. So redet der heilige Augustinus nicht [sondern]: Wenn wir gesund sind, dann ist am meisten [in uns] die scheußliche (insanus) Krankheit der bösen Lust; nicht [sage ich], daß man eine schlechte Gesundheit suchen sollte, sondern daß, wenn sie von Gott auferlegt ist, man sie als einen sehr guten Boten GOttes aufnehmen soll, weil man kein Gutes suchen soll ohne den Willen GOttes.

Fünftens. Gerade das hat GOtt auch sehr gut versehen um der Verstockten und der harten Nacken willen, daß einige Heilige für furchtbar und gleichsam für Rächer gehalten werden, wie St. Antonius und etliche andere. Aber wir lehren die Christen, welche eine gute Zuversicht haben, daß sie einst der Gesellschaft der Heiligen als Genossen eingefügt werden.

Sechstens. Ich weiß die neue Meinung derer, welche sagen, daß die Heiligen, gleichwie sie im Leben mit verschiedenen Gaben des Geistes beschenkt waren, so auch jetzt im Himmel verschiedene Gnadengaben haben, mit denen sie helfen können. Aber ich sehe nicht, wie sie dies beweisen können, es sei denn durch einen Schluß, der hergenommen ist von der Aehnlichkeit (a simili). Diese Art der Beweisung (argumentum) ist aber in den Dingen, die den Glauben anbetreffen, überaus schädlich. Denn was thut der Teufel, da er sich in einen Engel des Lichtes verstellt, anders, als daß er aus der Aehnlichkeit eine Beweisung macht? oder wie viele Irrthümer sind über die Kirche gefluthet unter dem Namen Christi und seiner Heiligen, unter dem Schein der Heiligkeit, unter dem Vorgeben (wie man sagt) der guten Meinung und des Eifers für den christlichen (catholicae) Glauben, das heißt, durch das Schließen aus der Aehnlichkeit (a simili)! ja, wie viele Wahrheiten werden verworfen als Aergerniß, unter dem Vorgeben, man habe ein zu zartes Gewissen, als daß man solches hören könne!

Das sehe ich, daß die Kirche nicht sagt: Heiliger Paulus, heile meine Wunden, sondern: Bitte für mich. Deshalb muß man wissen, daß alle Heiligen allerlei vermögen, und daß dir durch sie von GOtt so viel gegeben werde, als du glaubst, daß du empfangen werdest. Deshalb mögest du dem alle deine Angelegenheiten befehlen, zu dem du die größte Zuneigung hast; erstlich deine Sünden und die Noth der Seele, wie St. Cäcilia bei ihrer Hochzeit that, die alle Heiligen der Reihe nach anrief, daß sie ihr ihre Keuschheit bei GOtt bewahren möchten.

Freilich verwerfe ich das nicht, was auch St. Augustinus in einer gewissen Epistel als seine Meinung ausspricht: Daß GOtt durch einen Heiligen, ja, an einem Orte thut, was er an einem andern nicht thut. Wer weiß den Rath GOttes? Wenn er selbst austheilt, so ist es recht ausgetheilt. Aber darum ist es nicht heilsam für dich, wenn die Geschäfte auf die Heiligen vertheilt werden nach deiner Voraussetzung (sumptione). Daher hat er nicht gesagt, daß die Märtyrer in der Herrlichkeit immer andere und wieder andere Gaben empfangen, sondern er sagt: Er wollte nicht, daß bei allen Gedächtnissen der Heiligen das geschehe, was bei jenen geschehe etc. Wiewohl nun, wie ich gesagt habe, GOtt die Schwachen so trägt, welche allein solche [leiblichen] Dinge bei den Heiligen suchen, indem er ihnen nach seinem gnädigen Willen das schenkt, was sie begehren, so kommt es uns doch zu, die Christen zu unterweisen und zur Vollkommenheit zu leiten, und sie das zu lehren, was heilsam und verdienstlich ist, damit sie nicht etwa in den Irrthum gerathen, daß sie sich aus den Heiligen GOttes Abgötter erdichten. Denn es ist nicht ein Verdienst, daß der Leib gesund gemacht, und das Geld gemehrt wird, sondern wenn du krank bist und arm: das ist dann etwas Schönes, wenn du dieses Verdienst in der rechten Weise duldest.

Dies sei gesagt von der ersten Ehre der Heiligen, das heißt, der leiblichen und äußerlichen.

\quellseite{593}
\footnote{Löscher: „Das Folgende hat Luther in einer Predigt gesagt am Tage der heiligen Anna [26. Juli] 1516, nachdem er Weniges über die köstliche Perle vorausgeschickt hatte, die im Acker gefunden war.“ — „Das Exordium besitzen wir nicht mehr.“ (Weim. Ausg.)} Zweitens werden die Heiligen in Wahrheit und innerlich verehrt. Diese Verehrung aber besteht erstlich im Lobe des Herzens, darnach auswendig in Wort und Werk, das heißt, wenn GOtt gelobt wird in seinen Heiligen und die Heiligen in GOtt gelobt werden, wie geschrieben steht [Ps. 150, 1. 2. Vulg.]: „Lobet den HErrn in seinen Heiligen, lobet ihn nach der Größe seiner Macht.“ So schreibt Paulus Gal. 1, 24., daß die ersten Christen gethan haben, da er sagt: „Und priesen GOtt über mir.“ So thut der 111. Psalm [V. 1. 2. Vulg.]: „Ich will dir danken, HErr, von meinem ganzen Herzen, im Rath der Gerechten und in der Gemeine: Groß sind die Werke des HErrn und köstlich nach all seinem Willen.“ Denn der HErr will, daß seine Ehre erzählt werde und das Werk seiner Hände verkündigt werde, wie er spricht [Ps. 102, 22.]: „Auf daß sie zu Zion predigen den Namen des HErrn, und sein Lob in Jerusalem.“ Und wiederum [Ps. 9, 2.]: „Ich will alle deine Wunder erzählen“, desgleichen [V. 15.]: „Auf daß ich erzähle alle deinen Preis in den Thoren der Tochter Zion.“ Und viele andere Sprüche im Psalter. Denn um deswillen wirkt er so viele und so große Dinge in den Heiligen, damit er gepriesen werde. Wiederum heißt es [Ps. 34, 3.]: „Meine Seele soll sich rühmen des HErrn“, und anderswo [Ps. 105, 3.]: „Rühmet seinen heiligen Namen“, und Ps. 20, 6. [Vulg.]: „Wir wollen uns freuen in deinem Heil, und im Namen des HErrn, unseres GOttes, werden wir erhöhet.“ Denn so lehrt der Apostel [1 Cor. 1, 31.]: „Wer sich rühmet, der rühme sich des HErrn.“ Deshalb ist kein Zweifel, daß die Heiligen nicht in sich selbst, sondern in dem HErrn gelobt und verherrlicht werden wollen, und der HErr in ihnen. Und diese Weise der Verehrung wird nicht allein Gott angenehm und den Heiligen wohlgefällig sein, sondern auch dir heilsam und unvergleichlich nützlicher, als wenn du sie um deinetwillen und deiner Angelegenheiten willen verehrst, ja, allein diese Weise ist heilsam, und irgend eine andere ist schädlich. Darum wollen wir für die Einfältigen eine Weise anzeigen: Nämlich derjenige verehrt die Heiligen GOttes recht in Gott, welcher die Werke und die Gnade GOttes ansieht und durch ihre Betrachtung bewegt und erfüllt wird mit einer süßen Liebe gegen GOtt, daß er sich herbeigelassen hat, ihnen so große und herrliche Gaben zu verleihen, als ob er sagen wollte: Ich lobe dich und danke dir, du frommer GOtt, denn deine Barmherzigkeit währet ewiglich. Denn du hast dich dieses Heiligen erbarmt und dir aus dem Haufen derer heraus, die in der Sünde und Verdammniß liegen, ein so herrliches Gefäß der Ehre zugerichtet. So hast du GOtt in seinem Heiligen gelobt. Wiederum wirst du auch den Heiligen selbst in dem HErrn so loben: O Heiliger GOttes, N., gebenedeiet bist du unter den Menschenkindern, daß dich GOtt für würdig erachtet hat, daß du mit den Gaben einer so unaussprechlichen Gnade und Herrlichkeit Gottes geschmücket werdest, wie die heilige Elisabeth die heilige Jungfrau selig sprach, da sie sagte [Luc. 1, 42. 45.]: „Gebenedeiet bist du unter den Weibern, selig bist du, die du geglaubet hast; denn es wird vollendet werden, was dir gesagt ist von dem HErrn.“ Denn du mußt den Geruch der göttlichen Gnade, den er über seine Heiligen ausgegossen hat, wahrnehmen und die Tröpflein seiner Güte vorher kosten, damit auch du zur Liebe gegen GOtt entzündet werdest. Denn was sind die Heiligen anders als gleichsam Thautröpflein oder Nachttropfen in den Todten und auf dem Haupte des Bräutigams? wie er spricht im Hohenliede, Cap. 5, 2.: „Thue mir auf, liebe Freundin, meine Schwester, denn mein Haupt ist voll Thaues und meine Locken voll Nachttröpflein.“ Denn so nennt er Micha 5, 6. die Heiligen einen Thau, da er sagt: „Es werden die Uebrigen aus Jakob unter vielen Völkern sein, wie ein Thau vom HErrn“, denn alles, was sie sind und thun, ist Thau und himmlische Gnade, welche die überaus liebliche Barmherzigkeit GOttes über die Kinder der Menschen haucht. So, ja so muß man alle Ehre der Heiligen Gotte heimbringen, der ihnen diese Gewalt gegeben hat. Siehe die Kirche an, welche von der heiligen Jungfrau singt: Gebenedeiet seist du, o Tochter, (aber) vom HErrn, weil wir, nicht von dir aber durch dich, theilhaben an der Frucht des Lebens; und an den Festen aller Heiligen richtet die Kirche ihr Gebet nicht an die Heiligen, sondern zu Gott mit den Namen

\quellseite{594}
der Heiligen, indem sie bezeugt, daß die Verdienste derselben aus GOtt gekommen sind, und sodann durch diese Verdienste GOtte ihr Gebet befiehlt. Denn erst nachdem dies Opfer des Lobes in der rechten Weise vollbracht ist, können wir unser Gebet daran hängen, indem wir zuerst um eben solche Gnade GOttes bitten, und zuletzt um zeitliche Dinge, wenn ja ein Mensch um zeitliche Dinge beten kann, der nach den geistlichen dürstet.

Aber dieses Anrufen und diese Verehrung der Heiligen ist so völlig vernachlässigt, daß man nicht allein die Werke der Barmherzigkeit GOttes an seinen Heiligen nicht bedenkt, sondern auch nicht einmal ihre Thaten kennt, ja, auch nicht ihre Namen, wie jener, der die heilige Dreifaltigkeit „drei Falten“ nannte. So gar betrachten wir nicht, was für Gutes GOtt an den Heiligen gethan hat, damit er verherrlicht würde, sondern was für Gutes sie uns thun, damit wir erfreut werden. Aber der Heilige Geist und die Kirche Christi haben bei aller Verehrung und den Festen der Heiligen sicherlich dies im Auge, nämlich daß wir unser Herz zu GOtt und den göttlichen Werken erheben sollen, und wir aus der Betrachtung der Gnaden, die ihnen verliehen sind, eine Zuversicht [zu GOtt] erlangen, wie Christus sagt Matth. 5, 16.: „Also lasset eure guten Werke leuchten.“ Wozu? Daß mir mein Bein nicht wehtue? Nein, sondern „daß sie euren Vater im Himmel preisen“. Hier sagt St. Augustinus: sie sollen nicht verzweifeln; es könne dies auch ihnen widerfahren, wenn sie nur wollen. Und es haben freilich die Heiligen nach dieser Regel und zu dem Zwecke gelebt und Gutes gethan, damit ihr Licht leuchte zur Ehre GOttes, und es geschieht durch Wirkung des Teufels, daß ihr Leben und ihre Werke nicht zu dem Zwecke gepredigt oder gefeiert werden. Sie haben ihr ganzes Lebenlang nur das gesucht, was GOttes ist, wir aber suchen nicht allein in unserem Leben, sondern auch in dem Leben der Heiligen nur das, was unser ist, und kümmern uns nicht um die Ehre GOttes. Und überall wird nur gepredigt, was und wie viel sie vermögen zu unserem Vortheil, aber nicht, wie viel Barmherzigkeit GOtt in ihnen bewiesen hat. Daher versteigt man sich bis zu dem äußersten Frevel und festen Streite, welcher Heilige im Himmel der höchste sei vor GOtt, gleich als ob sie selbst irgend etwas gethan hätten und nicht vielmehr GOtt; wir haben unsere Augen ganz und gar ins Fleisch versenkt und von GOtt abgewendet.

Ja, zu unserer Zeit (was noch abscheulicher ist) ist es mit der Verehrung der Heiligen dahin gekommen, daß es besser wäre, daß man weder die Feste der Heiligen hielte, noch ihre Namen wüßte. Um das zu erkennen, so durchlaufe den Aberglauben des thörichten großen Haufens und siehe, wie verschiedene Schutzheilige die verschiedenen Handwerker haben. Die Goldschmiede haben den heiligen Eligium\footnote{So in der Weimarischen, Erlanger und bei Löscher. Im Manuscript, in der Baseler Sammlung und in der Loffherschen Ausgabe: Elogium; allein einen solchen Heiligen gibt es nicht. Wittenberger und Jenaer: Eulogium.}, die Schuster den Crispin und Crispinianus, die Walker den heiligen Severus, die Maler den heiligen Lucas, die Aerzte den Cosmas und den Damianus, die Juristen St. Ivo, die sich mit den freien Künsten beschäftigen (artistae) St. Catharina und bisweilen auch den Aristoteles zu Schutzheiligen; andere haben wieder andere, und die Kranken haben ihren Kilian; eine jegliche Nation ihren Heiligen. Aber nun siehe, wie sie diese ehren. Erstlich kümmern sie sich nicht um ihre Werke und Exempel. Zweitens, wenn es mit der Ehre aufs höchste kommt, so halten sie früh Morgens eine Messe und feiern diesen ganzen Tag. Doch zeigen sie dies durch kein anderes Zeichen als durch ihre Kleider und Müßiggehen, indem sie eine Feier mehr erheucheln als halten, da sie im Uebrigen die Zeit todtschlagen in Ausgelassenheit, mit Singen und Springen (indem ihr Herz niemals so fern ist von GOtt als an diesem Tage). Ja, sie spotten der Heiligen, als wenn sie Bachanalien oder Saturnalien feierten, und dies setzen sie fort bis in den andern Tag mit Trinken, Spielen und, um es kurz zu sagen, mit einer solchen Ausgelassenheit, daß es bisweilen zu Mord und Todtschlag kommt. Denn dies ist die Herrlichkeit der Festtage zu diesen bösen Zeiten, daß auch ein Anubis\footnote{Anubis, eine egyptische Gottheit mit einem Hundskopfe.} oder irgend ein Greuel der heidnischen Götter nicht schändlicher entehrt worden ist, ja, daß auch dein Schwein nicht dulden würde, daß man es so ehre. Aber GOttes Zorn ist so gewaltig geworden, und er hat sie in verstockten Sinn dahin gegeben, daß alles ganz verkehrt sein muß.

\quellseite{595}
Wie viel richtiger wäre es, wenn man mit demüthigem und ruhigem Herzen und im stillen Hause sich mäßiglich erquickte und die Armen herbeirief und ihnen ein Mahl zurichtete; ja, es wäre besser, man arbeitete und thäte auch alle niedrige Arbeit, als daß man so feierte. Denn nicht durch dieses unordentliche Wesen mit Prassen und Schwelgen, sondern durch ruhiges Verhalten (darum nennt man es auch den Sabbath feiern [sabbatissare]) und Liebe gegen den Nächsten dient man den Heiligen. Wenn du es nun selbst als etwas äußerst Schändliches verurtheilen würdest, so jemand am Osterfeste oder am Christtage dergleichen beginge, warum ehrst du denn damit deinen Schutzheiligen, den du gerade dann am meisten zu ehren dir vorgenommen hast? Oder warum wird er, der das ganze Jahr ruhig und sicher war vor deinen Spöttereien und Schmähungen, gerade an diesem Tage damit gequält, da es am wenigsten sein sollte? Warum fluchen wir den Türken, von denen gesagt wird, daß sie die Bilder Christi und der Heiligen schändlich beschimpeln? Oder verstehen wir noch nicht, daß GOtt uns durch deren Vorbild (figura) anzeigt, wie sehr es ihm mißfalle, daß wir Christen seine Heiligen in Wahrheit noch schändlicher beschimpeln? Welch ein schöner Gottesdienst und Frömmigkeit ist doch das! Die Heiligen GOttes lästert niemand schändlicher als gerade der, dessen Patron er ist, und zwar an seinem anderen Tage mehr als an dem, der seiner Verehrung geweiht ist. Stelle dir vor, ob du eine solche Verehrung, das heißt, Verspottung, wohl leiden könntest, auch von dem geringsten Menschen.

Hier sollten daher die Bischöfe aufs höchste wachsam sein; ja, es wäre sehr heilsam, wenn alle Festtage [der Heiligen] abgethan würden und mit sehr wenige beibehalten, an denen man früh Morgens das Wort GOttes hörte und Messe läse. Denn unsere Feiertage sind nichts Anderes als eine Schmach für die Kirche, und nicht eine Gottesfurcht, sondern eine Gottlosigkeit gegen GOtt, wie die Erfahrung genug und übergenug lehrt. Und GOtt kann zu uns mit viel größerem Rechte sagen als zu den Juden [Amos 5, 21.]: „Ich bin euren Feiertagen gram und verachte sie“, und das Wort Jes. 1, 14. 13.: „Meine Seele ist feind euren Neumonden und Jahrzeiten. Der Neumonden und Sabbathe, deren mag ich nicht.“

\footnote{Löcher: „Das Folgende hat Luther gepredigt im Jahre 1516 am 10. Sonntage nach Trinitatis [den 27. Juli]. — Die Exordien hiezu finden sich bei Walch, St. Louiser Ausg., Bd. XII, 1726. Bd. X, 1284 und Bd. XIX, 736.} Nun ist noch ein zwiefacher Mißbrauch in der Verehrung der Heiligen übrig.

Der erste besteht in der Hofsucht und dem Mißbrauch durch eitle Ehre, in welcher etliche auch unter einander zanken über die Heiligen, da sich jeder bemüht, den Heiligen, der sein Geschäft betrieben hat (sui studii) oder zu seinem Orden gehört, so berühmt zu machen, daß er ihn\footnote{Neumärsche: cum statt eines Druckfehlers.} allen anderen vorzieht, nicht sowohl, weil sie wollen, daß die Heiligen GOttes geehrt werden, sondern sie selbst, welche so große Leute zu Patronen haben. Wiewohl dieses nichtige Vornehmen (vanitas) klar zu Tage liegt und sehr viele damit umgehen, wage ich es doch nicht, weiter darauf einzugehen, eines theils, weil sie, mit dem schönsten Scheine geschmückt, sich leicht unterstehen möchten, mich als einen frevelhaften Richter auszuschreien, und sich auf das reinste zu entschuldigen, andererseits, weil diese Sache auch die Vornehmen und Großen in der Kirche anbetrifft (tangit), welche viel lieber leiden können [, daß man sie antaste], als der Himmel selbst. Denn auch gerade hierdurch hat es seinen Anfang genommen, daß ein so großer Schwarm von Brüderschaften\footnote{Gegenhard Pessinger, Rath des Churfürsten Friedrich des Weisen, gehörte zur Zeit seines Todes (1519) zu sechs und dreißig Brüderschaften. Das Verzeichniß derselben findet sich bei Kolde, Friedrich der Weise, S. 74 f.} aufgekommen ist, und ich weiß nicht, ob sie nicht durch so viele Brüderschaften es endlich dahin bringen, daß sie weder Brüder noch Schwestern behalten. Aber hierüber mögen andere urtheilen. Der Schein ist gut; man darf nicht allzuscharf strafen, bis daß sich diese Angelegenheit selbst klarer an den Tag gibt.

Ich möchte sagen, daß diejenigen hieher gehören, welche wollen, daß das Fest der Empfängniß der heiligen Jungfrau das erste und höchste von allen sei, und die, welche ihre ganzen Predigten füllen mit der Würdigkeit, den Vorrechten und dem Ablaß neuer Feste. Die Nachwelt wird sehen, daß auch wir nicht vergebens gelebt haben, denn wir haben mit neuen Festen es ausgerichtet, daß die alten in Wahrheit veraltet sind; wenn wir nichts Anderes ausgerichtet haben, so ist dies Eine genugsam ewigen Gedächtnisses werth.

\quellseite{596}
Ich hoffe, daß es dazu kommen wird, daß auch noch ein Fest Abrahams gefeiert werden wird, weil in seinem Glauben der Anfang des Heils war und der verheißene Christus. Denn wenn GOtt dem Abraham nicht die Verheißung gegeben hätte, so wäre Christus nicht geboren worden. Und vielleicht ist es verziehen, daß auch der Tag Adams und Evas gefeiert werden muß. Und so sind die Tage gekommen, an denen wir alles gar gerne lehren und hören, außer dem, was alte und bewährte Wahrheit ist, wie der Apostel sagt [2 Tim. 4, 4. 3.]: „Sie werden die Ohren von der Wahrheit wenden, werden sich aber zu den Fabeln kehren, indem sie sich selbst Lehrer aufladen, nach dem ihnen die Ohren jücken.“ Nicht als ob ich eine gottseelige Neigung zu neuen Festen verwerfe, sondern daß ich sehe, daß nicht wenige derselben gehalten und hochgehoben werden, daß man sogar sucht, andere verhaßt zu machen. Die Heiligen bedürfen unserer Hoffahrt und Zwietracht nicht; ja, nicht sie, sondern wir machen Abgötter aus uns selbst und ehren uns selbst unter dem Namen der Heiligen.

Der andere Mißbrauch kommt vom Geize her. Es werden überall Kirchen errichtet und sind schon errichtet über Bildern, die nicht geweiht sind, besonders da, wo der leicht verführbare Haufe des gemeinen Volks zuströmt. Und auch die Kirchen selbst werden nicht geweiht; in denselben wird nichts als Gewinn gesucht, nur daß der Schein und der Name der Gottseeligkeit da ist. Denn wenn der Teufel unter dem Namen Christi falsche Christus findet, wie viel mehr kann er unter dem Namen und dem Bilde der Maria oder eines andern Heiligen Irrthum und Täuschung wirken! Und zwar um so freier, je nachlässiger wir sind in der Beachtung des Rathes Pauli [1 Thess. 5, 21.]: „Prüfet alles, und das Gute behaltet“, und dessen, was Johannes sagt [1 Joh. 4, 1.]: „Prüfet die Geister, ob sie von GOtt sind.“ Wir aber begnügen uns mit einem schönen Schein, und laufen sofort hinzu, ohne zu prüfen. Dies Hin- und Herlaufen richtet der Teufel an: erstlich, damit er das Volk von den geweihten Stätten abziehe und ihr Herz dem Heiligthum entfremde und auf weltliche Dinge lenke, indem er sie immer mit neuen und fremden Dingen reizt und neugierig macht (denn das haben diejenigen verdient, welche ihrer Kirche, die sie zu Hause haben, überdrüssig sind). Zweitens, damit sie dadurch gleichsam eine gerechte Ursache bekommen, das Wort GOttes und den Gottesdienst in den Pfarrkirchen zu vernachlässigen. Ich will davon nichts sagen, daß sie durch die Reisen die Zeit vergeuden und ihr Vermögen verschwenden, und sich mit vielen Sünden beladen, indem sie viele unnütze Dinge schwatzen, oder hören, oder sehen. Aber ihr Haus, ihr Geschäft (rem), ihre Familie vernachlässigen sie, das ist, vergebens mühen sie sich ab und begehen viel Böses. Daß aber solch Zusammenlaufen vom Teufel sei, zeigt außer dem, was bereits gesagt ist, das ungestüme Zusammenlaufen des Pöbels selbst; denn man erzählt, daß Frauen und Dienstboten, während sie arbeiteten oder ihren Geschäften nachgingen, plötzlich entzückt (raptos) worden seien und gleichsam gezwungen, ihre Geschäfte zu verlassen, was unmöglich durch den Heiligen Geist geschehen kann, sondern ein ganz offenbarer Betrug des Teufels ist. Denn der heilige Geist ist nicht ein Geist, der zu frevelhaftem Verhalten und ungestümem Wesen treibt, sondern ein Geist des Rathes. Auch widerruft er nicht sein Gebot, in welchem er dem Weibe befohlen hat, daß es unter der Gewalt des Mannes sein solle, und daß ohne seine Einwilligung ihre Gelübde nichtig sind; und den Knechten, daß sie ihren Herren treulich dienen sollen und nichts veruntreuen: wie sollte er denn nun ebendieselben zum Ungehorsam hinwegreißen, den er so strenge verboten hat? Daher habe ich gehört, daß, wo solche Kirchen geweiht und zum ordentlichen Gottesdienste verordnet worden waren, ein solches Zusammenlaufen aufgehört habe. Aber man sagt, daß auch jetzt noch etliche darum nicht geweiht werden, damit der Zulauf nicht aufhöre und der Gewinn versiege. O über die Christen, welche durch solche Gottlosigkeit Gewinn suchen! Ist denn das aus GOtt, was das Kreuz und den Segen nicht leiden kann? Wer hätte je glauben können, daß Christen selbst durch den listigsten Teufel in diesen Abgrund des Irrthums gestürzt werden könnten!

Aber auch ich habe einen Menschen gekannt, welcher anfing, ein derartiges Heiligthum, das zum Betrüge dient, aufzurichten unter dem Namen der Jungfrau Maria, die ihm im Schlafe erschienen war (so gar heilig und würdig meint die Hoffahrt zu sein); aber als ihm dies von den Vorgesetzten der Kirche gewehret wurde, verhielt
\quellseite{597}
sich der Teufel an dem Orte alsbald still, und er wandelte anderswohin und setzte es durch, daß er dort sein Hurenhaus aufrichtete. Aber dies sind in Wahrheit Bilder und Götzen der Hurerei, die der HErr durch den Propheten [Hos. 9, 1. Ps. 73, 27.] straft, in welchen die gläubigen Seelen Christum verlassen durch ihre Hurerei. Denn er hat verboten, denen zu glauben, die uns auf gewisse Stätten hinweisen und uns predigen [Matth. 24, 23.\footnote{Weimarsche ebenso wie Walch: Matth. 24, 26.} 25.]: „Siehe, hier ist Christus, oder da.“ Glaubet ihnen nicht, sagt er, „siehe, ich habe es euch zuvor gesagt“. Darum soll es dir billig verdächtig sein, wie die Pestilenz, die im Finstern schleicht [Ps. 91, 6.], einen wie köstlichen Schein das auch haben mag, was an diesen ungeweihten Oertern vorgenommen wird. Denn nicht vergebens hat die Kirche verordnet, daß ihre Kirchen und Stätten geheiligt werden sollen durch das Wort und Gebet. Wenn nun jemandes Frau oder Knecht ihm sagen sollte, sie hätten einen solchen Geist zum Wallfahrten überkommen, so höre meinen Rath: so nimm auch du ein gutes eichnes Marterholz (crucem) und heilige ihren Rücken mit etlichen kräftigen Schlägen, und du wirst sehen, daß durch diesen Finger GOttes jener Teufel ausgetrieben wird. Denn da die Weiber leicht zu verführen sind, so pflegt er sie am häufigsten durch diese Possen zu betrügen.

Drittens. Wenn er diese Gewalt nicht an ungeweihten Stätten hat (denn er hat sie in Wahrheit nicht überall), so thut er dies an heiligen Orten. Dies thut er aber, wenn er dazu reizt, die Mutterkirche zu verlassen, und zu Wallfahrten zu Kirchweihen und Heiligenbegängnissen (patrocinia) anderer Kirchen, die entweder an wüsten Orten oder in Dörfern gelegen sind, wo du aber so viele Schenken und unordentliche Häuser sehen kannst, daß du glauben möchtest, es sei ein Babylon (denn in solcher Weise hält man heutzutage die Kirchweihfeste und die Feste der Schutzheiligen); aber am Abend kehren sie nach Hause zurück voll von Ablaß, das heißt, voll von Bier, Schwelgerei und anderen greulichen Sünden, denen sie dort gedient haben, wenn sie überhaupt nach Hause gehen; denn es kommt fast regelmäßig vor, daß auf den Kirchweihen etliche getödtet oder wenigstens schwer verwundet werden, in solchem Grade, daß auch ein Sprüchwort daraus geworden ist: Der Ablaß auf den Kirchweihen wird spät Abends ausgetheilt, weil sie sich dort todtschlagen oder wenigstens verletzen. Solche überaus jammervolle Schmach mehrt der Teufel fort und fort: und wir können noch in der Meinung stehen, daß es besser sei, diese Kirchweihen würden nicht abgethan, sondern sollten statthaben? Gebe doch der HErr endlich einmal unseren Bischöfen die Gnade, daß sie diesen Ablaß sammt den Kirchweihen widerrufen möchten, und nicht zuließen, daß der Teufel uns, da wir ermordet werden, noch verspottet, denn es wäre ihm damit schon allzuviel gestattet, daß er unser spottete durch Irrthum.

Aber hier werden mir die Wallfahrten entgegengehalten, welche [von der Kirche] angenommen worden sind (receptae), zu St. Peter in Rom, zu St. Jakob, nach Jerusalem, nach Trier, ja, auch nach verschiedenen Orten, um die Reliquien der Heiligen zu ehren und Ablaß zu erlangen.

Ich antwortete: Der würde nicht sündigen, welcher alle diese aufstehen ließe und zu Hause bliebe; sie sind nicht geboten, sondern man übernimmt sie freiwillig. Doch widersetze ich mich nicht dem, was nun durch langen Brauch und Gewohnheit eingerissen ist; diejenigen, welche wollen, mögen immerhin gehen, nur sollen sie das lernen, daß sie mit denselben Kosten unterdessen unvergleichlich besser GOtte dienen und ihre Seligkeit schaffen könnten, wenn sie zu Hause den Armen dienen, ja, denen, die ihnen befohlen sind, als, ihrem Weibe, Kindern, Gesinde, Oberherrn, sodann auch die Züchtigungen GOttes und irgendwelche Widerwärtigkeiten geduldig trügen. Es ist aber, als ob wir daheim nicht genug mit dem Teufel zu streiten hätten, wenn wir nicht wie Dina [1 Mos. 34, 1. 2.] herausgehen, und mehr Gelegenheit suchen, zu Schaden zu werden. Denn wie viele widerstehen den spitzfindigen Reden und Sorgen, mit denen er unsere Herzen in der heimathlichen Kirche anficht? Ferner mit dem Saufen, dem Spiel, dem Müßiggange plagt er uns alle Tage, und versucht alles, damit der Gottesdienst bei uns nicht rein bleibe, und unsere Seligkeit Schaden leide.

Auch das hat unsere Geldgier neulich erfunden, daß sie sieben goldene Messen (nämlich sie sind so genannt von den Goldmünzen, mit denen sie
\quellseite{598}
bezahlt werden) mit gewissen Wachskerzen und anderen Gebräuchen gehalten werden, ich weiß nicht für was für ein ganz gewiß eintretendes Ereigniß. So hat uns der schändliche Feind auch nicht einmal dieses Opfer unverletzt gelassen.

Wie? wenn nun auch dieser bodenlose und unerträgliche Geiz des Geizes hieher gehört, da man zum Kirchbau opfert, bettelt, und ein Schatz ohne Ende zusammengehäuft wird, und man zu diesem Zwecke die Reliquien und den Ablaß aufs listigste gebraucht: so daß, wo das Geld zusammengebracht ist, es nicht ausgegeben wird, weder für den Kirchbau noch für die Armen, sondern für Kriege oder Rechtshändel, die schon da sind, oder aufbewahrt wird für solche, die noch kommen werden? Denn so haben sie auch die Nägel Christi zu Dienern des Geizes gemacht, doch unter einem überaus heiligen Schein der Gottseligkeit, daß sie die Gold- und Silbermünzen [mit denselben] durchbohren. Warum durchbohren sie damit nicht lieber Papier oder ein Brett, wenn nur gesucht wird, daß man eine ehrerbietige Handlung (devotio) vornehme? Aber was unterstehe ich mich, den Pfuhl des Geizes in der Kirche zu beschreiben, da der Prophet selbst daran verzweifelt, und alles in Ein Stück zusammenfaßt und spricht [Jer. 6, 13.]: Alle folgen dem Geize nach vom Größten bis zum Kleinsten.\footnote{Löscher bemerkt: „Dies ist wider Tezel und seinen Ablaßkram gesagt. Darnach hat Luther insbesondere (ex professo) über den Ablaß geredet, was nachher aus dem Manuscripte mitgetheilt werden wird.“ In unserer Ausgabe findet sich dieser Passus Bd. XIX, 736 ff.}

Dieß soll schließlich die Summa von allem seyn:

Alle diese Verehrer der Heiligen sündigen wider das erste Gebot, weil sie nicht das, was GOttes ist, sondern das Ihre suchen, auch an GOtt und seinen Heiligen, und sich selbst bei diesem Werke der letzte Zweck (wie man sagt) und ein Abgott sind, indem sie GOttes gebrauchen und sich selbst dienen. Von diesem kann das Wort des 78. Psalms [V. 34. ff.] gesagt werden: „Wenn er sie erwürgte, suchten sie ihn und kehrten sich frühe zu GOtt, und gedachten, daß GOtt ihr Hort ist, und GOtt der Höchste ihr Erlöser ist (aber nun folgt, aus was für einem Herzen sie dies gethan haben), und heuchelten ihm mit ihrem Munde, und logen ihm mit ihrer Zunge; aber ihr Herz war nicht fest an ihm und hielten nicht treulich an seinem Bunde“; und jenes Wort Christi Joh. 6, 26.: „Wahrlich, wahrlich, ich sage euch, ihr suchet mich nicht darum, daß ihr Zeichen gesehen habt, sondern daß ihr gegessen habt und seid satt geworden.“ Aus derselben Ursache trieb er auch den von sich, der da sprach [Matth. 8, 19.]: „Meister, ich will dir folgen, wo du hin gehest.“ Er sagte aber [V. 20.]: „Die Füchse haben Gruben, und die Vögel unter dem Himmel haben Nester; aber des Menschen Sohn hat nicht, da er sein Haupt hinlege“, womit er freilich anzeigt, daß jener nicht allein um JEsu willen, sondern auch um seiner Habsucht (cupiditatem) willen dies gesagt hätte.

Es steht also das Urtheil fest [Matth. 16, 24.]: „Will mir jemand nachfolgen, der verleugne (nicht Geld, nicht Kleidung, ja, auch nicht Gesundheit, sondern) sich selbst, und nehme sein Kreuz auf sich und folge mir.“ Desgleichen [Luc. 14, 33.]: „Wer nicht absagt allem, das er hat, kann nicht mein Jünger sein.“ Desgleichen [Luc. 14, 26.]: „So jemand nicht hasset Vater und Mutter, auch dazu sein eigen Leben, der kann nicht mein Jünger sein.“

Ferner, da jene auch ihren Leib suchen und was des Leibes ist, so ist es gewiß, daß sie Christi nicht werth sind; jedoch sind sie einiger Gaben werth, wie die Lohnknechte, gleichwie die Kinder Abrahams von der Hettura [1 Mos. 25, 6.], aber nicht des Erbes mit Isaac. Denn sie aber Christi nicht werth sind, so ist es klar, daß sie seinen Willen nicht thun, auch sein Gebot nicht halten.

Dies sage ich nicht, daß ich wollte, alle diejenigen seien verdammt, welche solche Leute sind, sondern allein, daß sie lernen ihr Elend erkennen, und sich nicht darauf verlassen, daß sie um deswillen Christen seien, weil sie GOtt und die Heiligen in solcher Weise ehren und anbeten, vielmehr daß sie ihre Unvollkommenheit sehen, und verstehen, was mit dem Seufzen gemeint sei, von dem Christus uns gelehrt hat, daß wir es bekennen sollen, und darnach uns befohlen hat zu beten [Matth. 6, 12. 10.]: „Vergib uns unsere Schulden“ und: „Dein Wille geschehe“ etc., weil wir nämlich solche Leute nicht sind, wiewohl wir doch so sein sollten. Darum, wenn man den Rath GOttes fassen kann, so gibt er diese [leiblichen] Güter nicht deshalb, daß er mit denselben
\quellseite{599}
sättigen wollte, sondern mehr reizen durch diese kleinen Dinge, daß entweder die Unvollkommenen oder wenigstens seine Kinder der größeren begehren, denn er solches nicht gibt. Und dies sei genug von dieser Sache.

Aber damit die unseligen Ketzer, die Picarden, nicht denken, daß ich ihrer Sache Vorschub leiste, welche aus großer Unwissenheit, entrüstet in hochmüthigem Dünkel, uns Deutschen auflegen, daß wir die Heiligen GOttes verehren und Abgötter aus ihnen machen, und deshalb wider uns einen großen Haufen von Schriftprüchen sammeln, in denen verboten wird, mehr als Einen GOtt anzubeten, und, damit sie uns bei den Ihrigen, wie es scheint, mit gutem Grunde verhaßt machen, auf das hinterlistigste nicht anzeigen, daß geschrieben steht, der König David und Salomo und viele andere seien angebetet worden, zugleich als gottlose Verführer der Schrift und hinterlistige Verleumder unseres Christenthums (pietatis) (denn so lehren uns diese Bauern endlich, daß man nur Einen GOtt anbeten müsse, und rühmen sich, als ob wir gerade das jemals geleugnet hätten, während sie doch nicht leugnen können, daß häufig auch die Beamten am Hofe eines Königs geehrt und gleichsam angebetet werden, um desto leichter zum Könige zu gelangen): so sage ich um des überaus großen und ungereimten Unverstandes dieser Leute willen: daß man jedenfalls seine Zuflucht nehmen solle zu den Fürbitten der Heiligen, wie es bei Hiob [Cap. 5, 1.] heißt, und sich an irgend einen der Heiligen wenden, und gleichwie Salomo seinen Vater anführt [Ps. 132, 1. Vulg.]: „Gedenke, HErr, an David, und an alle seine Sanftmuth.“ Aber auch der Patriarch Jakob sagte von Ephraim und Manasse [1 Mos. 48, 16. Vulg.]: „Mein Name soll angerufen werden über diese Knaben, und die Namen meiner Väter Abraham und Isaak.“ Aber hiervon will ich seiner Zeit reden, so Gott will.

\footnote{„Dies ist sicher ein späterer Zusatz aus der Zeit, als Luther den Druck vorbereitete: leider schweigt hier Löscher über die alte Handschrift.“ (Weim. Ausg.)} Hier hätte auch das Vertrauen auf den Ablaß gestraft werden sollen, wenn dasselbe nicht schon mehr als zu viel auf mancherlei Weise von mir gestraft worden wäre.

\footnote{„Was hier zunächst folgt, scheint am 1. August 1516 gepredigt zu sein.“ (Weim. Ausg.) — Das Exordium findet sich Walch, St. Louiser Ausg., Bd. XII, 1726 ff.} Jetzt muß ich von der andern Art der Uebertretung sagen, nämlich, welche sich bei denen findet, die anstatt des wahren GOttes den Götzen ihrer Weisheit und Gerechtigkeit anbeten. Dies sind „die Stolzen“, deren König, wie GOtt, Hiob 41, 25., sagt, der Behemoth, der Teufel, ist. Bei diesen findet sich das Vertrauen auf den freien Willen [die Meinung], daß es nicht übel um sie stehe, daß sie eines Arztes nicht bedürften. Allein diesen ist Christus vergeblich gestorben, weil sie ohne ihn ein gutes Leben führen können. In solcher Weise sind vor Alters besonders die Pelagianer, sodann auch alle Ketzer verloren gegangen. Ferner gibt es auch heutzutage Leute, die zwar nicht Pelagianer sind, aber doch dieselbe Gesinnung oder noch eine ärgere haben als die Pelagianer. Denn die Pelagianer schrieben Gotte wenigstens das zu, daß er durch die Lehre des Gesetzes den freien Willen unterwiese; diese aber haben die Vernunft, die sie recht leitet,\footnote{Vgl. Luthers ausführliche Erklärung des Galaterbriefes (Cap. 2, § 347. Walch, St. Louiser Ausg., Bd. IX, 237, § 347.} welche nach dem Besten strebt, so daß sie der Unterweisung des Gesetzes nicht gebrauchen aus Noth, sondern nur, weil es ihnen bequem ist (pro commoditate). Doch darin weichen sie von den Pelagianern ab, daß sie zugeben, daß man ohne Christum nicht ein gutes Leben führen könne in der Weise, daß es verdienstlich sei, und so ist Christus gestorben, nicht um der Sünde willen, sondern weil sonst kein Verdienst da wäre. Es wäre nicht vonnöthen gewesen, daß er stürbe, um die Höllenstrafen zu bezahlen, sondern nur, damit man den Himmel verdienen könnte. Denn stelle dir vor, wenn ein Kind aufwüchse ohne eine wirkliche Todsünde (denn das, behaupten sie, sei zwar schwer, aber doch möglich), so ist ihm Christus nicht vonnöthen als ein Erlöser von der Gewalt des Teufels, sondern nur als ein Helfer zum Himmel, denn wenn ein solcher Mensch stürbe, käme er weder in die Hölle, noch in den Himmel. Durch diese Theologie scheinen sie mir zu behaupten, daß niemand als ein Kind des Zorns und ein Kind des Teufels geboren werde, es sei denn, man wollte sagen, ein Kind des Zorns und des Teufels sei der, welcher kein Verdienst habe. Und da der Priester bei der Taufe der kleinen Kinder den Satan ausbläst, daß er Raum gebe dem Heiligen Geist, so versteht man darunter, daß er aus einem solchen Gefäße, welches sein

\quellseite{600}
Verdienst habe, nur ein solches mache, das Verdienst habe. Denn erst jetzt wird der Apostel dahin gezwungen, daß er in solcher Weise „die Gefäße des Zorns“ [Röm. 9, 22.] verstanden habe; und das Sacrament der Taufe gibt Gnade, nicht zur Vergebung der Sünde, sondern nur zu einem Anfange des Verdienstes, und das Wort des [nicänischen] Glaubensbekenntnisses: „Ich bekenne eine einige Taufe zur Vergebung der Sünden“ muß nach der neuen Glosse nicht verstanden werden von den Sünden, die zur Verdammniß führen, sondern von den Sünden, die nicht verdienstlich sind. Und so hat die Erbsünde keinen Schaden in die Welt gebracht, sondern nur des Verdienstes und des Schauens GOttes beraubt, und dies Beraubtsein des Schauens (das ist, aller Güter) ist nicht ein Uebel [in der That], sondern nur dem Worte nach, nämlich privative [es zeigt nur das Beraubtsein an]. Denn in solcher Weise treiben wir heutzutage Theologie. Und [Gal. 3, 13.]: „Christus ward ein Fluch für uns“ und [2 Cor. 5, 21.] „ist für uns zur Sünde gemacht“, damit er uns vom Fluch und von der Sünde befreie, das wird nicht anders verstanden, als daß er gemacht ist zu etwas Unverdienstlichem (non meritorium), um uns von dem Unverdienstlichen zu befreien. Schließlich wird das Wörtlein „Sünde“ in der ganzen Schrift eine neue Bedeutung empfangen, das heißt, die Schrift wird ein ganz verwirrtes Buch werden (induct confusionem). Doch wir wollen diese Finsterniß fahren lassen und für eine andere Zeit aufsparen, weil wir auch anderswo sehr oft von diesen Dingen geredet und gehandelt haben.

Nun, wie wir oben gesagt haben, da wir, wie St. Augustinus spricht, Ein verderbter Haufe (Una massa perditionis) sind, so wird niemand anders gerechtfertigt als durch die lautere Gnade Christi, ohne alles Verdienst. Und darum muß nothwendiger Weise derjenige hochmüthig sein und viel von sich halten, welcher noch nicht verzweifelt hat an sich selbst, seiner Vernunft, seiner guten Meinung, an seinen innerlichen und äußerlichen Werken, und zitternd und elend zum Thron der Gnade hinzugetreten ist. Denn alle jene sittlichen Tugenden und Wissenschaften, die in menschlicher Erkenntniß stehen (scientiae speculativae), stürzen den Menschen in ein zweifaches Uebel: erstlich, daß sie nicht wahre Tugenden und Wissenschaften sind, sondern Sünden und Irrthümer, weil sie mit Nothwendigkeit aus einem bösen Herzen herkommen, das noch nicht durch die Gnade gesund gemacht ist, welches in allen Dingen nichts Anderes suchen kann als sich selbst, und so machen sie den Menschen leer und eitel (inanem). Zweitens, daß sie den Menschen betrügen und aufblähen machen, so daß er sich dünken läßt, er sei weder böse noch thöricht. Daher kommt es, daß er sicher ist, und sich nichts weniger angelegen sein läßt, als das Suchen der Gnade, da er ja weiß, daß er nichts Böses thue, wiewohl er zugibt, daß er in so vielen guten Werken nichts Verdienstliches thue. Daher kommt noch weiteres Uebel, daß er alle verachtet, richtet, verdammt, beißt, wider sie eifert, straft, herabsetzt, die nicht sind wie er, denn er ist wie jener Pharisäer [Luc. 18, 11.] „nicht wie andere Leute“, noch einer von ihnen. Diese sind es, die ihr Haus auf den Sand bauen [Matth. 7, 26. 27.], das in erschrecklicher Gefahr steht und einen Fall thun wird ins ewige Verderben.

Aber die rechten Diener der Gerechtigkeit wissen und bekennen, daß sie ganz und gar Sünde sind, und wollen, daß alles Gute, was sie haben, nicht in ihnen liege, sondern außer ihnen in Gotte und in seiner Barmherzigkeit, weil der Gerechte seines Glaubens leben wird, aber vor GOtt kein Lebendiger gerecht ist [Röm. 1, 17. Ps. 143, 2.]. Darum richten sie, niedergedrückt durch ihre Sünde, niemanden als sich selbst, setzen niemanden herab, verachten niemanden als sich selbst und halten sich nach dieser Erklärung (diffinitionem) von Gerechtigkeit: Gerechtigkeit besteht darin, daß man sich selbst zuerst anklage, und ein Gerechter klagt zuerst sich selbst an. Diese haben in Wahrheit Einen GOtt, aus dem, durch den und zu dem (in quem) sie gerecht werden, und sind ohne Sünde durch die Barmherzigkeit GOttes, der ihnen die Sünde vergibt, nicht durch ihre thätige Gerechtigkeit. So wird GOtt gepriesen und in Wahrheit verehrt, da seinen Werken alles zugeschrieben wird, was wir sind, ja, da unsere Werke nicht unser sind, sondern GOttes, wie Christus sagt [Joh. 7, 16.]: „Meine Lehre ist nicht mein.“ So bleibt uns nichts übrig als Sünde, Thorheit, Bosheit, Verderben und Schande, und darum können wir uns in keinem Stück gefallen oder einen Abgott aus uns machen, da wir wieder zu dem Nichts geworden sind, aus dem wir hergekommen
\quellseite{601}
sind, und GOtt allein alles in allen bleibt.

\footnote{Vielleicht gehört dies nachfolgende Stück zu der Predigt vom 3. August 1516, deren Exordium sich Walsh, St. Louiser Ausgabe, Bd. X, 1288 ff. findet; nur wird es für den Druck umgearbeitet sein. Man vergleiche u. a. Luthers Urtheil hier und dort über die Erklärung der spes (Hoffnung) [l. c. fol. 1290, z 4] beim Magister sententiarum. (Weimarische Ausgabe.)} Aber wenn wir dies sagen, so wirft man uns vor: Also braucht man nicht gute Werke zu thun! gleichwie dem Apostel, der ebendaselbe lehrte, auch vorgeworfen wurde [Röm. 3, 8.]: „Lasset uns Uebeles thun, auf daß Gutes daraus komme; welcher Verdammniß ist ganz recht.“

Darum muß diesen geantwortet werden: Erstens, daß es nicht gelehrt, sondern aufs eifrigste gelehrt wird, daß man Gutes thun solle. Der Streit beruht nur auf dem Worte „Gutes“, denn [Hebr. 11, 6.]: „Ohne Glauben ist es unmöglich GOtt gefallen“, und [Röm. 14, 23.]: „Alles, was nicht aus dem Glauben gehet, das ist Sünde.“ Freilich ist ihr Gutes, was sie so nennen, aus der Natur und aus einem sittlichen Bestreben (moraliter) gethan, darum böse vor GOtt, weil es nicht gethan ist aus GOtt und um GOttes willen (das heißt, es wird nicht erkannt [aus GOtt] und nicht bezogen auf GOtt), sondern der Mensch legt es sich bei und gefällt sich darin, als wäre es sein, wiewohl es nicht sein ist. Denn dieser Zusatz des Leviathan [des Teufels] macht, daß alles vor GOtt böse ist, wie großen Schein es auch vor den Menschen habe, daher ruft Paulus sowohl über die Heiden als auch über die Juden frei öffentlich aus [Röm. 3, 10. 12.]: „Da ist keiner, der gerecht sei, da ist keiner, der Gutes thue, sie sind alle abgewichen, und allesammt untüchtig worden.“ Dazu aber, daß der Mensch erkenne, daß seine Werke allein GOttes seien, erleuchtet ihn die Gnade; die Natur lehrt ihn das nicht. Daher kommt es, daß ein solcher Mensch geduldig und sanftmüthig ist, mögen seine Worte und Werke getadelt oder gelobt werden, indem er weiß, daß er nichts verliert, weil es nicht sein ist. Aber die Natur mühet, wenn sie getadelt wird, und ist unsinnig, wenn sie gelobt wird. Die Ursache von dem allen ist, daß der Glaube und die Hoffnung keine andere Grundlage hat als allein GOtt, Jesum Christum, den rechten Felsen, und durchaus keine Creatur. Die Vernunft aber und die Natur haben einen sandigen Grund, ihre guten Werke und ihre Weisheit. Daher kommt es, daß ein Mensch zürnt, wenn seine guten Werke dahinfallen durch üble Nachrede, fallen sie aber durch die Offenbarung des Gewissens und des göttlichen Gerichtes, so verzweifelt er ewiglich. Aber der Gerechte, wenn er auch fällt, so steht er wieder auf, und wider sein Gewissen stützt er sich auf GOttes Barmherzigkeit, wie geschrieben steht [Ps. 21, 8.]: „Denn der König hofft auf den HErrn, und wird durch die Güte des Höchsten fest bleiben.“

Daher ist die Beschreibung der Hoffnung bei dem Magister Sententiarum [Petrus Lombardus] entweder falsch, oder sie ist dunkel und nicht in verständlicher Weise gegeben, denn die Hoffnung kommt nicht her von den Verdiensten, sondern umgekehrt, die Verdienste kommen her aus der Hoffnung; denn gleichwie der Glaube, so müssen auch die Hoffnung und die Liebe nothwendiger Weise eher da sein als irgend ein verdienstliches Werk. Denn es kommt die Hoffnung nur davon her, daß GOtt sich erbarmt und sie eingießt, und hat keinen andern Gegenstand (objectum) oder Stoff (materiam) oder Grundlage als die bloße Barmherzigkeit GOttes, keineswegs unsere Werke, welche vielmehr der Gegenstand (objectum) und der Ursprung der Verzweiflung sind. Christus hat für uns gelebt und ist unser Verdienst, wenn wir an ihn glauben. Die aber an ihn glauben, leben hinfort nicht für sich selbst, häufen auch nicht ihre Verdienste auf, sondern dienen Christo, und wiederum, machen sich um Christum verdient. Denn so lehrt der Apostel Röm. 14, 8.: „Leben wir, so leben wir dem HErrn, sterben wir, so sterben wir dem HErrn.“ Er ist unsere Hoffnung, und wir (daß ich so sage) sind seine Hoffnung oder die, auf welche er wartet (expectatio ejus), wie ein Sohn die Hoffnung seines Vaters ist, und der Vater die Hoffnung des Sohnes, für den er alles vorsorglich beschafft und Schätze gesammelt hat.

Wie ich aber von der Hoffahrt in geistlichen Dingen gesagt habe, so mag es auch von der Hoffahrt in leiblichen Dingen verstanden werden, als da sind Reichthum, Schönheit, Ehre, Macht, Gunst, Geschlecht, Wohlleben und dergleichen. Denn wer darauf vertraut, und darauf stolz ist, und andere, die ihm nicht gleichstehen, verachtet, bei dem ist es offenbar, daß er diese Dinge zu seinem Eigenthum macht, sie
\quellseite{602}
Gotte raubt, und sie sich als Götzen aufrichtet. Wenn er erkennte, daß es allein Gotte zugehörte, so würde er keinen Unterschied machen zwischen sich und irgend einem noch so verworfenen Menschen, und würde bekennen, daß er nichts mehr habe, nur daß Jesus Christus allein seine Zuversicht ist. So würde es geschehen, daß, wenn dies Gebot gehalten würde, keine Hoffart da sein würde, das heißt, keine Wurzel, kein Anfang der Sünde, und infolge dessen würden keine Sünden da sein, sondern Friede, Liebe, Sanftmuth, Mildigkeit, Geduld und alle Tugenden auf einen Haufen, was aber in diesem Leben nicht zu hoffen ist. Darum bleiben wir immer Sünder und Uebertreter dieses Gebots, werden aber allein durch dieses Opfer erhalten, daß wir diese Uebertretung wohl erkennen und nicht leugnen, auch nicht mit gottlosen, boshaften Worten uns entschuldigen in unseren Sünden, sondern dieselben bekennen, und nach der Hülfe der Gnade seufzen, und daß das künftige Reich bald komme. Durch diese Demuth verdienen wir, daß uns in den Dingen, die wir zu wenig thun, Vergebung zu Theil werde, wie St. Augustinus sagt: Alle Gebote werden erfüllt, wenn uns das vergeben wird, was wir nicht thun. Es wird aber denen vergeben, die es bekennen, denn den Demüthigen gibt GOtt Gnade [1 Petr. 5, 5.].

Daher muß man sich hüten vor der schädlichen und schmeichlerischen Gloße, welche die Demuth von Grund aus gänzlich austreibt und zunichte macht, da man sagt: GOtt verlangt nicht, daß dieses Gebot in diesem Leben erfüllt werde; denn diese Gloße ruft die schädlichste Sicherheit hervor, sie macht die Hände schlaff und die Kniee der kräftigen Streiter wankend, und ist eins der Kissen und Pfuhle, welche der HErr in dem Propheten Hesekiel [Cap. 13, 18.] verdammt.

Es sei denn, man verstehe dies so, daß GOtt es nicht verlange von denen, welche selbst von sich eine solche Erfüllung fordern und Leid tragen, und bekennen, daß sie es nicht erfüllen, und deshalb nach einem baldigen Tode trachten, und wünschen abzuscheiden [Phil. 1, 23.], damit sie nicht länger in Sünden und im Ungehorsam gegen dieses Gebot liegen, sondern bei Christo seien und in Gerechtigkeit und in völligem Gehorsam gegen dieses Gebot. Diesen, sage ich, rechnet er es nicht zu, verlangt es auch von ihnen nicht. Aber von denjenigen, welche es nicht selbst von sich verlangen und fordern, sondern sich Kissen unter die Arme legen und ohne Furcht einhergehen, sicher, daß es nicht verlangt werde, wird er es gewißlich fordern bis auf den letzten Heller. Von diesen sagt der 10. Psalm [V. 13.]: „Warum soll der Gottlose GOtt lästern? Denn er spricht in seinem Herzen: Du fragst nicht darnach.“ Denn wer nicht erkennt, daß er dies Gebot schuldig sei, wie sollte der erkennen, daß er ein Sünder sei? Wer nicht erkennt, daß er ein Sünder sei, wie sollte der GOtt und sein Gericht fürchten? Wer sich aber nicht fürchtet, wie sollte der demüthig sein? Wer aber nicht demüthig ist, wie sollte der Gnade erlangen? Wer die Gnade nicht erlangt, wie sollte der gerecht werden? Wer nicht gerecht wird, wie sollte der selig werden?\footnote{Am St. Laurentiiustage, den 10. Aug. 1516, schloß Luther die Predigten über das erste Gebot. (Weim. Ausg., Bd. I, S. 74.) — „Hier hat sich vermutlich die Predigt vom 10. August 1516 angeschlossen.“ (Weim. Ausg., Bd. I, S. 430, Anm. 2.) Diese Predigt findet sich nach Walch, St. Louiser Ausgabe, Bd. X, 202 ff. Der erste Paragraph derselben ist noch einmal l. c. Bd. XII, 1728.}

Es folgt das andere Gebot:

\section[Das andere Gebot]{Das andere Gebot.\footnote{Dieses Gebot ist eingeleitet durch das Exordium, welches sich nach Walch, St. Louiser Ausgabe, Bd. XII, 1728 ff. findet, dessen Schluß [Col. 1733, § 6] zeigt, daß Luther am Tage der Himmelfahrt Mariä, den 15. August 1516, zur Auslegung des zweiten Gebots überging.}}

\begin{quote}
Du sollst den Namen des HErrn, deines GOttes nicht unnützlich führen.
\end{quote}

Dieses Gebot (wie auch alle anderen) fließt aus dem ersten Gebote. Denn wenn das erste gehalten ist, so wird auch dieses leicht gehalten, ja, es ist schon gehalten. Denn es ist um deswillen gegeben, daß unserer Blindheit unsere Thorheit offenbar gemacht würde, da dies nicht vonnöthen gewesen wäre, wenn Adam festgestanden hätte; denn dann hätten wir das alles gewußt.

\quellseite{603}
Aber jetzt wissen wir nicht einmal das, was nothwendiger Weise in äußerlichen Dingen geschehen muß, geschweige, daß wir das Innerliche wissen sollten. Deshalb wird im ersten Gebote das Herz und der inwendige Mensch unterwiesen, wie er sich gegen GOtt halten solle, in diesem Gebote aber der Mund. Denn wir sündigen mit drei Dingen wider GOtt, mit dem Herzen, mit dem Munde, mit Werken; darum ist für jedes einzelne ein besonderes Gebot, und diese Gebote sind alle drei so beschaffen, daß sie eine Verneinung (negativa) oder ein Verbot in sich enthalten (prohibitiva). Und es ist unter allen zehn Geboten nur eines, welches gebietet (affirmativum), nämlich: „Du sollst Vater und Mutter ehren“ etc. Denn das dritte: „Du sollst den Feiertag heiligen“, legt der HErr selbst durch ein Verbot (negative) aus und sagt [2 Mos. 20, 10.]: „Da sollst du kein Werk thun.“ Das zeigt auch der Name „Sabbath“ an, das heißt Ruhe, das ist, sich enthalten von Werken. Denn in demselben wird kein Werk geboten, daher sie sich an dem Tage nur mit dem Lesen des Gesetzes beschäftigen. Also verbietet er in allen Geboten Böses, deshalb hat auch kein anderes Gebot eine Verheißung als das eine, nämlich dasjenige, welches die Ehre gegen die Eltern gebietet (affirmativum). Die anderen drohen Strafe, weil sie Böses verbieten.

Darum Gleichwie der, welcher mit dem Herzen nicht sündigt, auch nicht sündigt mit dem Munde und Werk, So kann der, welcher mit dem Herzen sündigt, weder mit dem Munde noch mit einem Werke recht handeln: Wer das erste Gebot nicht hält, der hält auch das zweite und dritte nicht, Wer aber das erste hält, der hält auch das zweite und dritte.

Denn es steht nicht zu befürchten, daß derjenige, welcher auf GOtt seine Zuversicht gesetzt hat, GOtt über alle Dinge liebt, auf ihn hofft, ihn sucht, nach ihm dürstet, seinen Namen unnützlich und unehrerbietig führen sollte, gleichwie nicht zu fürchten ist, daß ein Unterthan, der seinen Fürsten liebt, nicht auch wollen sollte, daß der Name desselben überall durch ihn herrlich gemacht und gehalten werden sollte als ein Heiligthum.\footnote{Wie Löscher anmerkt, hat das Manuscript: haberi ut sanctum, was wir aufgenommen haben statt habere.} Daher haben wir bei dem ersten Gebote diejenigen schon namhaft gemacht, welche wider dasselbe sündigen, nämlich Leute zweierlei Art, ganz grobe und ganz thörichte. Hier muß in allen Dingen auf gleiche Weise geredet werden. Denn in diesem Gebote wird kein anderes Fluchen verboten, als wenn es bei dem Namen GOttes geschieht, aber auch keine Lästerung gegen einen Menschen, sie geschehe denn im Namen GOttes, sondern nachher [, wo es heißt]: „Du sollst nicht tödten.“

Das ist aber vor allen Dingen zu bemerken, daß dies Gebot in gebietender Weise (affirmative) auch das zu verstehen gibt, daß man den Namen GOttes in rechter Weise im Munde führen soll, oder wo es zur Seligkeit vonnöthen ist, und daß hier allein verboten wird, man solle ihn nicht unnützlich führen. Daher wird in der Schrift häufig geboten, daß man den Namen GOttes anrufen soll, loben, danken und preisen. So heißt es Röm. 10, 13.: „Wer den Namen des HErrn wird anrufen, soll selig werden.“

Daher wollen wir zuerst die groben [Uebertreter] ansehen, unter denen uns zuerst die Abergläubischen entgegentreten, die Wahrsager und diejenigen, welche aus jüdischen Fabeln abergläubische Dinge von dem tetragrammaton erdichten.\footnote{Anstatt des vorhergehenden Satzes steht, wie Löscher berichtet, Folgendes in der Handschrift: Erstlich schänden also den Namen GOttes die Aberglaubischen, die Wahrsager und andere, von denen oben gesagt ist. Zweitens etliche abergläubische Leute aus den Juden, welche unnützes Geschwätz machen von dem Namen tetragrammaton; doch dies gehört zum ersten Gebote. Drittens diejenigen, welche des christlichen Namens unwürdig leben, welche recht eigentlich keinen Namen schänden. Viertens Meineidige, Lügner, Lästerer (dolosi), Betrüger und alle, welche die Wahrheit verlassen, wo sie GOtt zum Zeugen anrufen und seinen Namen, [Anmerkung: Statt nomini divino ist zu lesen: nomen ejus immo. Diese und gleich folgende Correctur ergibt sich aus der Parallelstelle Col. 1196.] ja auch bei GOtt fluchen; desgleichen die Bischöfe, welche [mit Unrecht] [Anmerkung: injuste fehlt bei Löscher und in der Weimarschen.] bannen und absolvieren. Fünftens die [mit dem Namen GOttes] Scherz treiben. Sechstens die wider GOtt lästern. Dies sind die größten. Siebentens die Hetzer und Heuchler und die Leute, die einen hoffärtigen Sinn haben, Zänker, die wegen der Wahrheit streiten; ja, die ihre Gerechtigkeit anrufen unter dem Namen GOttes. Diese wollen wir jetzt einzeln vornehmen. Erstens: Daß diese etc.} Daß diese den Namen des HErrn unnützlich führen, ist offenbar, denn sie führen den Namen GOttes weder zur Seelen Seligkeit, noch zur Ehre GOttes, sondern zu ihrem Vorwitz, zum Bunde mit den Teufeln, in Zeichen, Worten und Geberden, wie oben [beim ersten Gebote, Col. 1139] gesagt ist. Denn es liegt
\quellseite{604}
ihnen nicht daran, daß sie dadurch die Seligkeit der Seele erlangen, viel weniger, ob GOtt dadurch geehrt werde, sondern nur, daß sie ihre Lust büßen. Wie sie daher GOtt verworfen haben aus ihrem Herzen, und ihn in sich unnütz gemacht, wider das erste Gebot, so schänden sie auch jetzt seinen Namen mit ihrem Munde, und führen ihn unnützlich.

„Vergebens“ und „unnütz“ aber wird das genannt, was ohne Noth und rechte Ursache geschieht. Es ist aber eine zwiefache Noth oder Nutzen, dazu man GOttes Namen führen soll, nämlich unsere Seligkeit und die Ehre GOttes; ja, allein GOttes Ehre, weil wir auch nicht unsere Seligkeit durch seinen Namen suchen sollen, es gereiche denn zu seiner Ehre. „Führen“ (assumere = [in den Mund] nehmen) aber versteht man so, z. B. wenn jemand etwas vor sich nimmt oder beabsichtiget zu thun oder zu sagen, und dazu den Namen GOttes gebraucht; gleichwie er im Herzen die Gnade GOttes durch den Glauben annimmt (assumit) zu seiner Seligkeit, so nimmt er (assumit) in seinen Mund mit Ehrerbietung den Namen GOttes zu seiner und anderer Leute Seligkeit. Davon ist nun beim ersten Gebote genug gesagt worden, daß man sich [des unnützen Gebrauchs des göttlichen Namens] entschlage. Sie können sich auch damit nicht entschuldigen, daß sie heilige Worte und Gebete anwenden, da sie gerade dadurch schändlicher sündigen als in anderen derartigen Dingen.

Würdest du den nicht für einen Rasenden halten, der in den heiligen Gewändern, die zum Gottesdienst bestimmt sind, zu Tanz, ins Theater, oder in den Krieg ginge, und einen solchen Scherz aus einem so ernsten Dinge machte? Wenn du nun diesen mit allem Ernste strafen würdest, warum strafst du denn dich nicht viel mehr als jenen, der du das, was heiliger ist als alle Kleider, ja, das, wodurch die Kleider und alles geheiligt wird, nicht zum Tanz, sondern zur Hurerei, zu den Werken des Teufels mißbrauchst, oder wenigstens zu leiblicher Hülfe, und der Seele nicht abstehst? Derartig ist die Entschuldigung jener, die dessen, der in solcher Weise die heiligen Kleider mißbraucht, und wenn er gestraft wird, sagt: Ei, es ist doch eine heilige und gute Sache. Denn was soll man einem solchen anders sagen als dies: [Du redest,] als ob ich nicht wüßte, daß es gut und heilig sei. Dies lehrst du mich jetzt, während ich dich um deswillen strafe, daß du des Heiligen in solcher Weise mißbrauchst. Aber wir sind jetzt solche Christen, die da Holzschuhe kaufen, damit wir unsere neuen Schuhe nicht besudeln, und schleppen gute Kleider gewiß nicht in den Koth, auch thun wir unseren goldenen Gefäßen nicht die Schmach an, daß wir Harn hineinthäten, ja selbst nicht unreines Wasser, aber den heiligen Namen GOttes ziehen wir zu irgend welchen unwürdigen Dingen, ohne alle Scham. Da dies nun niemand mit dem, was sein ist, so macht (stelle dir nur vor, daß jemand so rasend wäre, der solches an dem Seinigen thäte), wie wird es denn stehen mit dem, was anderen gehört, ja, was dessen ist, der da will, daß das Seine rein gehalten werde? Wie beliebt würde sich wohl der machen, der eines Fürsten Titel, Namen und Wappen in den Koth zöge und darin umkehrte, nämlich wenn er es selbst sähe und verböte, ja vielmehr beföhle, es hoch aufzurichten, wie es sich gebührt?\footnote{Wir haben die Lesart der Weimarischen Ausgabe angenommen: in sublime loco statui statt sublimi in den andern Ausgaben.} Aber wenn wir hören, daß der Türke die Kirchen entweiht, die Altäre und alle Heiligthümer besudelt, so kommen wir ganz außer uns, entbrennen vor Zorn, und denken daran, diese Schmach durch Krieg zu rächen, und beklagen uns, daß die Fürsten nicht Krieg führen wider den Türken. Aber greife nur in deinen Busen, dann hast du einen Türken ergriffen. Unterdessen gibt GOtt zu desto größerer Strafe die Fürsten dahin in Kriege unter einander, damit er uns mehr strafe als die Türken, weil wir die Heiligthümer schlimmer besudeln als sie.

Zweitens treten uns entgegen Meineidige, Lügner, Ränkemacher, Betrüger und alle, welche die Wahrheit verlassen, wo sie GOtt zum Zeugen anrufen und seinen Namen, ja, auch bei GOtt fluchen, diejenigen, welche ungerecht bannen, auch die [mit dem Namen GOttes] Scherz treiben und wider GOtt lästern, von denen oben schon etwas gesagt ist und nachher noch mehr gesagt werden wird. Hier aber wollen wir sehen, welche es sind, die mit Schwören wider den Namen des HErrn sündigen.

Es ist daher zu merken,\footnote{Statt des Abschnitts vom Anfang des vorigen Absatzes bis hieher bringt das Manuscript hier nur: „Zweitens die durch Schwören wider den Namen GOttes sündigen. Hier ist zu merken“ etc. Diese Auslassung hier hat ihren Grund darin, daß dasselbe schon vorher im Manuscript gesagt worden war. Vgl. Col. 1193, Anm. 2.} daß es ein zwiefaches
\quellseite{605}
Schwören gibt [ein gutes und ein böses]. Von dem guten heißt es 5 Mos. 6, 13. und 10, 20.: „Du sollst den HErrn, deinen GOtt, fürchten und ihm allein dienen, und bei seinem Namen schwören.“ Denn so schwört GOtt selbst öfters in den Propheten, indem er spricht [Ezech. 33, 11.]: „So wahr ich lebe, spricht der HErr“, und Ps. 110, 4.: „Der HErr hat geschworen, und wird ihn nicht gereuen“, und Ps. 132, 11.: „Der HErr hat David einen wahren Eid geschworen.“ So schwört Christus im Evangelio [Matth. 26, 63. f.]; die Patriarchen, die Propheten, die Apostel, und alle Heiligen schwören und haben geschworen. Und dies ist ein verdienstliches Werk, weil es in den Dingen geschieht, welche die Seligkeit antreffen. So oft daher jemand etwas redet oder thut, was zur Seligkeit dient, so soll er schwören, damit man es glaube, wenn zu fürchten steht, man glaube es nicht. So schwört der Apostel im Briefe an die Römer [Cap. 1, 9. f.], daß er sich oft vorgenommen habe, zu ihnen zu kommen. Welchen Grund hatte er denn anders für sein Schwören, als daß es ihnen nützte zu ihrer Seligkeit, daß sie das glaubten? damit sie erkenneten, er sei ein treuer Hirte und nicht ein gewinnsüchtiger Mensch, und die Zuversicht hätten, daß die rechte Liebe in ihm sei. So heißt es im 63. Psalm, V. 12.: „Wer bei ihm schwört, wird gerühmt werden.“ Die Ursache, warum Gotte dies Schwören gefällt, ist dies, daß durch dasselbe seine Wahrheit angerufen wird, und an ihr geglaubt wird, und um seinetwillen Friede und Einigkeit hergestellt wird unter denen, die da schwören. Deshalb wird er in diesem Werke heiliglich verehrt, weil das Werk des Teufels zerstört wird, nämlich Streit und Hader. Denn wer dir nicht glaubt und sich nicht zufrieden gibt, der glaubt dir um des Namens GOttes willen, den du anrufst, und gibt sich zufrieden; darum sind wir diese Ehre allein Gotte schuldig. Daher sagt der Apostel Hebr. 6, 16.: „Der Eid macht ein Ende alles Haderes.“

Das andere Schwören ist böse, und dies ist ein zwiefaches. Das eine ist das, welches aus Gewohnheit geschieht, weil so durch das beständige Gebrauchen des Namens Gottes die Ehrbietung gegen denselben verringert wird. Wie wenn man im Deutschen sagt: „Bei GOtt, wahrlich, fürwahr, so (als) wahr als GOtt im Himmel ist, bei GOtt und allen Heiligen“, bei dem lebendigen GOtt, bei dem heiligen GOtt, „auf mein theures Kreuz, auf mein theures Blut, bei dem Kreuz GOttes, es ist Marter groß, Marter schön“\footnote{Im Original: „Martyr schön“.}, das heißt, bei dem Leiden Christi ist es so. Desgleichen, wenn man so flucht: „Der Teufel hol mich, brech mir den Hals, GOtt helfe mir nimmermehr, als mir GOtt helfe, auf meine Seele“, auf mein Gewissen, „bei meiner Treu und Ehr, bei meiner Priesterschaft“. Dies und viel Aehnliches ist, leider, vielen so häufig im Munde, daß fast das zweite Wort ein derartiger Schwur ist. Hier sollten Hausväter und Mütter wachen über Kinder und Knechte und Gesinde, denn es ist eine überaus böse Gewohnheit und wider dies Gebot.

Aber um der Einfältigen willen möchte man hier fragen: Hat denn das etwas mit dem Namen GOttes zu thun, wenn man so schwört: Wahrlich, ich sage die Wahrheit? Ich antworte: Ja, denn der Name GOttes ist die Wahrheit, weil Er die Wahrheit ist. Denn so schwört der Apostel Röm. 9, 1. bei der Wahrheit und seinem Gewissen. Und wiewohl es scheinen könnte, als ob dies deshalb verboten sei, weil das Gewissen nicht der Name GOttes ist, da es ebensowohl Sünde ist, wenn man unnöthig schwört, als wenn man nicht bei dem Namen GOttes schwört oder bei etwas anderem als bei GOtt, wenn man auch die Wahrheit sagt, so bezieht dennoch Christus Matth. 5, 34–36. dies alles auf GOtt, da er sagt: „Du sollst nicht schwören weder bei dem Himmel, denn er ist GOttes Stuhl, noch bei der Erde, denn sie ist seiner Füße Schemel, noch bei Jerusalem, denn sie ist eines großen Königes Stadt. Auch sollst du nicht bei deinem Haupte schwören, denn du vermagst nicht ein einiges Haar weiß oder schwarz zu machen.“ Hieraus ist offenbar, daß derjenige, welcher bei dem schwört, was GOttes ist, oder was ihm zugeeignet wird, auch bei GOtt schwört, von dem, in dem und zu dem alle Dinge sind [Röm. 11, 36.].

Das andere ist das falsche Schwören, da jemand vorsätzlich und wissentlich falsch schwört mit den ebengenannten Eidesformeln; dies ist eine sehr schwere Sünde, und selten bleibt diese Sünde auch in diesem Leben ungestraft. Wenn du das bezweifeln solltest, so frage nur die Meineidigen darnach. Aber diese Sünde ist über die
\quellseite{606}
Maßen schwer, wenn dies geschieht Angesichts der Kirche in feierlicher Weise mit vorgeschriebenen Ceremonien (forma), wie die thun, welche sich heimlich mit einander verloben, und zur Strafe ihrer Sünde darnach einander hassen, und sich scheiden durch einen Meineid: eine erschreckliche Schuld, aber auch eine erschreckliche Strafe. Daher ist unter allen Geboten allein diesem eine Drohung der Strafe angefügt, indem es heißt [2 Mos. 20, 7.]: „Denn der HErr wird den nicht ungestraft lassen, der den Namen des HErrn, seines GOttes, mißbraucht.“ Doch hat auch das erste Gebot eine Drohung, daß der Herr die Mißethat derer, die ihn hassen, heimsuchen will bis ins vierte Glied [2 Mos. 20, 5.]. Aber weil der Apostel [Eph. 6, 2.] die Verheißung, welche in diesem [ersten] Gebote enthalten ist, nicht aufzählt, sondern nur die im vierten, von der er sagt, daß sie die erste sei, so scheint es, daß auch diese Drohung des ersten Gebotes nicht gezählt werden solle.

Und wahrlich, der Meineid ist eine gewaltige Sünde, das erhellt erstlich aus dem Zeugniß, welches die Schrift darüber ablegt, denn der HErr gebietet im Gesetze [3 Mos. 19, 12. Matth. 5, 33.]: „Du sollst keinen falschen Eid thun“, und die Propheten strafen den Meineid gar sehr. Daher sagt Hoseas, Cap. 4, 1. 2. [Vulg.]: „Es ist keine Wahrheit, es ist keine Barmherzigkeit, es ist keine Erkenntniß GOttes auf Erden. Fluchen und Lügen und Mord und Diebstahl und Ehebruch hat überhand genommen und eine Blutschuld kommt an die andere.“ Ja, so unverletzlich (sanctum) hat GOtt den Eid gehalten wissen wollen, daß er es auch rächte, daß Zedekias dem Könige von Babylon seinen Eid nicht gehalten hatte [2 Kön. 24, 17. bis 25, 7.]. Desgleichen durften sie auch den Gibeonitern [Jos. 9, 19.] nicht Schaden thun, wiewohl sie von ihnen betrogen worden waren, da sie ihnen den Eid leisteten. Auch wagten sie nicht den Kindern Benjamin ihre Töchter zu geben um ihres Eides willen [Rich. 21, 1. ff.], und ließen lieber geschehen, daß heidnisches Geblüt mit dem Blute ihrer Brüder vermischt würde, als daß sie ihren Eid gebrochen hätten.

Zweitens [erhellt das, wie große Sünde der Meineid sei,] aus der Sache selbst. Denn schwören ist GOtt zum Zeugen anrufen und ihn als ein Mittel gebrauchen, durch welches sich Zwieträchtige vereinigen. Aber so setzt der, welcher falsch schwört, GOtt dem Gespötte aus; ja, was erschrecklich ist, wenn er sagt: So wahr GOtt ist, lebt, wahrhaftig ist, sagt er im Herzen, daß dies nicht wahr sei. Und dadurch verleugnet er GOtt mit dem Munde, ja, auch mit dem Herzen, mit dem er der Verleugnung des Mundes beistimmt. Und dies ist die ärgste Verleugnung, weil er weiß, daß GOtt ist, und dies mit dem Herzen glaubt, und dennoch stimmt er dem zugleich auch mit dem Herzen bei, daß er ihn verleugne; denn er würde ihn äußerlich nicht verleugnen, wenn er ihn nicht verleugnen wollte, und doch verleugnet er ihn wider sein Herz und Gewissen. Darum ist dies die allerschwerste Lüge und geradezu wider GOtt, weil er wollte, daß kein GOtt wäre, damit seine Lüge bestehe.

Drittens erhellt dies aus einem Gleichniß. Wenn du einen Fürsten dazu bewegtest, daß er dir, damit man dir glaube in deiner Sache, Brief und Siegel gebe, daß man dir so um seines Namens willen Glauben schenkte, und du nachher das schändlicher Weise in Abrede nähmest und den Fürsten in Schande stecken ließest: was meinst du wohl, daß er von dir denken würde? Meinst du, daß du an ihm einen Bürgen haben würdest? Oder daß er dir zu Dienst stehen würde, wenn du ihn bätest, er möchte bei deinem Vertrage, den du mit einem andern machtest, zugegen sein, und du, wenn er käme, vor seinem Angesichte leugnen würdest, daß er ein Fürst sei, und ihn verwürfest? Aber eine wie große Majestät ist doch GOtt! Er ist allein der HErr.

Um dies klarer zu verstehen, ist zu merken, daß wider dies Gebot in zwiefacher Weise gehandelt wird bei dem Eide. Erstens durch Schwören, zweitens durch Brechen des Schwures. Durch Schwören, zum Beispiel: wenn jemand wissentlich in einer Sache einen Meineid thut, indem er betrüglich sagt, er wolle etwas thun, was er doch nicht thun will, um frei zu werden, oder wenn jemand fälschlich leugnet was er weiß, wie es geschieht im Handel, bei anvertrautem Gut, in Gerichtssachen. Diese sündigen schwerer als die Abergläubischen, denn dieselben führen den Namen GOttes unnützlich, diese aber zur Lüge, was viel schwerer ist; und dies geschieht hauptsächlich durch Ableugnungen (negationibus), wie es im 24. Psalm, V. 4. heißt: „Der seinem Nächsten nicht fälschlich geschworen hat.“ Durch Brechen des Eides, zum Beispiel: wenn jemand
\quellseite{607}
zwar ohne Betrug geschworen hat, jedoch nachher anderen Sinnes wird, und den Eid nicht hält, was bei Zusagen (affirmationibus) geschieht, wie der 15. Psalm, V. 4., sagt: „Wer seinem Nächsten schwört, und ihn nicht täuscht.“ Aber wenn er es nicht halten kann, hat er darum noch keinen Meineid geschworen. Aber du sagst: Wie steht es denn? Wenn ich nicht schwöre, so muß ich Schaden oder Gefahr oder den Tod leiden; da sieht man doch, daß ich [von der Sünde des Meineids] loszusprechen sei; das ist als wenn ein Räuber oder ein Feind dich zwingt zu bekennen, was du hast. Darauf ist das die Antwort: Man darf keine Sünde thun um irgend einer Sache willen, auch nicht um des Ewigen willen, viel weniger um des Zeitlichen willen. Aber du sagst: Wer wird es mir wiedergeben? Antwort: GOtt, für den du die Wahrheit gesagt hast. Aber ich kann nicht alle Fälle dieses entsetzlichen Greuels aufzählen, weil er heutzutage in Wahrheit wie eine Fluth hereingebrochen ist, und für nichts geachtet wird. O, wenn wir doch bedächten, wie so wenige Menschen selig werden, und uns fürchteten, wie bald würde diese Sicherheit unserer Seligkeit halben aufhören!

Daher sagt der HErr Matth. 5, 33. f.: „Ihr habt gehört, daß zu den Alten gesagt ist, du sollst keinen falschen Eid thun, und sollst GOtt deinen Eid halten. Ich aber sage euch, daß ihr allerdings nicht schwören sollt.“ An diesem Spruche haben sich viele damit abgemüht, ob man schwören dürfe. Einige sagen, es sei vom Uebel der Strafe halben, wenn man schwört, andere, es sei vom Uebel des Schwörenden halben. Aber wenn man auf den Sinn Christi sieht, so will er schlechterdings, daß man durchaus nicht schwören soll, weil er sagt: „Ihr sollt allerdings nicht schwören.“ Es ist also der Sinn: Den Juden ist geboten, daß sie keinen falschen Eid thun sollten, das Schwören ist ihnen aber gestattet nach ihrem Belieben; auch aber befehle ich, daß ihr auf keine Weise schwören sollt, weder bei dem Himmel \&c. Damit will er schlechterdings dies, daß keiner jemals aus seinem eigenen Willen schwören soll. Nämlich die Lust und den eigenen Willen zu schwören hat er verboten, und so viel an ihm ist, darf niemand schwören. Deshalb, wenn er mehr sagt als: „Ja, ja, nein, nein“, so sündigt er, und es ist böse. Damit verbietet er jedoch nicht, daß ein jeglicher Gläubige nicht schwören solle, wenn er dazu gezwungen wird, sei es nun, sage ich, daß er gezwungen werde durch die Gewalt eines anderen oder durch die Noth des Bruders. Ja, es kann sowohl Freund als Feind mich dazu bringen, daß ich schwöre,\footnote{Die Weimarische hat hier nach der editio princeps: ut orem, doch haben wir die Lesart, welche sich, in allen andern Ausgaben findet: ut jurem, vorgezogen.} und alles thue, was er will, wenn es nur etwas Erlaubtes ist, wie er im Folgenden sagt [Matth. 5, 40. f.], daß man dem, der den Mantel nimmt, auch den Rock lassen soll, und sogar noch zwei Meilen mit ihm gehen. Denn dann schwört er nicht selbst, weil es nicht in seinem Willen steht, sondern er in Demuth einem fremden Willen dient, oder durch die Liebe sich einer fremden Noth annimmt. Wo nun dies nicht vorhanden ist, darf man auf keine Weise schwören, ohne daß es Sünde und vom Uebel ist, weil man dann aus bloßer Willkür schwört, ohne Noth. Denn Christus unterweist an dieser Stelle den inwendigen Menschen, der nicht allein den Meineid, sondern auch den Eid meiden soll. Den Juden wird verboten, daß sie keinen falschen Eid thun sollen, den Christen aber auch, daß sie die Wahrheit nicht beschwören sollen, weil die Ehre des göttlichen Namens unter dem neuen Gesetze vollkommener und unverletzter sein soll, als unter dem alten Gesetze. Zweitens, weil ein Christ die zeitlichen Dinge nicht lieben soll, so soll er um derselben willen auch nicht schwören. Denn da er nicht suchen darf, was sein ist, wie sollte ihm denn gestattet sein, um derselben willen zu schwören? Darum muß er den Namen GOttes höher achten, als wie man sich damals bei den Juden dazu stellte (dissimulatum est). Wo man aber einem anderen einen Dienst leisten kann, da ist ohne Bedenken, doch in der Furcht GOttes, ein Bruder gehalten, das zu thun, was des anderen Noth erfordert. Und die Ursache dieses Verbotes ist, daß er weiß, daß wir immer den Namen GOttes nicht anders als unnützlich führen [, und er wird nicht recht gebraucht], außer wo er in Liebe und Demuth für andere oder auch zur eigenen Seligkeit gebraucht wird, was die Grundlage für dieses Verbot ist, 1 Cor. 12, 3.: „Niemand kann JESUM einen HErrn heißen, ohne durch den Heiligen Geist“ \&c. Oder um es noch besser zu sagen: „Ihr sollt allerdings nicht schwören“, wird wider die gesagt, welche eine Ausflucht
\quellseite{608}
suchen, indem sie schwören, aber nicht bei dem Namen GOttes, sondern bei dem Himmel, bei der Erde, bei Jerusalem, bei ihrem Haupte, bei ihrer Seele, als ob sie um deswillen keine Schuld auf sich laden, wenn sie falsch schwören, weil sie nicht bei dem Namen GOttes schwören; und in solcher Weise setzen sie ein derartiges Schwören so an, als ob sie nicht geschworen hätten. Daß dies die Meinung Christi sei, sieht man daraus, was er gleich darauf erklärt, wie er nicht wolle, daß man überhaupt schwören solle, nämlich, weder bei Jerusalem etc.

\section{Das dritte Gebot}

\subsection{Gedenke, daß du den Feiertag heiligest}

„Das ist mein Gebot, daß ihr euch untereinander liebet“, Joh. 15, 12. „Die Liebe ist des Gesetzes Erfüllung“, sagt der Apostel [Röm. 13, 10.]. Wenn man die hat, so ist kein Gesetz vonnöthen, ohne dieselbe ist kein Gesetz genugsam. Warum wird in diesem dritten Gebote nicht ein Werk geboten, vielmehr Ruhe, auf daß GOtt nicht durch Werke beleidigt werde. Denn es wird kein Werk in demselben geboten. Deshalb bereiten diese drei Gebote den Menschen Gotte zu, gleichsam als einen reinen Stoff (materiam), daß er ruhen soll mit dem Herzen, mit dem Munde, mit Werken, das heißt, nach dem inwendigen, und auswendigen, und mittleren Menschen, welche da sind der sinnliche (sensualis), der vernünftige, und der geistliche Mensch; und es soll eine lautere Ruhe sein. Hier sind vor allem zwei Dinge zu merken.

Erstens daß der Sabbath den Juden zu einem Vorbilde geboten war, wie der Apostel ausdrücklich anzeigt Col. 2, 17.: „Welches ist der Schatten von dem, das zukünftig ist, aber der Körper selbst ist in Christo.“ Daher bedeutet Dienstbarkeit (opera servilia) oder Arbeit der Hände (manualia), wiewohl sie durchaus gut ist, sündliche und unerlaubte Werke des alten Menschen. Denn gleichwie etliche Thiere nur nach ihrer Bedeutung (significative) unrein waren, die an sich gut, rein und gesund waren, weil GOtt sie geschaffen hatte, und doch Unreines und Böses bedeuteten: so sind auch die leiblichen Werke, welche an sich gut sind, dennoch [am Sabbath] verboten, um künftige Dinge, das ist, die Offenbarung von Uebelthaten zu bedeuten. So bedeutet der Sabbath die geistliche Zeit, welche Christus, die Sonne der Gerechtigkeit, erleuchtet hat, welche keine Nacht hat. Daher sagt Jesaias Cap. 66, 23.: „Es wird ein Monat am andern sein, und ein Sabbath am andern.“ Und der Apostel straft die Galater [Cap. 4, 10.], daß sie Tage und Monate und Zeiten und Jahre hielten nach jüdischer Weise, das ist, nur äußerlich. Daher hat dies Gebot eigentlich aufgehört, ja, alle Gebote, sofern es die vollkommenen Christen betrifft, denn [1 Tim. 1, 9.] „dem Gerechten ist kein Gesetz gegeben“.

Aber die Kirche hat um der unvollkommenen willen Feste beibehalten, weil es nothwendig ist, GOttes Wort zu treiben, denn der wahrhaft Gerechte ist Gotte so gleichförmig, so daß, wie bei GOtt kein Unterschied ist der Tage, der Stätte, der Person, so auch bei ihm jeder Tag ein Feiertag ist. Aber die Schwachen, die noch nicht getödtet sind nach dem alten Menschen, haben diese Dinge vonnöthen, daß sie zu schaffen haben (occupentur) mit gewissen Gottesdiensten, Tagen und Weisen, mit Wachen, Fasten, Arbeiten, Gebeten, Uebungen (disciplinis) und dergleichen, wodurch sie zum Wachstum des inwendigen Menschen gelangen, damit, wo der Leib fastet und geknechtet worden ist, und die Leidenschaften getödtet, sie dann auch allmälig aufhören, und um so viel gemindert werden, als der inwendige Mensch zugenommen hat, so daß, wenn dieser vollkommen geworden ist, jene Dinge ganz und gar aufhören müssen. Dies ist es, was der Apostel sagt [Gal. 3, 24.]: „Das Gesetz ist unser Zuchtmeister gewesen auf Christum.“ Denn das Gesetz hat nichts zur Vollkommenheit gebracht; Johannes aber ist vor ihm hergegangen, um dem HErrn ein bereit Volk zuzurichten [Luc. 1, 17.].

Wer sich darum dünken läßt, daß er durch das Thun jener Werke etwas Großes oder gar alles ausgerichtet habe, und darauf sich gründet, darauf vertraut, sicher ist, und dadurch nichts Anderes sucht, der ist [Hos. 10, 11.], „das Kalb Ephraim, gewöhnt, daß es gerne drischt“, der
\quellseite{609}
da meint, daß sei ein guter Mensch, nämlich, der da gute Werke thue, während es doch nur dazu dient, daß der Mensch dazu bereitet werde, daß er gut sei. Denn es ist nur eine Taufe mit Wasser [Matth. 3, 11.]. Von solchen Leuten wird in der Schrift vieles, ja, fast alles gesagt. Denn die durch Werke ihre Gerechtigkeit suchen (justiciarii), sind Heuchler, welche auch die Schrift dahin verdrehen, und, wie der Apostel sagt [1 Tim. 1, 7.], „nicht verstehen, was sie sagen oder was sie setzen“. Gleichwie der Lehrling in einem jeglichen Handwerk das thut, was der Meister nur dann thut, wenn es ihm beliebt, aber dadurch zur Meisterschaft fortschreitet (wenn er aber nun dabei stehen bleiben wollte, als ob er alles wüßte, so würde er allen als ein Thor und lächerlich erscheinen):\footnote{Diese Klammern sind von uns gesetzt.} so werden wir hier durch die guten Werke nicht gut, sondern vorbereitet. Denn nicht durch Thun, sondern durch Leiden sind wir gut, wenn wir GOttes Werk (actiones) an uns leiden, indem wir stille sind. Darum hält die Kirche dies Gebot und viele anderen bildlichen Dinge (figuralia), um die Schwachen zu üben, damit sie zunehmen, nicht, damit sie auf demselben Flecke stehen bleiben, und sicher seien, wenn sie diese Dinge gethan haben.

Zweitens ist zu merken, daß St. Gregorius hier fragt: warum der HErr nur dies Eine Gebot [der Liebe] anziehe, da doch die Schrift voll ist von Geboten? Aber (wie ich gesagt habe) es fordern alle Gebote die Liebe, da ohne die Liebe, das heißt, ohne einen leichten, bereiten, fröhlichen, geneigten Willen die Gebote nicht erfüllt werden, wenn sie auch erfüllt werden. Denn es bleibt ein unwilliger und darum auch ein schuldbeladener (rea) Wille, wiewohl er das Werk mit der Hand und äußerlich thut. Denn es heißt [Jac. 2, 10.]: „Wer an einem sündigt, der ist aller Gebote schuldig geworden.“ So spricht der erste Psalm [V. 2.], daß der Mann glücklich ist, der Lust hat am Gesetze des HErrn (das heißt, der die Liebe hat und den Nächsten liebt). Und Ps. 110, 3. heißt es: „Dein williges Volk“, und der Apostel sagt [2 Cor. 9, 7.]: „Einen fröhlichen Geber hat GOtt lieb.“ Denn GOtt hat kein Gefallen an einem gezwungenen Dienste, sondern am freien und willigem, wie der 54. Psalm [V. 8. Augl.] sagt: „Williglich will ich dir opfern.“ Dieser Wille ist die Liebe, wie St. Augustinus sagt.

Warum hat der HErr denn nicht allein dies Eine Gebot gegeben, und sich daran genügen lassen? Ich antworte: Es wäre zwar ausreichend gewesen, aber es mußte dem Menschen gezeigt werden, in welchen Dingen er Liebe haben müßte, damit er sich nicht selbst verführe durch falschen Wahn und Vertrauen, und meinte, er habe die Liebe, während er sie doch nicht hat. Denn es geschieht gar leicht, daß wir von uns eine sehr gute Meinung haben und eine sehr schlechte von den Nächsten, während es doch umgekehrt bei uns sein sollte. Darum sind die heiligen zehn Gebote gleichsam ein Spiegel, in welchem der Mensch sich selbst betrachten soll, worin und wie weit es ihm mangele oder vorgeschritten sei. Denn bei allen Geboten soll man nicht sowohl Acht darauf haben, daß man sie thue, sondern vor allen Dingen, aus was für einem Herzen und Willen man sie thue. Ich sage, sie sind gegeben, damit wir unseren Willen daraus erkennen, sodann aber auch, damit dieser geübt werde; dies alles ist denen eine Last, welche unwillig und unlustig sind. Denn siehe, damit wir eben dies klar machen, das erste Gebot enthält alle anderen in sich. Denn wer dieses hält, der hält alle, und wer irgend ein anderes nicht hält, der hält dieses nicht, weil sein Herz irgend etwas Anderes im Auge hat, aber nicht allein GOtt.

Aber unser alter Adam ist in so verkehrter Weise hochmüthig, daß er öfters für weiser und heiliger angesehen sein will als alle Heiligen. Deshalb untersteht er sich, mit frecher Stirne zu sagen, er halte das erste Gebot, er sei sich seiner Uebertretung bewußt, er habe keinen Abgott, er liebe nichts so sehr als Gott; darum müßte er auf Werke und Worte verwiesen werden, welche ihm Zeugniß geben möchten über ihn, und ihn seiner hochmüthigen Frevelhaftigkeit oder der Lüge überführen. Zum Beispiel, wenn er sich dessen vermißt, daß er keinen andern Gott habe, so soll er seine Zunge fragen, ob sie schwöre, ob sie den Namen Gottes unnütz führe. Denn [Matth. 12, 34.]: „Wes das Herz voll ist, deß geht der Mund über.“ Leicht ist ein Irrthum in der Meinung des Herzens, und betrügerische Vermessenheit. Und man kann nicht leicht wissen, ob man das erste Gebot halte, wenn man die Werke und Worte nicht darum befragt. Darum,
\quellseite{610}
Wenn jemand schwört, oder den Namen GOttes unnützlich führt, der soll wissen, daß er lüge oder sich betrüge, wenn er sagt, er liebe GOtt und halte das erste Gebot. Er hat sein Herz nicht recht erforscht und geprüft.

Aber siehe, es möchte jemand meinen, er halte auch das zweite Gebot, und ist sich seiner Uebertretung desselben bewußt: der befrage sich über das dritte Gebot, ob er bereitwillig den Feiertag heilige, ob er begierig das Wort GOttes höre, bete, bedenke, lese, rede von GOtt, beim Gottesdienst zugegen sei. Wenn er dies nicht thut, oder mit weniger Lust thut, als seine anderen Geschäfte, so lügt er, wenn er sagt, er halte das erste und das zweite Gebot; denn seine Werke geben Zeugniß wider ihn. Oder wäre es möglich, daß jemand GOttes Namen nicht unnützlich führen sollte, der selbst unnütz ist? Wie sollte ein unnützer Mensch etwas Anderes als Unnützes thun? Er redet unnütz, er denkt unnütz, auch wenn er über GOtt denkt und redet. Ein unnützer Mensch ist aber der, welcher mit anderen Dingen umgeht, als die GOttes sind, und der das, was Gottes ist, unnöthig thut, im Vergleich zu dem, was sein ist. Und jeder Mensch ist ganz unnütz (vanitas).

Drittens, so er sich dünken lassen sollte, daß er auch dieses halte, so prüfe er sich an einem größeren Zeichen: Er sehe zu, wie er die Haushalter (vicarios) GOttes aufnehme, die er sieht, nämlich die Eltern, die Vorgesetzten; wenn er diese gering achtet und verachtet, die er sieht, wie kann er sich unterstehen zu sagen, er achte GOtt hoch und nichts höher als ihn, da er ihn doch nicht sieht?

Viertens, in den Dingen, die den Nächsten anbetreffen, da ist es schlechterdings der Wille GOttes, daß du denselben liebest. Dieser Wille ist im ersten Gebote immer mit eingeschlossen (intelligitur). Wenn du nun sagst, du liebest deinen Nächsten, und hassest ihn, heißest ihn, setzest ihn herab, verletzest ihn, verachtest ihn, tödtest ihn, fluchst ihm, schmäbest ihn: siehe, so gibt dein Werk Zeugniß wider dich. Desgleichen, wenn du ihn nicht verteidigst, so man ihm Unrecht thut, ihn nicht entschuldigst, nicht Gutes von ihm redest; denn das thut diese Liebe, wenn sie eine wahre Liebe ist. Siehe also wohl in diesen Spiegel.

Fünftens. Wenn du auch daraus dich nicht erkennst, wie es denn dadurch nicht genugsam erkannt werden kann, weil die böse Lust gar mancherlei Art ist, und niemand auch nur seine eigene genugsam erkennt, so mußt du diese Prüfung durch vieles, ja, durch alles hindurchführen. Also: Hast du deinem Nächsten keinen Schaden gethan an seinem Leibe und an seiner Person, so hast du es vielleicht gethan am Leibe seines Weibes, welche Ein Fleisch mit ihm ist. Hast du nicht etwa ihrer begehrt, sie unzüchtig angesehen, betastet, angeredet, oder einem anderen, der solches that, nicht gewehret? ihm nicht geholfen und ihn nicht wieder zurechtgebracht? Wie kannst du doch, ich bitte dich, den lieben wie dich selbst? Denn du wolltest doch gewiß, daß ein anderer so gegen dich handelte, denn in solcher Weise liebst du dich, daran ist kein Zweifel, und das gönnst du dir.

Sechstens. Gesetzt den Fall, daß du dir auch hier nichts bewußt wärest, und dein Adam meinen sollte, er sei ganz unverletzt, wie er denn nach seiner Blindheit zu thun pflegt, während er doch in niemandem rein ist, so denke, was du etwa an seinem äußerlichen Gute gesündigt habest, nämlich ob du gestohlen habest, das heißt, ob du durch irgendwelchen Betrug etwas an dich gebracht habest, mit falscher Waare, Geld, Maß, Wucher, List, Raub, Diebstahl, ob du mit Rath und Willen ihn schadest oder geschadet habest?

Siebentens. [Erforsche dich,] wie es stehe in den Dingen, die den guten Namen betreffen. Ob du [deinen Nächsten] herabsetzest, beißest, und ob du seiner Sünde mit der Zunge gegen ihn schuldig seiest? Ob du dich freuest, wenn andere das thun, was du selbst nicht thust.

Achtens. Wenn du in Werken und Worten untadelig (integer) bist (was sehr selten vorkommt), so gehe nun in dein Herz, erforsche seine Geheimnisse, und erwäge seine Begierden, und wenn du nicht lügen willst, so wirst du finden, daß du voll böser Lust bist. Auch kannst du zu dir nicht sagen: Ich habe keine böse Lust, denn die Schrift sagt dawider [Spr. 20, 9.]: „Wer kann sagen: Ich bin rein in meinem Herzen?“ Kein Mensch ist inwendig rein. Denn es ist nicht genug, daß du deine böse Lust im Zaume haltest, sondern es soll dich auch nichts gelüsten; darüber soll am gehörigen Orte weitläufiger gehandelt werden. Darum sind die beiden letzten Gebote auch die größten, welche der Apostel im Briefe an die Römer [Cap. 13, 9.] allein anzieht, indem er sagt: „Dich soll nichts gelüsten.“

\quellseite{611}
Deshalb, wer die Liebe hat, das heißt, einen gesunden, völligen Willen zum Gesetze GOttes, und dasselbe liebt, der erfüllt das Gesetz GOttes. Denn es steht nicht zu fürchten, daß derjenige, welcher GOtt liebt, seinen Namen verunehren oder einen anderen Gott suchen sollte, oder den Feiertag schänden. Und wer den Nächsten liebt, von dem steht nicht zu fürchten, daß er die Eltern verachte, tödte, die Ehe breche, stehle, falsch Zeugniß rede, oder auch böse Lust hege, da er vielmehr darnach trachtet, dem Nächsten Gutes zu thun. „Die Liebe“, sagt der Apostel [Röm. 13, 10.], „thut dem Nächsten nichts Böses“, weder in Wort noch Werk noch Wunsch.

Doch wir wollen wieder auf den eigentlichen Gegenstand (materiam) dieses dritten Gebotes zurückkehren. Ich habe gesagt, daß der Sabbath einst, wie er auch noch jetzt ist, zu einem Vorbilde und einer Ceremonie gedient habe um der Schwachen willen, jetzt aber sei jeder Tag ein Festtag, an derer keinem man sündigen solle. Denn da vorzeiten der gesteinigt worden ist, welcher Holz auflas am Sabbathstage, 4 Mos. 15, 32. ff., was wird denen widerfahren, die in der Zeit der Gnade entsetzliche Verbrechen begehen, noch dazu meistens an Festtagen? Dies ist geschrieben, damit wir dadurch erschrecken; denn Holz auflesen ist nicht eine so große Uebelthat, sondern bedeutet (significat) eine Uebelthat, wie auch andere Dienstbarkeit am Feiertage.

\footnote{Löscher: „Das Folgende hat Luther am 14. Sonntage nach Trinitatis [den 24. August 1516] gepredigt.“ Das Exordium dazu findet sich Walch, St. Louiser Ausg., Bd. XII, 1736 ff.} Die ersten sind nun diejenigen, welche in grober und roher Weise am Feiertage sündigen, nämlich mit Schwelgen, Saufen, Spielen, Zorn, Mord, Fluchen, Raub, Nachstellen (insidiis), Betrug, bösen Ränken und dergleichen, von denen Maleachi, Cap. 2, 3., sagt: „Ich will den Koth eurer Feiertage euch ins Gesicht werfen, und soll an euch kleben bleiben.“ Siehe, Koth nennt er die Feiertage solcher Leute, weil sie an denselben auch Koth von sich geben, wie die unreinen Säue. Und Amos, Cap. 5, 21.: „Ich bin euren Feiertagen gram und verachte sie.“

Denn diese Leute setzen ebenso wie die Juden (von denen er hier redet) die ganze Herrlichkeit des Festes in schöne Kleider, Aufwand in Essen und Trinken,\footnote{Statt victu, splendore in der Weimarischen haben wir mit den anderen Ausgaben (Löscher, Petri, Wittenberger, Jenaer und Erlanger) victus splendore angenommen, weil bei ersterer Lesart vier Dinge herauskommen. Letzterer Lesart entspricht lautus victus in der „heimlichen Deutung“.} und Aufhören zu arbeiten. Wenn diese drei Dinge da sind, ohne daß, was sie bedeuten, so machen sie eine jüdische Feier aus, und sind nichts als Neigungen zur Unzucht und Lastern. Denn wenn der Bauch voll ist und die Seele müßig, und dann noch Schmuck vor die Augen gestellt wird, was sollte dabei anders herauskommen, als daß dadurch die Leute zu Schwelgerei und bösen Begierden veranlaßt werden? Und diejenigen, welche so die Feiertage zubringen, denken nicht, daß, wenn das den Feiertag halten hieße, dies auch eine jegliche Hure, oder jeder Türke, oder Heide thun könnte, ja, ein Pferd und ein Schwein, nämlich fressen und saufen, ferner, müßig gehen und endlich, sich prachtvoll aufputzen.

Hier kann man etliche finden, die am Abend vor dem Feste (vigiliam) fasten, bisweilen bei Wasser und Brod, aber den Tag selbst mit Fressen und Schwelgen zubringen. Und es ist ganz allgemein geworden, daß man den Sonnabend heiliger hält als den Sonntag, weil man sich wenigstens der Fleischspeisen, des Müßiggangs und der Pracht enthält, was am Sonntage in vollem Schwange geht, und alle heiligen Abende (vigiliae) werden heiliger gehalten als die Festtage selbst; freilich ganz verkehrt.

\subsection{Heimliche Deutung (mysteria)}

Das leibliche Ruhen (ocium) bedeutet die geistliche Ruhe, daß man ablasse von bösen Lüsten und Gedanken, damit man geschickt sei für GOttes Wort, welches eine leere Seele haben will. Deshalb muß sich ein jeglicher Christ fragen: Weshalb bist du heute müßig? Und in solcher Weise soll er dies als ein Erinnerungszeichen und Denkzettel gebrauchen, daß er feiern müsse (vacandum), um GOttes Wort zu hören.

Der Schmuck der Kleider bedeutet die guten Werke, als, Messe hören, beten, Almosen geben, und andere Werke der Barmherzigkeit thun, wie Christus wider die Pharisäer gelehrt hat; davon hernach weiter.

Köstlich Essen und Trinken (lautus victus) bedeutet, daß man die Wohlthaten GOttes, das Leiden Christi, die eigenen Sünden bedenke, das Wort GOttes lese und immer wieder überdenke. Davon hernach.

\quellseite{612}
Die zweiten sind, welche Werke thun, die sonst erlaubt sind, als da sind die Werke der Noth, und die Werke, welche dazu dienen, das zu bereiten, was wir zu unserem Unterhalte bedürfen. Und diese sündigen weniger als die vorigen, und um so viel weniger, daß sie auch in vielen Fällen entschuldigt werden.

Der erste Fall dieser Art ist die Noth. Wie man aber Noth verstehen müsse, ist zweifelhaft, weil hier eine außerordentlich große Verschiedenheit stattfindet. Dies kann aber nicht besser bemessen werden als nach der Regel unseres Heilandes, welcher im Evangelio sagt [Luc. 14, 5.]: es sei nicht Sünde, einen Esel oder einen Ochsen aus der Grube zu ziehen oder zu tränken. Hieraus entnimmt man diese Regel: Wenn ein so großer Schade zu befürchten stände, oder auch ein geringerer, als ein Esel oder dessen Gesundheit geschätzt wird, so kann das Werk geschehen, durch welches ein solcher Schade verhütet wird. Dies muß jedoch so verstanden werden: wenn etwas Derartiges nicht aus eigener Wahl, sondern wegen eintretender Noth durch einen Unfall eintreten sollte. Denn dies muß ich darum hinzufügen, damit nicht die Geizigen meinen, ihnen sei Freiheit gegeben nach ihrem Belieben zu verkaufen oder zu kaufen, weil nämlich, wenn sie nicht verkaufen oder kauften, ein so großer Gewinn an dem Tage verloren ginge. Denn in solcher Weise könnte kein Festtag angesetzt werden, weil ein jeglicher [Festtag] um den Gewinn bringt, da er die Geschäfte hindert, die Gewinn bringen, und die Arbeit. Die Regel wird aber dadurch bewiesen: Weil es den Juden erlaubt war, und sie daran keine Sünde thaten, viel mehr ist es den Christen zugelassen, welche freier sind von gesetzlichen Ceremonien!

Hiernach mögen nun die besonderen Fälle leicht beurtheilt werden. Zum Beispiel, wenn ein Kaufmann an einem Festtage kauft oder verkauft, weil er an einem andern Tage oder zu einer andern Stunde nicht ohne großen Schaden kaufen oder verkaufen kann, als, weil der Käufer oder Verkäufer auf der Durchreise begriffen ist und nicht warten will. Dieser wird deswegen entschuldigt, weil er das Werk und den Gewinn nicht gesucht hat, sondern das, was ihm geboten wurde, annahm, um Schaden zu verhüten, es aber sonst nicht gethan haben würde, und sein Herz so gestanden hat, daß er lieber gewollt hätte, es wäre ihm zu einer andern Zeit angeboten worden. Denn in solchem Falle, wie sich das Geschäft unverhofft zugetragen hat, so kann er es auch bald abmachen und sich wieder zu der Feier wenden. Ebenso verhält es sich mit jedem anderen Schaden, auch geringerem, doch unter der Bedingung, daß der Gottesdienst, nämlich die Messe und das Wort GOttes nicht versäumt werde, es sei denn große Noth, welche nach dem Urtheile eines guten Mannes als Entschuldigung dienen möchte. Aber, wie ich gesagt habe, dies ist zu verstehen von dem Schaden, welcher nicht von dem Stillehalten mit dem Gewerbe herkommt, oder von dem Unterlassen des Geschäfts, sondern von ungefähr. Denn ein Schade durch Unterlassung des Geschäfts kann zu jeder Stunde geschehen.

Zweitens von denen, die da kochen und die Nahrung zubereiten. Denn dies ist nothwendig, und man muß hier nicht nach dem jüdischen Aberglauben sich unsinniger Weise richten, da die Juden am Sabbathtage nicht kochen, wiewohl sie fressen, saufen und schwelgen. So ist es auch erlaubt, Speise und Trank herzuschaffen und, um auf unsere Sitten zu kommen, so ist es nicht Sünde, ein Faß Bier in den Keller zu legen, aber das ist mehr Sünde, daß man es in so viel Stübchen (modiis) und Kannen\footnote{Ein Stübchen = eine Gallone; Kanne = Quart.} herausbringt, so daß sie auch so viel saufen, daß sie es wieder speien. So ist es eine größere Sünde, wenn man aus der Küche und dem Keller allzuviel herausholt, als wenn man es hineinbringt: um dieses machen sich die Juden keine Sorge, aber vor jenem hüten sie sich ängstlich (curant), und viele folgen ihnen. Es ist aber eine wunderliche Thorheit, daß sie [am Sabbath] keine Speise bereiten. Warum lassen sie denn nicht auch das Ankleiden anstehen und gehen nackt einher? denn es ist ja das Anziehen auch eine Arbeit und ein Werk. Ferner, warum spazieren und gehen und stehen sie? Warum sitzen und schlafen sie nicht? Warum essen und trinken sie? Sind denn dies nicht auch Werke? Wie halten sie also den Sabbath, das ist, wie ruhen sie denn? Ferner, es ist nicht eine geringe Arbeit, daß man stehen muß bei der Messe, bei der Predigt und den gottesdienstlichen Handlungen; das beweise ich durch alle, welche dies mit Unlust thun, deren es hier und anderswo viele gibt. Also sagen sie: Ich muß feiern und
\quellseite{613}
nicht arbeiten; und handeln nach dem Gebote wider das Gebot. Denn dahin ist es durch die jüdische Weisheit gekommen.

Denn in solcher Weise sind einige auch bis ins Kleinliche gegangen, daß sie sogar sagen, es sei verboten, Kränze zu machen. Wenn sie dies verstehen von denen, die sie verkaufen, so ist es wahr. Aber eine Jungfrau kann wohl früh Morgens einen Kranz machen, um sich damit zu schmücken. Denn sie kann ihn ausbeßern, wenn er zerrissen ist, also auch einen machen; ja, sie kann ihre Kleider ausziehen und sich schmücken, ihre Kleider ausschütteln und reinigen, um sie anzuziehen. Daß sie ihren Mantel (peplum) in Ordnung bringe und anderen Schmuck, ist nicht wider dies Gebot, es sei denn, daß dies geschehe nach der Weise etlicher, die nicht vom Spiegel wegkommen können und deshalb dadurch die Gottesdienste versäumen. Die sich aber aus Hoffahrt oder um böser Lust willen schmücken, die sündigen nicht wider dies Gebot, sondern sündigen um so schwerer, weil sie wider ein anderes Gebot GOttes sündigen an einem heiligen Tage.

Drittens, man darf das Vieh mit Futter und Streu versehen, melken, weiden, versorgen etc.

Viertens ist zu sagen von Weinschenken (tabernariis), Fleischern, Bäckern, Fischern und anderen Handwerkern, welche andere Leute mit Nahrung und Kleidung versehen. Diese sündigen nicht, wenn sie nicht Säufern und Spielern dadurch Vorschub leisten, ihnen helfen oder günstig sind.

Fünftens [mögen die Werke geschehen], welche zum gemeinen Nutzen dienen, als, zur Ausbesserung einer Brücke, zur Verhütung einer Wasserfluth, zur Abwehr eines Feindes, zum Löschen von Feuer und zur Abwendung anderer öffentlicher Noth.

Sechstens [mag Arbeit gethan werden] von Aerzten,\footnote{Weimarische: medicis, in den andern Ausgaben: medicinis.} Läufern, Boten, welche nothwendige Dinge zu befördern haben, von Schmieden, welche Pferde beschlagen, von Mühlen etc. Dies alles wird in der Noth entschuldigt, wenn nur dadurch der Gottesdienst nicht versäumt wird.

Siebentens, die auf den Jahrmarkt ziehen, weil sie es zu anderer Zeit nicht können. Und hiernach mag über ähnliche Fälle geurtheilt werden.

Aber hier erhebt sich die Frage: ob die Wittenberger entschuldigt sind, daß sie an einem Festtage ein Armbrutschießen anstellen, um einen Vogel von der Stange zu schießen, da sie dies ebenso gut an einem andern Tage thun können, und dabei weder eine Noth, noch die Liebe (pietas), noch die gebührliche Bescheidenheit (modicitas) vorhanden ist, und dieser halbe Festtag nach dem Essen ebensowohl Gotte gegeben werden könnte, als der Vormittag? Ein ähnliches Urtheil ist zu fällen über die Militärparade (armorum ostensione), als ob dies nicht ebenso gut an einem andern Tage geschehen könnte. Hier sollten die Herren Bürgermeister Acht haben.

Das andere ist, daß man Maß halte (modicitas) oder nur wenig thue. Darum mag man hier alle die Werke thun, die man an anderen Tagen thut; wenn es nur wenige sind. Als wenn ein Schuster Ein Paar oder zwei Paar Schuhe zuschnitte, ein Schneider einen Rock, ein Krämer etwas verkaufte, was er nicht zum Verkauf ausgelegt hat. So magst du in ähnlichen Dingen durch allerlei Verrichtungen (officia) und Werke laufen.

Das dritte ist die Liebe, da man den Armen und Dürftigen zu Hülfe kommt, desgleichen den Gemeinden (ecclesiis). So ist es auch etwas Heiliges, wenn man einen Kranken badet, einen Unreinen wäscht, seine Kleider, sein Bett reinigt, das Bett macht, ihm dient mit Nahrung und Handreichung, sonst würde es nicht erlaubt sein, um Almosen zu bitten oder dieselben zu geben, da doch die Bettelorden dies nie [in solchem Grade], wie an den Festtagen\footnote{Weimarische: nunquam ut diebus festis; Löscher, Wittenberger, Jenaer: nonnunquam in diebus festis; Petri: nunquam, nisi in diebus festis.} thun, gleicherweise auch alle anderen. So mag man ihnen [den Armen] auch Getreide, Holz, Speise und Trank bringen, das ist erlaubt, wenn es nur aus Liebe geschieht, und nicht aus Gewinn eines Gegendientes (retributionis) oder des Freßens halben, wie es jetzt überall der Brauch ist.

\footnote{Löscher: „Das folgende hat Luther am Tage St. Matthäi [21. Sept.] im Jahre 1516 gepredigt, wie aus der Handschrift erhellt.“ Das Exordium hierzu findet sich Walch, St. Louiser Ausg., Bd. XII, 1744.} Die dritten sind diejenigen, welche den Tag nach dem Buchstaben heilig halten, aber nicht geistlicher Weise, da sie zwar die Werke der äußerlichen Heiligung thun, aber nicht mit dem

\quellseite{614}
Herzen. Und wiewohl sie der Kirche ein Genüge thun und weniger sündigen als die Vorigen, so sind sie doch Leute wie Kain (Cainistae), die nicht die Person opfern, sondern das Werk der Person. Sie hüten sich nur vor Aergerniß. Und wollte doch GOtt, daß alle, die es ärger machen, doch nur solche Leute wären.

Um den Feiertag zu heiligen, muß man fünf Dinge thun, wie sich aus dem Decret [des geistlichen Rechts] ergibt, nämlich Messe hören, das Wort GOttes hören, beten, opfern (wie einige wollen), und Reue haben über die Sünden.

Erstlich die Messe hören. Hier deuten dies einige überaus thöricht auf den Canon, indem sie um desswillen wollen, die Messe solle mit erhobener Stimme gelesen werden, weil geboten sei, nicht, sie zu sehen, sondern zu hören, als ob nicht vielmehr geboten sei, den Canon der Messe, welcher allein eigentlich die Messe ist, mit verhaltener (pressa) Stimme zu lesen, damit er nicht gehört werde. Das Uebrige aber sind Gebete und Vorbereitungen (paramenta) auf den Canon der Messe. Wenn sie es nun um des Evangelii willen lauter gelesen wissen wollen, warum bedenken sie denn nicht, daß in Bezug auf das Wort GOttes ein anderes Gebot gegeben ist als über die Messe? Doch fort mit diesen unnützen Träumen! da in der ganzen Kirche niemand anders als der Priester allein die ganze Messe hört, nämlich, weil dies Geheimniß groß ist, daß Christus, der wahre Priester, wiewohl er in der Kirche vieles durch äußerliche Sacramente und sichtbare Ceremonien wirkt, in welchen die Gläubigen mit einstimmen und mitwirken, wie jetzt der Clerus thut gegen den Priester mit Gesang, Gebet und Antwort, dennoch das Meiste für uns bei GOtt in unsichtbarer Weise ausrichtet, was das Volk und seine Gemeinde nicht sieht noch gewahr wird.

Darnach ist es eine unentschiedene (anceps) Frage, ob es genug sei, wenn man die Frühmesse höre; desgleichen, ob es genüge, wenn man sie in den Klöstern höre. Hier muß man einen Unterschied machen: Wenn eine Noth da ist (worüber oben geredet ist), so ist genug, in der Frühe eine Stillmesse (privatam) zu hören, oder auch eine Messe im Kloster, denn Noth hat kein Gebot. Wo aber keine Noth da ist, da macht man fast immer diesen Unterschied, daß, wenn die Pfarrkirche nicht aus Verachtung verlassen wird, sondern das Kloster aufgesucht wird aus Andacht, es dann freisteht, sie im Kloster zu hören. Aber nach meinem Urtheile möchte ich sagen (ein anderer möge seine Meinung haben), in dem großen Jammer unserer Zeit\footnote{Auf das zunächst Folgende bezieht sich Luthers Bemerkung im Druckfehlerverzeichniß: „Das, was auf der ersten Seite von Blatt H gesagt ist, daß man nicht in frevelhafter Weise die Messen in den Klöstern hören solle, ist gesagt und bereits gedruckt gewesen, ehe die Erklärung Leo's des Zehnten ausgegangen war, darin er die Erlaubniß gibt, die Messe überall in den Klöstern zu hören, wiewohl es auch jetzt, wo diese Erlaubniß da ist, sicherer ist, seiner eigenen Andacht nicht zu glauben.“ — Die Verfügung Leo's X., von der Luther hier redet, ging aus unter dem 13. November 1517. (Weimarische Ausgabe, Bd. IV, 538, Anm. 2.)} ist es durchaus nicht sicher, daß man einen jeglichen, besonders einen einfältigen und unverständigen Menschen, seiner Andacht und seinem Ermessen nachgeben lasse wider die klaren Regeln (canones) [des geistlichen Rechts]. Die Ursache davon ist, daß darum die Canones der Väter gegeben sind, daß sie das Volk leiten, und es muß einem jeglichen seine Andacht immer verdächtig sein, und zwar um so mehr, je mehr sie ihm gefällt, obgleich sie wider den gemeinen Gebrauch oder die Bestimmung (decretum) der Väter geht, da wir wissen, daß der Satansengel, der sich in einen Engel des Lichts verstellt, sehr zu fürchten ist, und man sich vor ihm aufs äußerste in Acht nehmen muß. Denn unter dem Schein der Andacht möchte er leicht alle Decrete zunichte machen, und auch das Wort 5 Mos. 12, 8. und Nicht. 17, 6., wo geboten, ja vielmehr verboten und gestraft wird, daß nicht ein jeglicher thue, was ihn recht und gut dünkt. Sonst geht zuerst die Furcht vor Menschen verloren und die Ehrerbietung gegen GOttes Kirche und gegen das Gemeinwesen. Darnach wird die Sicherheit genährt [und verlangt], daß allmählig auch andere größere Dinge nachlassen werden sollen. Drittens kehrt die Zügellosigkeit des alten [Juden-] Volkes wieder, daß auf den Höhen und in den Hainen und außerhalb der Kirchen Gottesdienst gehalten wird, was aufs höchste verboten war. Denn in der That sehen wir, daß, wenn diese Willkür der Andacht erlaubt wird, nicht allein die Pfarrkirchen, sondern auch die Klöster der Stadt verlassen werden, und ihnen die Feld- und Dorfkirchen unter demselben erheuchelten Schein der Andacht vorgezogen werden, und daß in solcher Weise der Gottesdienst
\quellseite{615}
der Kirche endlich besteht in Besuch von Schenken und Handel auf dem Jahrmarkt unter dem Vorwande der Kirchweihe, der Schutzheiligen (patronorum), der Gelübde etc.

Deshalb sage ich, daß nach meiner Meinung, was auch der Grund sei andermal Messe zu hören, keinem Menschen erlaubt sei, gerade hierin seinem Urtheil, seiner Meinung, seiner Andacht zu folgen, sondern der Ordnung (autoritatem) der Kirche, und daß man seinen Sinn unter den Gehorsam Christi gefangen geben soll [2 Cor. 10, 5.], weil geschrieben steht [Spr. 28, 26.]: „Wer sich auf sein Herz verläßt, der ist ein Narr.“ Und die Kirche kann nicht irren, aber ein jeder kann in seiner Andacht irren, da auch viele Propheten, Heilige und Könige in ähnlicher Weise geirrt haben, uns zu einer Warnung und Schrecken. Glaube daher ja keiner Glosse über dieses Decret. Denn wer solche Glossen macht, der ist unwillig in Gottes Gesetz und hat noch nicht erkannt, was der Teufel und der Mensch sei, viel weniger, was GOtt sei. Und es ist eine wunderbar verkehrte Glosse, welche gestattet, daß man seiner Andacht nach seinem Belieben nachlaufen dürfe, während der Text dies doch ausdrücklich hat verhindern wollen. Ja, halte immer alles verdächtig, was dich auf dich selbst hinweist und dir etwas nachläßt, denn es macht, „daß du deinen Willen übest, wenn du fastest“, wie Jesaias [Cap. 58, 3.] sagt. So spottete der Prophet Micha des Ahab und sprach [1 Kön. 22, 15.]: „Ja, zeuch hinauf“ etc. Aber allem, was dich von dir selbst abzieht, dem folge, als, Matth. 5, 25.: „Sei willfertig deinem Widersacher“ etc. Micha 7, 5. [Zulg.]: „Hüte dich vor deinem Freunde“; desgleichen [Jer. 9, 4.]: „Ein Freund verräth den andern.“

Das zweite ist GOttes Wort hören, was wiederum unter dem Gebote steht, daß man es höre in seiner Pfarrkirche, ebenso wie die Messe, es sei denn, daß es unterlassen wird wegen der Kirchweihe anderer Kirchen oder Klöster. Daß man Nachmittags GOttes Wort in diesen hören möge, ist zwar ein Rath der Kirche, aber doch\footnote{In allen Ausgaben sed tamen, was die Weimarische, wie es uns scheint, nicht gut, in sed non verändert hat.} dem Gebote GOttes gemäß, der da will, daß man sein Wort überall mit aller Andacht hören soll. Und dies zweite ist nothwendiger als das erste, ja, das erste ist um dieses zweiten willen, wie aus Christi Worte erhellt, da er sagt [1 Cor. 11, 25.]: „Solches thut, so oft ihr es thut, zu meinem Gedächtniß“, als ob er sagen wollte: Haltet keine Messe, es sei denn, daß ihr auch das Evangelium predigt. Und der Apostel sagt [1 Cor. 11, 26.]: „Denn so oft ihr von diesem Brod esset und von diesem Kelche trinket, sollt ihr des HErrn Tod verkündigen“ etc. Darum darf man keine Messe halten ohne das Evangelium; die Privatmesse mit sonderlicher, die öffentliche Messe mit öffentlicher Verkündigung des Evangelii. Ja, in der ersten Kirche waren auch die Gebannten, die Büßenden, die Katechumenen, die Besessenen (energumeni)\footnote{Vgl. Walch, St. Louiser Ausg., Bd. XVIII, 140.} bei der Messe zugegen bis nach dem Evangelium; dies wurde die Messe der Katechumenen genannt, wie man aus Dionysius sieht. Wollte GOtt, daß dies noch heutzutage geschähe, daß die Gebannten wenigstens das Wort GOttes hörten, damit sie [ihre Sünde] desto leichter erkennen und einen Stachel im Gewissen fühlen möchten. Denn es wäre ein überaus grausamer Bann, der auch das Hören des Wortes GOttes verbieten sollte, da es doch geboten ist, daß es allen Creaturen gepredigt werden soll.

Hier erhebt sich wieder die Frage, ob es genug sei, wenn man das Evangelium höre in der Stillmesse. Darauf ist das die Antwort, daß es genug sei, wenn er es verstehet und zugleich gezwungen ist, eine Messe anzutreten. Sonst, wenn er es auch verstehet, so ist er doch gehalten, das Wort GOttes zu hören, wenn es öffentlich (solemniter) mit lauter Stimme (vocale) verkündigt wird. Sonst würde es einem jeglichen freistehen, die öffentliche Predigt zu verlassen, und sie würde in Verachtung gerathen. Ja, wenn er von Noth wegen reisen muß, so ist er entschuldigt, wenn er es nicht verstehet, ja, auch nicht hört. Jedoch muß ein jeglicher bedenken, nicht allein, aus wie großer Noth er gezwungen sei, abwesend zu sein, da dies eine leibliche Noth ist, sondern auch, aus wie großer Noth man das Wort GOttes haben müsse, da dies eine geistliche Noth ist. Darum muß ein gottseeliges Verhalten in himmlischen Dingen (pietas coelestis) immer der irdischen Begierde den Rang ablaufen.

\quellseite{616}
\footnote{Die erste Hälfte dieses Absatzes kehrt fast von Wort zu Wort wieder in „Luthers Vorlesung über das Buch der Richter“, Cap. 1, 8. (Weim. Ausg., Bd. IV, S. 538.) Es liegt die Vermuthung nahe, daß dieser Passus dort auch auf Rechnung des Schreibers zu setzen sei, der höchst wahrscheinlich auch andere Lücken seines Manuscripts auf eigene Hand ergänzt hat. (Weim. Ausg., Bd. IV, Nachträge.)} Hier ist jedoch zu merken: Wenn nun das Volk mit so großer Strenge verbunden ist, GOttes Wort zu hören, Lieber, mit wie viel größerer Strenge sind die Priester dazu verbunden, daß sie das Wort GOttes auch verkündigen, und Sorge dafür tragen, daß es verkündigt werde! Da dies nun über alle Gebote aufs dringendste befohlen ist, so weiß ich nicht, wo der Jammer herkommt, daß man sich heutzutage nichts weniger angelegen sein läßt. An vielen Orten läßt man es ganz und gar anstehen, an mehreren Orten aber predigt man solche Dinge, die man lieber sollte anstehen lassen, und an anderen Orten werden menschliche Dinge, Rechtslehren und Weltweisheit gepredigt; das Evangelium aber, das ist, das Gedächtniß Christi, ist so selten, daß einem ein Grauen ankommen muß, wenn man das ansehen muß. Das Volk wird gezwungen zur Pfarrkirche, und gehorcht nicht, weil die Priester gezwungen werden, das Evangelium zu predigen, und es nicht thun. Was ist’s also Wunder, wenn das Volk herumstreicht zu anderen Kirchen in abergläubischer Zügellosigkeit, da die Priester in großer Unwissenheit mit Fabeln und nichtigen Lehren Narrenwerk treiben? Du predigst das Evangelium nicht, und hältst es doch für eine große Unthat, daß das Volk hin und wieder läuft zu den Höhen der Götzen? Da du dich der Kirche in deinem Amte und dem Worte [GOttes] entziehst, was kann anders daraus folgen, als daß das Volk sich auch dem Gehorsam und der Predigt (anditu) entzieht? Du bist buchstäblich inwendig ein unnützer Mensch, mit dem Geiste stehst du draußen (foris = außerhalb der Kirche). Darum wandert auch das Volk gleicherweise hinaus nach dem Leibe und nach dem Geiste, und sündigt dadurch um so viel weniger als du, je bereitwilliger es ist, das Evangelium zu hören, wenn du es predigest; ja, wenn du es auch nur gelernt hättest! Du gibst also den Anlaß dazu, daß die Kirche leer wird, und du reichst den Schafen nicht ihre Weide und Nahrung, während du doch das Vieh in den Stall treibst. Deshalb werden die Schafe Christi auf den Bergen zerstreut [Matth. 9, 36.], weil kein Hirte da war [Sach. 10, 2.], spricht der HErr.

Das dritte ist beten. Vom Gebete ist vieles zu sagen, aber hier muß es vom mündlichen Gebete verstanden werden, dies ist nicht zu leugnen; aber doch soll es das Geringere und ein Kleines sein gegen das geistliche Gebet. Denn das Gebet ist die Erhebung des Gemüthes zu GOtt. Die Worte sind die Leiter, aber die Stimme ist das Angeben (apparatus) der Leiter. Deshalb habe ich gesagt, daß diejenigen den Tag nach dem Buchstaben und schlecht heiligen, welche nur mit der Stimme beten; ja, sie beten nicht, sondern tragen nur die Gebete anderer vor, nicht anders, als wie ein Buch eines Anderen Schrift und Worte, das selbst todt ist und nichts versteht, oder wie eine Orgel den Gesang der Geistlichkeit (cleri) zu Gehör bringt, selbst aber davon nichts empfindet oder versteht. Hierin liegt heutzutage ein großes Geheimniß (mysterium = geistliche Deutung): nämlich wie die Orgel ist, so ist auch die Geistlichkeit, das heißt, ein Schreier ohne Verstand. Eine Stimme ohne Worte bringt, gleichwie die Orgel, nur einen gleichsam stummen und sprachlosen Laut und Ton hervor. Von dem geistlichen Gebete aber werden wir gleich nachher im fünften [Abschnitte]\footnote{Löscher, Wittenberger, Jenaer und Erlanger haben falsch: in quinto praecepto.} sagen.

Das vierte ist (wie einige annehmen) opfern; doch dies ist abgeschafft, wiewohl in der ersten Kirche dies Darbringen (oblatio) von Brod und anderen Speisen ungemein geschah; davon wurde geopfert, und was übrig blieb, wurde den Armen gegeben. Aber jetzt sind anstatt dessen die Kirchen, Klöster, Hospitäler etc. mit Stiftungen für die Armen versehen.

Fünftens, was das Gröbste und erste von allen ist, nämlich, daß man sich mit GOtt versöhne durch Prüfung des Gewissens und Reue über die Sünden. Diese Reue ist aber so anzustellen, daß sie nicht so sehr aus Haß [gegen die Sünde] als aus Liebe [gegen GOtt] hervorgehe. Aus Liebe wird sie aber dann hervorgehen (um davon den Einfältigen ein Exempel zu geben), wenn ein Mensch wieder und immer wieder die Wohlthaten überdenkt, die GOtt ihm in seinem ganzen Leben erwiesen hat. Dann
\quellseite{617}
möge er dagegen seine Undankbarkeit und seine Uebelthaten halten, die er begangen hat durch Mißbrauch solcher Güter, als, daß er uns fünf Sinne gegeben, uns in so vielen Gefahren bewahrt hat, so viele Güter des Leibes, der Seele und allerlei Dinge (rerum) gegeben hat. Dies alles hat er nicht den Thieren, den Steinen, ja, auch sehr vielen Menschen nicht gegeben. Darnach steige man zu den geistlichen Wohlthaten empor, daß er seinen Sohn ins Fleisch gesendet hat, daß er ihn hat kreuzigen und sterben lassen, und andere Gaben unermesslicher Größe, endlich, daß er uns die ewigen Güter verheißen hat. Wenn dies alles wohl erwogen und gehalten wird gegen die eigenen Sünden, so erweckt dies einen außerordentlichen Haß und Abscheu gegen uns selbst, aber Liebe und Lob gegen GOTT. Dies ist eine ganz wahrhaftige Reue, die lebendig und kräftig ist, während diejenige, welche hervorgebracht wird durch die Furcht vor der Hölle und durch die Schmach der Sünde, nur eine Reue ist dem Buchstaben nach, erdichtet und von kurzer Dauer, weil sie nicht gewurzelt ist in der Liebe, sondern nur aufgebracht durch Furcht.

So liest man von einem Menschen, der sein ganzes Leben lang sehr übel gelebt hatte, aber dennoch durch wunderbare Wohlthaten und Gaben GOttes emporgehoben wurde, endlich sogar zum Bischoffamte gelangte. Da schlug er in sich, hielt bewundernd GOttes Güte gegen seine Unwürdigkeit und rief aus: HErr, du hast gesiegt, und hast mich durch deine übergroße Güte gefangen; siehe, von nun an will ich dir ganz und gar dienen; und so that er auch. Und man soll dieses grobe Aufsteigen von den niedrigen (primis), leiblichen Gütern etc. nicht verachten. Denn man liest in den Geschichten, daß zur Zeit des Conciliums zu Costnitz zwei Cardinäle, die zum Concilium ritten, einen Hirten auf dem Felde weinen sahen. Aber einer derselben, dessen Herz von Mitleid sehr bewegt ward, sagte: Ich kann nicht vorübergehen, ohne diesen Menschen zu hören und zu trösten, ritt zu ihm hinan und forschte, warum er weine. Als nun der Hirte die Ursache seiner Traurigkeit\footnote{Hier haben wir die Lesart des Manuscriptes: tristitiam angenommen statt tristis in den Ausgaben.} nicht sagen wollte, bestand der Cardinal darauf [daß er es wissen wollte] und zwang ihn, es zu bekennen. Endlich sagte er: Siehe dieses Thier an (und sah auf eine Kröte): mich hat GOTT zu einer so herrlichen Creatur gemacht, zu einem aufrechten, schönen Menschen, der Vernunft hat, und ich habe ihm noch nie dafür gedankt, daß er mich nicht auch zu einer so ungestalten Kröte gemacht hat; das ist die Ursache, weshalb ich bitterlich weine. Da nun der Cardinal sich entsetzte über die ganz ungewöhnlich große Reue dieses Menschen und plötzlich inne wurde, daß er sich gegen viel größere Wohlthaten, die ihm zutheil geworden waren, noch viel undankbarer bewiesen hätte, fiel er wie todt von seinem Maulthier. Da er nun endlich durch seine Diener aufgehoben und in die Stadt gebracht wurde, schrie er immer: O heiliger Augustin, wie richtig hast du gesagt: Die Ungelehrten stehen auf und reißen den Himmel zu sich, und wir mit unserer Gelehrsamkeit wälzen uns in Fleisch und Blut!

Dies sind die Werke des wahren Heiligens, welche durch alle äußerlichen Werke bedeutet werden. So heißt es in den Klageliedern Cap. 3, 20. [Vulg.]: „Ich werde daran gedenken, und meine Seele wird in mir verschmachten.“ Ohne diese zwei Stücke wird das Gemüth und das Herz nicht entzündet; wenn aber das Herz nicht entbrennt, so geschieht alles andere kalt, nichtig und schädlich, weil es ohne Frucht abgeht. Und diese zwei Stücke müssen aus dem Evangelio gehört werden und das Gehörte immer wieder wohl überdacht werden (ruminari). Denn es sind diese zwei Stücke, nämlich Gottes Guthaten (bona) und unsere Uebelthaten (mala), gerade die Leiter zu GOTT, durch welche Menschen uns herniedersteigen und zu GOTT aufsteigen, wie es abgebildet ist 1 Mos. 28, 12. Es sind auch diese zwei Stücke Opfer des Lobes und des Bekenntnisses, von denen es heißt [Ps. 50, 23. Vulg.]: „Das Opfer des Lobes wird mich ehren“; und es folgt, was das Aufsteigen zu GOTT sei: „Und das ist der Weg, daß ich ihm zeige das Heil GOttes“, das heißt, ich will ihm die ewige Seligkeit offenbaren. Von dem anderen, nämlich [dem Opfer] des Bekenntnisses [sagt Ps. 51, 19.]: „Die Opfer, die GOTT gefallen, sind ein geängstigeter Geist; ein geängstigetes und zerschlagenes Herz wirst du, GOTT, nicht verachten.“ Und diese zwei muß man zugleich opfern, damit sie desto vollkommener seien.

\quellseite{618}
\section{Das vierte Gebot}

\begin{quote}
\footnote{„Vermutlich begann Luther über das vierte Gebot am 5. October [den 20. Sonntag nach Trinitatis] 1516 zu predigen, wir wissen sonst das Exordium (Walch, St. Louiser Ausgabe, Bd. XII, 1748) nicht passend einzufügen; am 12. October scheint er kein Stück des Decalogs behandelt zu haben.“ (Weim. Ausg. I, 447.) — In unserer Ausgabe ist l. c. irrtümlich das Evangelium Matth. 9, 1–8. über die Auslegung gesetzt, es sollte Joh. 4, 46 ff. heißen, denn „nach der consuetudo ecclesiae Romanae ist dies die Perikope des 20. Sonntags nach Trinitatis“. (Weim. Ausg. I, 87.) Die Auslegung handelt von diesem letzteren Evangelium. — Am 12. October [den 21. Sonntag nach Trinitatis] 1516 gab Luther nach dem Exordium die Erklärung eines Theiles des Vater-Unsers. Beides findet sich Walch, St. Louiser Ausgabe, Bd. XII, 1752. Löscher, Reformations-Acta, Bd. I, S. 293, merkt an, daß der größte Theil von Luthers Auslegung des Vater-Unsers verloren gegangen sei.} Du sollst deinen Vater und deine Mutter ehren, auf daß [es dir wohlgehe, und] du lange lebest auf Erden.
\end{quote}

Die Juden haben dies Gebot kraftlos (irritum) gemacht durch ihre Satzungen, gleichwie auch alle anderen [Gebote], weil sie diese Ehre allein auf Worte und Zeichen gedeutet haben und so die Worte des Gesetzes entleert haben (irritatores facti). Denn weil „ehren“ bedeutet, ein Zeichen der Ehrerbietung geben, so haben sie es nach der äußeren Schale dieser Bedeutung genommen, so daß sie lehrten, man solle sie allein mit äußerlichen Zeichen ehren. Aber, wie ich öfters gesagt habe, das Gesetz des HErrn ist geistlich, ohne Wandel, erquickt\footnote{In der Vulgata: convertens (Ps. 19, 8.).} die Seelen und unterweist den inwendigen Menschen. Deswegen legt es dem Herzen zuerst seine Verpflichtungen auf (obligat) und fördert das Herz. Darum wollen wir fortfahren nach diesem zweifachen Verstande, nämlich des Buchstaben und des Geistes.

Die geistliche Ehre gegen die Eltern ist ein Hingeben (exhibitio) des Herzens, und Ehrerbietung, die aus einem bereiten Willen herfließt. Hierin sind zwei Stücke begriffen: erstens der bereite Gehorsam und Wille zu allem, was die Eltern wollen; zweitens gebührendes Werthhalten und hohe Meinung von den Eltern. Denn die Ehre, welche das Herz erweist, ist ein großer Gottesdienst (pietas), und größer als die Liebe gegen den Nächsten, weil die Liebe nur so liebt, daß sie bereit ist, dem Nächsten zu dienen, die Ehre aber flößt ihm überdas eine heilige Furcht ein\footnote{Die Worte: honor vero etiam ultra hoc ei timorem exhibet sind nach der von Löscher gegebenen Anmerkung aus dem Manuscript hier einzufügen; dies ist in der Weimarischen Ausgabe geschehen. In allen andern Ausgaben fehlen sie. Den unverständlichen Zusatz: minus pretiosum existimans hat die Weimarische Ausgabe mit Recht weggelassen.} und hält ihn [den Nächsten] überaus hoch. Deshalb ist dies Gebot gleich nach den Geboten der ersten Tafel gesetzt, weil es von denen handelt, welche Statthalter GOttes sind; deshalb, gleichwie GOtt zu verehren ist mit Ehre und heiliger Furcht, so muß auch sein Statthalter in solcher Weise verehrt werden.

Es ist also offenbar, daß, wer die Eltern in geistlicher Weise ehrt, ihnen erstlich einen bereitwilligen Gehorsam erzeigt und seinen Willen ihnen unterordnet, und um deswillen ihnen nicht allein in Nöthen zu Hülfe kommt, sondern auch in allen Dingen nach ihren Wünschen sich richtet. Und wer seinen Eltern nicht ein solches Herz entgegenbringt, der soll wissen, daß er dies Gebot nicht erfülle, wiewohl auch jetzt sehr viele blind sind in diesem [geistlichen] Verständnis und nur an der bloßen Oberfläche der Worte hängen bleiben. Es heißt also: „Du sollst ehren“, nämlich von Herzensgrunde und aus voller Neigung, „deine Eltern“ etc. Denn wie sollte der den Eltern ungehorsam sein, der ihnen nicht allein sein Herz darbietet, sondern auch so darbietet, daß es mit Ehrerbietung (honore) gegen sie erfüllt ist?

Zweitens (was noch größer ist), er erweist Ehrerbietung (reverentiam) (nicht auswendig nach dem Buchstaben, sondern inwendig im Geiste), das heißt, daß er in seinem Herzen eine hohe Meinung von ihnen hat, und dies wird hauptsächlich in diesem Gebote anbefohlen, wodurch es sich auch (wie ich gesagt habe) von der Liebe unterscheidet, durch welche ein jeglicher dem andern gehorsam und dienen soll, (Gal. 5, 13.)

Aber woher entsteht denn diese hohe Meinung von den Eltern in den Kindern, zumal da sie bisweilen arm sind, niedrigen Standes, krank, ungestalt oder sonst gar geringen Ansehens? Denn an solchen Dingen ärgern sich viele, so
\quellseite{619}
daß sie vertrießt, solche Eltern zu haben, und sie sich derselben schämen, und es ihnen leid thut, daß sie nicht von Reichen, Gewaltigen, Edelleuten etc. geboren sind. O große Blindheit! Diesen Leuten muß diese Regel gegeben werden: Die hohe Meinung von den Eltern kommt her durch Betrachtung (respectu) GOttes und seines Willens. Denn ein frommes Kind wird so denken: Siehe, wenn es der höchsten Majestät nicht unwürdig gedäucht hat, daß ein solcher Mensch mein Vater oder meine Mutter sein sollte, warum sollte es mir unwürdig erscheinen? Wenn es dem höchsten Gotte gefiel, mich in ihnen zu wirken und zu schaffen, warum sollte es mir mißfallen, in ihnen zu entstehen oder geschaffen zu sein? Denn was thun doch die Eltern, um dem Kinde seine Gestalt zu geben, welche nie auch nur daran denken, ob es ein Knabe oder ein Mädchen sein werde, ob es Augen, Ohren, Füße haben werde, ob es lebendig oder todt, weiß oder schwarz sein werde, ob es Haare haben werde oder nicht? Also will ich die Werkstatt meines GOttes verehren, der mich gebildet hat. Ich will meine Augen nicht richten auf die Armuth, häßliche Gestalt, niedrige Herkunft meiner Eltern, sondern auf GOtt, den Werkmeister. So, ja so fließt die Ehrerbietung von GOtt auf die Eltern, kommt nicht von den Eltern her, indem sie höher angesehen werden, als wie sie Fleisch und Blut sind, vielmehr als eine Werkstatt der höchsten Majestät. Denn wer wird die Werkstatt des Urhebers (principis) verachten? Welches künstlich verfertigte Werk haßt seine Werkstatt, in der es bereitet ist? Da nun GOtt will, daß die Eltern geehrt werden sollen, und eine überaus gerechte Ursache dafür da ist, nämlich, weil hier auch eine Werkstatt ist, die GOttes Stelle vertritt, so unterscheidet sich dieses Gebot von den [drei] vorhergehenden allein dadurch, daß GOtt in jenen geehrt wird als in sich selbst, in diesem [Gebote] aber als in einem anderen, nämlich in den Vorgesetzten (regentibus), welche sein Stuhl, seine Werkstatt, sein Altar, sein Gnadenthron sind. Dies ist also die rechte Ehre.

Die Ehre der Eltern nach dem Buchstaben geschieht mit unlustigem Herzen und aus einem Gemüthe, welches sie verachtet, wenn man nämlich allein aus Furcht vor der Strafe äußerlich mit Geberden (signo) die Eltern ehrt. Denn wer sie allein aus Furcht vor der Strafe ehrt, nämlich damit er vor den Menschen nicht zu Schanden werde, oder nicht hier oder im künftigen Leben Strafe erleiden müsse, der ehrt sie so lange, als er sich fürchtet und an die Strafe gedenkt, verachtet sie aber immer inwendig im Herzen; vergibt er aber der Strafe, so verachtet er sie auch äußerlich mit dem Werke, wie durch viele Exempel bewiesen wird und man auch täglich erfährt.

\footnote{Löscher: „Das Folgende hat Luther am 22. Sonntage nach Trinitatis [den 19. October] 1516 gepredigt.“ — (Ein Exordium zu dieser Predigt besitzen wir nicht.“ (Weim. Ausg.)} Es ist gesagt worden, wie sich die Kinder gegen die Eltern verhalten sollen. Aber es ist vonnöthen, daß diese sich wiederum als Eltern erweisen, daß die Eltern ihres Namens würdig seien, damit die Ehre der Kinder nicht Götzen und todten Bildern der Eltern erzeigt werde. Ja, bisweilen sind sie etwa ärger als Räuber, weil sie des Fleisches der Kinder pflegen, aber ihre Seelen tödten. Denn es geschieht gar oft, daß ein Kind an seinem Vater einen noch schlimmeren Mörder hat, und es besser wäre, daß das Kind von Räubern getödtet als vom Vater aufgezogen würde. Oder weshalb meint ihr, daß die Pest hauptsächlich unter den kleinen Kindern wüthet? Aus keinem andern Grunde, als weil GOtt an ihnen das erfüllt, was das Buch der Weisheit (Sap. 4, 11. 14. sagt: „Und wird hingerückt, daß die Bosheit seinen Verstand nicht verkehre, noch falsche Lehre seine Seele betrüge; darum eilet er mit ihm aus dem bösen Leben.“ Und hernach [V. 16.]: „Denn es verdammet der verstorbene Gerechte die lebendigen Gottlosen, und ein Junger, der bald vollkommen wird, das lange Leben des Ungerechten“ etc. Aber davon nachher ausführlicher. Nun wollen wir sehen, was sie ihnen schuldig sind, damit sie Eltern seien.

Nachdem der Apostel Eph. 6, 1. ff. den Kindern befohlen hatte, daß sie die Eltern ehren sollten, und ihnen dies Gebot ans Herz gelegt und erklärt hatte, indem er sagte:\footnote{Statt dicit in den Ausgaben (auch in der Weimarschen) muß mit dem Manuscripte dicens gelesen werden, weil statim de parentibus etc. den Nachsatz bildet zu: postquam filios jussit. Bei der Lesart dicit müßte Luthern ein Fehler in der Reihenfolge zur Last gelegt werden, der sich nicht bei ihm findet.} „Ihr Kinder, seid gehorsam euren Eltern in dem HErrn“, und schlechterdings den Willen kund gibt, daß
\quellseite{620}
sie [die Eltern] nicht bloß nach dem Fleisch Eltern sein sollen, wie die Heiden, sondern „in dem HErrn“, und die Kinder ihnen auch gehorchen sollen „in dem HErrn“: so sagt er alsobald von den Eltern [V. 4.]: „Ihr Väter, reitzet eure Kinder nicht zum Zorn (was er Col. 3, 21. auslegt: „Auf daß sie nicht scheu werden“), sondern ziehet sie auf in der Zucht und Vermahnung zum HErrn.“

Das erste, was er ihnen [den Eltern] aufgelegt wissen wollte, das geht freilich die Seele an (denn er hat nichts von ihrem Leibe gepredigt), daß sie die Kinder nicht reizen sollen zum Zorn, und sie nicht kleinmüthig machen. Dies geht gegen die, welche die Kinder mit Wüthen regieren, und da das Herz noch weich ist, wird es ganz und gar mit Furcht und Scheu erfüllt, in solcher Weise, daß sie aus Haß gegen die Eltern endlich davonlaufen, und thun, was sie sonst nie gethan haben würden. Denn was steht von einem Menschen zu hoffen, der in Mißtrauen und Haß gegen die Eltern an ihnen verzweifelt? Denn das ist nicht die Meinung des Apostels, daß man die Kinder nicht erzürnen oder ihnen nicht Schläge geben solle, sondern man soll sie aus Liebe züchtigen, aber nicht, daß die Eltern ihren Zornmuth an ihnen auslassen sollen, wobei ihnen nichts daran liegt, wie sie den Fehl der Kinder bessern. Denn wenn jemand kleinmüthig geworden ist, der wird zu allem unnütz, und verzagt an allem, ist furchtsam bei allem, es zu thun und anzugreifen, und, was noch schlimmer ist, solche Furchtsamkeit, die in zartem Alter eingepflanzt wird, kann nachher kaum jemals ausgerottet werden. Denn weil sie gelernt haben, bei jedem Wort der Eltern in Schrecken zu gerathen, so fürchten sie sich später selbst vor einem rauschenden Baumblatt. Aber auch den Andern soll man wehren, welche den Kindern bange machen (illudunt) durch schreckliche Gestalten und Geberden, besonders bei Eintritt der Nacht (de nocte). Mit allem Fleiße muß ein Kind zu einer rechten Furcht herangezogen werden, daß es das fürchte, was gefürchtet werden muß, aber nicht, daß es furchtsam sei, wie es etlichen Eltern genug ist, wenn sie die Kinder nur furchtsam haben, was für das spätere Leben überaus schädlich ist.

So viel von dem Abweichen zur Linken. Jetzt aber, damit sie auch nicht nach der andern Seite hin sich verfehlen, sagt er:

Zweitens: „Sondern ziehet sie auf in der Zucht und Vermahnung zum HErrn“ [Eph. 6, 4.], das ist, lehret sie, was gelehrt werden muß, und strafet sie, wenn sie es nicht thun. Denn beides ist für sie nothwendig, daß man sie unterweise in dem, was sie nicht wissen von GOtt, und daß man sie strafe, wenn sie dem, was sie erkannt haben, nicht folgen. Deshalb siehe zu, daß du dafür sorgest, daß die Kinder vor allem in geistlichen Dingen unterwiesen werden, daß sie eher GOtte übergeben werden als weltlichen Geschäften. Aber ach, wie ist dies alles jetzt verkehrt! Und es ist auch kein Wunder, da auch die Eltern selbst nicht unterwiesen sind, noch auch zurechtgebracht (probati) durch die Strafe. Es war aber noch Eine Hoffnung übrig, die beruhte auf dem Lehren der Schulen, daß sie wenigstens da unter dem Einfluß der Wissenschaften in der Gotteligkeit unterwiesen werden möchten, aber auch diese ist dahingefallen. Alle Heiden, besonders aber die Juden unterweisen ihre Kinder mit größerem Fleiße, als die Christen. Deshalb steht es auch überaus schlecht um die Kirche, weil ihre ganze Kraft beruht auf den Nachkommen (successoribus); diese aber werden in ihrer Jugend vernachlässigt, wie ein Garten zur Frühlingszeit.

Darum müssen sie aufgebracht werden in der Unterweisung zum HErrn. Es ist aber die Unterweisung zum HErrn diejenige, durch welche man sie den HErrn JEsum Christum erkennen lehrt, ihnen lebendig ins Gedächtniß prägt, wie er für uns gelitten hat, und was er gethan hat, was er verheißen hat, wie den Kindern Israel geboten war [5 Mos. 11, 19.], daß sie ihren Kindern kund thun sollten die Wunder, welche der HErr an ihren Vätern in Egypten gethan hatte, damit es auch das künftige Geschlecht wissen möchte [Ps. 78, 4.]. Und wenn sie diese Erkenntniß haben, aber doch nicht lernen, den HErrn zu lieben und ihm zu danken und zu beten und Christo zu folgen, dann muß man die Vermahnung (correctionem) zum HErrn anwenden, das heißt, den Schrecken des Gerichtes GOttes und die Drohungen des Zornes GOttes über die Bösen.\footnote{Zu dem vorhergehenden Abschnitt bemerkt Luther im Druckfehlerverzeichniß: „Was über die Unterweisung und Zucht des HErrn an den Kindern gesagt ist, würde besser genommen von der Unterweisung und Zucht, welche sie [die Eltern] nach der Weise GOttes, nicht nach der Weise der Menschen und ihren Leidenschaften, ausüben sollen, damit die Kinder wissen, daß sie durch die Eltern mehr von GOtt unterwiesen und gestraft werden [als durch die Eltern].} Denn wer dies von frühester Jugend
\quellseite{621}
an gelernt hat, nämlich die Wohlthaten GOttes und seine Verheißungen, wegen derer man ihn lieben soll, sodann die Strafen und Drohungen, wegen derer man den HErrn fürchten soll, der wird das leicht festhalten, wenn er älter wird. Denn in diesen beiden Stücken will der HErr verehrt werden, daß man ihn liebe als einen Vater wegen der vergangenen, gegenwärtigen und künftigen Wohlthaten, und daß man ihn fürchte wegen der vergangenen, gegenwärtigen und künftigen Strafen, wie er spricht Mal. 1, 6.: „Bin ich nun Vater, wo ist meine Liebe? Bin ich HErr, wo fürchtet man mich?“ In solcher Weise sollen nun die Kinder lernen, dem HErrn zu singen von Gnade und Recht [Ps. 101, 1.]. Denn diese zwei Dinge hat der Apostel im Auge, da er sagt: „in der Zucht und Vermahnung zum HErrn“. Zu der Zucht gehört, daß du ihm [dem Kinde] saget, wie GOtt alle Dinge geschaffen hat, ihm die Sinne, das Leben, die Seele gegeben hat, täglich mit allem Guten dem beisteht, was er geschaffen hat; daß er für das Kind gelitten und Wunder gethan hat, gepredigt, und noch größere Dinge verheißen hat: damit es dafür GOtte dankbar sei, und ihn als seinen Vater erkenne und liebe. Zur Strafe gehört, daß du ihm sagest, mit wie großen Plagen er einst die Egypter geschlagen habe, die Heiden, die Sodomiter, die Kinder Israel, ja, auch alle Menschen in Adam. Desgleichen, wie er noch vor unsern Augen schlägt mit Pestilenz, Galgen, Schwert, Wasser, Feuer, wilden Thieren, Krankheiten, und wie große Strafen er zukünftig droht in der Hölle mit den Teufeln. Und er will, daß man sich vielmehr dessen gebrauche als der Strafe der Menschen, das ist der eigenen Strafe, und zwar nicht ohne Ursache, nämlich damit sie lernen immer aufzublicken zu GOtt, und sich nicht vor Menschen, sondern vor GOtt zu fürchten, damit sie nicht, wenn sie den Ernst (rugas) der Eltern fürchten, auch endlich in dem, was GOttes ist, Menschen fürchten, das heißt, damit sie nicht kleinmüthig werden. Denn so muß man mit den Kindern handeln, nicht daß sie die Eltern fürchten, sondern daß sie wissen, daß sie GOtt beleidigen, wenn sie die Eltern nicht fürchten. Denn so werden sie nicht kleinmüthig werden, weil sie, wo sie gelernt haben, auch ihre eigenen Eltern nicht anders als in der Furcht GOttes zu fürchten, auch hernach in Widerwärtigkeiten und im Wohlgehen nicht aus irgend einer Furcht oder Liebe von GOtt meiden. Daß dies aber ein GOtt angenehmes Opfer sei, wenn man die Kinder so unterweist, ist klar aus dem ersten Buch Mosis [Cap. 18, 17–19.], wo der HErr dem Abraham nicht verbergen konnte, was er thun wollte. Allein hierdurch veranlaßt sagt er: „Denn ich weiß, daß er seine Kinder lehren wird, daß sie den HErrn fürchten sollen.“ Siehe, wie er sagt, daß die Strafe des HErrn über Sodom durch Abraham dazu dienen würde, daß seine Kinder unterwiesen würden und in [rechter] Furcht ständen. In gleicher Weise wird Jonadab, der Vater der Rechabiter, herrlich gelobt und in seinen Kindern [Jer. 35, 19.] gesegnet, weil er sie fromm und gottseelig aufgebracht hatte in der Furcht des HErrn. So sind auch Tobias und das Weib des Tobias, Susanna, auferzogen worden. Dagegen sieht man im ersten Buche Samuelis Cap. 3, 13., was Eli damit verdient hat, daß er seine Kinder vernachlässigte. Desgleichen, wie ist es doch mit dem fünfjährigen Knaben bei St. Gregorius gegangen?\footnote{In der Leipziger Ausgabe, Bd. IV, S. 38, steht (mehr erklärend als übersetzend): „So besteht auch St. Gregorius, daß der böse Feind einem Vater sein Kind, das fünfjährig war, von dem Schooße hinwegführe; denn er hatte es übel gezogen, ließ es schwören und fluchen etc.“} Und wir sind in der That überaus thöricht. Siehe, an dem eigenen Kindern kann man entweder das Leben oder den Tod verdienen, und hat doch nicht Acht darauf. Denn was nützt es dir, so du durch deine Schuld böse Kinder auferziehst, wenn du auch an dir selbst untadelhaft bist? Viele bemühen sich in ganz verkehrter Weise, GOtte einen schweren Dienst zu leisten in solchen Dingen, von denen er nichts wissen will (alienis studiis), mit Fasten, mit Tragen harter Kleider und anderen selbsterwählten Auflagen (affectatis observantiis), und über dies gottlose Verhalten gehen diese blinden Leute nicht weniger [leichtfertig] hinweg, als einst die Juden, welche den Tempel GOttes verließen und auf den Höhen opferten. Zuerst solltest du doch darauf Acht haben, was der HErr von dir verlange, die Ausrichtung welches Amtes er dir aufgelegt hat, wie der Apostel sagt [1 Cor. 7, 20.]: „Ein jeglicher bleibe in dem Beruf, darinnen er berufen ist.“ Denn dafür soll man vor allen
\quellseite{622}
Dingen Sorge tragen und damit umgehen vor allem Ablaß, Gebet, Besuchen von Kirchen, oder irgend welchen Gelübden.

\footnote{Bücher: „Das Folgende hat Luther gepredigt am zweiten Sonntage (Domin. post Octavam) nach Allerheiligen [9. Nov.] im Jahre 1516.“} Es ist gesagt worden, wie sich die Eltern gegen ihre Kinder halten sollen, damit sie des Namens „Eltern“ werth seien und nicht aus den Eltern Verderber der Kinder werden, die zugleich mit den Kindern verloren gehen.\footnote{Im Lateinischen ist hier ein Wortspiel, welches sich im Deutschen nicht wohl wiedergeben läßt: parentes — perdentes — pereuntes.} Nun wollen wir sehen, wer denn diese Verderber und Zugrundrichter der Kinder sind.

Die ersten sind nämlich die, welche ihre Kinder versäumen und sie wissentlich aufwachsen lassen ohne Zucht und Vermahnung zum HErrn; wiewohl sie ihnen durch kein böses Beispiel schaden, so verderben sie sie doch durch Nachgiebigkeit (permissione), da sie dieselben nach dem Fleische allzusehr lieben und verzärteln, indem sie sagen: Es sind Kinder, sie verstehen nicht, was sie thun. Und sie sagen die Wahrheit. Aber auch weder ein Hund, noch ein Pferd, noch ein Maulthier versteht, was es thut. Aber siehe, wie sie gehen, kommen, thun und lassen lernen, was sie nicht verstehen. Aber auch ein Stein und ein Holz versteht nicht, ein wie ungeschicktes Ding es sei zu einem Hause, aber es wird von dem Werkmeister zugerichtet und versteht es alsbald; wie viel mehr ein Mensch? Oder wenn es die Kinder anderer Leute verstehen, warum verstehen es denn deine Kinder nicht? Diese Eltern werden daher die Sünden ihrer Kinder tragen, weil sie dieselben zu den ihrigen machen. Oder hat denn jener fünfjährige Knabe, von dem St. Gregorius sagt, es nicht auch gethan, ohne es zu verstehen?

Die zweiten sind, die ihnen mit unkeuschen Worten und Flüchen, dann auch mit bösen Beispielen und Geberden vorgehen. Es hat Leute gegeben (und wollte GOtt, daß jetzt nirgend mehr solche wären), wie ich weiß, welche von dem Hurenlohn ihrer Töchter lebten, und auf Geheiß der Eltern trieben die Töchter und die Frau Hurerei. Wer kann daran zweifeln, daß es für diese Töchter besser gewesen wäre, wenn sie Räuber zu Eltern gehabt hätten? Dann gibt es auch so thörichte Leute, denen es außerordentlich gefällt, wenn sie kriegerische Söhne haben, die leicht zuschlagen, als ob sie um deswillen des höchsten Lobes werth wären, weil sie niemand fürchten. Der elendliche Lohn für die Thorheit dieser Leute pflegt Trauer und Leid zu sein, welches ihnen zustößt durch die plötzliche Ermordung ihrer Söhne; und mit Recht. Nun sind die Kinder (wie denn die Jugend hitzig ist) geneigt zu böser Lust und zum Zorn: darum müssen sich die Eltern hüten vor [bösen] Exempeln in Worten und Geberden. Denn was sollte der Sohn eines Fluchers und unkeuschen Menschen anders [von ihm] lernen als fluchen und schamlos reden?

Die dritten sind, die ihren Kindern Gelegenheit geben, die Welt lieb zu gewinnen, da sie um nichts Anderes bekümmert sind, als wie sie zierlich gehen, Eigentänze tanzen lernen, sich putzen, den Menschen gefallen, die bösen Künste reizen, der Welt gleichförmig werden. Wo sind zu unserer Zeit die Leute, welche mit so großer Sorgfalt darauf aus sind, daß ihre Kinder GOtt haben, als sie dafür sorgen, daß dieselben Kleider, Pracht, Reichthum, Ehrenstellen etc. haben? Hier treffen wir auf ein ganzes Meer voll großes und kleines Gewürms.

Siehe, welche ungebührliche Kleidung (portenta vestium) diese Zügellosigkeit bei beiden Geschlechtern aufgebracht hat, vornehmlich bei den Weibern! GOtt behüte nur die Männer, daß sie nicht ein leinenes Obergewand (peplum) tragen, weil die Weiber auf den bloßen Haaren Barette tragen wie die Männer. Vielleicht aus Rache, weil die Männer angefangen haben, Haarnetze zu tragen, die den Weibern zukamen, so haben deswegen die Weiber von den Männern die Barette angenommen. Und weshalb meint die Welt, daß dieses wohl anstehe, als weil ihr eine neue Mode (mos) wohlgefällt? Aber bedenke du, was das sei, daß Haupt eines Mannes in ein Weibernetz hüllen, und zugleich das Geheimniß, welches der Apostel [1 Cor. 11, 7.] berührt hat, da er nicht wollte, daß das Haupt des Mannes bedeckt werde.

Aber endlich wollen wir einen Unterschied machen. Entweder geschieht ein solcher Schmuck mit Willen und auf Geheiß der Männer: dann können die Weiber sicherlich entschuldigt werden, da sie es doch wider ihren Willen thun, nur um den Männern zu gefallen; ich sage aber davon, wenn die Männer vernünftig und wohlhabend
\quellseite{623}
sind; wie St. Cäcilia auswendig in einem goldenen Kleide einherging, aber darunter am Fleisch mit einem härenen Hemde angethan war, und gleicherweise Anastasia; desgleichen Esther im alten Bunde (lege). Ja, St. Augustinus strafte einst ernstlich eine gewisse Editia,\footnote{Im Manuscript: Aegyptiam.} daß sie wider den Willen ihres Mannes ihre Frauenkleider abgelegt und ein Wittwengewand angezogen hatte.

Indem wir nun hiervon abbrechen, weil dem Manne in den Dingen Gehorsam geleistet werden muß, welche ohne Sünde geschehen können, wollen wir zusehen, wie es mit den Weibern stehe, welche wider den Willen ihrer Männer diesen Aufstand fordern, und einen Zwang auf ihre Männer ausüben. Darüber höret nicht mich, sondern den ersten Apostel Petrus, welcher zu einem treuen Hirten über alle eingesetzt ist, und 1 Petr. 3, 1–6. sagt: „Desgleichen sollen die Weiber ihren Männern unterthan sein, auf daß auch die, so nicht glauben an das Wort, durch der Weiber Wandel ohne Wort gewonnen werden, wenn sie ansehen euren keuschen Wandel in der Furcht (geistlicher Weise und im Herzen sollen sie unterthan sein, das heißt, in Furcht oder Ehrerbietung unterthan sein, daß sie eine hohe Meinung von ihnen haben). Welcher Schmuck soll nicht auswendig sein mit Haarflechten und Goldumhängen oder Kleideranlegen; sondern der verborgene Mensch des Herzens unverrückt, mit sanftem und stillem Geist, das ist köstlich vor GOtt. Denn also haben sich auch vorzeiten die heiligen Weiber geschmückt, die ihre Hoffnung auf GOtt setzten und ihren Männern unterthan waren; wie die Sara Abraham gehorsam war, und hieß ihn Herr, welcher Töchter ihr worden seid, so ihr wohl thut und nicht so schüchtern seid.“

Diese Lection wird euch gelesen, ihr Weiber, dies Latein wird euch (als den Kindern) aufgegeben; das lernt und habt Acht darauf. Euer Bischof, Papst, Apostel, Vater, Hirte sagt dies und legt es euch zu halten auf. Hernach wollen wir auch hören, was er den Männern auflege.

Und erstlich, damit niemand meine, daß man in diesem Gebote nichts von den Weibern sagen solle, soll er wissen, daß hier geboten wird, alle menschliche Obergewalt zu ehren, wie ich auch zuvor gesagt habe. Der Mann ist aber das Haupt und der Herr des Weibes und ihr Vorgesetzter. Ebenso aber wird hier wiederum alle Obrigkeit (potestas) erinnert, daß sie sich der Ehre würdig erzeigen soll. Deshalb ist, nachdem die Kinder ihre Unterweisung empfangen haben, ja, zugleich mit ihrer Unterweisung, ganz passend und eng damit verbunden (proxime), auch von dem Verhalten der Eheleute unter einander zu reden.

Siehe aber, wie lieblich und bescheiden, wie ruhig und gelind der heilige Apostel im Lehren verfährt. Er ist nicht stürmisch, er donnert nicht, er droht nicht, sondern lehrt sanft und lieblich, so daß schon die angenehme Weise, in der er dies sagt, dazu bewegen sollte, das zu thun, was er lehrt, geschweige denn der Nutzen, den die Dinge haben, welche er lehrt und allein zu lehren beabsichtigt.

Erstlich will er, daß die Weiber ihren Männern unterthan sein sollen. Damit stimmt Paulus Col. 3, 18.: „Ihr Weiber, seid unterthan euren Männern in dem HErrn, wie sich’s gebührt“; und Eph. 5, 22. ff.: „Die Weiber seien unterthan ihren Männern, als dem HErrn. Denn der Mann ist des Weibes Haupt; gleichwie auch Christus das Haupt ist der Gemeinde, und er ist seines Leibes Heiland“, das ist, nicht ein Verderber. „Aber wie die Gemeinde ist Christo unterthan, also auch die Weiber ihren Männern in allen Dingen.“ Worin aber diese Unterthänigkeit bestehe, erklärt er [1 Petr. 3, 2.] im Folgenden genugsam, da er ihren Wandel gegen die Männer in Ehrerbietung und Furcht preist als nützlich zur Seligkeit der Männer.\footnote{Löscher: Hier ist eine Lücke [?], welche aus dem Manuscript so zu ergänzen ist: „Und so erweist die rechte Ehe, geistlich verstanden, in schöner Weise gerade die beiden Dinge, die oben genannt sind, nach dem Gehorsam und nach der Ehrerbietung. St. Bernhard.“} Also, daß sie unterthan seien, das heißt, daß sie den Mann ehrerbietig behandeln und ehren, ihm in allen Dingen gehorchen. Denn deshalb führt er das Exempel der Sara ein, daß sie ihrem Manne unterthan gewesen sei, weil sie ihn Herr genannt habe, und dadurch ihre Furcht und Ehrerbietung kund gethan habe, und daß sie bereit sei, ihm in allen Dingen zu gehorchen und nachzugeben. Da nun Petrus gebietet, den ungläubigen Männern so unterthan zu sein, wie viel mehr soll man dies gegen die gläubigen thun, wie auch Paulus wiederum sagt, Eph. 5, 33.: „Das Weib fürchte ihren Mann“, das ist, sie
\quellseite{624}
verhalte sich ehrerbietig gegen ihn. Hingegen handeln nun die Weiber, welche ihre Männer verachten, habsüchtig und halsstarrig sind, immer das letzte Wort haben wollen, oder ohne Wissen und Willen der Männer handeln mit Geldausgeben \&c. Diese alle sind dieses Gebotes schuldig, weil sie ihrem Oberen nicht unterthan sind. Denn der Apostel sagt 1 Cor. 14, 34., das Gesetz sage, daß die Weiber den Männern unterthan sein sollen. Das Gesetz Mosis hat dies zwar an vielen Orten, aber das erste Gesetz ist 1 Mos. 3, 16. das Gebot GOttes an die Eva: „Dein Wille soll deinem Manne unterworfen sein, und er soll dein Herr sein.“ Hier ist die ausdrückliche Bestimmung, daß der Mann die Herrschaft über das Weib haben solle. Wie es sich nun nicht geziemt für einen Sohn, daß er hadere und sich empöre wider seinen Vater, ein Knecht wider seinen Herrn, das Volk wider seinen Fürsten oder Priester, so geziemt es sich auch für ein Weib nicht, daß sie das thue wider ihren Mann. Und dies merket zuerst, denn es ist etwas Großes und wird sehr selten gehalten. So that aber die heilige Monica, die schwieg still gegen ihren Mann und gab nach; aber zu anderen Weibern, die von ihren Männern geschlagen wurden, sagte sie, daß sie ihrer Stellung eingedenk sein sollten, nicht hochmüthig sein gegen ihre Herren, da sie aus der Verlesung des Ehecontracts wissen müßten, daß sie Dienerinnen (ancillas) geworden seien. Deshalb soll ein Weib entweder unterthan sein, oder nicht heirathen. Wenn sie nicht einen Herrn haben will, so stehe sie davon ab, einen Mann zu nehmen, denn so befiehlt es GOtt und verordnet es durch die Apostel und seine heilige Schrift.

\subsection{Schlußfolgerung}

Das Weib, welches einen wunderlichen (morosum) Mann hat, nehme es nicht in den Sinn, Gleiches mit Gleichem zu vergelten und Gewalt mit Gewalt zu vertreiben, sondern suche durch einen gerade entgegengesetzten Wandel und ruhiges Wesen eine Aenderung hervorzubringen. Denn so lehrt hier der Apostel [1 Petr. 3, 1.]: „Auf daß die Männer durch den guten Wandel der Weiber zum Glauben bekehrt werden.“ Was aber des Glaubens halber gethan wird, das soll man auch wegen irgend eines anderen geringeren Gutes thun. Denn so sagt auch Paulus 1 Cor. 7, 14.: „Denn der ungläubige Mann ist geheiliget durch das gläubige Weib.“ Denn ein Uebel wird durch ein anderes Uebel nicht überwunden, sondern vermehret und gestärket, es wird aber durch Gutes überwunden.

Zweitens. Ihr Schmuck soll nicht auswendig sein mit Haarflechten oder Gold oder Kleidern. Aber dies ist jetzt vornehmlich der Schmuck der christlichen Weiber und heutzutage bis aufs äußerste (ad monstra) gebracht. Davon sagt der Apostel, daß dieser äußerliche Schmuck sich für heidnische Weiber zieme, aber christliche Weiber sollen einen andern Schmuck haben. Wer kann aber schweigen, da der Mißbrauch so hoch gestiegen ist?

Erstens. Ich weiß etliche Weiber (denn nicht alle trachten nach dieser thörichten Eitelkeit), welche gezwungen werden, sich wider ihren Willen mit diesen Kostbarkeiten zu schmücken, und einen großen Abscheu dagegen haben, wie die Königin Esther, die Königin St. Elisabeth; und häufig sind dies gerade solche Weiber, welche im Ueberflusse leben, so daß sie sich wohl schmücken könnten. Denn diejenigen trachten um so mehr darnach, welche es nicht vermögen, wie es jetzt geht bei den Edelleuten und Bürgern: je ärmer sie sind in Vergleich zu ihren Vorfahren, desto verschwenderischer sind sie im Aufwand. Daher wollen wir von diesen sagen.

Wiewohl die heilige Agatha von edler Herkunft war, ging sie nicht allein mit einem Schleier, wie ein Weib, sondern sogar im Sclavenkleide, und da sie deshalb getadelt wurde, rühmte sie sich, daß sie eine Magd Christi sei, wie aus ihrer Legende erhellt.

Und erstlich weiß ich nicht, ob ich die Männer oder die Frauen mehr strafen soll, oder wenigstens die Obrigkeiten der Städte, daß sie den Weibern solches zulassen, wodurch ihr Vermögen mit Wissen und Willen geschädigt wird. Denn warum gibt die Obrigkeit [hier] nicht Gesetze, wie es anderswo geschieht? Oder warum hält der Mann nicht sein Weib in Schranken, da er doch ihr Herr ist? Oder wenn er dieses williglich erträgt, wie kann er sich wundern über die Auflagen der Fürsten? Und sehet wohl zu, wenn ihr euch nicht ändert, so wird die Schabung nicht aufhören, sondern zunehmen. Das will ich euch versichern, daß, wenn diese beendigt ist, eine andere, größere kommen wird, weil ihr GOtt so herausgefordert habt und noch herausfordert. Muß sich eine ehrbare Frau nicht schämen
\quellseite{625}
, daß sie so kostbare Schleier und Hüllen (calyptras) hat, daß ein mäßiger Bürger sein ganzes Hauswesen damit erhalten könnte? Was soll ich aber sagen von den Kleidern und dem überaus kostbaren Mißbrauch der Seide, des Purpurs, des Goldes und des Silbers? Und es ist freilich ein erträglicherer Schade, daß man daraus silberne Hefteln, Ketten und Ringe macht, als daß man aus Gold Hauben, Netze und Franzen macht, aus Seide Gewebe, Stickereien und Damast; denn diese Dinge werden durch den Gebrauch verzehrt, jene bleiben brauchbar, wenngleich sie zur Zeit nicht nützlich sind.

Aber ich weiß nicht, was ich von den neuen Gewinn- oder Schaden-Meistern sagen soll, welche theure Maschinen mit Räderwerk bauen, um Schleier und anderes Leinenzeug zu preßen (plananda), so daß die Schleier auf achtzehn bis zwanzig Gulden zu stehen kommen, die Preßmaschine aber auf vierzig bis fünfzig: konnte der Blitz die für die heilige Catharina bereiteten Räder zerschmettern und kann nicht auch diese ganz eitlen und weltlichen Maschinen treffen?

Aber dies alles sind Zeichen, daß uns der inwendige Schmuck mißfällt, der uns nicht so mißfallen würde, wenn uns der äußerliche Schmuck nicht so gefiele. Leicht verachtet eine ehrbare Frau den äußerlichen Schmuck, wenn sie dem inwendigen Schmuck nachtrachtet; um beiderlei Schmuck kann man sich nicht auf gleiche Weise annehmen.

Daher ist jetzt die Zeit der Synagoge wiedergekommen, welche Jesaias Cap. 3, 16–24. mit diesen Worten beschreibt: „Darum, daß die Töchter Zions stolz sind, und gehen mit aufgerichtetem Halse, mit geschminkten Angesichten, treten einher und schwänzen, und haben köstliche Schuhe an ihren Füßen; so wird der HErr den Scheitel der Töchter Zions kahl machen, und der HErr wird ihr Geschmeide wegnehmen. Zu der Zeit wird der HErr den Schmuck an den köstlichen Schuhen wegnehmen, und die Hefte, die Spangen, die Kettlein, die Armspangen, die Hauben, die Flittern, die Gebräme, die Schnürlein, die Biesemäpfel, die Ohrenspangen, die Ringe, die Haarbänder, die Feierkleider, die Mäntel, die Schleier, die Beutel, die Spiegel, die Köcher, die Borten, die Kittel; und wird Gestank für guten Geruch sein, und ein loses Band für einen Gürtel, und eine Glatze für ein kraus Haar, und für einen weiten Mantel ein enger Sack“ etc. Siehe, wie viele und wie große Pracht der Weiber hat er zusammengebracht; einundvierzig Schmuckgegenstände sind für ein Weib nicht genug! Wer sollte sich nicht scheuen ein Weib zu nehmen? Nennen nicht sowohl die Griechen als auch die Lateiner mit einem sehr guten Worte den Weiberschmuck die Welt?\footnote{κόσμος und mundus heißen sowohl Schmuck als Welt.} Denn diese Sammlung von allerlei Dingen macht eine Art Welt aus. Lies auch den Theophrastus, St. Hieronymus gegen Jovinian, und Plautus in der Aulularia, und du wirst sehen, daß ein Weib, wenn ihr der Wille gelassen wird, eine unersättliche Begierde nach Schmuck hat; deshalb muß dieselbe nicht befriedigt, sondern mit allen Kräften eingeschränkt werden.

Und was thut eine ehrbare Frau oder Jungfrau, daß sie mit so großem Eifer wünscht, den Huren gleichförmig zu werden? Oder was suchst du mit so großem Schmuck? Willst du denn die Herzen fremder Männer an dich locken? Oder willst du kund geben, daß du die Keuschheit nicht gern hast? Wünschest du etwa eine Hure zu werden? Warum suchst du einem anderen zu gefallen als deinem Manne? Wer sollte diesen Wunsch, zu gefallen, nicht verdächtig halten? Ein Weib gefällt genug und übergenug, wenn sie so geschmückt ist, daß sie ihrem Manne gefällt: denn ihm allein soll sie gefallen; ihm ist sie gegeben und unterworfen.

Ebenso auch die Jungfrauen. Warum trachten sie in thörichter Weise darnach, den jungen Männern zu gefallen? Denkst du nicht, daß er dadurch abgeschreckt wird, dich zu heirathen, wenn er daran denkt, daß er gezwungen sein sollte, dich, wenn er dich geheirathet hat, mit so großen Kosten zu putzen? Willst du einen jungen Mann durch Liebe fesseln, so höre einen sehr guten Rath: Sei züchtig, wenig geputzt, rede wenige Worte und starre ihm nicht etwa ins Gesicht. Die größte Zierde einer Jungfrau und eines Weibes ist eine züchtige Zurückhaltung, denn diese reizt und fesselt die Herzen der Männer mehr als aller Schmuck. Außerdem macht ein solches Verhalten die Ehe fest, während die fleischliche Liebe, welche durch den Schmuck erweckt worden ist, bald einen Ueberdruß nach der Hochzeit erzeugt, weil die Ehe sich nicht auf gute Sitten gründet, sondern auf den eitlen Schmuck. Thue so, und du wirst (auch mit Gottes Segen)
\quellseite{626}
eher einen Mann bekommen als die leichtfertigsten Verschmenderinnen der Güter, die den Wölfen und Huren ganz ähnlich sind, von denen der 144. Psalm, V. 12., sagt: „Ihre Töchter sind wie die ausgehauenen Erker, gleichwie die Paläste“, das heißt, vor anderen herrlich geziert, wie ein Palast vor anderen Häusern hervorragt, geziert ist und gesehen wird etc.

\footnote{Löscher: „Das Folgende hat Luther am Tage des heiligen Martin [den 11. November] im Jahre 1516 gepredigt.“ Das Exordium findet sich Walch, St. Louiser Ausgabe, Bd. XII, 1774.} Es ist gesagt worden, wie ein Weib sich gegen den Mann verhalten solle, damit sie ihn nach diesem Gebote als ihren Oberen ehre, nämlich durch zwiefachen Dienst (cultu), nämlich des Gehorsams und der Ehrerbietung, daß sie ein bereites Herz habe, sich nach seinem Willen zu richten, sodann ihn auch von Herzen ehre, indem sie hoch von ihm hält. Wenn sie aber fragt, wie sie einen Mann in Ehren halten könne, der ungestalt ist, unmürrisch (morosus = dem man nichts recht machen kann) und ganz und gar auch nicht einmal der Liebe werth, geschweige denn der Ehre, so antworte ich, wie ich von den Eltern gesagt habe, daß man GOtt in ihm ansehen soll, weil es Gott gefallen hat, daß der Mann sein (ejus = GOttes) sei und seine (ejus) Stelle vertrete, wie der Apostel sagt [1 Cor. 11, 7.]: „Der Mann ist GOttes Bild und GOttes Ehre, das Weib aber ist des Mannes Ehre.“ Deshalb muß er ihr nicht mißfallen, weil sie sieht, daß er Gott gefallen habe und gefalle. Es ist eine große Würde des Mannes, daß er Theil hat an GOttes Namen und Amt; darum soll er von dem Weibe geehrt werden, denn GOtt ist der Mann der ganzen Kirche und der Seele. Dieses überaus große Geheimniß ist auch die Ursache, weshalb das Weib den Mann ehren, fürchten und hören soll. Wende also dein Auge vom Fleische und erkenne und diene (cole) GOtt in deinem Manne, dann wird es dir auch leicht werden, ihn zu ehren.

Aber dagegen muß jetzt davon gesagt werden, daß die Männer sich auch des Namens eines Mannes und eines Oberen würdig erweisen sollen. St. Petrus fährt nun an der oben erwähnten Stelle [1 Petr. 3, 7.] fort: „Desgleichen, ihr Männer, wohnet bei ihnen mit Vernunft, und gebet dem weiblichen, als dem schwächsten Werkzeuge, seine Ehre, als auch Miterben der Gnade des Lebens, auf daß euer Gebet nicht verhindert werde.“

Erstlich lehrt er die Männer, daß sie nicht bloß bei ihren Weibern wohnen sollen, sondern mit Vernunft, das heißt, nicht nach der Leidenschaft und nur nach der Sinnlichkeit der Wohllust halben, indem sie keinen Unterschied machen zwischen einem Eheweibe und einer Hure; denn so wohnen auch die Thiere und die Heiden bei einander. Er unterweist deshalb das Auge des Mannes, wie es sein Weib ansehen soll; nämlich nicht mit dem Auge der bösen Lust, sondern der Vernunft und der Erkenntniß.\footnote{Hier sind wir der Lesart des Manuscripts gefolgt: cognitionis statt cogitationis in den Ausgaben.} Durch diese Vernunft soll der Mann an seinem Weibe nach GOttes Geheiß zwei Dinge wissen und erkennen. Erstens sagt er: „als dem schwächsten Werkzeuge“ (vasculo), zweitens: „als auch Miterben“ sollen die Weiber von ihnen gehalten werden. Und so erklärt er sich selbst, was er unter „dem Wohnen bei ihnen mit Vernunft“ verstanden wissen wolle, nämlich, daß der Mann das Weib erkenne als ein gar schwach Gefäß, und sie dennoch als eine Miterbin der Gnade in Ehren halte, und nicht als ob er sie habe zum Büßen seiner Lust (in passione desiderii). Dies wollen wir weiter erklären.

Das sollt ihr Männer wissen, daß das Weib ein schwächeres Gefäß (vasculum) ist als der Mann, darum sollt ihr nach dieser Erkenntniß (scientiam) bei ihnen wohnen. „Gefäß“ nennt er sie nach dem Leibe und nach der Seele. Denn auch der Mann ist nach derselben Weise ein Gefäß, aber ein stärkeres. Denn wie das Weib einen zarteren Leib hat, so auch ein zarteres Gemüth; darum soll ihm von den stärkeren Männern Ehre erzeigt werden. „Gefäß“ werden wir alle genannt, weil die Seele im Leibe als in einem Gefäße weilt.

Es ist aber diese Ehre des Werkzeuges (vasculi) die Ehre, die dem Leibe des Weibes gegeben wird, welche in nichts Anderem besteht, als daß man sein Weib mit einer keuschen Liebe lieb habe, wie es 1 Thess. 4, 4. 5. heißt: „Daß ein jeglicher unter euch wisse, sein Faß zu behalten in Heiligung und Ehren, nicht in der Lustseuche (passione desiderii), wie die Heiden, die von GOtt nichts wissen“, das heißt, man soll sich von seinem Weibe entziehen und der
\quellseite{627}
bösen Lust nicht genugthun. Denn die Heiligung ist diese Keuschheit.

Aber indem der Apostel Petrus dies lehrt, hat er in verdeckter Weise den Männern ihre Weiber ans Herz legen (commendare) wollen, daß sie Geduld mit ihnen haben sollen. Gleichsam mit einem ganz kurzen Wörtlein (epitheto) nimmt er die Geduld der Männer in Anspruch (tangit), indem er sagt: „dem weiblichen, als dem schwächsten Werkzeuge“,\footnote{„In der Erlanger Ausgabe [und nur in dieser] ist hier der Text so wiedergegeben, wie ihn Luther im Druckfehlerverzeichniß der ersten Ausgabe ausdrücklich verworfen hat.“ (Weim. Ausg.)} als ob er sagen wollte: Sie ist ein Weib, sie ist schwach, sie ist ein Werkzeug, indem sie eine gar edle Seele hat. Da ihr dies nun wisset, so gebet ihnen Ehre um ihretwillen, sodann auch um euretwillen, damit euer Gebet nicht verhindert werde. Denn wen läßt das Trachten (passiones) nach böser Lust zum Gebet kommen, ja, auch nur irgend etwas Ehrbares vornehmen? Denn das Gebet erfordert aufs stärkste, daß das Herz sich wende von den Lüsten des Fleisches. Aber da das Weib ein gar schwaches Werkzeug ist, so muß man sie auch tragen und ihre Schuld übersehen und sie nicht immer scharf strafen wegen geringfügiger Dinge, wie denn etliche Männer, die allzu ungestüm sind, bei der kleinsten Ursache ein Trauerspiel (tragediam) aufführen. Wie der Psalmist [Bl. 130, 3.] zum HErrn sagt: „So du willst, HErr, Sünde zurechnen, HErr, wer wird bestehen?“ So ist es unmöglich, daß in einem Hauswesen Friede gehalten werde zwischen Mann und Weib, wenn einer dem andern nichts zugute hält und übersieht, sondern alles bis auf das Kleinste beachtet. Denn wer gibt nicht bisweilen Anstoß? Deshalb muß man vieles übersehen, sehr vieles nicht beachten, damit ein friedliches Verhältnis bestehen könne. Darum soll das den Mann dazu bewegen, sein Weib feind zu erhalten, weil sie das schwächere Werkzeug ist, nämlich wegen des zarteren Leibes und des weicheren Gemüthes; sodann aber viel mehr, weil sie die Mithelferin einer so großen Gnade ist; drittens am allermeisten, damit nicht auch die Seligkeit der Männer selbst verhindert werde. Siehe, welch eine kräftige und gewaltige Vermahnung!

In ähnlicher Weise sagt der Apostel Paulus Eph. 5, 25. 33.: „Ihr Männer, liebet eure Weiber, gleichwie Christus auch geliebet hat die Gemeine“, und: „Ein jeglicher habe lieb sein Weib, als sich selbst; das Weib aber fürchte den Mann“, das heißt, sie soll ihn ehren und eine ehrerbietige Scheu vor ihm haben (reverentia). Gar schön hat er unterschieden: „der Mann liebe“, „das Weib fürchte“, beides in keuscher Weise; die keusche Liebe soll sich herablassen, die keusche Furcht soll sich aufschwingen; diese ist der Untergebene dem Oberen schuldig, jene der Obere dem Untergebenen. Und Col. 3, 19. sagt er: „Ihr Männer, liebet eure Weiber und seid nicht bitter gegen sie“, das heißt, wie ich oben gesagt habe, sie sollen nicht allzusehr hinter ihnen her sein und nicht jede Gelegenheit wahrnehmen, um sie zu schelten.

Wie hat denn Christus die Gemeine geliebt? Schlechterdings geistlicher Weise. Wiewohl er an ihr nichts finden konnte, das er hätte lieben können, sondern nur was er an ihr hassen mußte, so hat er sich dennoch selbst für sie dargegeben, und hat sie geheiligt etc. So muß ein Mann, wenn er an seinem Weibe etwas findet, was nicht gar liebenswürdig ist, sich selbst ändern und sich dargeben, das heißt, er lasse seine Wünsche fahren, die nicht erfüllt werden wegen des Fehls seines Weibes, so daß er in solcher Weise nur das an ihr ansehe, was liebenswerth ist, und sie ihm selbst so reinige, daß er sie sich darstelle ohne einen Makel. Dies geschieht nicht, wie ich gesagt habe, wenn er nicht das Böse mit Gutem überwindet, und lernt, auch das zu lieben, was nicht liebenswerth ist, so daß dies liebenswerth wird, indem dasselbe anders wird, oder er selbst sich ändert. Denn da alle Dinge Mitteldinge (indifferentes) sind, so können sie sowohl geliebt als auch gehaßt werden, je nachdem sich das Gemüth ändert; also, die Ehre der Welt, wiewohl sie dieselbe bleibt, wird von demselben Menschen geliebt und gehaßt, je nachdem er sich von ihr abwendet oder sich ihr zuwendet. Und dies ist die beste Art der Aenderung, wo nicht die Sache und der Gegenstand, sondern die Willensrichtung (potentia) geändert wird, denn das ist etwas Beständiges. Bei den Christen findet sich diese Aenderung, und zur Rechten des Allerhöchsten, aber bei weltlichen Leuten ist allein das Suchen nach Aenderung der Gegenstände.

Doch hiervon ist dies genug. Jetzt von dem Hausgesinde.

„Du sollst deinen Vater und deine Mutter ehren.“ Daß unter dem Vater auch der Haus-
\quellseite{628}
Herr verstanden werde, sieht man aus 2 Kön. 5, 13., wo die Knechte Naamans ihn Vater nennen. Daher kommt nun der Name Hausvater. Darum müssen die Knechte ihre Herren ehren nach derselben Weise, wie sie oben angezeigt worden ist, nämlich mit geistlicher Ehre inwendig [im Herzen], bereit, ihnen zu gehorchen und hoch von ihnen zu halten um des HErrn willen, dessen Namen und Amt sie in ihnen leuchten sehen, dessen ausgesprochener Wille dies auch ist, wie es 1 Petr. 2, 18. heißt: „Ihr Knechte, seid unterthan (siehe, das ist der Gehorsam) mit aller Furcht (siehe, das ist das ehrerbietige Verhalten) den Herren, nicht allein den gütigen und gelinden, sondern auch den wunderlichen“, das ist, denen, mit welchen schwer umzugehen ist, schlechten, mürrischen, unbilligen, harten etc. Wiewohl der Apostel hier von leibeigenen Knechten redet, nicht von gemietheten, so geht sie doch dasselbe Gebot an. Darum sind die Knechte, welche die Güter ihrer Herren mit Vorsatz und Willen veruntreuen, nicht anders anzusehen als Diebe, und werden angehalten zur Wiedererstattung, denn sie bekommen ihren Lohn, damit sie treu seien. Aber nun hat die ganze Welt große Klage über das Gesinde. Darum sollen sie dies Gebot mit allem Fleiß merken; denn das heißt nicht die Herren ehren, wenn man sie Herren nennt, sondern wenn man ihnen von Herzen gehorsam ist und sie fürchtet.

Und gar fein tröstet der Apostel mit diesem Worte das Gesinde, welches sich über seine Herren beklagt, daß sie böse und wunderlich sind, das heißt, daß mit ihnen nicht auszukommen ist (difficiles), indem er sagt [1 Petr. 2, 19. 20.]: „Denn das ist Gnade, so jemand um des Gewissens willen zu GOtt das Uebel verträgt und leidet das Unrecht. Denn was ist das für ein Ruhm, so ihr um Missethat willen Streiche leidet?“ Und Paulus sagt Col. 3, 22–24.: „Ihr Knechte, seid gehorsam in allen Dingen euren leiblichen Herren, nicht mit Dienst vor Augen, als den Menschen zu gefallen, sondern mit Einfältigkeit des Herzens und mit Gottesfurcht. Alles, was ihr thut, das thut von Herzen, als dem HErrn, und nicht den Menschen; und wisset, daß ihr von dem HErrn empfangen werdet die Vergeltung des Erbes; denn ihr dienet dem HErrn Christo.“ Siehe, wie schön zieht der Apostel dieses Gebot auf den geistlichen Verstand, daß die Herren von Herzen geehrt werden sollen, weil sie nicht den Herren, sondern dem HErrn dienen in ihnen. Wenn ich dies so dürr heraus sagte, so würde es niemand glauben. Aber es folgt auch der Trost für die Knechte, wie bei Petrus, indem er sagt [Col. 3, 25.]: „Wer aber Unrecht thut, der wird empfangen, was er Unrecht gethan hat (siehe, wie er sie zur Geduld zu bewegen sucht), und gilt vor GOtt kein Ansehen der Person.“ Das heißt, lasset euch nicht bekümmern, daß jene Herren sind, ihr aber Knechte; das ist einerlei bei GOtt; wer Gutes oder Böses thut, der wird seinen Lohn empfangen. Das selbige sagt er Eph. 6, 5–8.: „Ihr Knechte, seid gehorsam euren leiblichen Herren mit Furcht und Zittern, in Einfältigkeit eures Herzens, als Christo; nicht mit Dienst allein vor Augen, als den Menschen zu gefallen, sondern als die Knechte Christi (das heißt, weil ihr so\footnote{sie von uns aufgenommen nach dem Manuscript. Daß es zu setzen ist, erhellt aus dem Parallelismus mit dem in der folgenden Parenthese Gesagten.} Christo dient), daß ihr solchen Willen GOttes thut von Herzen, mit gutem Willen. Lasset euch dünken, daß ihr dem HErrn dienet, und nicht den Menschen (denn darin thut ihr nicht den Willen der Menschen, sondern den Willen GOttes und dienet Gott); und wisset, was ein jeglicher Gutes thun wird, das wird er von dem Herrn empfangen, er sei ein Knecht oder ein Freier.“ Und Tit. 2, 9. 10.: „Den Knechten [stelle dich selbst zum Vorbilde], daß sie ihren Herren unterthänig seien, in allen Dingen zu Gefallen thun, nicht widerbellen, nicht veruntreuen, sondern alle gute Treue erzeigen, auf daß sie die Lehre GOttes, unseres Heilandes, zieren in allen Stücken.“

Dagegen sollen die Herren sich auch nicht als Tyrannen erzeigen, weil es, wie oben gesagt ist, unmöglich ist, daß ein Knecht nicht auch einmal etwas unterlasse oder begehe; man muß an einem guten Knechte vieles übersehen. Denn wenn sie Gehorsam und Ehrerbietung erzeigen sollen, so müssen schlechterdings die Herren auch Gütigkeit und Wohlwollen beweisen. Daher heißt es Eph. 6, 9.: „Und wisset, daß auch euer HErr im Himmel ist, und ist bei ihm kein Ansehen der Person.“ Denn wer kann einem solchen Menschen dienen, von dem man nichts hört als schelten und schimpfen, und der immer poltert?

Hiezu gehören auch noch andere Lohnarbeiter
\quellseite{629}
und Handwerker, als Handarbeiter, Maurer, Zimmerleute und andere derartige Leute, über welche und über deren Untreue überall große Klage ist; um theuern Lohn werden sie gedungen, aber sie sind faul bei der Arbeit, und größtentheils betrügen sie die Herren, denen sie dienen, und, leider, machen sie sich kein Gewissen darüber, als ob sie nicht öffentlich Räuber wären und Diebe fremden Gutes, welches sie unrechtmäßiger Weise empfangen, während sie doch nicht genug Arbeit dafür leisten.

Endlich gehören hieher auch die Bischöfe und die Fürsten, die Pfarrer auf den Dörfern (plebani) und die Rathsherren [in den Städten], überhaupt die Obrigkeit in geistlichem und weltlichem Regiment (utrumque dominium), welche man gar sehr fürchten soll, weil sie das Schwert des HErrn tragen, denen zu widersprechen sich jedermann hüten soll, weil das Gesetz sagt [2 Mos. 22, 28.]: „Den Obersten in deinem Volk sollst du nicht lästern“, und wiederum: „Den Göttern sollst du nicht fluchen.“ Und, leider, man ist nirgends mehr bei der Hand mit herabsetzenden Reden als bei diesen; da ist keine Furcht vor GOtt, welcher in ihnen das Regiment ausübt [Röm. 13, 1.]: „Denn alle Gewalt ist von GOtt“, und wenn sie auch schaben und die Bauern oder ihre Bürger oder die Priester bedrücken, so muß man sie doch tragen, nicht anders als die Hand des HErrn, welche uns wegen unserer Sünde und Fehle drückt. Weil diese Leute wegen der [geringen] Gewalt, die sie haben, mit Werken nicht schaden können, so schaden sie mit Worten; darum gebietet das Gesetz, daß man nicht verkleinerlich von ihnen reden noch sie lästern soll, denn das hieße seinen Vorgesetzten richten, was durchaus der Ordnung widerstreitet. Aber je mehr ihr [der Regenten] Leben klar zu Tage liegt, desto mehr wird es gerichtet und getadelt, aber nicht in christlicher Weise (catholice), ja, auch im Widerspruch zu der rechten Lehre ([minus] orthodoxe).

\subsection{Zum Beschluß}

Die Juden haben dies Gebot aufgehoben, da sie, um ihrem Geize zu fröhnen, lehrten [Matth. 15, 5.]: „Das Opfer (munus), das von mir dargebracht wird, ist dir nütze“, das heißt, wenn ich die Gabe für dich Gotte opfere, so ist sie dir nützer, als wenn ich sie dir gäbe. Doch auch jetzt herrscht jene Thorheit überall, indem derselbe Geiz der Priester im Schwange geht, so daß die Leute mehr an die todten Steine und Holz hängen und mit willigerem Herzen, als sie den eigenen Eltern geben und den lebendigen Nächsten. Warum (wie ich öfters gesagt habe), damit wir nicht beschuldiget werden mögen, als hätten wir in ganz nichtiger jüdischer Heuchelei des Gesetzes GOttes vergessen, so wollen wir zuerst das Gebot GOttes ehren, und in solcher Weise ehren, daß wir es geistlich vollbringen und verstehen, das heißt, daß wir vor allen Dingen mit diesen äußerlichen Gütern, sei es mit dem Leibe oder der Habe (denn GOtt bedarf derselben nicht, hat auch in den ersten Geboten nicht ein Wort davon geredet, daß man sie ihm darbringen solle, wie er bei Jeremia Cap. 7, 22. sagt), den Eltern dienen, wenn sie arm sind, darnach den Verwandten, sodann auch den Armen und Dürftigen; endlich sollen wir allen Nächsten und Gläubigen, die es vonnöthen haben, leihen oder umsonst geben, aber uns nicht so verhalten, daß wir an diesen vorübergehen (wie jetzt dies unsinnige Verhalten gang und gäbe ist), und Kirchen und in denselben Bilder, Decken, Geräthe aufs kostbarste schmücken. Dies sollte das letzte sein und mit Maßen geschehen, weil es nicht nothwendig ist. Denn an Nothwendigem leiden wir Mangel, weil wir überflüssige Dinge beschaffen. Denn es ist nicht nothwendig, daß wir auf Ueberflüssiges bedacht seien, da so viele Dinge zur Seligkeit nothwendig sind, daß niemand im Stande ist, sie zu vollbringen. Sonst, wenn diese Dinge nothwendig wären zur Seligkeit, wie sollten die Armen selig werden, welche dieselben nicht haben?

Aber jetzt in dieser ganz elenden Zeit besteht man so auf diesen Sachen, als ob sie nothwendig wären und allein zum Gottesdienst gehörten, da sie doch den Menschen zum Troste gegeben sind, und GOtt ihrer nicht bedarf, sondern der Mensch: dennoch werden sie in wunderbarer Blindheit nicht den Menschen gegeben, sondern Gotte. Und es gibt Leute, welche sagen: Wir wissen, daß GOtt dieser Dinge nicht bedarf; aber sie antworten nicht, wenn du zu ihnen sagst: Warum gibst du denn GOtt diese Dinge, der ihrer nicht begehrt, und entziehst sie dem Bruder, der ihrer bedarf, wider den Willen GOttes, der da befiehlt [dem Bruder zu helfen]? Und das sei genug von diesem Gebote.

\quellseite{630}
\section{Das fünfte Gebot}

\begin{schriftzitat}
Du sollst nicht tödten.
\end{schriftzitat}

\footnote{Löscher: „Das Folgende hat Luther am 24. Sonntage nach Trinitatis 1516 gepredigt.“ Diese Angabe wird nicht richtig sein. Der 24. Sonntag nach Trinitatis ist der 2. November. Nun steht aber dies Stück zwischen den Predigten vom 11. und vom 30. November; daher ist es wohl auf den 26. Sonntag nach Trinitatis, den 16. November, zu setzen. (Weim. Ausg.)} Es ist von den Geboten gesagt worden, welche sich auf GOtt beziehen, sowohl an sich selbst, als auch auf die, welche an seiner Statt sind. Nun beziehen sich die sechs folgenden auf den Nächsten. Und siehe die bewunderungswürdige und passende Ordnung an. Denn das Verboten beginnt mit dem Größten und schreitet fort bis zum Kleinsten. Denn der größte Schade ist, wenn ein Mensch getödtet wird; darnach ist der nächste, wenn ihm sein Ehegemahl entehret wird (violatio); der dritte ist, wenn ihm seine Güter genommen werden. Welche nun in diesen Dingen nicht Schaden thun können, schaden wenigstens mit der Zunge; darum ist das vierte das Schädethun an seinem guten Namen. Wenn sie sich nun in allen diesen Dingen nichts zu Schulden kommen lassen, so treten sie wenigstens im Herzen dem Nächsten zu nahe, indem sie begehren, was sein ist. Darin besteht auch eigentlich der Neid; diese wollen wir jetzt vor uns nehmen. Ein solcher Fortschritt ist auch in den ersten vier Geboten. Denn das Größte ist, wenn man GOtt im Herzen, mit dem Munde und mit Werken verwirft; darnach, wenn man wenigstens mit dem Munde dessen Namen unnützlich führt, den man mit dem Herzen und mit dem Munde bekennt; endlich wenn man den, welchen man mit Herz und Mund verehrt, mit dem Werke nicht ehrt, und schließlich, wenn man ihn in seinen Stellvertretern verachtet. Wiewohl dies das geringste ist unter den vier ersten, so ist es doch größer als die anderen, welche darnach folgen, um deswillen, weil [durch] die Uebertretung dieser letzteren nicht geradezu (proprie) GOtt betroffen wird, sondern der Mensch und der Nächste. Dies Gebot aber, gleichwie auch alle anderen, haben die Pharisäer gezogen auf den Buchstaben, als ob niemand wider dasselbe thäte, der nicht mit eigener Hand einen anderen tödtete. So lehrten sie und so lebten sie auch. Daß sie aber so lehrten, das erhellt aus allen Propheten. Daher auch der Psalmist, der da ansieht, daß sie nur die äußerlichen Werke lehrten, nicht aber von Herzensgrunde, von Unwillen darüber bewegt, im 19. Psalm [V. 2.] sagt: „Die Himmel erzählen die Ehre GOttes und die Feste verkündigt seiner Hände Werk“, als ob er sagen wollte: Diese Leute lehren nur Menschenwerke, nicht die Werke GOttes. Daher kommt es, daß sie das Gesetz GOttes verderben, welches ohne Wandel ist und die Seelen erquickt [Ps. 19, 8.], aber jene bekehren nur die Hände. Und der Grund davon ist, daß das Gesetz geistlich ist. Und hier ist zu merken, daß, wenn das Gesetz „geistlich“ genannt wird, dies nicht davon verstanden wird, daß es nach geistlicher Deutung (mystice) zu verstehen sei, wie die Figuren und Geheimnisse (mysteria). Denn „geistlich“ ist etwas Anderes, je nachdem man es von geistlicher Deutung (mysticum) oder von geistlichen Dingen (spirituale) versteht. Es wird aber [das Gesetz] geistlich (spiritualis) genannt, weil es allein durch den Geist erfüllt wird und den Geist erfordert, das heißt, wenn es nicht von Herzen mit fröhlichem Willen erfüllt wird, so wird es nicht erfüllt. Aber ein solcher Geist ist nicht in uns, sondern wird durch die Gnade des Heiligen Geistes gegeben, welche Leute macht, die Lust haben zum Gesetze des HErrn. Wenn man daher irgend ein Gesetz hörte, welches dies oder jenes gebietet, so muß man immer denken und mit dabei verstehen, daß es besteht, mit Lust solches zu thun, das heißt frei, ohne Furcht vor der Strafe und mit fröhlichem Herzen. Da nun das in uns nicht ist, so wird sofort verstanden, daß das Gesetz uns zwingt, zu der Gnade zu gehen, damit es erfüllt werde. So heißt es im ersten Psalm, V. 2.: „Sondern hat Lust zum Gesetz des HErrn.“ Welche aber „geistlich“ so verstehen, daß „nach geistlicher Deutung“ (mysticum) und „geistlich“ (spirituale) dasselbe seien, die müssen auch die heiligen zehn Gebote nicht geistlich (spiritualem) verstehen, sondern nur die Ceremonien des Gesetzes, da kein Gebot mit geistlicher Deutung

\quellseite{631}
(mysticum) in dem Decalog ist. Und dadurch widerstreben sie dem Apostel, welcher Röm. 7, 14. sagt: „Das Gesetz ist geistlich“, und er redet doch von dem Gesetze in den heiligen zehn Geboten, wie dies der Text ausdrücklich anzeigt. Daher unterscheiden sich geistliche Deutung (mysticum) und Geistliches (spirituale), wie Höheres und Niederes. Alles Geistliche (spirituale) ist geheimnißvoll (mysticum), aber nicht umgekehrt. Denn das Geistliche (spirituale) ist nichts Anderes, als was ein lebendiger Wille ist, etwas zu thun oder zu lassen, oder denselben bedeutet. Das Geheimnißvolle (mysticum) aber bedeutet oft äußere Handlungen (operationes), wie das Opfern eines Sperlings die Kasteiung des Fleisches bedeutet. Aber diese gehört noch dem Buchstaben an, wiewohl sie eine heimliche Deutung (mystica) ist, denn sie kann geschehen ohne Geist und Willigkeit, weil öfters auch die Gebote des Evangelii nach dem Buchstaben gehalten werden aus Furcht, und nicht williglich, wiewohl das rechte Leben nach dem Evangelio darin besteht, daß man die Gebote mit Lust und gern thue, und wer anders lebt, der lebt nicht dem Evangelio gemäß, sondern nach dem Buchstaben des Gesetzes.

Deshalb sind diejenigen, welche unter geistlichem Verstande und Lehre des Geistes das Evangelium selbst verstehen, und zwar nach den äußerlichen Werken, fast Pelagianer, denn solche Werke geschehen auch von Gottlosen, aber die Werke des Evangeliums haben ihren Namen nicht von ihrem Aeußeren, sondern von ihrem Inneren. Darum sind sie verborgene Werke; wiewohl sie nach außen leuchten, so ist doch der Wille verborgen. Aus derselben Ursache werden sie auch GOttes Werke genannt, weil sie Werke der Gnade und des Geistes sind, da der Mensch diesen Willen nicht aus sich hat, und um deswillen auch nicht die Werke selbst. Dieser Werke geschieht nun in der Schrift häufig Erwähnung, und die Werke der Menschen werden scharf gestraft, weil sie den rechten Werken zwar ähnlich sind, aber nicht rechte Werke sind, weil sie nicht mit Lust und umsonst gethan werden, sondern immer entweder die Furcht vor der Strafe oder Gesuch des Nutzens vor Augen haben (pro fine). Denn so lebten und lehrten die Juden, wie ich gesagt habe. Denn da sie Christum zu Pilato führten, wollten sie ihn nicht tödten, damit sie nicht sündigten, sondern überantworteten ihn, daß er getödtet würde. So sagte auch Saul [1 Sam. 18, 17.] über David: „Meine Hand soll nicht an ihm sein“, und überantwortete ihn in die Hände der Philister, als ob er darum unschuldig wäre, weil er ihn nicht mit der Hand tödtete, sondern mit dem Herzen.

Darüber sagt Christus Matth. 5, 20.: „Es sei denn eure Gerechtigkeit besser, denn der Schriftgelehrten und Pharisäer, so werdet ihr nicht in das Himmelreich kommen.“ Und warum das? Weil sie das Gesetz nahmen nach dem Worte und nach dem Laut der Buchstaben, das heißt, nicht nach der Absicht des Gesetzgebers. Darum sagt er [V. 21.]: „Ihr habt gehört, daß zu den Alten gesagt ist“, als wollte er sagen: Es ist „gesagt“, aber nicht verstanden; sie haben nur die Worte gehört, aber den Verstand haben sie nicht erkannt; darum haben sie nach dem Buchstaben gelebt, und dies alles darum, weil sie alte, fleischliche Menschen gewesen sind, Adamskinder, Irdische. „Ich aber sage euch“ [V. 22.]. Welchen „euch“? Schlechterdings nicht „den Alten“, sondern den Neuen, Himmlischen, den Christen, den Geistlichen sage ich geistlicher Weise. Darum wollen wir hören, wie Christus dies Gesetz geistlicher Weise auslegt, und wir werden sehen, wie viele Arten von Menschen wider dies Gebot sündigen.

Die ersten sind nun diejenigen, welche ihren Zorn mit der That beweisen, welche tödten nach dem Buchstaben und mit dem Werke, von denen der HErr sagt [Matth. 5, 21. Vulg.]: „Wer aber tödtet, der ist dem Gerichte (judicio) schuldig“, oder vielmehr, der ist des Gerichts schuldig, oder „dem Gerichte“, im Dativ (in dandi casu), dafür, daß er straffällig ist dem Gerichte gegenüber. Dies verstehen alle leicht, nur [muß bemerkt werden,] daß unter diesem Werke auch alle die mitbegriffen werden, welche dazu aufgefordert haben, durch Rath, Hülfe oder auf irgend eine andere Weise Ursache zum Morde gegeben haben, oder Genossen und Theilnehmer gewesen sind.

Die zweiten sind, die ihren Zorn mit Worten beweisen, von denen der HErr sagt [Matth. 5, 22.]: „Wer zu seinem Bruder sagt: Du Narr, der ist des höllischen Feuers schuldig.“ Es werden aber unter dem Worte: „Du Narr“ alle die Flüche, Schmähreden und Lästerungen verstanden, welche aus einem wohlüberlegten Zorn hervorgehen. Denn wiewohl solche Leute nicht

Luther's Werke. Bd. III.
\quellseite{632}
mit dem Werke tödten, so thun sie es doch mit Worten; sie reden solche Dinge über ihn, daß er jedenfalls sterben würde, wenn das geschähe. Denn wie sollte er leben, wenn ihn der Teufel holen würde, wenn ihn die Pestilenz tödtete, wenn ihn das Donnerwetter (fulmen) erschlüge? Ja, bisweilen geht es sogar so weit, daß sie ihm den Tod der Seele an den Hals wünschen, indem sie sagen: „Daß dich das höllische Feuer verzehre!“ Christus nimmt aber ein gelinderes Wort, nämlich: „Du Narr“, damit wir verstehen, was es mit denen auf sich habe, welche härter fluchen, und zwar aus einem heftigen (plena) Zorn, der den Untergang dessen wünscht, dem er flucht.

Die dritten sind, die mit Geberden ihren Zorn auslassen. Von denen sagt der HErr [Matth. 5, 22.]: „Wer zu seinem Bruder sagt: Racha, der ist des Raths schuldig.“ Darunter werden alle anderen Zeichen verstanden, welche aus demselben tödlichen Zorn hervorgehen, als da sind, ihm den Rücken kehren, ihn nicht grüßen oder ihm nicht danken, sich mit freudigen Worten über sein Unglück aussprechen, Bedauern ausdrücken über sein Wohlergehen, woraus der Philosoph [Aristoteles] im zweiten Buch der Ethik die Göttin (virtutem) Nemesis gemacht hat. Denn Christus gebraucht hier ein ganz geringes Zeichen, damit er Anlaß gäbe, über das Vorhergesagte ein rechtes Urtheil zu fällen, ja, weil der Geist die Bosheit der Menschen nur schwach (parce) und mit Zurückhaltung ausdrückt, aber sie gar tief verstanden wissen will. Wenigstens das Wort „Racha“, wenn ich mein Urtheil sagen darf, scheint mir sogar in allen Sprachen gebräuchlich zu sein; denn auch wir, wenn wir entrüstet sind, bringen mit einem gewissen Schnarren der Kehle gleichsam [das Knurren] eines Hundes: „arch“ oder „rach“ hervor. Denn ein solches Schnarren und rauhen Laut der Kehle, meine ich, habe der HErr durch diesen Ausdruck bezeichnen wollen, und es sei das, was andere als einen Ausruf eines Entrüsteten auslegen. Denn, wie bekannt ist, gibt es kein Ausrufungswort, welches die Entrüstung, und zwar heftiger, ausdrückt, als „rach“ oder „rarch“. Wiewohl daher ein solcher Mensch weder mit dem Worte noch mit dem Werke tödtet, so drückt er dies doch, durch Zorn bewegt, durch ein Zeichen aus, durch welches er seinen Wunsch ausdrückt, daß er vernichtet werde (non esse); also tödtet er mit einem Zeichen.

Die vierten sind, die im Herzen (affectu) zürnen; und diese Stufe ist die Ursache, weshalb das, was die [drei] Genannten thun, Sünde ist; wenn das nicht statthätte, so wäre es nicht Sünde, wenn es auch geschähe und geschehen könnte (denn die Bewegung des Zorns ist das Haupt und das Leben des Zorns, der in Worten, Zeichen und Werken zu Tage tritt; ohne diese wäre es nicht Zorn); von diesen sagt der HErr [Matth. 5, 22.]: „Wer mit seinem Bruder zürnet, der ist des Gerichts schuldig.“ Denn wiewohl sie nicht mit Werken, Worten und Zeichen tödten, so thun sie es doch mit dem Herzen. Von diesen sagt Johannes in seiner ersten Epistel (Catholica) [Cap. 3, 15.]: „Wer seinen Bruder hasset, der ist ein Todtschläger.“ Weshalb? Weil er ihm das Leben nicht gönnt, sondern den Tod; darum ist er, was das Herz anbetrifft, ein Todtschläger des Bruders vor Gotte.

Daher ist dies Gebot ein sehr tiefes, und es ist niemand, der es ohne die Gnade erfüllen kann, weil niemand ist, der nicht jemanden habe, dem er zürne, wenn er sich selbst recht prüft. Sonst wäre das Vater-Unser falsch, welches sagt: „Wie wir vergeben unsern Schuldigern.“ Denn ohne die Liebe ist es unmöglich, daß dies Gebot nicht übertreten werde, denn ein Mensch liebt entweder seinen Nächsten, oder er haßt ihn. Doch es sind jetzt etliche Leute aufgestanden, die hierin eine Mittelstellung einnehmen wollen (neutrales), welche sagen, man könne so sein, daß man weder liebe noch hasse; aber ich rathe dir, daß du diesen Leuten nicht trauest.

Darum, wer dies Gebot erfüllen will, der bemühe sich auf alle Weise, jeden Menschen mit einer süßen Zuneigung zu lieben, ja, er bitte den HErrn um die Gnade der Liebe, und schmeichle sich nicht vergeblich damit, daß er niemanden hasse; denn das muß ein sehr sanftmüthiger und demüthiger Mensch sein, der in diesem Gebote nicht sündigt. Denn der Zorn im Geiste (spiritualis), welcher bei dieser Gemüthsbewegung (affectu) verboten ist, ist so tief, daß er, wenn er auch nicht auswendig erscheint in Zeichen oder Wort, inwendig aufs allerstärkste leben mag. Von diesen sagt Jesaias (Cap. 29), 15. [Vulg.]: „Wehe euch, die ihr einen Hinterhalt habt im Herzen (profundi corde) und deren Werke im Finstern sind.“ Es ist aber auch noch ein anderer geistlicher
\quellseite{633}
Zorn, und zwar ein guter, welcher „der Eifer um GOtt“ genannt wird. Denn wenn der HErr nicht auch hätte anzeigen wollen, daß es einen guten Eifer gebe, so hätte er ohne irgend eine nähere Bestimmung (absolute) gesagt: Wer da zürnt; nun aber sagt er: „Wer mit seinem Bruder zürnet“, damit er sich selbst nicht widerspreche, da er sagt [Luc. 14, 26.]: „Wer [nicht]\footnote{non steht in der Vulgata, fehlt aber in allen Ausgaben.} sein eigenes Leben haßt“ in dieser Welt etc. Darum ist Zorn und Haß, den man gegen sich selbst hat, etwas überaus Edeles; davon anderswo.

Es ist also für uns diese Erklärung des HErrn über dieses Gebot genugsam; nur daß es viele bewegt, warum der HErr diese Ordnung inne gehalten habe, daß er den einfachen Zorn dem Gerichte überwiesen habe, „Racha“ dem Rathe, „du Narr“ dem höllischen Feuer. Denn das wird leicht erkannt, daß unter „Gericht“, da das Gesetz sagt: „Wer aber tödtet, der soll des Gerichts schuldig sein“, diese Dinge alle drei im Geiste begriffen werden, wie der HErr selbst das Gesetz auslegt; so daß, gleichwie in dem Worte: „Du sollst nicht tödten“ das Verbieten alles Zorns verfaßt ist, so unter dem Worte „Gericht“ alle Strafe des Zorns. Warum unterscheidet er denn jetzt „Gericht“ von „Rath“ und „Hölle“? Erstlich hält Lyra dafür, daß der Text umgestellt worden sei, da es scheint, daß „Rath“ hätte vor der Strafe stehen sollen. St. Augustinus sagt im ersten Buche „über die Bergpredigt“ (de sermone domini in monte): daß eine Ordnung der Schwere (gravitas) hier angezeigt wird in der Strafe, wie sie auch in der Ordnung der Schuld ist; aber wie das unsichtbarer Weise in den Seelen vor sich gehe, könne man nicht sagen. So läßt er es unerörtert (incognitum), und behauptet nur dies, daß durch jene Stufen eine immer schwerere Strafe ausgedrückt ist. Denn im Gerichte wird noch Gelegenheit zur Vertheidigung gegeben, im Rathe aber wird über den Schuldigen befundenen verhandelt, mit welcher Strafe er belegt werden solle, in der Hölle aber ist eine gewisse und festbeschlossene Verdammniß des Schuldigen. Durch diesen gerichtlichen Gang wird die Schwere der verschiedenen Verbrechen ausgedrückt. Da aber St. Augustinus nicht angegeben hat, wie diese gestraft werden, auch sagt, daß dies nicht angegeben werden könne, so will auch ich davon abstehen, dies ausdrücklich anzuzeigen. Denn wenn der HErr durch diese Stufen verschiedene Strafen hätte bezeichnet wissen wollen, so daß der des Gerichts Schuldige nicht schuldig wäre des Rathes oder der Hölle, so folgt, daß tödlicher Haß des höllischen Feuers nicht schuldig wäre, ebenso wenig Fluchen und Herabziehen (detractio) [des Nächsten], was doch ganz falsch ist. Ja, es würde folgen, daß derjenige, welcher mit dem Wert tödtet, eine schlimmere Strafe verdiene als das höllische Feuer, weil dies das ärgste und schwerste von den dreien ist, während es doch keine schlimmere Strafe gibt.

Deshalb — möge es mir vergönnt sein, meine Gedanken ohne Vermessenheit darzulegen — gibt es für diese vier Stufen allzumal, oder mögen es noch mehr Grade des Todtschlags sein, auch im Werke (was gewißlich wahr ist, weil der einen schwereren Mord verübt; der seinen Vater und seine Mutter oder sein Weib tödtet als einen anderen Menschen), nur Eine Strafe, nämlich die ewige Verdammniß und das höllische Feuer, die aber in verschiedenen Graden aufgelegt wird, was durch diese Stufenunterschiede angezeigt wird. Denn wie der ferner ist von der Strafe und langsamer zur Strafe kommt, der nur in zeitlichen Dingen vor Gericht gestellt wird, so wird der langsamer, das ist, weniger scharf und leichter gestraft in der Hölle, welcher nur zornig gewesen ist. Denn was in der Zeit Langsamkeit oder Schnelligkeit ist, das ist in der Ewigkeit Verschärfung oder Milderung. Ja, die Langsamkeit und Schnelligkeit in leiblichen Dingen entspringt aus der Anspannung und dem Nachlassen der geistlichen Dinge, wie man sieht an feurigen und trägen Seelen,\footnote{Es scheint uns, daß statt animalibus (in allen Ausgaben) gelesen werden sollte: animabus. Denn sogleich folgt anima. Es ist auch schwer zu glauben, daß Luther den „leiblichen Dingen“ (corporalibus) die Thiere als „geistliche Dinge“ (spiritualia) sollte entgegengestellt haben.} von denen die einen Vieles in kurzer Zeit, die andern gar Weniges in langer Zeit zuwege bringen. Denn eine schlaffe Seele thut wenig, während eine angespannte Seele sehr viel thut. So auch, gleichwie der, welcher des Raths schuldig ist, seiner Strafe (malo) näher ist als der, welcher des Gerichts schuldig ist, und dennoch ferner als der, welcher der Hölle schuldig ist, so wird dieser
\quellseite{634}
[des Raths Schuldige] schärfer in der Hölle gestraft als jener [des Gerichts Schuldige]; der aber der Hölle schuldig ist, der hat keinen Aufschub mehr, das heißt, er wird überaus scharf gestraft, und nicht gelinde. Wer aber tödtet, und schlimmer als tödtet, der wird ohne Zweifel am schärfsten gestraft werden. Dies will ich so gesagt haben, daß ich nichts als gewisse Wahrheit behauptet habe.

Aber es sind etliche heimlich hoffärtig, die sich sicher einbilden, sie seien dieses Gebotes nicht schuldig, weil sie nicht tödten, nicht zürnen, sich auch keines Hasses bewußt sind; das ist möglich und kann wahr sein. Aber deshalb müssen sie nicht meinen, daß sie gerechtfertigt seien und dies Gebot bis auf den Grund erschöpft und erfüllt haben. Denn das Gesetz ist ohne Wandel [Ps. 19, 8.] und im Feuer durchläutert und bewähret siebenmal [Ps. 12, 7.]. Deshalb soll sich niemand eher rühmen, er sündige nicht durch Zorn wider seinen Bruder, als bis er sich so sanftmüthig und friedfertig fühlt, daß, wenn ihm auch alles und sein Leben dazu genommen würde, er nicht zürnen wolle gegen den, der es ihm nehme, und keinen Zorn hege, wenn es geschehen wäre. O wie ein großes und tiefes Wort ist dies! Es ist leicht, nicht zu zürnen gegen die Friedfertigen und gegen die, welche dir kein Leid thun; das thun auch die Heiden, das thun auch die Thiere, ja, auch die wilden Thiere. Welch ein herrlicher Ruhm der Gerechtigkeit ist also das, daß du nicht zürnst gegen friedfertige Leute! Da ist es ja so weit mit dir gekommen, daß du eine Sanftmuth hast wie die Schlangen und die Löwen. Aber wenn dir Geld, Haus, Vater, Mutter, Bruder, Schwester, Sohn, Acker, die Glieder des Leibes und selbst das Leben genommen wird, da siehe zu, ob du da könnest nicht allein nicht zürnen, sondern auch fröhlich sein und GOtt preisen und denen, die solches thun, wohlthun und günstig sein. Wenn du das kannst, dann danke GOtt, weil du GOttes Gnade hast. Wenn du dich aber so befindest, daß du es nicht thun kannst, so wisse, daß der alte eingewurzelte Zorn, der geistliche Wille zu tödten, noch in dir verborgen sei. Denn diesen verbietet das Gesetz und bekämpft ihn; darum ist hier die Gnade vonnöthen, welche zu suchen das Gesetz dich zwingt, nachdem du durch dasselbe in solcher Weise die Unmöglichkeit erkannt hast [dir selbst zu helfen].

Nun sagst du mir: Muß denn der Mensch so rein sein, daß er nicht einmal den Willen haben soll, zu zürnen bei allem Uebel, das man ihm zufügt? Ich antworte: Schlechterdings so rein, weil nichts Unreines in das Himmelreich eingehen wird [Offenb. 21, 27.]. Denn so lange dieser Wille in uns bleibt, welcher in Zorn ausbricht, sobald sich ein Anlaß bietet, sind wir unrein. Darum muß der Mensch so rein werden, daß er (wie ich gesagt habe) nicht allein nicht zürne, sondern auch [Matth. 5, 44.] segne, die ihm fluchen, wohlthue denen, die ihn verfolgen, GOtt danke in Widerwärtigkeit, und sich noch mehr Widerwärtigkeit wünsche, so daß er nur die Sünde haßt und nicht die Strafe, wie es heißt [Gesang der drei Männer, V. 70.]: „Licht und Finsterniß, lobet den HErrn“, und [Ps. 34, 2.]: „Ich will den HErrn loben allezeit, sein Lob soll immerdar in meinem Munde sein.“ Darum, daß die [falschen] Heiligen der Juden solche Leute nicht gewesen sind, ist von ihnen im 8. Psalm, V. 3., gesagt worden: „Vertilge den Feind und den Rachgierigen.“ Darum habe ich gesagt, daß das Gesetz geistlich ist; denn wer dies Gebot anders versteht, der versteht es in jüdischer Weise und nach dem Buchstaben, und seine Gerechtigkeit ist nicht besser als die Gerechtigkeit der Pharisäer.

Daher setzt St. Augustinus im ersten Buche „über die Bergpredigt“ sechs Stufen in diesem Gebote. Die erste Stufe ist: Leid thun, wenn man dir auch nichts zu Leide gethan hat. Das ist die höchste Ungerechtigkeit. Die zweite ist: Ein schwereres Leid zufügen, als man dir gethan hat, was dem Gesetz Mosis [2. Buch, 21, 12. f.] nahe kommt. Die dritte ist: Soviel beschädigen, als du beschädigt worden bist. So weit ist das Gesetz Mosis gegangen, welches sagt [2 Mos. 21, 24.]: „Zahn um Zahn, Auge um Auge.“ Denn dieses Gesetz läßt etwas zu (donat), weil nach allen Rechten der Schädiger einer größern Strafe werth ist als derjenige, welcher keinen Schaden gethan hat (innocens). Die vierte ist: Wenn man dich beschädigt hat, und du nicht so viel Schaden thust, als du erlitten hast. Die fünfte ist: daß man überhaupt nicht wieder schädige, wenn man geschädigt worden ist. Dies kommt dem Gebot des HErrn nahe, aber noch nicht völlig. Die sechste ist: daß du, wenn man dir Schaden gethan hat, bereit seiest, dich noch mehr schädigen zu lassen; das heißt
\quellseite{635}
[Matth. 5, 39.]: „den andern Backen auch darbieten“. Siehe also, wie er den Zorn von Grund aus ausgerottet wissen will, so daß der Mensch nicht allein nicht zürne, sondern auch wünscht, daß man ihm Leid thue; denn ein solcher Mensch ist rein. Denkst du, der du dies hörst, nicht an das Wort des Psalms [Ps. 119, 96. Vulg.]: „Dein Gebot erstrecket sich überaus weit“? Wenn nun jemand nicht von einem anderen geschädigt wird, so ist es vonnöthen, daß er sich selbst schädige, und seinen Zorn gegen sich selbst auslasse durch viele Reue. Darum ist nichts Besseres, als daß ein jeglicher, gleichwie geboten ist, den Nächsten zu lieben wie sich selbst, so auch den Nächsten hasse wie sich selbst, damit so alles in die rechte Gestalt gebracht werde.

Deshalb ist die Neigung zum Zorn (vis irascibilis), welche durch dies Gebot getödtet und gänzlich ausgerottet wird (mortificatur), selbst eine Feindin des Kreuzes Christi. Deshalb muß sie auch, weil sie eine Mörderin ist und tödtet, nicht bloß verwundet, sondern gänzlich getödtet und ausgerottet werden, und es muß das Recht der Wiedervergeltung an ihr ausgeübt werden. Denn wer tödtet, der soll getödtet werden. Sie wird aber getödtet durch dieses Gesetz, das heißt, es wird angezeigt, daß sie getödtet werden soll, aber allein die Gnade tödtet sie in der That. Wenn sie aber getödtet ist, so wird der Mensch gütig, geduldig, sanft, auch bereit, irgend welche Feinde zu lieben und ihnen wohlzuthun, fröhlich zu sein in Trübsalen. Dies alles hindert jene ungeduldige und rohe Neigung zum Zorn. Ach wie gar wenige haben Acht darauf; sie sind sicher und verlassen sich auf ihre anderen guten Werke, und wissen nicht, wie gar geistlich und unbefleckt das Gesetz des HErrn ist.

Also derjenige, welcher das Kreuz des HErrn liebt und sich desselben rühmt, der ist’s, der nicht tödtet, nicht zürnt. Denn dies ist der rechte Unterschied zwischen dem guten und dem bösen geistlichen Zorne, daß der böse Zorn nur zürnt wegen des Uebels der Strafe, der gute Zorn aber allein wegen des Uebels der Schuld, und sich deshalb so sehr mißfällt und sich selbst haßt wegen der Sünden, alle anderen Menschen liebt, und niemanden haßt wegen der Güter, die man ihm genommen hat.

Aber, sagst du, wenn dies so wäre, so würde niemand selig, oder doch nur überaus wenige. Ich antworte: Deshalb habe ich gesagt, daß das Gesetz sehr geistlich ist, und deshalb muß man sich bemühen, es in solcher Weise zu erfüllen. Es ist auch kein Wunder, daß wenige selig werden; doch deshalb soll man nicht verzweifeln, denn was uns unmöglich ist, das ist Gott möglich, wie der HErr sagt Luc. 18, 27.: „Was bei den Menschen unmöglich ist, das ist bei GOtt möglich.“ Das ist es, was ich gesagt habe, das Gesetz, recht und geistlich verstanden, schreckt den Menschen und zwingt ihn, nachdem er gedemüthigt worden ist, zu der Gnade seine Zuflucht zu nehmen, und, da er an sich selbst verzweifelt, seine Augen zum Himmel zu erheben und zu sprechen [Ps. 123, 1.]: „Ich hebe meine Augen auf zu dir, der du im Himmel sitzest“ etc., und wiederum [Ps. 121, 1.]: „Ich hebe meine Augen auf zu den Bergen, von welchen mir Hülfe kommt.“ Denn hier gibt es kein Entfliehen; denn dahin muß es mit dem Menschen kommen, daß er so rein sei, mag dies nun im Leben geschehen oder im Tode oder im Fegfeuer. Darum sei nicht hoffärtig [und meine nicht], daß du nicht sündigest, wenn du noch empfindest, daß du zürnest um zeitlicher Dinge willen, weil du noch nicht im Geiste bist, gerecht mit Christo, sondern im Fleische, und mit mosaischen Dingen zu thun hast.

Was ich aber von diesem Gebote gesagt habe, ist von allen in gleicher Weise zu verstehen. Denn ein jegliches Gebot ist geistlich, ohne Wandel, zuverlässig, bewährt, geläutert siebenmal, wiewohl es selten so ausgelegt wird, wie sich der 12. Psalm\footnote{Weimarsche Ausgabe (ebenso wie Walch): „Ps. 12, 7.“ Dies Citat ist falsch, denn es handelt zwar der ganze Psalm von der angegebenen Sache, nur der siebente Vers nicht.} über derartige Leute beklagt. Daher hat St. Hieronymus mit Recht gesagt, daß das [ἐκ] was im Griechischen steht: „Wer mit seinem Bruder zürnet ohne Ursache“ getilgt werden müsse, denn es wird, wie er sagt, jeder Anlaß zum Zorn in diesem Gebote aufgehoben, und jenes Wort ist ein Zusatz, denn in den rechten Büchern findet es sich nicht.

Aber es wird eingewendet: Du redest von den Vollkommenen, es werden aber nicht alle Unvollkommenen verdammt. Ich antworte: Das ist wahr. Es sind Stufen in diesem und jedem anderen Gebote, wie der HErr in diesem Evangelio vier Stufen unterschieden hat, und St. Augustinus, wie wir gesehen haben, sechs. Denn das erste ist, daß du nicht tödtest mit dem Werke,
\quellseite{636}
darnach\footnote{tamen alle Ausgaben, vielleicht aber tum zu lesen. (Weim. Ausg.)} mußt du weiterkommen, daß du es auch nicht thust mit Worten, und ferner nicht mit Zeichen, endlich auch nicht mit dem Herzen. Wenn du nun auch verdienstlich leben (mereri) willst, so mußt du auch gern Widerwärtigkeit herbeiwünschen und dieselbe mit Freuden tragen. Ich weiß, daß man nicht auf einmal vollkommen sein muß, sondern daß man allmälig hinaufsteige. Aber endlich muß man vollkommen werden, und immer darnach streben, daß man weiter komme, und nicht stillstehen, und zufrieden sein mit irgend einer der zuvor genannten Stufen, wie viele thun, welche ärger damit sündigen, daß sie nicht fortschreiten, als andere damit, daß sie nicht anfangen, darum, weil jene stolz sind wegen ihrer angefangenen Gerechtigkeit.

Deshalb kommt es öfters gar dahin, daß die ärgsten Sünder und Huren in tiefer Reue alsbald so vollkommen werden, daß sie auffliegen zur höchsten Stufe, bereit, alles zu verlieren und auszuhalten mit Freuden, während andere, die sich viele Jahre abmühen, noch nicht die zweite Stufe erreichen. Wenn du nun sagen solltest: Das ist nicht zu verwundern, daß große Sünder so tiefe Reue haben [, so sage ich dagegen]: Siehe, siehe, hier höre ich dich, du Pharisäer; du bist nicht ein großer Sünder, du bist nicht wie andere Leute, du hast deine Wege für gut befunden, und deine Blindheit ist allzu tief und weit, als ob du durch diese Hoffahrt nicht ärger wärest als diese alle.

Deshalb muß man immer darauf bedacht sein, daß nicht bloß der Zorn, sondern auch selbst der Zunder des Zorns und der ganze Adam getödtet werde, der Baum mit den Früchten und der Wurzel. „Denn des Menschen Zorn thut nicht, was vor GOtt recht ist“ [Jac. 1, 20.]. Denn die Gebote GOttes sollen nicht bloß so gelernt werden, daß du wissest deine Beichte zu thun, denn dieser Zunder gehört nicht zur Beichte, sondern nur die Werke; er wird auch nicht durch die Beichte weggenommen, weil er ja auch nicht durch die Taufe entfernt wird; sondern [die Gebote sollen gelernt werden], damit du deine Sünden erkennest, und was du thun sollest, und was du von GOtt erbitten sollst, „denn durch das Gesetz kommt Erkenntniß der Sünde“ [Röm. 3, 23.], weil das Gebot GOttes mehr verlangt, als wir vermögen. Aber hier schreien sie: Sollte er also Unmögliches befehlen? Das ist eine Lästerung. Daraus würde folgen, daß er ungerechter Weise verdammte! Antwort: Nein; sondern er gebietet es um deswillen, damit er sich erbarme und die Menschen demüthige, auf daß sie [seine] Gnade suchen. Denn wer die Gebote nur um deswillen lernt, damit er beichten könne, der geht dahin in der Vermessenheit, daß er das thue, was er gelernt hat, und fällt immer ärger, indem er nichts Anderes thut, als daß er sein Gewissen mehrt und groß macht.

\subsection{Folgerungssatz (Corollarium)}

Daß der Mensch nicht vollkommen ist, wie er sein soll, ist eine Sünde, aber einigen wird sie zugerechnet, anderen wird sie nicht zugerechnet.

Denen wird sie zugerechnet, welche stehen, weil solche Leute ihre Sünde nicht hassen, die in ihnen ist, nämlich den Zunder, und darum lieben sie dieselbe und machen ein Bündniß mit den Samaritern. Denen wird sie nicht zugerechnet, welche mit beständigem Eifer dawider streiten; denn wegen dieses eifrigen Bemühens und ihres Fortschreitens wird sie ihnen nicht zugerechnet, wiewohl sie in ihnen ist, weil sie, wiewohl sie in ihnen ist, doch nicht in ihrem Willen ist, sondern wider ihren Willen im Fleische. Darum ist das nicht in ihnen, was in ihnen ist. Daher sagt der Apostel [Röm. 7, 18.]: „Ich weiß, daß in mir, das ist in meinem Fleische, wohnet nichts Gutes“; es wohnt [das Böse] in ihnen, aber es herrscht nicht in ihnen, wie die Jebusei\footnote{In der Weimarschen, bei Löscher und in der Erlanger: Jesubei statt: Jebusei.} unter den Kindern Israel [wohnten], aber nicht herrschten [1 Kön. 9, 21.]. Und Röm. 8, 1.: „So ist nun nichts Verdammliches an denen, die in Christo Jesu sind, die nicht nach dem Fleisch wandeln.“ Warum? Weil die, welche Christo angehören, ihr Fleisch gekreuzigt haben mit seinen Lüsten [Gal. 5, 21.]. Denn sie wandeln beständig wider das Fleisch, nicht nach dem Fleische, das heißt, sie streiten wider den Zunder, welchen jene in tiefem Schlafe (sterntes) in Sicherheit unbeachtet lassen, und sich an ihren guten Werken genügen lassen. Denn was, meinst du, bedeutet das Leiden und der Tod Christi anders, als den Tod des alten Menschen und des
\quellseite{637}
ganzen Adam? Wiewohl derselbe nicht auf einmal getödtet wird, so muß er doch mit so viel Backenstreichen, Faustschlägen, Geißeln, Dornen, ja, auch mit Nägeln gemartert und durchbohrt werden, bis daß er das Haupt neigt und verschiedet. Denn das Haupt ist dieser Zunder, die innerste Wurzel des Zorns und der bösen Lust, welcher nicht getödtet wird, wenn er nicht, durch viele Widerwärtigkeiten gebrochen, endlich ruht, wie ein Todter. Dasselbe ist dadurch bedeutet, daß vorzeiten die Kinder Israel alle gestorben sind in der Wüste außer Josua und Caleb, und nur ein anderes Geschlecht in das Land der Verheißung gekommen ist, welches die zweite Beschneidung im Geiste abbildete; so daß [das erste Geschlecht]\footnote{Ohne eine derartige Ergänzung, wie wir sie gemacht haben, scheint uns dieser Satz völlig unverständlich zu sein.} nicht allein um die zeitlichen, sondern auch um die geistlichen Güter gekommen ist.

Nun wollen wir sehen, was unsere Theologen von diesen Geboten halten. Sie sagen, es sei das nicht die Absicht des Gesetzgebers gewesen, daß sie in der Liebe gehalten werden sollten, es sei denn, daß jemand [durch das Halten derselben] ins Leben eingehen wolle, nach dem Worte [Matth. 19, 17.]: „Willst du zum Leben eingehen, so halte die Gebote.“ Sonst, sagen sie, würde folgen, daß derjenige sündigte, welcher nicht tödtete, und dies Unterlassen [des Tödtens] von ihm ohne Liebe geschähe. Siehe, so reden die, welche zu den Vornehmsten unter den Theologen gehören. Was sagen sie aber anders, als was auch die Juden halten und sagen? nämlich sie machen einen Buchstaben und eine Satzung aus dem allerg geistlichsten Gebote Gottes. Denn sie legen [das Gebot]: „Du sollst nicht tödten“ so aus: Das heißt, du sollst das Werk des Todtschlags nicht vollbringen. Da sich dies nun überaus leicht thun läßt, so fallen sie alsbald zu, und sagen, dies Gebot sei gehalten, wenn es auch etwa nicht in der Liebe geschehen sei; aber dies sei nicht Sünde, auch nicht eine Uebertretung dieses Gebots. Was sollten wir diesen pharisäischen Auslegern gegenüber machen, wenn der HErr nicht selbst dies Gebot ausgelegt hätte? denn sie wären im Stande zu sagen, daß auch der Zorn im Herzen in demselben nicht verboten sei, ausgenommen, wenn ein Mensch, der nicht beleidigt worden sei (non irritatus), aus Bosheit zürnte; doch weiß ich nicht, ob dies nach ihrer Meinung hieher gehöre. Wenn sie nun wüßten, daß dieser zu Grunde liegende (radicalis) boshafte Zorn (malignitas irae) hier von Grund aus verboten wird, so würden sie gewiß nicht leugnen, daß geboten sei, dies Gebot in der Liebe zu erfüllen, es sei denn, sie wollten sagen, es könne jemand aus eigenen Kräften den Zorn in Sanftmuth verwandeln, ohne die Gnade.

Zweitens. Auch das träumen sie, oder wollen dazu zwingen, daß man es träume: daß, „Du sollst nicht tödten“ etc. sich beziehe auf die Zeit, wo man in Ruhe gelassen werde; wenn jemand dann zürnen würde, so würde er sündigen. Das ist natürlich die Meinung dessen, der die Glosse gemacht hat: „ohne Ursache“, als ob er sagen wollte: [Wer mit seinem Bruder zürnet,] wenn er nicht durch ein Unrecht gereizt sein sollte etc.; als ob er nicht sündigte, wenn er wegen einer Beschädigung entrüstet wäre über den, der ihm den Schaden gethan hat. Denn sie sind nicht so weit in die Erkenntniß dieses Gebotes eingedrungen, daß sie sehen sollten, daß, weil sich alles bezieht auf die Zeit der Anfechtung und der Prüfung, der Mensch dann den Geist der Gebote und seine Gleichlichkeit erkenne, da er inne wird, er werde mit Gewalt hingerissen zur Sünde, und nichts Gutes in seinem Fleische findet, damit er seufze nach der Gnade, die ihn gesund mache. Denn hier kann der Mensch ohne die Gnade seinen Zorn und Ungeduld nicht ablegen, wie die Erfahrung zeigt. Sonst würde den Menschen hier gelehrt, daß er seinen anderen Zorn ablegen sollte als den teuflischen Zorn, der darin besteht, daß man, auch wenn man nicht beleidigt ist, zürne und Schaden thue, den doch kaum ein Thier hat, oder doch nur sehr wenige. In gleicher Weise träumen sie von der Handlung (actu), daß man GOtt liebe über alles, nämlich in der Zeit der Ruhe, und denken nicht daran, wie unmöglich dies sei in Anfechtung und Leiden. Dies alles verbergen sie gar fein, wenn sie nur von den Werken der Gebote disputiren nach dem Laut der Silben; denn dann finden sie die Greuel der jüdischen Blindheit, nämlich daß die Gnade nicht nothwendig ist, auch nicht gemeint (intentata) ist in den Geboten, es sei denn bedingungsweise, und zwar nicht wegen des Gebrechens (defectum) bei den Menschen, sondern wegen der Absicht dessen, der die Gebote gegeben hat (praecipientis
\quellseite{638}
); so daß die Gnade allen Menschen verhaßt wird, da vielmehr gleichsam ein neuer Zwang aufgelegt wird, der über das Gesetz hinausgeht, und die Seligkeit unmöglicher macht als das Gesetz selbst, und als ob die Gnade nicht vielmehr denen ein Beistand und Trost sein solle, welche das Gesetz erfüllen sollen. Aber in solcher Weise treiben die Theologen dazu, daß man sich Greuel der Unwissenheit und Hirngespinnste (prodigia) der Finsterniß einbilden muß. Daß daher die Gebote in der Gnade erfüllt werden müssen, das erfordert nicht allein die Absicht des Gebietenden, sondern auch die Noth (indigentia) dessen, der sie thun soll. Denn es ist wahr, daß derjenige sündigt, welcher nicht tödtet, während er außer der Gnade steht, er sündigt aber nicht damit, daß er nicht tödtet nach den Buchstaben, sondern weil er tödtet inwendig im Herzen, im Geiste, da er ohne die Gnade nicht ohne Zorn sein kann und ohne den Zunder der Entrüstung, wie klar ist aus der Erfahrung, wenn Anfechtung da ist, und Anlaß gegeben wird durch Widerwärtigkeiten; denn dieses verborgene Uebel des Zorns verursacht, daß er nicht anders kann als tödten; also ist er dieses Gebotes schuldig nach dem Geiste, wiewohl er ebendesselben Gebotes nicht schuldig ist nach dem Buchstaben; es ist also offenbar, daß sie von dem Gesetze reden nach dem Buchstaben.

Drittens machen sie aus ebenderselben Blindheit die Schlußfolgerung, daß etliche Gebote nur verbieten (negativa), andere nur gebieten (affirmativa). Dies ist wiederum nicht wahr, es sei denn, du behieltest nur die äußere Schale, die Silben. Denn durch dies Gebot „Du sollst nicht tödten“ drückt er das allerstärkste Gebieten (affirmativam) aus, nämlich dies: Du sollst sanftmüthig und von Herzen gütig sein und geduldig und stille und friedfertig. Denn eine Verneinung wird um desswillen in der Schrift gesetzt, weil eine Bejahung nicht so stark ist; also, an der Stelle [Luc. 10, 42.]: „Maria hat das gute Theil erwählet, das soll nicht von ihr genommen werden“, das heißt, in alle Ewigkeit soll sie damit überschüttet werden. Und im 110. Psalm, V. 4.: „Der HErr hat geschworen, und wird ihn nicht gereuen“, das heißt, es wird ihm aufs allerhöchste und ewiglich gefallen. Und das Wort im ersten Psalm [V. 1. Vulg.]: „Wohl dem Manne, der nicht weggegangen ist [in den Rath der Gottlosen]“, das heißt, dieser hat überaus wohl zugenommen unter den Gottseligen. Und so in vielen anderen Dingen. Denn auch der Apostel legt Röm. 7, 19. das Wort [Röm. 7, 7.]: „Laß dich nicht gelüsten“ in bejahender Weise aus, indem er sagt: „Das Gute, das ich will“, das ist das Gegentheil der bösen Lust, nämlich Keuschheit und Liebe, „das thue ich nicht“, denn „thun“ bezeichnet in bejahender Weise (positive). So ist nun dies Gebot zwar verbietend nach den Buchstaben, aber ein überaus starkes Gebieten nach dem Geiste, weil der HErr fordert, daß sie sanftmüthig und friedfertig sein sollen. Jene aber belassen den innerlichen Unflat, den der HErr der Synagoge aufrückt, indem er klaget. 1, 9. sagt: „Ihr Unflat klebet an ihrem Saum“, und 5 Mos. 32, 5. [Vulg.]: „Sie haben an ihm gesündigt, und sind nicht seine Kinder in ihrem Unflat.“ So ist auch das Wort: „Du sollst nicht ehebrechen“ ein Verbot, aber nur nach den Buchstaben; es enthält aber das überaus geistliche Gebot: Du sollst keusch, enthaltsam, mäßig, nüchtern, bescheiden sein, und zwar von Grund deines Herzens, willig und gern. Dazu ist aber schlechterdings die Gnade nöthig, nicht um GOttes willen, der dies bei seinem Gebieten im Auge hat, sondern um unseres fleischlichen Wesens willen, welches durch sein Thun nichts auszurichten vermag. So ist auch das Wort: „Du sollst nicht stehlen“ ein verbietender Buchstabe, aber ein gebietender Geist, nämlich: Du sollst arm sein im Geiste, du sollst dir genügen lassen an dem, was du hast, und bescheiden sein etc. Ja, auch das erste: „Du sollst nicht andere Götter haben“ ist ein Gebot (affirmativum), nämlich: Du sollst den HErrn, deinen GOtt, lieben und ehren, allein (unum), über dich und alle Dinge, sonst wäre er nicht der einige GOtt (unus), wenn du etwas Anderes zugleich mit ihm liebtest. Ebenso auch das zweite: „Du sollst den Namen des HErrn, deines GOttes, nicht unnützlich führen.“ Dies ist der Buchstabe, aber der Geist sagt [Ps. 111, 9.]: „Heilig und hehr ist sein Name“, und wiederum [Joel 3, 5.]: „Wer den Namen des HErrn anrufen wird, der soll errettet werden.“ Es will also [dies Gebot], daß der Name GOttes gebraucht werden soll in stetiger Ehrerbietung, und angerufen werden soll mit Furcht; mehr aber noch soll er gerühmt und gepriesen werden durch ein demüthiges Bekenntniß der eigenen Schande. Und das dritte: „Du sollst den Feiertag heiligen
\quellseite{639}
“, das heißt nach dem Buchstaben: Du sollst am Feiertage kein Werk thun; der Buchstabe ist ein Verbot, während darin die edelsten und größten Werke GOttes geboten werden mit Hören, Lehren und Betrachten des Wortes Gottes inwendig im Geiste. Auswendig ist Ruhe, inwendig ist die größte Thätigkeit GOttes, der da wirkt. Und wenn du diese drei Gebote geistlich auslegst, das heißt, von den innerlichen Werken, so wirst du finden, daß das erste die Liebe gegen GOtt sei, das zweite der Glaube, das dritte die Hoffnung. Denn keinen anderen Gott haben ist GOtt über alle Dinge lieben. Den Namen GOttes nicht unnützlich führen ist an seinen Namen glauben und ihn im Geiste des Glaubens inwendig anrufen und preisen. Denn [Röm. 10, 14. 13.]: „Wie sollen sie anrufen, an den sie nicht glauben? Denn wer den Namen des HErrn wird anrufen, soll selig werden.“ Ferner, eine heilige Feier halten (sanctificare quietem) ist Gotte stille halten (sese passibilem praestare), damit GOtt allein in ihm wirke; hier ist Geduld und Hoffnung vonnöthen, denn hier tritt man in die Finsterniß, wo der Mensch nicht wirkt, sondern wunderbarlich auf dem Wege des Leidens geführt wird. So oft du nun leidest, so oft wirkst du nicht, sondern feierst (quiescis), und GOtt wirkt in dir, aber du weißt nicht was, weil du dich leidend verhältst und ein bloßer Stoff bist [, an dem GOtt arbeitet]. Dies ist es, was gesagt ist [Ps. 46, 11.]: „Seid stille und erkennet, daß ich GOtt bin.“ Hier herrscht allein die Hoffnung in der Liebe durch den Glauben. Darum sagt St. Augustinus mit Recht, daß die Verehrung und der Dienst GOttes bestehe in Glauben, Hoffnung und Liebe. Aber diese sind die allerstärksten Gebote (affirmativissima) und die drei theologischen Tugenden, welche unter dem Buchstaben dieser drei Worte als Verbote (negativa) ausgedrückt sind. Es bleibt allein Ein Gebot (affirmativum): „Du sollst deinen Vater und deine Mutter ehren“, aber wie auch dieses ein geistliches Gebot sei, ist oben gesagt worden. Wiewohl der Buchstabe nur die Ehre in äußerlichen Gebärden (signi) zu lehren scheint, so erfordert das Gebot doch darunter eine herzliche Furcht und einen bereitwilligen Gehorsam, der ohne die Gnade unmöglich statthaben kann, wie man erkennt in der Prüfung der Anfechtung, wo die Gelegenheit an einen jeglichen herantritt, es zu halten, wenn es wider seinen Wunsch geht. Nun bei dem achten: „Du sollst nicht falsch Zeugniß reden“, ist es genugsam bekannt, daß der Buchstabe zwar ein Verbot ist, aber daß es nach dem Geiste ein Gebot ist, daß du von Herzen ein Anwalt (paracletus) deines Nächsten sein sollst mit Entschuldigen, Vertheidigen, und indem du alles thust, daß du wolltest, daß man dir thäte in Gefahr des guten Namens und [in Nöthen] des Gewissens. Von den noch übrigen beiden Geboten ist gleicherweise bekannt, daß „sich nicht gelüsten lassen“ nach dem Buchstaben sei, dem Nächsten von Herzen alle Gute gönnen und wünschen, daß ihm kein Böses widerfahre.

Aus diesem allem ist offenbar, daß die zehn Gebote, wie sie nach dem Buchstaben lauten, richtig im Geiste von Christo und den Aposteln ausgelegt werden, da sie den Glauben lehren, Hoffnung, Liebe, Gehorsam, Ehrerbietung, Demuth, Sanftmuth, Frieden, Geduld, Bescheidenheit, Keuschheit, Armuth, Gütigkeit, Mildigkeit, Fröhlichkeit, Wohlwollen unter einander \&c. Dieses sind zwar Gebote, aber durchaus eine geistliche Predigt der heiligen zehn Gebote.

Nun wollen wir wieder zur Sache kommen. Es beginnt aber [das Gebot] mit der Sanftmuth, und lieblichen Wesen des Geistes, den Menschen zu unterrichten, daß er sich bestreben solle, Gott ähnlich zu sein in seinem Verhalten gegen den Nächsten, gleich wie er wollte, daß GOtt sich gegen ihn verhielte. Denn durch Gütigkeit wird GOtt aufs höchste gepriesen; deshalb, damit du ein sanftmüthiger Mensch seiest, gütig, lieblich, freundlich, milde, umgänglich, nachgiebig, gut, das heißt, Gotte überaus ähnlich, sagt er: „Du sollst nicht tödten“, du sollst nicht ein Todtschläger sein, du sollst nicht roh, scharf, hart, bitter, unverträglich, entrüstet sein. Dies aber wirkt der Geist der Gnade, den zu suchen und zu erbitten dies Gebot zwingt. Ganz dasselbe sagt der HErr Matth. 5, 5.: „Selig sind die Sanftmüthigen, denn sie werden das Erdreich besitzen.“

\subsection{Folgerungssatz}

Vergebens schmeicheln sich diejenigen, welche nicht tödten, und doch nicht von Herzen darnach trachten, daß sie gütig seien, weil sie dadurch nur den äußeren Schein des Gebotes und den todten Buchstaben haben, wodurch gleichsam als durch Schattenbilder der inwendige Schatten
\quellseite{640}
des Behemoth bedeckt wird. Und wir wollen ein Gleichniß geben:

Der Kalk scheint kalt zu sein und ohne Wärme, aber wenn Wasser darüber gegossen wird, braust er in Hitze auf, aber nur durch Oel wird er wieder gelöscht, daß er ganz stille wird. So hat der Mensch, entzündet im Feuerofen der angeerbten (maternae) bösen Lust, die verborgene Wurzel des Zornes, welche hervorbricht, wenn eine Versuchung an ihn herantritt, und wird nicht anders geheilt, als durch das Oel der Gnade.

Ein anderes Gleichniß: Wenn ein Räuber oder ein Mörder oder ein anderer Wegelagerer auch nicht mit der That tödtet oder verwundet, ist denn etwa der Fürst damit zufrieden, wenn er die Waffen nicht niederlegt? denn er wird den Verdacht hegen, daß er morden wird, sobald sich nur eine Gelegenheit bietet. So hat auch der Mensch, wenn er gleich nicht in der That zürnt, dennoch inwendig das Schwert des Zornes, wiewohl es in der Scheide verborgen ist, aber er gefällt GOtte nicht, wenn er das Schwert nicht niederlegt und von sich wirft. Es ist auch nicht genug, daß es in der Scheide ist; denn es ist zu besorgen, es möchte bei künftiger Gelegenheit gebraucht werden. Denn dies ist das Schwert, mit dem der Sohn GOttes getödtet worden ist. Deswegen bist du, auch wenn du nicht tödtest, dennoch vor GOtt ein Todtschläger, weil du in den Waffen erfunden wirst, mit denen sein Sohn ermordet wurde. Denn wer könnte sicher vor dem Angesichte eines Fürsten das Schwert tragen, von dem es feststeht, daß damit der Sohn des Fürsten getödtet worden sei? Deshalb muß der Zorn so weit hinweggethan und aus dem innersten Herzen entfernt werden, daß es nicht möglich scheint, daß er noch wieder hervorbrechen könne.

So auch ein Hund, wenn er auch in der That nicht beißt, ist er etwa um deswillen nicht bissig? Ja, er ist um so ärger, je tückischer er ist. Aber ach, wie wenige haben Acht auf diese Tücken unseres Herzens! Daher sagt der 5. Psalm, V. 10.: „Ihr Inwendiges ist Herzeleid“ (insidiae), wo wir [in der Vulgata] haben: „Ihr Herz ist eitel“ (vanum).

Daher ist die erste Stufe, die denen eigen ist, welche anheben, daß sie sich des Werkes des Todtschlags enthalten, sei es nun durch eigene That, oder durch Bewilligen, oder Rathen, oder Helfen, wie es die falschen Zungen (tercia lingua), die Ohrenbläser, die Anstifter des Mordes zu thun pflegen, daß sie Zorn und Zwietracht anstiften. Zank, Feindschaft, Mord, wie der Apostel (Gal. 5, 20). lehrt. Diese Sünde scheint leicht und gering zu sein, aber wenn man sie recht ansieht, ist sie gar groß, und wenige sind derselben nicht schuldig, weil wenige ihre Zunge so im Zaume halten, daß sie nicht Zwietracht säen und dadurch, so viel an ihnen ist, Todtschlag veranlassen. Daß seine Mutter von diesem großen und häufigen Uebel durch besondere Gnade GOttes frei gewesen sei, sagt Augustinus im 9. Buche seiner „Bekenntnisse“. Und es halte niemand dafür, daß er leicht davon frei sei, es sei denn, daß er friedliebend und ein guter Vermittler sei, der nur das redet, was erzürnte und zwieträchtige Leute versöhnen kann und ihnen gefällt, oder vielmehr, was Einträchtige in Eintracht erhalten kann, dadurch, daß er verschweigt, was böse ist, und das Gute, was jeder an sich hat, bei beiden rühmt. Es liest man ein Exempel von einem klugen Knechte, der als Bote abgeschickt ward von seinem Herrn zu einem Manne, damit er ihm den Befehl überbrächte, aus seinem Hause zu ziehen. Er aber sagte ihm, daß er ihn von seinem Herrn grüßen sollte. Und da er zu seinem Herrn zurückkehrte, sagte er, daß der Mann ausziehen wolle. Endlich wurde der Herr durch wunderliche Schickung (davon ich jetzt nicht weiter sagen will) begütigt, und erkannte die Klugheit des Knechtes. Darum hat man recht gesagt, daß ein Ohrenbläser oder ein Verleumder ein dreifältiger Mörder ist, denn er tödtet drei Leute mit einem Schlage: erstens sich selbst, zweitens den, welchen er in die Ohren bläst, drittens den, welchen er verleumdet, weil der Mund, welcher lügt, die Seele tödtet [Weish. 1, 11]. Aber sie entschuldigen sich und sagen, daß sie die Wahrheit reden; aber es ist eine schlechte Entschuldigung, denn die Wahrheit muß auch in der rechten Weise gesagt werden nach dem Sprüchwort: Was recht ist, mußt du auch recht ausrichten. Denn es wird auch die Wahrheit zur Lüge, wenn sie nicht gesagt wird, wo, wie, wann, zu wem, und so viel gesagt wird, als gesagt werden sollte. Daher sagt auch St. Bernhardus, daß der Verleumder den Teufel in der Zunge hat; der, welcher ihm zuhört, im Ohre. Der zweite, der getödtet wird, ist der, welcher dem Verleumder zuhört, weil er in gleicher Weise
\quellseite{641}
sündigt, außer dem, daß er ihn zum Zorn, das ist, zum Tode veranlaßt. Der dritte [, der getödtet wird,] ist, der dies erleidet, welcher, wenn er dies Böse erfährt, gleicherweise in Zorn, das ist, in den Tod stürzt; denn wenn er nicht dahinfällt, so ist es nicht die Wohltat des Verleumders, sondern GOttes; der Verleumder hat gethan, so viel er vermochte.

Von anderen Uebeln, die ein Verleumder anrichtet, anderswo. Es müßte die Epistel Jacobi gehandelt werden, um klar zu machen, wie weit sich die erste Stufe dieses Gebotes erstrecke, und wie wenige auf derselben stehen. Und kurz wollen wir das thun.

Er sagt im ersten Capitel [Jac. 1, 19. 20.]: „Darum, lieben Brüder, ein jeglicher Mensch sei schnell zu hören, langsam aber zu reden, und langsam zum Zorn. Denn des Menschen Zorn thut nicht, was vor GOtt recht ist.“ Mit diesen Worten reißt er zur Sanftmuth, deshalb folgt [V. 21.]: „Darum so leget ab alle Unsauberkeit und alle Bosheit, und nehmet das Wort an mit Sanftmuth, das in euch gepflanzet ist“ etc., als ob er sagen wollte: Vollbringet ja nicht alle Unsauberkeit und überaus viele Bosheit, das heißt, gebt euch der Unsauberkeit und der Bosheit nicht hin, daß ihr sie alle ausrichtet, sondern es ist genug, daß, wo ihr Widerstand leistet und dieselbe von euch werfet, euch doch noch Unsauberkeit und Bosheit anhängen wird. Er fügt hinzu: „alle“ und „überaus viele“ (abundantiam malitiae in der Vulg.), um anzuzeigen, daß diejenigen, welche dagegen kämpfen, dennoch nicht ohne Unreinigkeit und Bosheit sind, wie viel mehr haben die auch alle Unreinigkeit, welche nicht [dawider] kämpfen, sondern ihr gehorchen, wie der Psalm [Psalm 36, 5. Vulg.] sagt: „Er stehet fest auf allem bösen Wege“, denn es folgt: „aber die Bosheit hat er nicht gehaßt.“ Ferner nennt er die böse Lust eine „Unreinigkeit“ oder einen Unflat, denn sie ist ein Unflat der Seele und des Geistes, welcher die Ursache des Zornes ist unter den Menschen; denn deshalb zürnen sie einer gegen den andern, weil der eine die böse Lust des andern hindert; darum lehrt er, daß wir diese, welche gleichsam das Ganze und das Haupt der Unreinigkeit ist, ablegen sollen. „Bosheit“ aber nennt er den Zorn selbst, welcher auch etwas ist, das im Geiste zu viel und überflüssig ist, wie die ekelhaften Aussonderungen (foeda) des Leibes etwas Ueberflüssiges sind, als ob er sagen wollte: den Mist und die überflüssige Bosheit, das heißt, die Bosheit, welche ein Mist und eine schnöde Aussonderung (obscœnum) der Seele ist, leget ab und heget ihn nicht. Dies geschieht aber, wenn jemand schnell ist zu hören und langsam zu reden“. Denn zornige Leute sind nicht bloß schnell zu reden und langsam zu hören, sondern auch gar nicht im Stande zu schweigen und zu hören; denn wer schnell ist zum Zorn, der ist auch schnell zu reden, und wiederum. Darum ist das beste Mittel, den Zorn zu brechen, daß der Mensch sich Gewalt anthue und gewöhne, erstens, daß er langsam sei zu reden, wenn er erzürnt sein sollte; zweitens, und dies ist ein vollkommenerer Grad, „langsam zum Zorn“, so daß er sich im Herzen vornehme, daß er sich nicht zum Zorn bewegen lassen und endlich einmal in Erfahrung bringen wolle, ob er nicht schweigen könne, wenn er beleidigt wird.

Aber es entschuldigen sich einige damit, daß sie um GOttes und der Gerechtigkeit willen zürnen und reden. Denen tritt er entgegen und spricht [Jac. 1, 20.]: „Des Menschen Zorn thut nicht, was vor GOtt recht ist.“ Es ist unmöglich, daß der ein GOtt wohlgefälliges Werk thue, der erzürnt ist und aus Zorn etwas thut, da GOtt ein GOtt des Friedens und des stillen Wesens (suavitatis) ist, dem kein ungestümes Verhalten (turbulentum) gefällt. Ein merckenswerther Ausspruch, den viele allzusehr außer Acht lassen, welche gegen unrechtes Verhalten anderer ihre Entrüstung an den Tag legen, ihr Unrecht aber beschönend (palpant) und bisweilen sich dessen sogar rühmen. So sagt er auch hernach [Jac. 1; 26.]: „So sich jemand unter euch läßt dünken, er diene GOtt, und hält seine Zunge nicht im Zaum, sondern verführet sein Herz, deß Gottesdienst ist eitel.“ O ein großes Wort! Wie viele machen gar viele Worte, sind aber auch (und von diesen redet er hier eigentlich) bitter, zornmüthig und überaus scharf, und wollen doch für Gottesdiener und Heilige angesehen werden! Es ist also ein und dasselbe, daß der Gottesdienst eines geschwätzigen Menschen eitel ist, und daß der Zorn des Menschen nicht thut, was vor GOtt recht ist. So macht der Zorn, durch die Zunge, den Gottesdienst und alle Gerechtigkeit zunichte, das heißt, von außen ist er [der Gottesdienst] schön, aber inwendig nichtig; darum „verführet er“ auch, nicht andere, sondern
\quellseite{642}
„sein eigen Herz“, indem er meint, er sei ein Gottesdiener, während er es doch nicht ist.

Und im dritten Capitel [V. 1.] sagt er, da er die Gefahren und Uebel der Zunge beschreibt: „Lieben Brüder, unterwinde sich nicht jedermann Lehrer zu sein, und wisset, daß wir desto mehr Urtheil empfangen werden.“ Er sagt: Wolle doch nicht eine große Anzahl [plures in der Vulg.] (das heißt, alle) Lehrer sein, und rasch zu reden und zu lehren, und langsam zu hören, sondern ein jeglicher wünsche, daß er ein Schüler werden könne, und lasse einen andern seinen Lehrer sein, damit nicht ein ganz wüstes, zwieträchtiges Wesen entstehe, so daß da viele seien, welche lehren, und niemand, der da höre, und ein jeglicher wider den andern, so daß sie alles mit Zwietracht erfüllen, wie die Ketzer in der Kirche gethan haben. Lehrer werden ist nicht böse, aber daß viele es werden wollen, das heißt, daß sie verschiedene Meinungen hegen wollen (esse diversi), indem keiner den andern hört, das ist es, daß sie desto mehr Urtheil empfangen werden. Und davor warnt er durch Anzeigung solcher Gefahr, weil wir auch ohne dies größere Urtheil, welches solche Leute empfangen werden, uns schon genug zu Schulden kommen lassen. Es ist nicht unnöthig, Anlaß zur Sünde zu suchen und mehr Urtheil zu empfangen, denn [Jac. 3, 2.]: „Wir fehlen alle mannigfaltiglich.“ Wir alle sind solche Leute, daß wir fehlen, und zwar nicht bloß mit Worten, sondern in vielen Dingen, nämlich mit den Augen, Ohren, Händen und allen Sinnen, sondern auch mit Gedanken. Und diese Fehler sollten uns genug sein, daß wir uns nicht noch mit größeren beladen. Sodann vermahnt der Apostel auch dazu, daß nicht jemand in hochmüthiger Weise sich vermesse, eines Andern Lehrer zu sein, als ob er selbst in keinem Stücke fehle, als ob er sagen wollte: Was tadelt ihr euch unter einander heftig, und warum bemüht sich ein jeglicher, des Andern Meister zu sein? Wisset ihr denn nicht, daß wir alle (ohne Ausnahme) nicht bloß in Einem Stücke, sondern mannigfaltiglich fehlen? Wenn jemand nun auch in wenigen Dingen oder im Werke nicht fehlt, wie wird es aber mit dem Worte stehen? [Jac. 3, 2.] „Wer aber auch in seinem Worte fehlet, der ist ein vollkommener Mann“, als ob er sagen wollte: Der mag von sich rühmen, daß er vollkommen sei, aber auch nicht eher, als bis er auch im Worte keinen Fehl begeht, als wollte er sagen: Solcher Mensch wird unter euch nicht zu finden sein, weil wir alle mannigfaltiglich fehlen, viel mehr auch in Worten. Was seid ihr denn hoffärtig und erhebet euch einer über den andern, und vermahnet nicht vielmehr einer den andern in Sanftmuth und traget euch unter einander? „Der ist“, sage ich, „ein vollkommener Mann,\footnote{Hier sind wir der Interpunction der Vulgata und der Jenaer Ausgabe gefolgt. Die anderen uns vorliegenden Ausgaben, Löcher, Petri, Wittenberger, Erlanger und Weimarische bieten: Hic, inquam, perfectus vir. Potest etc.} und kann auch den ganzen Leib im Zaum halten“, das heißt, wenn er seine Zunge zähmen kann, so kann er auch seinen ganzen Leib vor der Sünde bewahren, denn wer mit der Zunge nicht sündigt, der hat ohne Zweifel auch seine anderen Glieder schon gezähmt und hat sie im Zaume.

Danach legt er drei Gleichnisse vor, deren eines er schon begonnen hatte [Jac. 3, 3.]: „Siehe, wir legen den Pferden Zäume in das Maul, daß sie uns gehorchen, und lenken den ganzen Leib.“ Zweites [V. 4.]: „Siehe, die Schiffe, ob sie wohl so groß sind, und von starken Winden getrieben werden, werden sie doch gelenkt mit einem kleinen Ruder, wo der hin will, der es regiert. Also ist auch die Zunge ein kleines Glied und richtet große Dinge an“, oder: wirft sie hin und her, als ob er sagen wollte: Die Pferde können wir im Zaum halten und lenken, und die Schiffe, ja (wie hernach kommt [V. 7.]), auch die wilden Thiere, dagegen aber die Zunge, da sie nur ein kleines Glied ist, sollte am meisten im Zaum gehalten und gelenkt werden. Aber sie schleudert vielmehr und bewegt große Dinge, welche sonst wohl regiert werden, bewegt alles und wirft es umher, wie der Wind die Schiffe, gerade entgegengesetzt dem Ruder, welches die Schiffe im Zaum hält. An Größe gleicht sie der Ruderpinne der Schiffe, und ist doch in der Wirkung derselben entgegengesetzt, vielmehr ähnlich den Winden und Wirbelstürmen. Das dritte Gleichniß [V. 5.]: „Siehe, ein kleines Feuer, welch einen Wald zündet es an!“ das heißt, es trägt sich zu, daß ein klein Feuer viel Material (materiam) verzehrt, nämlich einen ganzen Wald. [V. 6.] „Und die Zunge ist auch ein Feuer, eine Welt voll Bosheit“ (mundus iniquitatis), griechisch: ὁ κόσμος τῆς ἀδικίας (mundus injustitiae = eine Welt der Ungerechtigkeit
\quellseite{643}
), als ob er sagen wollte: Sie richtet so viele Ungerechtigkeit an, daß man meinen könnte, sie wäre gleichsam eine Welt, weil sie ein Feuer ist, und vieles in Brand setzt.

Mit allen diesen und den folgenden Worten ist er darauf aus, daß er uns erinnere und lehre, weil unsere eigene Zunge unser größter Feind ist, vor der wir uns nicht weniger in Acht nehmen müssen, als vor unserem größten Feinde, und je kleiner sie ist, desto schädlicher ist sie, da sie sehr klein ist der Größe halben (in re), aber sehr groß in ihrer Kraft. Denn nichts thut dem Gewissen leichter Schaden als die Zunge. Denn deshalb ist auch in den Klöstern das Schweigen auferlegt, wiewohl wider dasselbe der Teufel auch dort die Logik eingeführt hat, die überaus geschwätzige Lehrerin der Geschwätzigkeit.

Aus derselben Ursache, und damit er das Schnellsein zu reden dämpfe, droht der HErr, daß die Menschen von einem jeglichen unnützen Worte Rechenschaft geben müssen [Matth. 12, 36.]; ein hartes aber notwendiges Urtheil wider ein so schlüpfriges Uebel und überaus wirksames Gift.

Diese haben auch die Philosophen gesehen, daß dem Menschen um deswillen zwei Ohren gegeben worden seien, aber nur Eine Zunge; und die Ohren selbst sind nicht nur weit und offen nach außen, sondern auch mit den Ohrmuscheln versehen, die sie umgeben, damit das Ohr leichter höre; endlich gewunden und gewissermaßen gegen das Wort gerichtet, damit nicht leicht ein Wort vorüberrausche. Die Zunge aber ist nicht bloß verschlossen und verborgen, sondern auch mit einer beinernen Mauer umschlossen, nämlich den Zähnen, sodann mit einem Bollwerk von Fleisch, das ist, den Lippen, außer dem, daß sie so weit entfernt ist von der Zunge und von dem Herzen, und so vieler Werkzeuge bedarf, während die Ohren keines bedürfen; dennoch ist sie schneller als alle Sinne.

Siehe, die fünf Sinne schwimmen gleichsam wie fünf Schiffe in diesen leiblichen Dingen und bringen in das Gedächtniß unzählige Dinge, nämlich durch Sehen, Hören, Fühlen, Riechen und Schmecken. Dennoch können sie nicht so viel hineinbringen, daß die Zunge nicht noch mehr hinausbringe, indem sie sehr oft Vermuthungen redet, Argwohn, Meinungen, vermessene Träume, leichtfertige Erdichtungen, die weder sind, noch waren, noch sein werden.

Desgleichen, deshalb war alsbald die zweite Plage Egyptens das Senden der Frösche, welche die unnütze Geschwätzigkeit bedeuten; denn die erste Plage war das Blut, welche bedeutet, daß, wenn der Mensch Fleisch und Blut wird, er ganz zur Zunge wird und nur eitle Dinge redet, und je eitler er ist, desto mehr will er reden. Und hierdurch magst du als an einem zuverlässigen Zeichen einen eitlen Menschen erkennen, wenn er geschwätzig ist; und umgekehrt: wie tüchtig und gehaltvoll (gravior) ein Mensch sei, magst du an diesem Zeichen erkennen, wenn er karg ist mit Worten, reich an Werken. Denn solche Leute haben einen Widerwillen gegen das Reden [und reden nicht], wenn sie nicht durch Gehorsam oder Noth dazu gezwungen werden, wie es heißt im 39. Psalm, V. 3.: „Ich bin verstummet und still, und schweige der Freuden“, denn wie sie ihr Leben anstellt, so auch das Reden. Wiederum [Spr. 10, 19.]: „Wo viel Worte sind, da geht es ohne Sünde nicht ab“; und ebendaselbst [Spr. 14, 23.]: „Wo man mit Worten umgehet, da ist Mangel.“ Hiob 11, 2. [Vulg.]: „Wird denn der, welcher viel redet, gerecht sein können?“ und Ps. 140, 12.: „Ein böses Maul wird kein Glück haben auf Erden.“

Aber dies ist am allerbößesten, wenn der Mensch vom Zorn getrieben (ira magistra) redet, daher sagt man, Aesopus habe Zungen gekauft, als er den Auftrag erhielt, das beste Fleisch zu kaufen, und ebenso, als er darnach auch das schlechteste Fleisch kaufen sollte. Ja, viele haben sich selbst das Leben genommen aus Furcht vor übler Nachrede (infamiae), und die Dichter sagen, daß Lycambes durch die [beißenden] iambischen Verse des Archilochus dahin gebracht worden sei, sich selbst zu erhängen. Wenn nun noch jemand daran zweifeln sollte, der frage einen Mann, der ein böses, zänkisches und geschwätziges Weib hat.

Ferner, was für die wilden Thiere die Zähne, die Klauen, Schnäbel, Hörner, Stachel, Gift oder sonst irgendwelche Waffen sind, ja, was für böse Menschen so mancherlei Geräthe und Waffen, das ist für die Seelen die Zunge, daher der 57. Psalm [V. 5.] sagt: „Die Menschenkinder, ihre Zähne sind Spieße und Pfeile.“ Desgleichen der 75. Psalm [V. 6. Vulg.]: „Erhebet euer Horn nicht so hoch.“ Daher ist [ihre Zunge] „ein scharfes Schwert“.\footnote{Wir haben nach Ps. 57, 5. [Vulg.] die Conjectur der Weimarischen Ausgabe angenommen: acutus gladius statt: arcus, gladius in den Ausgaben.}
\quellseite{644}
Deshalb wird sie mit Recht „eine Welt voll Ungerechtigkeit“ genannt und „ein kleines Feuer, welches einen großen Wald anzündet“.

Und es folgt [Jac. 3, 6.]: „Also ist die Zunge unter unsern Gliedern, und befleckt den ganzen Leib“, als ob er sagen wollte: Ist es nicht zu verwundern, daß das befleckt, und den ganzen Leib befleckt, was ganz und gar nicht außer uns, vielmehr so völlig in uns ist, daß es unter unseren eigenen Gliedern ist? Ist dies nicht einer wunderbaren Begebenheit (monstri) ganz ähnlich, daß ein Glied nicht bloß einem Gliede, sondern dem ganzen Leibe schadet, wiewohl sich doch die Glieder unter einander mit großer Sorge annehmen? Nur dies Eine Glied befleckt den ganzen Leib, das heißt, es macht ihn der Sünden schuldig, so daß um der Zunge willen der ganze Mensch bei GOtt und den Menschen verhaßt wird, weil ein Mensch mit einer bösen Zunge (linguosus) so schändlich und greulich ist, daß, wenn er auch am ganzen Leibe schön wäre, geschickt und reich, dennoch alles an ihm von den Menschen als ein Greuel geflohen wird. Denn wer geht gern mit geschwätzigen und plauderhaften Leuten um, und hat nicht vielmehr einen Abscheu und Ekel vor ihnen, und spricht: „Ei, wie wäßig ist der!“

[Jac. 3, 6. Vulg.]: „Und entzündet das Rad unserer Geburt und wird entzündet von der Hölle.“ Es scheint der Apostel unter „Geburt“ das zu verstehen, als was wir geboren sind, wie oben Cap. 1, 23. [Vulg.]: „Der ist gleich einem Manne, der das Angesicht seiner Geburt im Spiegel beschauet“, und wiederum [Jac. 1, 18. Vulg.]: „Er hat uns gezeuget nach seinem Willen durch das Wort der Wahrheit, daß wir eine Art Anfang (initium aliquod)\footnote{initium, welches in der Vulgata steht, fehlt in den Ausgaben. Petri, die Jenaer und die Wittenberger haben aliquid.} seiner Creatur wären.“ Also ist „Geburt“ die Zeugung oder Schöpfung und Geschöpf, zu dem wir geboren, hervorgebracht, geschaffen sind. Also wird „das Rad der Geburt“ genannt das ganze Wesen, und alles, als was wir geschaffen sind, gleichwie ein ganzer Leib. Und dieses „Rad der Geburt“, und den ganzen Inbegriff (circulum) unseres Wesens oder dessen, was wir sind, zündet die Zunge an. Denn sie ist ja ein Feuer und eine Welt voll Ungerechtigkeit; darum entzündet sie nicht einen Theil von uns, sondern alles, was wir sind und leben. Denn er hat sie vielleicht auch deshalb ein „Rad“ genannt, nicht bloß wegen unseres völligen Seins, sondern auch wegen unseres völligen Lebens und Wirkens in dieser Zeit; denn unser Leben ist ein flüchtiges Rad. Er fügt aber hinzu: „entzündet von der Hölle“, damit man nicht verstehe, daß hier eine gute Flamme entzünde, da sie von höllischem Feuer entzündet ist, das heißt, durch Zorn und Haß entzündet sie alles, was wir sind, mit Zorn und Haß, da sie durch Zorn reizt und bewegt, wie man an zornigen Leuten sehen kann, welche inwendig und auswendig zittern und glühen.

Das vierte Gleichniß [Jac. 3, 7. 8.]: „Denn alle Natur der Thiere, und der Vögel, und der Schlangen und der Meerwunder, werden gezähmt, und sind gezähmt von der menschlichen Natur; aber die Zunge kann kein Mensch zähmen.“ Siehe, mit wie gewaltigen Worten und in wie vielerlei Hinsicht (sensibus) er die Zunge angreift, um sie dem höllischen Feuer gleich zu machen und herunterzusetzen unter die wilden Thiere. Wer sollte sich nicht vor sich selbst fürchten, da er ein schändlicheres und grausameres wildes Thier bei sich hat, als die ganze Welt? Wer sollte seine Zunge nicht hassen? Wer sollte sich nicht scheuen zu reden, da er hört, daß seine Zunge durch so viele greuliche Namen geschmäht wird?

Es folgt [Jac. 3, 8.]: „Das unruhige Uebel voll tödtliches Gifts.“ Es ist [dem Apostel] zu wenig, daß sie [die Zunge] unruhig sein sollte, was an sich schon ein Uebel wäre, sondern „ein unruhiges Uebel“, das heißt, welches immer schadet und zu schaden nicht aufhört. Ferner „voll tödtliches Gifts“, das heißt, welches tödtet, nicht bloß verderbt, das ist, voll eines rasch wirkenden Giftes, nicht eines langsamen oder unkräftigen (tardo). Hierdurch wird freilich klar, wie tief das Gift der Ohrenbläser eindringt, so daß denen, die sich gegenseitig vergiftet haben, kaum im ganzen Leben ein Heilmittel verschafft werden kann, da der Neid und der Haß bis ins Mark gedrungen ist; sehr oft sterben sie lieber mit ihrem Hasse, als daß sie sich versöhnen sollten. Daher sagt er sehr nachdrucksvoll und überaus stark „voll tödtliches Gifts“ nicht bloß „verderbliches Giftes“.

[Jac. 3, 9.]: „Durch sie loben wir GOtt den Vater, und durch sie fluchen wir den Menschen,
\quellseite{645}
nach dem Bilde GOttes gemacht.“ Als ob er sagen wollte: Man kann nicht beides zugleich thun, wenn es nicht zum Scheine geschieht. Darum betrüge sich niemand, daß er Gott liebe; wenn er einem Menschen flucht, so lobt er auch GOtt nicht, dessen Werk der Mensch ist. Denn wer könnte ein Gemälde oder ein Bildwerk tadeln oder fluchen, daß ein solcher Fluch nicht auf den Maler und Meister zurückfiele? So sagt Jeremias, Cap. 15, 10., von denen, die den Werken des HErrn fluchen: „Alle fluchen sie mir.“ Und Jesaias, Cap. 8, 21.: „Wenn sie aber Hunger leiden, werden sie zürnen, und fluchen ihrem Könige und ihrem GOtt.“ Darum war es vornehmlich im Gesetze verboten, und ist es noch [2 Mos. 22, 28.]: „Den Göttern sollst du nicht fluchen, und den Obersten in deinem Volk sollst du nicht lästern.“ Und der Apostel sagt [1 Cor. 6, 10.]: „Die Lästerer werden das Reich GOttes nicht ererben.“

Wenn nun die Menschen dies beachteten, daß sie durch Fluchen der Menschen auch Gotte fluchen, dessen Werk und Geschöpf sie fluchen, so würden sie gewiß nicht so bereit sein zu fluchen, wie es heißt, Spr. 14, 31. [Vulg.]: „Wer den Armen verleumdet, der lästert desselben Schöpfer.“ Denn was thut der, welcher lästert, anders, als daß er sagt: Ih, daß du nicht anders bist und anders gebildet! während er sich doch dessen erbarmen sollte, beten und sprechen: O HErr, mache ihn anders! gleichwie man einen Werkmeister erinnern und bitten darf, aber nicht lästern, wie es ebendaselbst heißt, Spr. 14, 31.: „Der ehrt ihn aber“, nämlich seinen Schöpfer, „wer sich des Armen erbarmt.“

Was nun St. Petrus verbietet [1. Ep. 3, 9.], daß man nicht Scheltwort mit Scheltwort, noch Böses mit Bösem vergelten soll, das lehrt gleicherweise auch Christus [Matth. 5, 44.]: „Bittet für die, so euch beleidigen und verfolgen“, und Paulus Röm. 12, 14.: „Segnet die euch fluchen,\footnote{In der Vulgata steht Röm. 12, 14., ebenso wie in unserer Bibel: „verfolgen“, deshalb hätte Luc. 6, 28. angezogen werden sollen.} segnet und fluchet nicht.“ Was wird denen widerfahren, welche denen fluchen, welche ihnen nicht fluchen? Darum sagt der 34. Psalm, V. 14. f.: „Behüte deine Zunge vor Bösem, und deine Lippen, daß sie nicht falsch reden. Laß vom Bösen und thue Gutes; suche Frieden und jage ihm nach“, das heißt, „vom Bösen“, nicht allein mit Erdulden, sondern daß du es auch nicht nachtragest oder wiedervergeltest, und „das Gute“ nicht allein, daß du es vergeltest, sondern es auch erweisest für Böses, welches dir widerfahren ist; sonst wird nicht der Friede gesucht oder ihm nachgejagt, wenn du Böses mit Bösem vergiltst.

Dieser Grad ist also der zweite, der vom HErrn verboten ist, daß jemand seinem Nächsten nicht fluche, ihn nicht verletze und verleumde. Und gleichwie im ersten Grade mehr Leute schuldig sind, als sie selbst glauben, so auch in diesem. Denn es werden durch dies Gebot die drei größten Vergehen der Zunge verboten, nämlich Fluchen, üble Nachrede und Ohrenbläserei; aber dies dritte gehört zu dem ersten Grade. Und so werden alle Uebel, welche die Zunge anrichtet, verboten durch das dritte, fünfte, sechste und siebente Gebot. Durch das dritte Meineid und Lästerung und irgend welche andere Beleidigungen, die wider GOtt gehen, aber durch das fünfte die Beleidigungen gegen den Nächsten, durch das sechste unkeusche Worte, durch das siebente Worte, durch welche man täuscht und betrügt in Geschäften oder vor Gericht. Denn auch im siebenten Gebote ist üble Nachrede, aber öffentlich und vor Gericht; hier aber ist eine andere.

Es ist nun zwischen Fluchen, und übler Nachrede, und Ohrenbläserei dieser Unterschied, daß das Fluchen geschieht gegen einen, der selbst gegenwärtig ist, oder gleichsam gegenwärtig ist, und ins Angesicht. Hierunter sind auch begriffen alle Vergehen, die man sich zu Schulden kommen läßt durch Schmähnen, Lästern, Anschuldigen, Tadeln, Geben von Beinamen. Üble Nachrede geschieht, wo die Person selbst abwesend ist, zu einem anderen, und folgt auf das Fluchen. Denn wenn jemand Böses wider jemanden aufgebracht hat, der entweder äußerlich gegenwärtig ist, oder den er sich inwendig vor das Gedächtnis stellt, und dies einem anderen mittheilt, so ist aus dem Fluchen schon eine üble Nachrede geworden. Darum ist keiner, der übel nachredet, der nicht auch ein Flucher sei, wenigstens bei sich selbst. Und diese beiden Stücke hat der Spruch [Ps. 34, 14.] schon ausgedrückt: „Behüte deine Zunge vor Bösem (das heißt, daß du niemandem fluchest, sei es nun mit Worten oder mit wortlosem [tacito] Gemurmel), und deine Lippen
\quellseite{646}
, daß sie nicht falsch reden“ (das heißt, hüte dich, daß du nicht übel nachredest), denn es ist eine Falschheit, wenn jemand hinter dem Rücken redet, da es der nicht weiß, dem er schadet. Aber umgekehrt ist nicht jeder Flucher auch ein Mensch, der übel nachredet, wenigstens nicht gegen seinen Mitmenschen, wiewohl er Gotte übel nachredet, wenn er den Menschen flucht. Wenn er nun aber auch das Böse, welches zwei wider einander gethan haben, beiden zuträgt, oder auch nur von einem zum anderen, als der Freund des einen und der Feind des andern, dann ist er ein zweizüngiger Ohrenbläser- und ein Stifter der Zwietracht, der ärger ist als die beiden anderen; deshalb ist er nicht in diesem Grade, sondern auf die erste Stufe zu stellen mit den Todtschlägern.

So wollen wir nun diesen Grad abschließen. Gleichwie niemand die erste Stufe der Erfüllung dieses Gebotes erreicht, als wer nicht nur nicht tödtet, oder verletzt, oder Zwietracht säet (denn auch durch Unterlassen kann man dies begehen), sondern auch so sanftmüthig ist, daß er vielmehr lebendig mache und erquicke durch Werke der Barmherzigkeit und des Wohlwollens, und so friedfertig, daß er sich auch bemühe, mit sanften und lieblichen Worten die Zwieträchtigen zu versöhnen, so viel er vermag, und ins Mittel zu treten, daß sie nicht uneinig werden (denn diese werden die Worte hören [Matth. 5, 5. f.]: „Selig sind die Sanftmüthigen, selig sind die Barmherzigen, selig sind die Friedfertigen; denn sie werden das Erdreich besitzen, Barmherzigkeit erlangen, GOttes Kinder genannt werden“: siehe, diese drei Seligpreisungen sind in diesem Gebote begriffen): so erreicht die zweite Stufe der Erfüllung niemand, als wer nicht allein nicht flucht, übel nachredet oder tadelt, sondern auch in allen Dingen die segnet, die ihn verfolgen, und die Angeklagten entschuldigt, für alle bittet, und denen widersteht, die üble Nachrede machen, soviel er vermag. Siehe, das ist der geistliche Verstand dieses Gebotes.

\footnote{Löscher: „Dies hat Luther als Predigt gehalten am St. Andreastage [den 30.] November im Jahre 1516.“ Das Exordium zu dieser Predigt findet sich Walch, St. Louiser Ausgabe, Bd. XII, 1778.} Jetzt auch gleicherweise vom dritten Grade.

Der dritte Grad des Zorns ist, daß jemand zu seinem Bruder sagt „Racha“ oder „Archa“, indem man einen rauhen Kehllaut und gleichsam ein Rasseln der Luftröhre hervorbringt. Ich habe gesagt, daß darin alle Zeichen des Zorns begriffen sind, außer den Worten, ja, auch die Worte außer den Flüchen und üblen Nachreden, als da sind Spott (ironiae) in Worten, Aufwerfen der Lippen, Runzeln der Stirn, Rümpfen der Nase, Glühen des Angesichts, sodann auch Schweigen, das Gesicht abwenden, grimmiger Blick: das alles sind Zeichen eines erzürnten Herzens, das da erfüllt ist mit Haß. Dazu gehört auch noch Schreien, Knirschen mit den Zähnen etc.

Hier sagen einige, es sei zwar vonnöthen, daß man den Groll fahren lasse, aber nicht die Zeichen des Grolles. Wiewohl dieser Ausspruch in einiger Hinsicht wahr ist, ist es doch nicht durchgängig der Fall, und es würde viel besser so ausgedrückt: Es ist bisweilen nützlich, die Zeichen des Grolles anzuwenden, nicht als ob sie Zeichen des Grolles wären, sondern den Zeichen des Grolles ähnlich; sonst sind sie nicht Zeichen des Grolles in der That, oder man darf sie auf keine Weise behalten, wenn sie es in der That sind. Daher würde noch am besten gesagt: daß man die Zeichen der Liebe bisweilen anstehen lassen müsse, und Zeichen des Zorns zeigen. Denn wenn man die Zeichen des Grolles nicht ablegen darf: warum sollte man sie denn wegen irgend welches Nutzens (fructum) nicht vielmehr zeigen? damit nämlich nicht bloß eine Gelegenheit gesucht werde, das Gute zu unterlassen.

\footnote{Um diesen etwas schwierigen Absatz zu verstehen, vergleiche man den vorletzten Absatz von den nachfolgenden „Folgerungssätzen“. Die Schwierigkeit hat hauptsächlich darin ihren Grund, daß die „Zeichen des Grolles“ so beschaffen sind, „daß dein Bruder nicht wissen kann, wie er mit dir daran ist“.} Es ist also besser, daß man [die Zeichen des Zorns] zeige, als daß man [dieselben] nicht fahren lasse; denn wer sie zeigt, der thut es nicht aus Zorn, sondern aus einem guten Eifer; darum bleibt er des Bösen unverdächtig. Wer aber dieselben nicht aufgibt, der ist der Beschuldigung ausgesetzt (arguitur), daß er zuvor solche Zeichen aus Zorn gethan habe, und darum läßt er den Bruder, wenn er dieselben nicht ablegt, in dem Verdachte des Zorns, und hindert dadurch wahren und völligen Frieden und Eintracht, da jener nicht glauben kann, daß sein Herz versöhnt sei, und geärgert wird [und veranlaßt], wieder dergleichen zu thun.

\quellseite{647}
Deshalb soll man solche Lehrer zu diesem Gebote der Natur weisen: Was du nicht willst, daß man dir thue, das thue auch einem anderen nicht, und: Was du willst, daß man dir thue, das thue einem anderen; und ein Heide spricht: Willst du geliebt werden, so liebe. Aber es ist kein Zweifel, daß alle wollten, daß ein Mensch, der von uns beleidigt worden und zornig auf uns gewesen ist, nicht bloß die Zeichen des Zornes ablege, sondern sogar die entgegengesetzten Zeichen zeige, damit wir seiner gewiß werden. Deshalb begehrt man stärkere Zeichen, als wenn er nicht beleidigt worden wäre. Denn in solcher Weise hat auch GOtt dem menschlichen Geschlechte nach dem Sündenfall (peccatum) die Wohlthaten in reicherem Maße erzeigt als vor der Sünde, dadurch, daß er seines eigenen Sohnes nicht verschont hat etc. [Röm. 8, 32.]. Wenn er nun diese so großen Zeichen nicht gegeben hätte, wer hätte glauben können oder dürfen, daß er versöhnt wäre? Darum sollen GOttes Kinder ihrem Vater nachahmen.

Durch diese Aussprüche wird uns gelehrt, wie schwierig es sei, das Gebot GOttes zu halten. Denn durch die Erfahrung ist es offenbar, wie schwer es den Menschen sei, diese bösen Zeichen abzulegen und gute Zeichen zu erweisen, wie sie sich krümmen und winden, ehe sie dazu bewegt werden können, daß sie es wenigstens zum Schein thun; denn in der That und von Herzen werden sie es nie thun, es sei denn, daß GOttes Gnade ihnen beistehe, denn sie sagen: Ich kann dir nicht gut sein, nicht freundlich gegen dich sein. Und diese niedrigen Menschen bedenken nicht, daß sie nicht wollten, daß ihnen solches geschähe, weder von GOtt noch von Menschen. Darum sieht man an diesem Greuel, was die Hölle sein werde, und was diese Welt sei, nämlich ein Reich des Teufels, in welchem in den Menschen ein so großes Uebel gefunden wird, welches sie hartnäckig verstockt gegen die Natur und das, was ihnen die Natur in ihrem eigenen Herzen vorschreibt (dictamen).

Deshalb ist zwischen Zorn und Gehässigkeit (invidia) kein anderer Unterschied als wie zwischen neuem Wein und altem Wein; denn Gehässigkeit oder Haß ist ein eingewurzelter (inveterata) Zorn, wie St. Augustin in seiner Regel sagt: damit nicht der Zorn aufmache zum Haße, und einen Balken mache aus dem Splitter. Denn der Zorn, welcher bald vorübergeht und nicht zum Haße wird, kann kaum ein Zorn genannt werden, denn er wird alsbald der Vergessenheit übergeben. Aber wenn er eingewurzelt ist, wird er ein überaus böses Uebel, weil er ein geistliches Uebel ist, so daß es auch ein Sprichwort ist: Nichts bleibt länger im Gedächtniß als ein Unrecht, und nichts wird schneller vergessen als eine Wohlthat; beides geschieht durch die Schuld der natürlichen Verderbtheit, da es umgekehrt sein sollte. Denn daran wird die gute Natur eines Menschen erkannt, wenn er eine Beleidigung rasch vergißt und einer Wohlthat immer gedenkt, so daß er sich nicht gern wohlthun läßt, indem er sieht, daß er mit einer großen Menge Verpflichtungen zu Dank überhäuft wird und zu einer Wiedervergeltung, die ihm unmöglich ist. Aber dagegen ist die Art einer bösen Natur, eine Beleidigung immer zu behalten, und lange Vergangenes immer wieder vorzuhalten und aufzurücken.

Deshalb wollen wir es mit Exempeln erläutern, wann man Zeichen des Zorns oder der Härte zeigen dürfe. Als, wenn jemand so halsstarrig böse ist, daß er, wenn er ein-, zwei- oder dreimal vermahnt ist, nicht ablassen will, so darf man nicht allein, sondern man muß Zeichen des Zornes gegen ihn beweisen, daher heißt es 1 Cor. 5, 11.: „So jemand ist, der sich läßt einen Bruder nennen, und ist ein Hurer, oder ein Geiziger, oder ein Abgöttischer, oder ein Lästerer, oder ein Trunkenbold, oder ein Räuber“, und 2 Thess. 3, 14. f.: „So aber jemand nicht gehorsam ist unserm Wort, den tadel (notate), und habet nichts mit ihm zu schaffen, auf daß er schamroth werde. Doch haltet ihn nicht als einen Feind, sondern vermahnet ihn als einen Bruder.“ Siehe, nicht als ein Feind, sondern als ein Bruder soll er gebessert werden. Das ist ein heiliger und guter Zorn, der in der Schrift sehr gelobt wird.

Deshalb ist hier ein Unterschied zu machen: Erstens, daß er ermahnt werde nach dem Worte Christi [Matth. 18, 15.]: „Sündiget dein Bruder an dir.“ Uebrigens soll man ihn nicht gleich das erste Mal von sich stoßen, sondern ihn ermahnen. Denn wenn er siebenzigmal siebenmal um Vergebung bittet, soll man sie ihm nicht abschlagen.

Zweitens, unterscheide, ob du um deswillen so thust, weil er dir etwas zu Leide gethan hat, oder weil er eine Sünde begangen hat. Denn
\quellseite{648}
wenn du GOtt mehr liebst als dich selbst, so mußt du mehr darüber zürnen, daß er GOtt beleidigt hat, als dich. Sodann mußt du dich das ebenso jammern lassen, daß er sich selbst innwendig ein ärgeres Leid zugefügt hat, als dir auswendig. Wenn du darauf Rücksicht genommen hast, wird man dich leicht lehren können, wie und mit was für einem Herzen du ihm die Zeichen der Liebe versagen, oder nicht versagen sollest. Aber damit dich dein Herz nicht betrüge, und du nicht sprichest: So thue ich; ich zürne um GOttes, und nicht um meinetwillen, und du vielleicht lügest, so prüfe dich auch selbst in Bezug darauf an dir selbst, und siehe zu, ob es dir ebenso Leid thut und bewegt, wenn er etwa einen anderen verletzte, als weil er dich verletzt. Wenn du nämlich in jenen Falle nicht bewegt wirst, so ist es ein Zeichen, daß du dich selbst suchst und rächt, aber [deinen Nächsten] nicht besserst.

Und was halte ich mich lange dabei auf? Es ist eine vergebliche Disputation, ob die Zeichen des Grolles abzulegen seien, denn man muß sie ganz und gar ablegen. Wenn aber eine Ursache da ist, weshalb sie gezeigt werden sollten, da muß man ganz von Herzen zürnen, nicht bloß die Zeichen sehen lassen. Die Zeichen des Grolles, wenn sie vom Grolle herrühren, die soll man ganz und gar ablegen. Sind sie aber nicht Zeichen des Grolles, und es ist nöthig, sie um der Besserung des Bruders willen zu zeigen, muß auch ein vollkommener Zorn hinzugethan werden, damit er sehe, daß sie ihm von Herzen erzeigt werden, und aus welcher Ursache. Und gehe ja nicht so schweigend und in falschem Schein (simulans) einher, daß dein Bruder nicht wissen kann, wie er mit dir daran sei. Eine solche Lehre ist böse und ganz und gar gegen die Liebe, denn sie lehrt Heuchelei.

Aber weitläufig und viel sind die Gebote, welche die Liebe alle leicht sehen und thun würde; wenn die nicht da ist, so kann nichts genugsam gelehrt werden. Denn sie lehrt, wie der Mensch selbst geliebt werden will, daß man Sorge für ihn trage, sich sein annehme, ihm wohlthue, nichts gegen ihn in verstellter Weise vornehme: in solcher Weise lehrt sie, dem Nächsten ebendaselbe zu erweisen. Denn ein Freund muß nicht betrüglich einhergehen, wie St. Hieronymus sagt: Eine wahre Freundschaft kann Verstellung nicht leiden. Ist Groll da, so zeige ihn; ist er nicht da, so zeige es wiederum, und zwar entweder mit Zeichen der Gütigkeit oder der Strenge.

\subsection{Folgerungssatz}

Es ist klar, daß auch dieser Grad nicht gehalten wird, wenn sich jemand nicht bemüht, nicht allein die Zeichen zu unterlassen, sondern auch aus einem wohlwollenden Herzen die entgegengesetzten Zeichen zu beweisen. Denn durch dieses Gebot bildet uns GOtt zu aller Gütigkeit heran, und beabsichtiget, in uns die Wurzel des Vermögens zu zürnen auszurotten, so daß da nur ein liebliches Herz zurückbliebe, welches in der Gnade wirkt, nicht durch die Natur.

Es folgt der vierte Grad.

Der vierte Grad ist, mit dem Herzen zürnen, und dieser ist der erste und die Wurzel aller anderen. Denn diesen meint auch das Gesetz vornehmlich, da es sagt: „Du sollst nicht tödten.“ Denn es fordert ein reines Herz, nicht bloß die Hand, wie die Heuchelei der Juden [vorgab], wie der 24. Psalm, V. 4., sagt: „Der unschuldige Hände hat, und reines Herzens ist.“ Daher sind fast alle Seligpreisungen Matth. 5, 3. f. in diesem Gebote begriffen und legen dasselbe aus. Denn: „Selig sind die Sanftmüthigen, selig sind die Barmherzigen, selig sind die Friedfertigen, selig sind die Geduldigen, selig sind, die reines Herzens sind“, werden in diesem Gebote begriffen. Daß nun aber auch: „Selig sind die Armen, selig sind, die da hungert nach der Gerechtigkeit, selig sind, die da Leid tragen“, die bezeichnen, welche sich selbst als solche erkennen, die dieses Gebot nicht erfüllen, und darum trauern und wünschen und hungern: wer wollte daran zweifeln, daß hier die Leute sind, oder wenigstens solche Leute sein müssen, daß sie dieses und jegliches andere Gebot erfüllen.

Darum sind in diesem Grade alle bösen Gedancken mit ihrer Wurzel, nämlich dem Vermögen zu zürnen, welches geneigt ist zur Rache, als da sind Zorn, Haß, böser Argwohn, Bitterkeit, Groll etc. Und gleichwie dies der erste Grad ist, so hat er auch weniger Arten als der nächste Grad nach ihm. Denn der Todschlag im Werke hat sehr viele Arten, als, Todtschlag, Brudermord, Gattenmord (uxoricidium), Vatermord, Mord einer geweihten Person (sacrilegium) etc.

Und hierdurch ist klar, daß in diesem Gebote vier Stufen sind und in den einzelnen Stufen noch viele Grade, und darum ist kein Mensch in
\quellseite{649}
diesem Leben, der nicht in einem derselben wäre. Selig aber ist der, welcher auf den höchsten Stufen steht, unselig, wer auf den untersten ist. Denn Gütigkeit ist der Himmel, Zorn die Hölle; in der Mitte aber von beiden ist diese Welt. Denn [Ps. 122, 4. Augl.] „dahin gingen die Stämme hinauf, die Stämme des HErrn“, hier [5 Mos. 16, 30.] sind sie lebendig in den Pfuhl und in die Hölle gefahren; um so näher der Sanftmuth, je näher sie dem Himmel sind, und wiederum.

Darum ist es für diese vierte Stufe nicht genug, wenn wir dem Bruder nicht zürnen, wenn wir nicht auch eine überaus gütige Gesinnung und ein ganz sanftes Herz gegen ihn haben. Da wir nun so nicht beschaffen sind, so müssen wir zu der Gnade Christi fliehen, daß er unterdessen aus seiner Fülle unserem Mangel abhelfe, bis daß wir auch vollkommen werden.

Es folgt das sechste Gebot.

\section{Das sechste Gebot}

\begin{schriftzitat}
Du sollst nicht ehebrechen.
\end{schriftzitat}

\footnote{Löscher: „Dies hat Luther am St. Nicolaustage [den 6. December 1516] gehalten.“ — Das Exordium hierzu ist nicht mehr vorhanden.} Du ein gar kurzes Wort, aber eine sehr weitgehende Meinung, darum bedarf es der Auslegung. Darum wollen wir Christum hören, welcher Matth. 5, 27. f. sagt: „Ihr habt gehört, daß zu den Alten gesagt ist: Du sollst nicht ehebrechen. Ich aber sage euch: Wer ein Weib ansiehet, ihrer zu begehren, der hat schon mit ihr die Ehe gebrochen in seinem Herzen.“ Auch dies redet er kurz, indem er es frei läßt, aus dem vorhergehenden Gebote dieselbe Unterscheidung der Stufen herzuleiten, und nur den setzt, der im Herzen ist. Daher wollen auch wir hier dieselben vier Stufen setzen, nämlich im Herzen, in Zeichen, in Worten, in Werken. Denn dies Gebot hat der HErr um deswillen ausgelegt, weil die Alten, das heißt die Juden, es nahmen nach der Bedeutung und Schärfe des Wortes, und so lehrten, und dies Gebot zunichte machten, ebenso wie alle anderen, daß nur der ein Ehebrecher wäre, welcher mit dem Werke darüber gethan hätte, sich aber nicht kümmerten um die Begierden, Zeichen und Worte; darum wurden sie unrein. Und in dieser Meinung standen sie so fest, daß Jeremias zu ihnen sagte [Cap. 5, 8.]: „Ein jeglicher wiehert nach seines Nächsten Weibe, wie die vollen müßigen Hengste.“ Daher redet er zu den Neuen, das ist zu uns, anders, damit alle Begierde hier ihren Raum verliere. Denn es sagt St. Augustinus an dieser Stelle über die Bergpredigt: daß unter dem Namen Ehebrecher, deren in diesem Capitel gedacht wird, alle fleischliche und unkeusche Begierde verstanden werde.

Die erste Stufe ist also das äußerliche Werk, welches ohne Zweifel aus der innerlichen Begierde hervorgeht, wie St. Augustinus ebendaselbst beweist. Wenn darin aber alle fleischliche Begierde verboten wird, so wird also auch jedes Werk verboten, welches aus derselben hervorgeht. Deshalb werden viele Unterschiede der Werke gesetzt.

1. Einfache Hurerei eines Ledigen mit einer Ledigen; davon unterscheiden andere 2. die Hurerei, weil jene nur mit Einem Weibe zu schaffen hat, diese aber eine unstäte Brunst der Unkeuschheit ist gegen viele. 3. Schändung, das heißt, Schwächen einer Jungfrau. 4. Entführung, das heißt, wenn eine Tochter oder ein Sohn geraubt wird, worin zugleich auch ein Diebstahl begangen wird, und zwar der allergrößte Diebstahl. 5. Ehebruch, der bisweilen einfach ist, bisweilen zweifach, je nachdem beide verheirathet sind, oder nur einer. 6. Blutschande, welche gegen die Eheverbietung geht, welche man der Blutsverwandtschaft schuldig ist. 7. Sacrilegium, das heißt, wenn eine geistliche Person (religiosus), ein Priester, eine Nonne und alle anderen, welche Gottes Keuschheit gelobt haben, Hurerei treiben. 8. Etliche fügen noch das Uebermaß in der Ehe hinzu. Diese Grade im Werke sind in solcher Weise
\quellseite{650}
verschieden und unter einander größer und geringer, je nachdem die Umstände sich ändern. Denn so sind auch im vorhergehenden Gebote beim Todtschlage vielerlei unterschiedliche Werke verboten worden, wie gesagt ist, je nach verschiedenen Umständen.

Wiewohl aber jedes Gebot GOttes gleichsam ein Licht ist, durch welches erkannt wird, was der Mensch, die Welt und das Fleisch sei, und wie weit wir von GOtt entfernt seien, so zeigt dies doch vornehmlich dieses Gebot. Denn der Apostel sagt [Röm. 3, 20.]: „Durch das Gesetz kommt Erkenntniß der Sünde.“ Also ist das Gesetz ein Licht, durch dessen Leuchten wir das zahllose Elend des menschlichen Geschlechtes sehen.

Denn wer ist so gar ein Stein, daß er sich auch nur des Seufzens und der Thränen enthalten könnte, wenn er diese Gebote GOttes hört, und sie ihn erleuchten, ja, die ganze Welt zeigen, wenn er nur Acht darauf haben wollte, eine wie unendlich große Zahl der Menschen ihr Leben so führen, daß sie nicht nach den Geboten wandeln? Wenn keine andere Ursache da wäre, weshalb wir niemals fröhlich sein, sondern immer trauern sollten, so wäre diese allein ausreichend, wenn wir uns das Leben der Menschen in Vergleich zu den Geboten GOttes vor Augen stellten. Denn dann würden wir sehen, wie das menschliche Geschlecht ein ganz elender verlorner Haufe ist, da man so selten Leute sieht, welche nach den Geboten GOttes leben und handeln, und besonders in diesem Gebote. Ist es nicht erschrecklich, daß es so viele Grade der Werke gibt, durch welche gegen dies Gebot gehandelt wird, und daß die Menschen ohne alle Furcht dieselben frech vollbringen (ruere) und verloren gehen? Und dennoch können wir über solche Dinge lachen! Deshalb sollen wir auf dies Gebot Acht haben als auf eine Leuchte, und durch dasselbe [erkennen], wie so unendlich viele Menschen zu Falle kommen. Dann werden wir sehen, daß diese Welt ein greulicher Pfuhl ist und der Hölle ganz nahe, und wenig Gerechtigkeit und Wahrheit in ihr; ja, die Schrift sagt [Ps. 39, 6. 12. 62, 10.], daß alle Menschen „gar nichts“ sind und „Lügner“ [Ps. 116, 11.], deshalb ist in Wahrheit die Hölle in der Welt.

Das erste Elend der Menschen ist also einfache Hurerei, in welcher nicht allein die sündigen, welche sie begehen, sondern auch die, welche sie begünstigen, dazu helfen, dazu rathen, nicht Widerstand leisten, lachen, kuppeln, vermitteln, und besonders die, welche Anlaß dazu geben.

Es „begünstigen“ sie aber diejenigen, welche sie in ihren Häusern beherbergen, oder auch, wenn Fürsten, Oberherren und Richter es nicht wehren, sondern zulassen. Diesen wird es nicht nützen, daß sie selbst nicht huren, sondern es wird zu ihnen gesagt werden [Ps. 50, 18. 21.]: „Wenn du einen Dieb siehest, so läufst du mit ihm, und hast Gemeinschaft mit den Ehebrechern. Das thust du, und ich schweige, da meinst du, ich werde sein ewig wie du. Aber ich will dich strafen und will dir’s unter Augen stellen.“ Und dieses Elend, ach! als eine wie große Pestilenz herrscht es, da jeder Seinesgleichen unterweist, verführt, verderbt, sowohl unter Jünglingen als auch Jungfrauen etc.

Es „vermitteln“ aber die, welche Zwischenträger und Kuppler sind. Darüber kannst du eine schöne Historie finden in der fünften Abtheilung, im 80. Capitel des „Spiegels der Grempel“. Denn was der Teufel allein nicht kann noch vermöchte, das richtet er durch diese Leute aus.\footnote{Vergleiche die Erzählung Walch, St. Louiser Ausgabe, Bd. VIII, 544, § 132, und Tischreden, Cap. 43, § 40. Walch, St. Louiser Ausgabe, Bd. XXII, 1146.}

Es „leisten aber nicht Widerstand“, die nicht ermahnen oder strafen nach dem Gebote Christi, weder Mitleid haben mit ihrer Seele, noch für sie beten, wie es Christen geziemt. Zudem machen sie einen Scherz daraus, mit Erzählen und Lachen über das Verderben ihres Nächsten, während es doch, wie ich gesagt habe, als ein Zustand eines großen Elendes angesehen werden sollte.

„Anlaß geben“ aber diejenigen, welche mit schlüpfrigen Geberden und unkeuschen Worten oder überaus herrlichem Putz die Lüste anderer reizen, besonders wenn sie es mit Absicht thun; denn der übermäßige Schmuck ist immer eine Ursache, wiewohl nur eine Nebenursache (per accidens) der bösen Lust, wie es geschieht beim Tanze, auf den Straßen, bei Gastmählern etc. Wehe aber denen, die mit unzüchtigen Worten die Herzen der Unschuldigen verwunden, wie es jetzt in erschrecklicher Weise im Schwange geht in den öffentlichen Herbergen, wo die Kaufleute zusammenkommen. Es ist, wie man sagt, nicht nöthig, Läuse in den Pelz zu setzen, denn sie machen von selbst darin.

Denn unser Fleisch ist durch die erste Sünde
\quellseite{651}
mit einer zwiefachen Wunde überaus schwer verletzt. Die erste ist die Empfindlichkeit gegen Beleidigungen, die andere die böse Lust; diese zwei Wunden werden uns durch die Gebote offenbart, aber durch die Gnade werden sie geheilt. Aber die schlimmere Wunde ist die böse Lust [, und um so gefährlicher], je weniger sie schmerzt und bekümmert macht, sondern in angenehmer Weise bewegt; darum wird sie von thörichten Leuten nicht für ein Uebel geachtet, wiewohl sie zwiefach schlimmer ist als die andere, so daß die Schrift sagt [Hiob 40, 11.]: „Seine Kraft ist in seinen Lenden, und sein Vermögen im Nabel seines Bauchs.“ Denn vor (ante) der Sünde war nichts, das beleidigte, nichts, darnach man gelüstete. Jetzt ist aber beides da.

Deshalb soll ein Christenmens[ch] das wissen, daß, wenn er fühlt, er werde zu Zorn oder zu böser Lust bewegt, er nicht anders denke, als daß er ein tödtliches Gift und ein überaus böses Geschwür fühle, und soll dafür zum HErrn seufzen, damit er geheilt werde. Denn es ist die gewisse Wahrheit, daß, wiewohl dieser Kitzel angenehm und lieblich ist, er doch ein überaus feuriger Biß der alten Schlange ist, welcher die erschrecklichste Gewissensangst hervorbringt und die äußerste Bitterkeit. Und dadurch wird er bewegt, die Gnade zu suchen, wie die Figur im vierten Buche Mosis [Cap. 21, 9.] das in sich enthält, daß die von den feurigen Schlangen Gebissenen zu der ehernen Schlange aufblickten und gesund wurden. So soll auch derjenige, welcher den Kitzel des Fleisches empfindet, zu Christo, dem Gekreuzigten, aufblicken. Denn hierdurch wird Rechenschaft dafür gegeben, warum der HErr, damit nicht zufrieden, daß der Mensch von sich selbst geneigt ist zum Bösen, noch den Teufel als einen Versucher und Antreiber gebraucht, nämlich, [erstens] da ja der Mensch nachlässig ist, seine Genesung bei GOtt zu suchen, darum drängt er ihn durch Anfechtungen, damit er ihn zwinge, die Barmherzigkeit und Gnade GOttes zu suchen; zweitens auch deswegen, damit er sich nicht etwa erhebe in den Gaben, die er empfangen hat. So hat er es mit dem Apostel Paulus gemacht, so mit St. Hieronymus, so mit vielen anderen bis an ihren Tod. Siehe in dem „Spiegel der Grempel“ unter dem Titel „Anfechtung“.

Es sind aber andere, welche von diesem Uebel sehr heftig angegriffen werden, andere aber gelinder, noch andere gar wenig. Dennoch bleibt niemand keusch (continet) ohne die Gnade GOttes, wie es Weish. 8, 21. heißt: „Da ich aber erfuhr, daß ich nicht anders könnte tüchtig sein, es gäbe mir’s denn Gott (und daßelbige war auch Klugheit, erkennen, wes solche Gnade ist); trat ich zum Herrn“ etc. Siehe, wie man die Heilung von dieser Krankheit durch das Gebet suchen muß, und es ist eine große Weisheit, daß man dies sein Uebel und den Ort der Arznei erkenne. So sagt auch der Märtyrer St. Chrysantus: Der irrt, der da meint, daß er sich aus eigenen Kräften in der Keuschheit erhalten könne. Daher sagt St. Augustin im 1. Buche „über die Bergpredigt“: Wer da fühlt, daß seine fleischliche Lust sich wider seinen rechten Willen auflehne durch die Gewohnheit der Sünden, durch die, wenn sie nicht gezügelt ist, er mit Gewalt in Gefangenschaft geführt wird, so gedenke er, so viel er vermag, einen wie großen und herrlichen Frieden er durch Sündigen verloren habe, und rufe aus [Röm. 7, 24. f.]: „Ich elender Mensch, wer wird mich erlösen von dem Leibe dieses Todes? Ich danke GOtt durch Jesum Christ, unsern HErrn.“ Denn so, da er sich als einen elenden Menschen ausbreitet, fleht er durch sein Trauern die Hülfe des Trösters an, und die Erkenntniß seiner Unseligkeit ist ein nicht geringer Zugang zur Seligkeit. Dies ist ein goldener Ausspruch. Denn er sagt nicht, daß man lachen, Gefallen haben und fröhlich sein soll im Kitzel, sondern Leid tragen, weil er das Gift des Todes in sich fühlt. Ferner, wo sind die, welche den freien Willen rühmen? Warum lassen sie nicht sofort ab, wie sie wollen? ja, warum wollen sie nicht, und können nicht wollen? Warum konnten St. Hieronymus und Paulus den Stachel des Fleisches nicht ablegen? Denn wenn sie gewollt hätten, so wäre es schlechterdings geschehen, aber sie konnten auch nicht genugsam wollen; darum sollten sie wider ihr Wollen. Ebenso steht es auch mit dem Zorn. Denn warum halten die, welche zürnen, damit nicht alsbald ein, wenn sie können, ja, wenn sie wollen? Aber sie können nicht, noch wollen sie, weil ihr Wille bereits anderswohin gerissen und gefangen geführt ist.

Drittens\footnote{Dieses „Drittens“ entspricht dem „[erstens]“ und „zweitens“ zu Ende des vorletzten Absatzes.} ist das unser Trost, daß die, welche hier Leid tragen, getröstet werden sollen [Matth.

\quellseite{652}
5, 4.]. Denn in solcher Weise führt uns die allerweiseste Barmherzigkeit Gottes in wunderbarer Weise heraus, indem sie mit Gift das Gift vertreibt, und uns aus der Unkeuschheit ein Pflaster der Keuschheit macht, da er dieselbe um deswillen einbrechen läßt, damit der Mensch sie fühle, seufze, Leid trage, die Gnade suche, und so desto größere Neigung zur Keuschheit erlange. Denn er ist um so keuscher, nicht, je ruhiger er ist, sondern je mehr er Leid darüber trägt, daß er nicht keusch sein kann, wie er gern möchte. Und hieraus entnehmen wir (wiewohl dies alles zur vierten Stufe gehört, davon hernach gehandelt werden wird, dennoch, damit es uns nicht entfalle, da sich gerade die Gelegenheit bietet,)\footnote{Diese Klammern sind von uns gesetzt.} daß diejenigen getröstet werden sollen, welche fühlen, daß sie von solchen Stacheln gequält werden, damit sie nicht aus einem thörichten Wahne alsbald die Hände sinken lassen und sich der Verzweiflung hingeben, als ob der nicht keusch sein könne, der in seinem innersten Mark und Gedankenbrunst leidet. Dieser Gedanke ist sehr gefährlich und überaus geneigt zu allen Sünden im Werk, wegen der Verzweiflung. Denn so sagt ein junger Mensch, der vorher die fleischlichen Gedanken nicht erfahren hat, und meint, er könne nicht keusch sein, wenn er nicht ohne dieselben sei, alsbald: Was zögere ich, das Werk zu thun? denn es ist einerlei, ob ich es auch thue, denn die Keuschheit des Herzens ist doch verloren. Nicht also, du elender Mensch, nicht also; die Keuschheit ist durch solche Gedanken nicht verloren. Vielmehr bist du niemals keuscher gewesen, wenn du nur das dabei empfindest, daß du sie wider deinen Willen habest. Denn die Keuschheit, gleichwie auch jede andere lebendige Tugend, ist so verborgen, daß sie von dem nicht erkannt werden kann, der sie hat, sondern allein von GOtt. Denn GOtt verbirgt das ganze Leben der Heiligen so tief, daß sie es selbst nicht wissen können, wie der erste Psalm, V. 6., sagt: „Denn der HErr kennet den Weg der Gerechten.“ Deshalb, gleichwie die lebendige und wahre Ehre des Gerechten in der Schande besteht, die wahre Weisheit in Thorheit, die wahre Ruhe in Trübsal, die wahre Freude in Trauer, die wahre Freiheit in Gefangenschaft, der wahre Reichthum in Armuth, so besteht auch die wahre Keuschheit in Unkeuschheit, und je abscheulicher die Unkeuschheit ist, desto herrlicher ist die Keuschheit. Ach, daß diese Weisheit den Menschen unbekannt ist, wie großen Schaden thut doch das! Ein rechtes Zeichen einer lebendigen Keuschheit ist aber dies, wenn der Mensch fühlt, daß er jetzt ein Mißfallen habe, nicht daß es ein völliges Mißfallen sei (denn sonst würde der Kübel nicht anhalten), sondern ein gemischtes: Bald will er, bald will er nicht, bald liegt er oben, bald unten, so daß er, gleichwie ein Rad, durch seine Gedanken im Koth gewälzt wird, und doch der Vorsatz der Keuschheit bleibt. Denn wenn da nicht die lebendige und wahre Keuschheit wäre, würde er nicht beharren noch widerstehen, sondern ließe seinen Gedanken den Zügel frei schießen, und würde keine Beschwerde fühlen. Denn weil er will und sich ergötzt an den Gedanken, und dennoch zugleich fühlt, daß er nicht will, werden sie ihm gemischt, süß und bitter, und dadurch beschwerlich, da es ihm nach seiner Seite frei ist zu gehen, sondern er in einem mittleren Zustande herumgerollt wird. Denn der geistliche Mensch, das ist die Neigung zur Keuschheit, bleibt und bewahrt die Keuschheit, wiewohl der äußere Mensch in den Gliedern und im Herzen wunderliche Stürme erregt. Darum geht es im Kampfe der Keuschheit nicht anders zu, als wenn ein Schifflein im Meere durch die Fluthen hin und her geworfen wird, während Christus darin schläft. Darum muß man nur dafür Sorge tragen, daß er aufgeweckt werde, und gebiete dem Meere, das ist dem Fleische, und dem Winde, das ist dem Teufel. Denn er gibt auch den Priestern oft solche Unreinigkeiten ein, um sie vom Altar abzuhalten; den soll man verachten, nicht allein in den bloßen schändlichen Gegenständen, die er in Gedanken erregt, sondern auch, wenn er den Willen dazu hinneigt und böse Neigungen hervorbringt; ja, desto mehr muß man zum Sacramente laufen, um Gnade und Hülfe zu erlangen. Aber ich kenne jetzt viele, die nicht haben hinzu gehen wollen, es sei denn, daß sie nicht allein Ruhe hätten vor solchen Neigungen, sondern auch frei wären von diesen schändlichen Gegenständen, als allzu thörichte gehorsame Diener des teuflischen Betruges und seiner Bosheit.

Also in solcher Weise gebraucht GOttes Erbarmen und Güte unseres Uebels, daß er durch dasselbe unser Bestes mehr befördere, so daß
\quellseite{653}
[Röm. 8, 28.] den Auserwählten alles zum Besten dienen muß, auch die Sünden; dies alles [thut GOtt], um das Leben und die Gnade, die in uns ist, zu verbergen und zu bedecken, damit wir nicht lauarm werden, wenn wir sie erkennen, und uns damit zufrieden geben; sondern vielmehr, da wir nicht davon wissen, gleichsam als ob wir sie nicht hätten, desto mehr darnach trachten und seufzen. Wie aber von der Keuschheit gesagt ist, so ist es auch mit allen anderen Tugenden, welche dann am meisten geschenkt werden, wenn sie versagt werden, und dann da sind, wenn man meint, daß sie am allerfernsten seien, und dann der Mensch nicht anders glauben kann, als daß sie nicht vorhanden seien, und er verloren sei: daher kommt dann das Seufzen, der Haß gegen sich selbst, das Verlangen nach den Tugenden, und das Flehen um dieselben. Dann gibt GOtt den Demüthigen seine Gnade, welche er nachher wiederum verbirgt und wegnimmt, um größere Gnade hinzuzufügen, indem er das Gegentheil unter dem Gegentheil verbirgt. Dies ist deshalb die Weisheit, welche aus dem Verborgenen gezogen wird und im Verborgenen ist.

Diese Keuschheit aber ist golden im Vergleich zu der, welche ohne böse Gedanken ist, gleichwie Geduld und Sanftmuth, ohne daß Anreizung [zum Zorn] da ist. Denn was wäre das für eine Sanftmuth? Denn eine solche Sanftmuth haben auch die wilden Thiere und wenigstens die Heiden. So ist auch die Keuschheit, welche ohne Geilheit und Gedanken in der Ruhe ist, selbst bei Huren und den schlimmsten Hurentreibern, die dann keusch sind, wenn das Fleisch, ersättigt, stille ist. Uebrigens gehört diese Ruhe dem zukünftigen Leben an. Darum ist die Keuschheit gefährlich, da sie zur Hoffahrt führen mag, wenn sie offenbar ist (das heißt, wenn sie ruhig ist); sie mag aber auch gefährlich sein zur Geilheit, wenn sie verborgen ist (das heißt, wenn sie angefochten wird). Da nun auf beiden Seiten Gefahr ist, so ist doch die Gefahr des Hochmuths größer als die der Unzucht.

Darum geschieht beides. Das Gebot wird geistlicher Weise nicht erfüllt, so lange noch böse Lust in uns ist; von dieser kommt es her, daß niemand vermessen sein kann wegen seiner Keuschheit, und sich überall nicht rühmen kann, daß er dieses Gebot erfüllt habe, sondern sich nur selbst anklagen und seine Sünde bekennen. Wiederum wird es dadurch erfüllt, daß der Geist diese böse Lust haßt. Deshalb ist er ganz rein von ihr und in der That ein Erfüller des Gebotes, weil er in allem dem Gesetze beistimmt, es gutheißt und liebt, indem er sein Elend haßt.

Andere aber, unzüchtige Leute, verbergen ihre Keuschheit nicht, sondern löschen sie vielmehr aus. Es ist also nichts, daß sie sich rühmen könnten, sie kennten ihre Keuschheit nicht. Sie gehören nicht hieher.

\footnote{Löscher: „Das Folgende hat Luther am Tage der Empfängniß Mariä [den 8. December 1516] gepredigt.“ — Das Exordium hierzu findet sich Walch, St. Louiser Ausgabe, Bd. XII, 1788.} Das zweite [unterschiedliche Werk] ist die gemeine Hurerei (meretricium), welche ein großes Uebel ist, da die Unkeuschheit übergeht in eine Lebensweise (habitum) und Gewohnheit, welche sehr schwer abgethan werden (curantur). Denn die Schrift sagt [Ps. 137, 9.]: „Wohl dem, der seine\footnote{In der Vulgata und in unserer Bibel: „deine“.} jungen Kinder nimmt, und zerschmettert sie an den Stein.“ Wer aber seine kleinen Kinder machen läßt zu großen Riesen, wer sie alt und stark werden läßt, wie kann er sie dann zerschmettern, und wird nicht vielmehr zerschmettert? Und diese überaus edle Lehre des Heiligen Geistes sollten die Eltern ihren Kindern einflößen, daß dieselben nicht allein mit Gewalt im Zaume gehalten, sondern auch williglich zur Keuschheit gezogen würden. Denn vergeblich ist der Kampf der Keuschheit, wenn nicht diese Regel des Heiligen Geistes gehalten wird; nämlich der Fels ist Christus, unsere jungen Kinder sind die bösen Regungen, die bösen Begierden. Wenn jemand die empfindet, was soll er thun? wie soll er sie überwinden? Höre: Nicht anders, als daß er sie an dem Felsen zerschmettere. Wie geschieht das? Gleichwie jemand mit der Hand ein Würmlein an einem Steine zerreibt, so wird, wenn der Wille und die Gedanken Christum berühren, sofort der böse Gedanke zermalmt. Versuche es nur, und du wirst sehen, wie lieblich es sei, so die Begierden zu zermalmen; denn sofort, wie das Herz Christum, den Gekreuzigten, berührt, weicht alles Uebel. So ist es einst im Gesetze, 4 Mos. 21, 9., abgebildet worden, da die von den feurigen Schlangen Gebissenen durch das Ansehen der ehernen Schlange, die am Holze hing, geheilt wurden. Ferner ist die feurige Schlange jene schlüpfrige

\quellseite{654}
und sanft schleichende Wollust, welche die Gedanken vorspiegeln: Fliehe, fliehe diese grausame Schlange, denn sie entzündet und macht das Fleisch außerordentlich geil und voller Brunst. So heißt es auch 1 Mos. 3, 15. [Vulg.]: „Sie [das Weib] wird dir den Kopf zertreten“, das heißt, die gläubige Seele zertritt den Anfang und die erste Regung; sonst, wenn sie groß geworden oder ganz hineingeschlüpft ist, vergiftet und tötet sie. O wenn wir doch nur zu der Zeit, wenn die Unkeuschheit uns anficht, an das Bild der giftigen, kriechenden Schlange denken könnten und die geistliche Deutung (mysterium) derselben in Anwendung bringen, dann würden wir mit unseren Gedanken leicht zu Christo fliehen.

Die dritte Art ist die Schändung, das ist das Schmähen einer Jungfrau, was früher im Gesetze milder bestraft wurde als die einfache oder die gemeine Hurerei, weil der Mann gehalten war, sie auszusteuern oder sie zu heirathen. Aber jetzt ist es etwas sehr Großes; denn damals war die Jungfrauschaft eine überaus große Schande, jetzt aber ein sehr hoher Ruhm, damals eine verdammliche Unehre, jetzt aber eine unvergleichliche Zierde, wenigstens für die, welche nicht den Vorsatz gehabt haben, zu heirathen. Aber diejenigen, welche Willens sind, sich in die Ehe zu begeben, haben in der That (actuale) eine so große Zierde nicht, aber könnten sie haben, wenn sie zur Jungfrauschaft angefeuert würden. Darüber haben viele gar herrliche Dinge geschrieben, weil in der That keiner Jungfrau eine volle Wiedererstattung zu Theil werden kann. Verstehe dies aber recht: Gegen eine Jungfrau, die eingewilligt hat, wird der Schänder nicht so hart zur Strafe gezogen (tenetur); aber wider ihren Willen kann ihr die Jungfrauschaft nicht genommen werden; sie kann aber dazu bewogen werden, und so kann sie ihr genommen werden.

\footnote{Löscher: „Das Folgende trug Luther vor am 3. Adventssonntage [den 14. December] des Jahres 1516.“ — Das Exordium hiezu findet sich Walch, St. Louiser Ausg., Bd. X, 1212.} Die vierte Art ist die Entführung. Diese wird in den Rechten mit Todesstrafe belegt. Sie betrifft nicht allein eine Jungfrau, Tochter oder Magd, sondern auch das Eheweib oder eine andere Person, die unter seiner Obhut steht; aber das größte Vergehen ist die Entführung eines Eheweibes. Es geschieht aber die Entführung bisweilen mit Einwilligung des Mädchens, indem sie nämlich durch Schmeicheleien bewogen wird; bisweilen wider ihren Willen; dies ist das schwerste Verbrechen. Aber man liest viele Fälle, wo es nicht wohl ablief. Guarinus von Verona hat ein elegisches Gedicht gemacht über einen kläglichen Fall, der sich mit einer gewissen Jungfrau Alba zugetragen hat. Diese verließ ihre Eltern, folgte ihrem Entführer, und ist endlich von ihm in einem Walde geschändet und verstoßen. Da sie sich nun schämte, zu ihren Eltern zurückzukehren, bat sie ihn, er möchte sie mit dem Schwerte erstechen. Dies that er auch und tödtete sie. Diese Sünde begreift auch einen Diebstahl in sich, und zwar den allergrößten.

Die fünfte Art ist der Ehebruch. Dies ist das Verwickelste und Schwerste, was die Strafe anbelangt; aber es wird [die Strafe] heutzutage allzu sehr vernachlässigt. Diese Sünde ist die Ursache vieler Fragen, welche das Eigenthumsrecht anbetreffen; ob die Frau dem Manne aufsagen solle, daß das Kind ein Bastard sei. Siehe den Angelus in der Summa,\footnote{Diese Summa angelica nennt Luther in der großen Auslegung des Galaterbriefs (Walch, St. Louiser Ausg., Bd. IX, 532 § 168) eine „teuflische“.} welcher Regeln darüber gibt. Aber wer wird wagen, den Regeln zu glauben, oder wer kann einen Fall in Schranken fassen? Denn die Gemüther der Männer sind verschieden; heute denken sie so, und werden nach einem Jahre vielleicht anders denken, selbst wenn sie in diesem Jahre Geduld hätten mit ihrem Weibe, die ihren Ehebruch bekennt.

Aber es ist sicherlich eine gegründete Ursache, darüber erbittert zu werden, daß der Mann seinen Leib für den Leib seines Weibes gegeben hat, und ein anderer den Leib gebraucht, für den er seinen eigenen Leib gegeben hat, um dessentwillen er sein Gewerbe ausübt, um dessentwillen er alles thut, leidet; ja, auch zum Sclaven hat sich der gemacht und sich zu vielen verbunden, der sich an ein Eheweib gebunden hat, und siehe, ein anderer untersteht sich, dies alles zu besudeln: wer könnte das wohl ruhig ertragen? Siehe auch das Eifergesetz [4 Mos. 5, 12. ff.] an, welches einst bestand; dessen sind nun die Weiber enthoben und nicht mehr damit belastet; darum sollten sie desto lieber an ihren Männern hangen.

\quellseite{655}
Die sechste Art ist incestus, was auf deutsch „Blutschande“ genannt wird, weil es eine Schande ist gegen die Blutsfreundschaft und die Verwandtschaft, der um so größere Ehre gebührt, je näher sie ist. Denn auch im Gesetze hat der HErr gewisse Grade verboten, und einst waren sieben Grade verboten. Denn im Gesetze wird fast nur der zweite Grad der Blutsfreundschaft und der Verwandtschaft verboten, weil noch nicht die Gnade da war, welche die böse Lust heilt, sondern das Gesetz herrschte, welches dieselbe reizt und vermehrt.

Die siebente Art ist das Sacrilegium, wo bereits nicht allein die Keuschheit verletzt wird, sondern auch das, was allein Gotte gewidmet war, aufgehoben und das Heilige entweiht wird. Aber dies ist bei den Priestern mehr nach einer Verordnung der Kirche als nach der Ordnung GOttes; aber bei denen im geistlichen Stande (religiosis = Mönchen und Nonnen) ist es eine überaus schwere Sünde, weil sie sich freiwillig dem HErrn geweiht haben, und sich ihm wiederum entziehen.

Als achte Art wird [von etlichen] hinzugefügt das Uebermaß der Eheleute, wie St. Ambrosius angeführt wird im zweiten Buche wider den Julianus des heiligen Augustinus: Ein allzu hitziger Liebhaber seiner eigenen Ehefrau ist ein Ehebrecher. Davon sagt der Apostel 1 Thess. 4, 3. ff.: „Das ist der Wille GOttes, eure Heiligung, daß ihr meidet die Hurerei“; dies ist das erste, danach: „Daß ein jeglicher wisse sein Faß (das ist sein Ehegemahl) zu behalten in Heiligung und Ehren, nicht in der Lustseuche, wie die Heiden, die von GOtt nichts wissen.“ Und 1 Petr. 3, 7.: Desjenigen gleichen, ihr Männer, wohnet bei euren Weibern mit Vernunft, und gebet dem weiblichen, als dem schwächsten Werkzeuge, seine Ehre“ etc., davon oben genug gesagt ist. Aber es ist schwer, denselben eine Regel zu geben, ausgenommen, daß es Lasterhaft ist, bei jeder Regung des Fleisches seinem Gelüsten und Begierde genug zu thun, denn das heißt in der Lustseuche sein.

Jetzt folgen drei andere Arten: Weichlichkeit, da sich jemand machend selbst befleckt; Sodomerie, da ein Mann des andern, ein Weib des andern mißbraucht, oder auch der Mann eines Weibes oder ein Weib des Mannes; dies alles ist zu schändlich, als daß man es sagen dürfte; viehisch Weisen; davon siehe 3 Mos. 18, 23. Das andere gehört in die Beichte.

Ueber das sind noch die nächtlichen Befleckungen, welche an sich nicht Sünde sind, wiewohl das Gesetz 5 Mos. 23, 10. f. sagt: Wenn unter euch ein Mensch ist, der des Nachts von Samen befleckt ist, so soll er aus dem Lager gehen, und nicht zurückkehren, bis daß er sich gegen den Abend mit Wasser bade, und nach Sonnenuntergang mag er in das Lager zurückkehren. Dies hat eine geistliche Bedeutung: ein so Befleckter (wiewohl er auch noch anderes bedeutet) soll sich enthalten der Gemeinschaft der Heiligen (das ist des Sacramentes), bis daß die Sonne (das ist die Hitze und Brunst, die von der Befleckung übriggeblieben ist) zur Ruhe gekommen ist, und so, nachdem er sich gewaschen hat durch demütige Beichte, hinzutreten. Diese [Befleckung] hat aber öfter eine Ursache, welche eine Sünde ist, als Völlerei, Müßiggang, Verkehr mit dem andern Geschlechte oder davon gedenken.

Endlich gibt es auch Befleckungen im Wachen und bei Tage, aber ohne ihren Willen, was vielen widerfährt, während sie mit etwas ganz anderem zu schaffen haben, als, sich kratzen, reiten, fahren, das heißt, ohne daß sie es veranlassen. Von diesen sagt die Schrift 3 Mos. 15, 2. ff.,\footnote{Hier hat die Weimarische Ausgabe, ebenso wie Walch, das falsche Citat 3 Mos. 15, 16.} daß sie einen Fluß haben.

\footnote{Löscher: „Dies hat Luther am 4. Adventssonntage [den 21. December] 1516 gepredigt.“ — Das hiezu gehörige Evangelium findet sich Walch, St. Louiser Ausgabe, Bd. XII, 1790. Dasselbst hätte als Text nicht Joh. 1, 19. ff., sondern „Luc. 3, 5.“ gesetzt werden sollen.} Es folgt nun die zweite Stufe des Ehebruchs, welche in Worten besteht, wie denn etliche überaus schändliche Reden führen von den fleischlichen Sünden, dawider der Apostel redet Eph. 5, 3. f.: „Hurerei aber und alle Unreinigkeit, oder Geiz, lasset nicht von euch gesagt werden, wie den Heiligen zustehet; auch schandbare Worte und Narrentheid inge, oder Scherz, welche euch nicht ziemen, sondern vielmehr Danksagung“, das heißt, laßt das Gerücht nicht aufkommen, daß ihr solche Leute seid.

Er hat aber drei Sünden der Zunge gesetzt, nämlich in den die böse Lust erregenden „schandbaren Worten“, das ist Unflätigkeit, Unzüchtigkeit, Schamlosigkeit in Worten. Darin wird zwiefach gesündigt: Erstens von denen, die aus diesen schandbaren Worten ein Gelächter zurichten und sich
\quellseite{656}
gegenseitig gleichsam Muth einsprechen (consolationem), als ob es eine liebliche Sache wäre, von der sie reden. Hier thut einer nichts Anderes, als daß er zeige, wie unflätig auch er sein könne, indem er mit aller Schandbarkeit die Glieder und die Werke der Glieder mit Namen nennt; ein anderer aber erzählt Geschichten, die hierzu gehören, in aller Unverschämtheit; dann lacht man, und es wird für etwas ganz Köstliches angesehen, bis man zuletzt auch die Mönche und die Priester mit ihren Sünden heranzieht. O wie erschrecklich herrscht doch dieser Unfug (morbus) heutzutage in den Gasthäusern und Herbergen, so daß es gefährlich oder wenigstens bedauernswerth ist, wenn ein Geistlicher oder auch sonst ein züchtiger Mann über Land reisen muß.

Ich bitte euch um GOttes willen, sehet doch, was diese Leute thun. Erstlich, wenn sie Christen sind, wie sie doch genannt werden wollen, so sollten sie ein züchtiges Wesen so lieben, daß sie auch nicht zuließen, daß solche Dinge unter ihnen genannt würden, wenn sie es hindern könnten. Zweitens sollte es ihnen wenigstens leid thun, wenn sie es nicht hindern könnten. Drittens sollten sie über die Erzählungen von Gefallenen, seien es nun Priester oder Mönche oder Fürsten oder obrigkeitliche Personen (denn dieser Leute Fehler werden um so lieber erzählt, je stärkeres Beispiel sie denen geben, die ihnen nachfolgen, und je zügelloser sie sich verhalten haben), sich nicht anders betrüben, als über ihren eigenen Fall (denn das fordert die Liebe). Ja, je höher ein Stand ist, desto mehr sollte man für sie seufzen und GOtt für sie bitten, denn das wäre christlich.

Aber jetzt sehe, wie weit diese ganz abgeschmackten Leute von christlichem Verhalten entfernt sind. Erstlich thun sie nicht allein dieses nicht, was sie doch sicherlich Christo schuldig sind, sondern zweitens ergötzen sie sich gegenseitig mit den allerbösesten Dingen. Drittens kitzeln und reizen sie sich selbst mit diesen feurigen Pfeilen zur bösen Lust. Viertens beflecken sie das Leben und den guten Ruf anderer Leute durch diese überaus schwere böse Nachrede, indem sie den Stand der Todten aufrühren, während sie doch die Todten begraben sollten. Fünftens verspotten sie ihr eigenes und anderer Leute Verderben und Elend und sind fröhlich über ihre und anderer Leute Verdammniß.

Ich bitte euch, bedenket, wenn der HErr sagt [Matth. 12, 36.], daß die Menschen am jüngsten Gericht Rechenschaft geben müssen von einem jeglichen unnützen Worte, das sie geredet haben, Lieber, wo wollen denn diese bleiben? Wehe, wehe, wie weit regieret diese Pest!

Diese Leute sind gleich, erstens, den Unsinnigen, ja, wer ist so unsinnig, daß er in seinem Elende und Plage lache und fröhlich sei? Ist das nicht ein Jammer über allen Jammer? Schmerzt ein Geschwür oder ein Auge, so hört das Klagen und Seufzen nicht auf, und jedem ist das Lachen so gänzlich vergangen, daß er auch die, welche lachen, als die allerbößesten Leute haßt und verabscheuet; man sucht nur einen, der Mitleiden hat und hilft oder guten Rath gibt. Aber siehe, wenn das Fleisch kitzelt, ja auch das Gewissen an den schwersten Schmerzen leidet, ferner deine eigene Seele oder die eines anderen, deines Nächsten, dahinfällt: da wird so gar nicht gesucht nach einem, der Mitleiden habe oder Rath gebe, so gar wird nicht Leid getragen und geseufzt, daß auch derjenige, welcher Mitleiden hat und rathen oder hindern will, verspottet, verachtet und gehaßt wird. Großer GOtt, wann werden mir die Augen öffnen, daß wir erkennen, was die Welt sei? Ist dies nicht eine Verkehrtheit über alle Maßen? Ist nicht die Unkeuschheit die schlimmste aller Wunden? Ist sie nicht ein Geschwür und eine Plage, die über alles Geschwür und Plage geht? Wie Jeremias, Cap. 10, 19., sagt: „Ach, meines Jammers und Herzeleid's! Ich denke aber, es ist meine Plage, ich muß sie leiden.“ Und Jesaias, Cap. 1, 6., sagt von dem Volke Israel: „Die Wunden, und Striemen, und Eiterbeulen sind nicht geheftet, noch verbunden, noch mit Oel gelindert“, und wiederum: „Von der Fußsohle bis zum Scheitel ist nichts Gesundes an ihm.“ Dies sind Klagen und Wunden, die nicht von einem Menschen, sondern vom Teufel beigebracht worden sind. Siehe also, eine Wunde, die ein Mensch macht, wird beklagt, aber eine solche, die der Teufel macht, wird belacht und für etwas gar Liebliches gehalten. Aber ach, daß es für etwas Liebliches gehalten wird! Selig aber ist der, der die Wunde erkennt und merkt, daß es eine gar bittere Wunde ist, wie sie denn in Wahrheit ist. Denn wiewohl ihre Bitterkeit nicht empfunden wird, so wird sie doch endlich im Gewissen außerordentlich
\quellseite{657}
empfunden werden. Daher heißt es. Spr. 5, 2. ff. [Bulg.]: „Laß dich nicht berücken durch den Trug eines Weibes (das ist, des Fleisches und der Sinnlichkeit), denn die Lippen der Hure sind süße wie Honigseim, und ihre Kehle ist glatter denn Oel; aber hernach bitter wie Wermuth und scharf wie ein zweischneidig Schwert.“ Und Cap. 23, 31. f.: „Er [der Wein] gehet glatt ein, aber darnach beißt er wie eine Schlange und sticht wie eine Otter.“ Darum sollen wir von GOtt lernen, daß, wenn man mit solchen Worten Scherz treibt, darüber lacht, sich daran ergötzt, dasselbe ist, als wenn man sich an seinem eigenen Tode und Jammer freute, was nur ein Unsinniger thut, oder einer, der mehr ist als unsinnig.

Zweitens sind sie dem gleich, von welchem Marci 5, 2. ff. und Luc. 8, 27. ff. geschrieben steht, daß er von einer Legion [Teufel] besessen war, der keine Kleider anhat, auch in keinem Hause blieb, sondern in den Gräbern, und da er mit Ketten gebunden ward und in Gewahrsam gethan, zerriß er die Ketten, und wurde vom Teufel in die Wüste getrieben, und war Tag und Nacht in den Gräbern und in den Bergen, schrie, und schlug sich mit Steinen. Ich hoffe, wenn die, welche schändliche Reden führen, einsehen, daß ihr Laster an diesem Menschen beschrieben werde, so würden sie sich mehr fürchten. Denn was ist es, daß er nackt war, als daß ein Mensch, der schandbare Worte redet, auch seine eigene Schande mit Worten entblößt, während er sie doch vielmehr bedecken sollte, damit er geheilt würde? Sodann bleibt er nicht im Hause, nämlich bei den Lebendigen, und hat mit den Angelegenheiten der Lebendigen zu schaffen, sondern in den Gräbern, und hat mit Todten zu thun, das heißt, er erzählt die Laster und bösen Thaten anderer, und lebt und webt ganz und gar darin. Denn ein Grab, geistlich gedeutet, ist das Gedächtniszeichen eines Todten, das heißt, der üble Ruf und die garstige Geschichte eines Sünders, nämlich eines Fürsten, Mönchs, Priesters etc. Desgleichen, er schlägt sich selbst mit Steinen und schreit, das heißt, er schadet sich selbst, da er meint, sich am meisten und in angenehmer Weise zu nützen. So heißt es Jes. 3, 9.: „Sie haben ihre Sünde nicht verhohlen, sondern haben sie gerühmt, wie die zu Sodom.“ Wiewohl dies von der pharisäischen Gerechtigkeit im Geiste gesagt ist, so gehört es doch auch dem Buchstaben nach hierher. Desgleichen, er hat eine Legion [Teufel], so daß dort niemand vorübergehen kann wegen ihres Wüthens, das heißt, daß sie das Werk vieler Teufel ausrichten, indem sie vielen Leuten Schaden thun. Und wo sie am meisten schmeicheln, da sind sie am grausamsten, so daß niemand ohne Gefahr mit ihnen umgehen kann, wie ich auch oben gesagt habe. Ich weiß aber, daß niemand so thöricht und unsinnig ist, daß er lachen könnte, wenn er einen solchen Menschen sähe, sondern er würde vielmehr so weit fliehen, als er könnte: und doch lacht der, und man lacht über ihn, dessen Figur jener Mensch nur gewesen ist. Endlich ist auch dies dort gesagt worden, daß er mit Ketten und Gewahrsam nicht gehalten werden konnte, ja, daß er die Bande zerriß, und in die Wüsten getrieben wurde; niemand konnte ihn bändigen. Dies stimmt fast mit dem Apostel Jacobus überein [Jac. 3, 8.]: „Die Zunge kann kein Mensch zähmen.“ So lassen sie sich durch keine Menschen verbieten, durch keine Worte des Evangelii im Zaume halten, daß sie nicht mit ungezähmter Zunge hervorbrechen sollten und umherschweifen, wo sie umherzuschweifen pflegen, nackt und zerrissen, und entblößen und zerreißen andere Leute in gleicher Weise. Und daß ich wider dieses ganz abscheuliche Laster auch grob heraus rede: Sind nicht die Schweine unreine Thiere, weil sie Menschenkoth fressen? Aber jene Leute wühlen mit Zunge und Zähnen in ihrem eigenen und anderer Leute Kothe. Ist denn etwa der Koth des Leibes ärger als der Koth der Seele? Denn man muß mit den Säuen säuisch reden, damit sie es verstehen. Wenn du nun siehst, daß ein Schwein Unflath frißt, so gedenke: Siehe, ein Mensch, der unfläthige und schandbare Worte redet, ist ein solches Vieh. Der Apostel sagt sehr recht [1 Cor. 12, 23.], daß wir den unehrlichsten Gliedern am meisten Ehre anthun; jene aber entehren auch die ehrlichen Glieder durch die unehrlichen. Aber könnte es aber leiden, daß ihm ein anderer mit Gelächter seinen Hintern oder seine Scham oder seine Schande aufdeckte? Aber welche Schande ist größer als der Unflath der Sünden an uns und dem Nächsten? Das Gesetz aber hat verboten [3 Mos. 18.], daß wir unseres Nächsten Scham nicht bloßen sollen.

Zweitens wird auch durch unflätige Reden gesündigt, wenn solche Dinge vor unschuldigen
\quellseite{658}
Knaben und Mädchen genannt und geredet werden. Sie sind aber aller der Sünden und alles Blutes (mortium) schuldig, das daraus erfolgt. Denn das zarte und in solchen Dingen unerfahrene Alter wird überaus leicht befleckt, und, was das Schlimmste ist, es behält solche schändlichen Dinge, die es gehört hat, sehr fest, gleichwie ein Flecken, der in ein reines und zartes Gewand kommt, fester anhaftet, als wenn er auf ein unreines oder gröberes Zeug kommt, was auch die Heiden aus der Erfahrung gelernt haben, wie Horaz sagt: Ein neuer Topf behält gar lange den Geruch dessen, was zuerst hineingethan wird, und Juvenal: Nimmst du etwas Schandbares vor, so hüte dich mit allem Fleiße, daß nicht ein Kind dadurch geärgert werde; verachte ja nicht seine Jugend.

Was thun aber die, welche sich solches zu Schulden kommen lassen? Erstlich, da es etwas überaus Gutes ist, daß man das zarte Alter wohl in Acht nehme, es bewahre und fördere in Keuschheit und schamhaftem Wesen, wie man sagt, daß es Gotte das angenehmste Opfer ist, daß man um die Seelen sich eifrig annehme, so sollten alle durchaus den größten Fleiß daran setzen und machen, daß sie die Knaben und Mädchen behüteten und bewahrten, daß sie ja nicht etwas lernten, sähen und hörten, was unzüchtig ist, da sie so schon im innersten Mark voller Brunst sind, und man muß nicht darauf bedacht sein, wie man nur Feuer hinzutrage und sie entzünde, sondern daß man vielmehr Wasser über sie schütte und lösche. Aber das thun sie nicht allein nicht, sondern richten auch des Teufels recht eigentliches Werk aus und verderben die unschuldigen Seelen durch ihre giftigen Schwätzereien von schandbaren Dingen. Denn der Teufel wird ein Verderber der Seelen genannt, er verderbt sie aber nicht anders als durch den Dienst der Zunge jener Leute, welche ihm anhängen und ihm nachfolgen. Denn wie kann ein Knabe oder ein Mädchen ein schandbares Wort wieder ausrotten, welches es einmal gehört hat? sondern der Same ist ausgestreut und wurzelt im Herzen auch wider seinen Willen; darnach wächst er zu wunderlichen und seltsamen (incognitas) Gedanken, die er nicht zu bekennen mag, und weiß doch nicht ihrer loszuwerden. Wehe aber dir, der du diesem einfältigen Gemüthe, das solche Dinge nicht kannte, diese Last, diese Gefahr und dieses Gift eingegossen hast! Dem Leibe hast du keine Gewalt angethan, aber die Seele, die unvergleichlich viel edler ist, hast du (so viel an dir ist) geschändet. Denn durch das Ohr hast du es ihm eingegossen und ihm die Seele mit der allerschädlichsten Frucht (foetu) geschwängert. Daher sagt Baptista von Mantua: Durch die Sinne, gleichsam als Bäche, kößt Venus ihr scharfes Gift ins Gemüth; wenn dies es unvorsichtig aufgenommen hat, verursacht es eine Wunde, die nur durch die Kraft (dextra) der Gottheit geheilt werden kann.

Daher ist die Schandbarkeit (obscoenitas) dieser Leute größer als derer, von denen vorher geredet ist, weil es bereits nicht mehr Schandbarkeit, sondern ein Mord ist unschuldiger Seelen, und sie gehören dem Geschlechte Herodis an, welcher viele Kindlein in Bethlehem Juda tödtete [Matth. 2, 16.]. Aber siehe, ist es nicht so, daß kein Vater leiden könnte, daß man seine Kinder nach dem Leibe tödtete? Da du nun nicht wolltest, daß es dir geschähe an dem Leibe deiner Kinder, warum unterstehst du dich denn, dies an den Seelen der Kinder eines anderen, ja, der Kinder GOttes zu begehen? Der heilige Ludwig, König von Frankreich, hat gar wohl (recte) von seiner Mutter gelernt, daß sie lieber sehen wollte, daß ihre Kinder leiblich getödtet würden, als daß sie eine Todsünde begingen. Denn in solcher Weise hatte seine Mutter auch zu ihm geredet.

Aber gar entsetzlich droht der HErr solchen Leuten Matth. 18, 6., indem er die furchtbaren Worte spricht: „Wer aber ärgert dieser Geringen einen, die an mich glauben, dem wäre besser, daß ein Mühlstein an seinen Hals gehängt, und er ersäuft würde im Meer, da es am tiefsten ist.“ Siehe, wie sich Christus mit großer Sorge der kleinen und unschuldigen Kinder annimmt, daß er auch eine neue Strafe für die bestimmt, welche ihnen Aergerniß geben; denn für niemand anders hat er eine solche Strafe aufgelegt, indem er nämlich seinen Willen aussprechen wollte, daß sie vor anderen mit einer neuen und sonderlichen Strafe gemartert werden sollten. Dann folgt [Matth. 18, 7.]: „Wehe der Welt der Aergerniß halber! Es muß ja Aergerniß kommen; doch wehe dem Menschen, durch welchen Aergerniß kommt.“ Und hernach [V. 10.]: „Sehet zu, daß ihr nicht jemand von diesen Kleinen verachtet. Denn ich sage euch: Ihre Engel im Himmel sehen allezeit das Angesicht meines
\quellseite{659}
Vaters im Himmel.“ Siehe, er will, daß man wenigstens aus Rücksicht auf die Engel der Kinder schöne und ihnen gegenüber Zurückhaltung beweise. Wenn nun jemand solche Leute um deswillen keusch nennen wollte, weil sie mit Worten spielen, aber in der That keusch seien, wie auch ihr Vater David lügt: Glaube mir, mein Leben ist ganz anders als meine Gedichte: Mein Leben ist ehrbar, meine Muse scherzt, so antwortet Christus und stopft ihnen den Mund [Matth. 12, 34.]: „Wes das Herz voll ist, deß geht der Mund über.“ Wie der Baum, so die Blätter.

Hierin liegt auch die Ursache des verderbten Zustandes der Kirche ganz und gar; denn wenn sie jemals wieder zur Blüthe kommen soll, so ist es nothwendig, daß man an der Unterweisung der Kinder anfange. Dies ist von den schaubaren Worten gesagt. Daß man hier aber auch von „Narrentheidingen“ und „Scherz“ [Eph. 5, 4.] rede, das gehört nicht hieher, wiewohl diese beiden Dinge jetzt bei den ehrbarsten und gottesfürchtigsten Leuten (religiosissimos) so sehr im Schwange gehen, daß sie gerade durch das für ehrbar angesehen sein wollen, wodurch sie zeigen, daß sie Narrenwerk und unnützen Scherz treiben.

„Narrentheidinge“ aber sind alle unnützen Fabeln, in welchen keine Unterweisung und Wissenswerthes (scientiae) enthalten ist, wie er es selbst auslegt und sagt [Eph. 5, 4. Vulg.]: „welche nicht zur Sache dienen“. Dasjenige nennt Christus „unnütze Worte“; darum soll man nur nützliche oder nothwendige Dinge reden. Derartig sind die Altweibermärlein und die läppischen Possen, die Historien des Pfaffen von Kalenberg, Dietrich von Bern und allerlei Gerede (disputationes) von Dingen, die uns fern liegen und außer unserem Bereiche sind. Der Art sind auch die Meinungen der Philosophen und die Glossen über Aristoteles, das Geschwätz nach der Dialectik, die Träume der Astrologen. Denn er nennt diese Dinge „Narrentheidinge“, weil es Worte, nicht der Weisheit, sondern der Thorheit sind.

Eutrapelia oder „Scherz“, das ist witzige (facetia) oder feine Unterhaltung (urbanitas), welche Aristoteles für eine Tugend hält, wenn nämlich nicht schändliche noch thörichte, sondern zugleich angenehme und lehrreiche Dinge gesagt werden, so daß unter dem lieblichen Verkehr auch Nutzen sich findet. Daher ist eine solche Gesprächigkeit oder Redseligkeit, oder seine Unterhaltung oder Widerrede, welche geschieht, um Traurige und Angefochtene im Geiste aufzurichten, etwas Lobenswerthes; aber diejenige, welche ohne Ursache geschieht, nur um der Eitelkeit und dem Fleische ein Genüge zu thun, zumal da hiervon immer nur die Leichtfertigkeit die Ursache ist, und allein die Lust zu schwätzen und zu lachen dazu bewegt, geziemt den Christen nicht. So sagt auch St. Augustinus in seinen „Bekenntnissen“, daß es gar nicht christlich sei, mit Freunden Possen zu reißen; daß der Christen Wesen und Worte mit wohlgegründeten und großen ernsten Dingen zu schaffen haben müssen, die einen vollen Ernst und Zuverlässigkeit (autoritatem) erfordern, wie der Psalm [Ps. 35, 18. Vulg.] sagt: „Unter einem ernsten Volke will ich dich rühmen.“ Ferner wird diese Possenreißerei und dieser Scherz um so höher, je heiliger die Dinge sind, in welchen er sich bewegt, wie denn einige mit der dreistesten Leichtfertigkeit der Sprüche der heiligen Schrift zu diesem schändlichen Treiben mißbrauchen, indem sie in lächerlicher und witziger Weise (jucunde) dieselbe dahin verdrehen, wo sie ihre Eitelkeit suchen, während doch diesen heiligen Worten immer und überall die höchste Ehrerbietung und Furcht gebührt.

\footnote{Löscher: „Das Folgende hat Luther am Tage des Thomas [den 21. December] im Jahre 1516 geredet.“ — Das Exordium hiezu findet sich Walch, St. Louiser Ausgabe, Bd. XII, 1794.} Die dritte Stufe des Ehebruchs ist die, welche in Gebärden besteht, deren eine der HErr ausgedrückt hat, da er spricht [Matth. 5, 28.]: „Wer ein Weib ansieht, ihrer zu begehren“ etc. Siehe, selbst das Ansehen, welches ein Zeichen der bösen Lust ist, ist Ehebruch. Ist dieses Gesetz Christi nicht sehr hart? Aber darum ist es überaus heilsam und unbefleckt und bekehrt die Seelen. In dieser Geberde werden aber alle anderen begriffen, da das Sehen das kleinste und feinste Zeichen ist. Daher haben etliche fünf Unterabtheilungen (lineas) der Liebe gemacht: Sehen, Reden, Betasten, dann Küssen und die That. Und hier wird auch ein anderes Reden als die schandbaren Reden verstanden, gleicherweise auch ein anderes Betasten, desgleichen auch Küssen. Doch von den einzelnen Stücken wollen wir weitläufiger handeln. Es sind aber noch mehr, als, Sehen, Hören, Reden,

\quellseite{660}
Briefe, Zwischenträgereien, Betasten, Umarmen, Lachen, Küsse, aber nicht als ob so viel Ehebrüche da wären, als jemand Zeichen gegeben hat; denn es kann jemand in einem Ehebruche alle diese Zeichen erweisen: sonst wäre das Werk selbst immer ein dreifacher Ehebruch, da es ohne Berühren und böse Lust nicht zugeht.

Das Ansehen ist also das erste. Dies ist aber das geschwindeste und häufigste von allen diesen Zeichen, weil es gar weit und breit und rasch sein Werk ausrichtet, und sehr viel und sehr deutlich wahrnimmt. Darum wird in der Schrift so viel davon gesagt, daß man Acht haben solle auf die Augen, da kein Sinn so leicht und durch so viele Gegenstände gereizt wird, auch kein anderer Sinn die Seele so greulich verwundet; denn es heftet [das Auge] eine Sache um so tiefer in das Herz, je klarer es dieselbe erkennt, so daß man auch das Sprichwort hat: Was das Auge sieht, das glaubt das Herz. Und Horaz sagt: Das, was durch das Ohr überliefert wird, reizt weniger, als das, was die zuverlässigen Augen sehen; und wiederum: Es erregt mehr etc., wie aus der Erfahrung offenbar ist, weil der, welcher vom Tode hört, nicht so bewegt wird, als der, welcher ihn sieht, auch der, welcher schändliche, fleischliche und weltliche Dinge hört, nicht so sehr bewegt wird als der, welcher sie sieht. Als Beispiel möge Alipius dienen, der Genosse des heiligen Augustinus, über den er im 6. Buche der „Bekenntnisse“ schreibt, daß er wider seinen Willen zum Schauspiel geführt wurde, und die Augen zumachte, da er es verachtete: dennoch, als er nachher die Augen öffnete, ist er so davon hingerissen worden, daß er selbst nachher auch viele hinzuführte.

Wie es daher bei jungen Männern und Jungfrauen der höchste Schmuck ist, wenn sie die Augen senken und niederschlagen, so ist es etwas überaus Tadelnswerthes, wenn sie dieselben umherschweifen lassen, und nicht bewahren, so daß etliche das als ein Zeichen der Gefallenen (corruptionis) angeben, daß sie die Augen umherschweifen lassen; und das ist wahr. Sirach 26, 12. [Vulg.] heißt es: „Ein hurrisch Weib erkennt man an dem frechen Aufheben ihrer Augen und an ihren Augenlidern“ etc. Siehe daselbst nach. Entweder es ist ein Zeichen eines Weibes, die schon gefallen ist, oder einer solchen, die sehr leicht zu Fall gebracht werden kann; denn ein Weib, das in Wahrheit keusch ist, starrt einem Manne nicht in das Gesicht, sondern wird verschämt, wie von Rebecca 1 Mos. 24, 64. geschrieben steht, daß sie, sobald sie ihren Bräutigam Isaak sah, ihren Schleier nahm und sich verhüllte.

Desgleichen liest man von einem heiligen Manne, der seine Augen schloß, als er mit dem heiligen Antonius nach Rom gekommen war, und von dieser so großen Stadt nichts sehen wollte als die Kirche des heiligen Petrus. Desgleichen liest man ebenso von Sulpannus, daß er immer mit verhüllten Augen einherging und auch die Sonne nicht hat ansehen wollen. Desgleichen sagt Hiob, Cap. 31, 1.: „Ich habe einen Bund gemacht mit meinen Augen, daß ich nicht achtete auf eine Jungfrau“, das heißt, auf ein junges Mädchen.

Dagegen da Dina herausging [1 Mos. 34, 1. ff.], die Weiber des Landes zu sehen, ist sie zu Schanden geworden. Und da Bernhard ein gewisses Weib starr angesehen hatte, strafte er sich selbst deshalb sehr hart. Und St. Benedictus wurde sehr schwer dadurch angefochten, daß er ein Weib gesehen hatte, als er noch in weltlichem Leben (saeculo) war. Desgleichen, was begegnete David, als er die Bathseba ansah [2 Sam. 11, 2.], und den Kindern Israel, als sie die Moabitinnen ansahen [4 Mos. 25, 1. ff.]? Desgleichen, da ein anderer durch das Ansehen eines Weibes angefochten war, nun aber hörte, daß sie todt sei, heilte er seine böse Lust durch den Gestank ihres Leichnams.

Daher erregt das Ansehen nicht bloß in der Gegenwart böse Lust, sondern es bleibt auch von vergangenen Dingen der Eindruck fest und lange im Gedächtniß, und verursacht wunderliche Anfechtungen. Darum sagt Jeremias in den Klageliedern, Cap. 3, 51.: „Mein Auge frißt mir das Leben weg“, und wiederum Jer. 9, 21.: „Der Tod ist zu unsern Fenstern hereingefallen“, und 2 Petr. 2, 14.: „Sie haben Augen voll Ehebruchs.“ Siehe, wie er den Ehebruch in die Augen setzt. Desgleichen ist Eva [1 Mos. 3, 6.] zuerst durch das Anschauen des Apfels bewegt worden, daß er lieblich anzusehen wäre etc. Und von den Kindern GOttes steht 1 Mos. 6, 2. geschrieben, daß sie durch das Sehen gefallen sind, da es heißt: „Da sahen die Kinder GOttes nach den Töchtern der Menschen, wie sie schön waren, und nahmen zu Weibern, welche sie wollten.“ Daher lehrt auch der weise Mann [Sir. 9, 5.]:

\quellseite{661}
Siehe nicht in das Gesicht einer Jungfrau, damit dich ihr schönes Aussehen nicht zu Schanden mache. So soll auch wiederum ein Mädchen nicht in das Gesicht eines jungen Mannes schauen, wegen ebenderselben Gefahr.

Darum sagt Johannes [1 Joh. 2, 16.]: „daß alles, was in der Welt ist, ist Fleisches Lust, und der Augen Lust, und hoffärtiges Leben“. Unter „der Augen Lust“ wird der Vorwitz verstanden, vieles durch die Sinne zu erfahren und kennen zu lernen.

Man liest von der heiligen Elisabeth, daß sie von dem HErrn plötzlich gestraft ward, als sie ihren Gatten in der Kirche einmal lüstern (jucunde) angesehen habe. Was wird denen widerfahren, die in der Kirche fast nichts Anderes thun, als was David sagt: Sie kommen, um zu sehen, sie kommen, daß sie selbst gesehen werden; nichts als sehen, gesehen werden, begehren, begehrt werden? Aber es ist freilich durch eine Art Natural-Anlage den Weibern eigen, daß sie gerne begehrt und geliebt werden wollen, der Männer Art aber, daß sie begehren und lieben, wiewohl es auch umgekehrt geschieht, daß ein Weib voller Weise begehrt und liebt, ein Mann aber just begehrt und geliebt zu werden.

Es ist daher ein sonderbares Elend, daß uns unsere eigenen Sinne so schädlich sind, so daß Christus mit Wahrheit sagt [Matth. 10, 36.]: „Des Menschen Feinde sind seine eigenen Hausgenossen“ (das ist, seine Sinne), und der 38. Psalm [V. 12.]: „Meine Lieben und Freunde stehen gegen mir.“ Denn was thut der Vorwitz der Augen anders, als daß er wider die Seele streitet? Daher nennt St. Augustinus auch die Sinnlichkeit eine Schlange, weil sie schlüpfrig ist und leicht zu bewegen durch viele Gegenstände, welche abgebildet ist durch die Schlange im Paradiese. Deshalb siehe, daß Christus im Evangelio nicht umsonst gelehrt hat, daß wir uns selbst hassen und verleugnen sollen, daß wir unsere Augen und alle Sinne und alle Glieder hassen sollen, weil sie die Seele gefangen nehmen in das Gesetz der Sünde, besonders die Augen. Daher sagt Sirach, Cap. 31, 15. [Vulg.]: „Was ist Schalkhaftigeres erschaffen, als das Auge?“

Aber es ist nicht allein dasjenige böse, was die böse Lust durch das Auge thut, sondern auch, daß es nicht gesättigt wird, wie der Prediger Salomo Cap. 1, 8. sagt: „Das Auge siehet sich nimmer satt“, und Spr. 27, 20.: „Die Hölle und Verderbniß werden nimmer voll, und der Menschen Augen sind auch unsättig.“ Was nützt es also, daß man die Begierde der Augen reizt und, wenn der Anfang des Sehens beendigt ist, ein endloses Begehren erzeugt, welches kein Ende nimmt? Es ist besser, daß du sofort von Anfang an das Auge (visum) im Zaume haltest, ehe es den Schlund seiner Begierde so weit aufthut, daß du ihn nicht füllen kannst.

Daher steht es so, wie St. Augustinus sagt in seiner Regel: Saget auch nicht, daß ihr keusche Herzen habet, wenn ihr unkeusche Augen habt. Denn ein unkeusches Auge ist die Anzeige (nuntius) eines unkeuschen Herzens. Es soll sich deshalb niemand vermessen, daß er nicht ein Ehebrecher sei, wenn er das Werk nicht vollbringt, aber dennoch mit dem Auge begehrt.

Das sei gesagt von dem ersten und größten Zeichen.

Das andere ist das Hören, durch welches die böse Lust angezeigt wird, die inwendig brennt. Dazu gehören diejenigen, welche entweder dazu reizen, daß schandbare Dinge gesagt oder gesungen werden (und das sind die Aer[g]sten), oder welche still und gern zuhören, wenn sie auch nicht dazu Anlaß geben; denn sie ergötzen sich an diesen Dingen, wenden sich auch nicht davon ab als von einer Unreinigkeit und einem schädlichen Dinge. Die dritten möchten noch die sein, welche dergleichen nicht verhindern, denjenigen nicht strafen, welcher es mit anhört, oder ihn mit besseren Dingen unterrichten, auch den Leuten nicht unter Augen treten und sie schelten, und ihren Abscheu gegen solche Dinge an den Tag legen, damit sie die schrecken, welche zuhören. Nämlich diese sündigen durch Unterlassen, weil sie nicht bekümmert sind um die Seligkeit des Nächsten, und sich nicht angelegen sein lassen, ihn aus seinem Verderben zu ziehen, und weniger sorgfältig sind, als sie sein sollten.

Das dritte ist Reden, nämlich mit einem Weibe oder mit einem Manne. Denn die böse Lust bringt das zuwege, daß sie mit einander schwatzen, einander Höflichkeiten erweisen (salutent), sich mit einander zu schaffen machen; wiewohl dies bisweilen ehrbare Dinge sind, geschieht es doch nicht aus ehrbarer Ursache. Denn daher kommt das Lachen, die Scherze, das Geschwätz, die Unterhaltung, nur damit sie sich gegenseitig gefallen und sich zu gegenseitiger
\quellseite{662}
Liebe reizen. Und an diesem Fehler leiden die Weiber gar sehr, wenn sie lieben. Aber bei einem Mädchen ist nicht allein dies Reden, sondern durchaus alle Redeligkeit sehr unanständig (turpissima). Es ist auch nicht ein Zeichen einer beständigen Jungfräulichkeit oder Keuschheit, wenn ein Mädchen geschwätzig und plauderhaft ist, besonders mit jungen Männern. Denn siehe, da die heilige Jungfrau von dem Engel mit einem so großen Texte [Luc. 1, 28.] angeredet ward, sagte sie doch weiter nichts als dies [V. 34.]: „Wie soll das zugehen?“ etc., aber, wie St. Ambrosius sagt, sie erwiederte nicht einmal seinen Gruß. Deshalb, was thun die Mädchen (quae), welche sich gern mit ihrem Gesange oder ihrem Reden hören lassen wollen, und wiederum das Singen oder Gespräch des anderen Geschlechtes hören, anders, als daß sie beweisen, daß sie von böser Lust zu diesem Geschlechte brennen? Wenn sie nicht brennten, so würden sie sich nichts aus der Unterhaltung machen, oder darnach trachten, sie zu sehen, wie sie sich um andere Dinge nicht kümmern, zu denen sie durch keine Neigung hingezogen werden. Nun aber trachten die meisten Mädchen mit großer (tota) Leichtfertigkeit darnach.

Das vierte ist das Betasten, als da ist das Geben der Hände und das Umfassen, was zwar dem Scheine nach auch ehrbar geschehen kann, jedoch indem die böse Lust inwendig sich regt. Hier wird nun das Feuer der Unkeuschheit heftig entzündet, daß auch St. Hieronymus sagt, daß die böse Lust durch das Berühren sofort das Geschlecht erkennt. Darum muß man das Berühren so sehr als möglich fliehen, wie man von einem Pabste liest (des Namens erinnere ich mich jetzt nicht), daß er sich die Hand abhauen ließ, als er, wie es Sitte ist, seine Hand zum Küssen gereicht hatte, und durch den Kuß eines Weibes einen Kitzel empfand.

Das fünfte ist das Küssen. Das wollen wir übergehen.

Diese Zeichen gehen aber nirgends mehr im Schwange und in einer gröberen Weise als bei den öffentlichen Tänzen. Man muß sich wundern, wie viele und wie große Sünden da geschehen, was da das Sehen und das Hören aufstache, was das Betasten und das Reden erzeuge. Kurz, Welt ist Welt, ja, unrein\footnote{mundus — immundus, ein Wortspiel, das sich im Deutschen nicht wiedergeben läßt.} und eine Feindin GOttes, und man darf nichts Gutes suchen, das Gotte gefalle in der Welt: da treibt und häuft man Sünde auf Sünde, da sieht man den über alle Maßen großen Zorn GOttes; doch lacht man, tanzt und ist fröhlich, als ob es ganz unschädliche Dinge wären, ja, gut und nützlich.

Das sechste und zwar nicht ein geringes Zeichen ist der ausländische und neue Schmuck, der jetzt durch so viele Erfindungen, Künste und Bemühungen eifrig zuwege gebracht wird, daß man nicht anders sagen kann als: die Welt sei verrückt, während doch keine andere Ursache da ist, als daß sie gefallen, und sich der Welt gleichförmig machen wollen. Denn die Noth oder der Nutzen, ja, auch die Ehrbarkeit, welche allein gerechte Ursachen sind [sich zu schmücken], erfordert solche Dinge nicht, sondern ist mit mäßigem [Schmucke] zufrieden. Wer kann aber genugsam von jenen Dingen reden, da es ganz unerhörte Arten des Aufputzes und der Kleider sind? Aber diese Dinge sind ein heftiges Uebel und ein Zunder der Unkeuschheit; wiewohl man sich darum nicht kümmert. Ich glaube, daß sie endlich noch nackt einhergehen werden, da sie jetzt schon halbnackt gehen. Aber laß gehen, unsere Sitten werden auch ihre Zeit haben: Der Zorn GOttes ist endlich über die Welt gekommen [1 Thess. 2, 16.], und wohl dem, der da erkennt, daß der Zorn GOttes in Wahrheit schon über jene Leute gekommen ist.

Ich übergehe hier auch der Liebenden, das ist der Unsinnigen,\footnote{amantium — amentium, ein Wortspiel.} Verrücktheit, welche mit nächtlichem Dienste ihren Liebhaberinnen aufwarten mit Spiel, Gesang und Musik, und alles anbieten, damit sie ihnen gefallen, denn darüber urtheilt auch die Welt, daß es überaus großes Narrenwerk ist; deshalb gehen wir darüber hinweg.

Die vierte Stufe endlich ist die innerliche böse Lust selbst, von der oben genug gesagt worden ist. Niemand ist in diesem Leben ohne dieselbe; doch werden wir dadurch errettet, daß wir nicht in dieselbe willigen, wie es Röm. 6, 13. 12. heißt: „Auch begebet nicht eure Glieder der Sünde zu dienen, ihr Gehorsam zu leisten in ihren Lüsten.“ Denn obgleich das ganz reine (immaculata) Gesetz GOttes fordert, daß durchaus keine böse Lust da sein soll, sondern eine völlige und ganz unbefleckte
\quellseite{663}
Reinigkeit, so trägt doch unterdessen die Barmherzigkeit GOttes, und läßt sich wohlgefallen, daß der neue Mensch inwendig darinnen kämpft, damit nicht alle Ungerechtigkeit über uns herrsche. Diese muß fort und fort getödtet werden durch Seufzen, Wachen, Arbeit, Gebet, Erniedrigung (humiliatione) und andere Stücke des Kreuzes, endlich auch durch den Tod.

Aus allen diesen Dingen ist offenbar, daß dies Gebot nicht allein verbietet, wie es sich nach den Worten und dem Buchstaben ansehen läßt, sondern aufs stärkste gebietet, nämlich weil es die Keuschheit und eine ganz reine Enthaltsamkeit gebietet, innerlich und äußerlich. Wer sollte aber so thöricht sein, und leugnen, daß die Keuschheit etwas für sich selbst Bestehendes (rem affirmativam) sei? Aber sie wird in verneinender Weise geboten, da die Unkeuschheit verboten wird, damit verstanden werde, daß sie ganz rein, ohne alle Beimischung von Unkeuschheit, geboten sei, wie der 12. Psalm, V. 7., sagt: „Die Rede des HErrn ist lauter wie durchläutert Silber, im irdenen Tiegel bewähret siebenmal.“

Aber hierdurch wird nun klar, wie durch das Gesetz Erkenntniß der Sünde kommt [Röm. 3, 20.], denn ich wußte nicht, daß die böse Lust Sünde sei, „wo das Gesetz nicht hätte gesagt: Laß dich nicht gelüsten“. [Röm. 7, 7.], nämlich weil man allein das Werk für Sünde hielt. Darum ist es nicht genug, daß du mit dem Werke die Ehe nicht brichst, nicht genug, daß du in Worten keusch bist, und mit Anziehen oder irgend einer anderen Geberde; ja, auch das ist nicht genug, daß du mit dem Herzen nicht einwilligst, wenn du nicht auch endlich rein wirst von allen bösen Regungen, die inwendig aufwallen. Darum bleibt nichts Anderes übrig, als daß du seufzest nach der Gnade GOttes mit dem Apostel, Röm. 7, 24.: „Ich elender Mensch, wer wird mich erlösen von dem Leibe dieses Todes?“\footnote{Es sollte hier mit unserer Bibel und mit der Vulgata gelesen werden: de corpore mortis hujus. Dagegen bieten alle Ausgaben: de morte corporis hujus, und selbst die alte Uebersetzung: „von dem Tode dieses Leibes“.}

Es folgt das siebente Gebot.

\section{Das siebente Gebot}

\begin{schriftzitat}
Du sollst nicht stehlen.
\end{schriftzitat}

\footnote{Löscher: „Das Folgende hat Luther am dritten Weihnachtstage [Fer. III. Nativ., den 27. December] 1516 gepredigt.“ Das Exordium hiezu findet sich in der St. Louiser Ausgabe, Bd. X, 198 und (Duplicat) Bd. XII, 1800, an beiden Stellen mit falscher Zeitbestimmung: „Am zweiten Weihnachtstage“. Der Text ist der des dritten Weihnachtstages. (Weim. Ausg. I, 115.)} Wir haben gehört, daß das fünfte Gebot: „Du sollst nicht tödten“ die Schädigung (offensionem) des Nächsten an seiner eigenen Person verbietet, das sechste aber: „Du sollst nicht ehebrechen“, daß man ihn nicht schädige an der Person, die mit ihm verbunden und sein Fleisch mit ihm ist, und beides, sowohl im Werke als auch im Herzen, sowohl in Gedanken (affectu) als auch mit der That (effectu), so daß ganz offenbar ist, daß der geistliche Verstand dieser Verbote das Gebot ist, daß wir sanftmüthig und keusch sein sollen, inwendig und auswendig, ganz und gar ohne alle Bitterkeit des Zorns und ohne die angenehme Empfindung (dulcedine) der bösen Lust. Diese beiden Stücke kann man durchaus nicht haben ohne die Liebe des Geistes, denn sie sind Früchte des Geistes, 1 Cor. 13, 1. ff. und Gal. 5, 14. Darum ist die Liebe des Gesetzes Erfüllung [Röm. 13, 10.]. Denn dieselbe ist sanftmüthig und thut niemand Schaden, sondern nützt allen und wünscht allen dienstlich zu sein. Sie ist auch keusch und heilig; darum befleckt sie niemanden.

Jetzt folgt das Verbot, daß man dem Nächsten an seinen Gütern keinen Schaden thun soll, was gleicherweise ohne die Liebe nicht geschehen kann, wenigstens im Geiste. Darum ist auch hier zu merken, daß dies Gebot in zwiefacher Weise verstanden wird: Erstens dem Buchstaben und Wortlaut nach. So wird hier nur das Werk des Diebstahls selbst verboten, denn es sagt: „Du sollst nicht stehlen.“ (Es sagt schlechterdings nicht: Du sollst [den Diebstahl] nicht lehren, nicht dazu anleiten, nicht wünschen, sondern, du sollst ihn nicht thun. So verstanden es die Juden, und sind hoffärtig geworden
\quellseite{664}
in ihrer eigenen Gerechtigkeit, als ob sie keiner Uebertretung schuldig wären. Als daher Christus selbst Luc. 16, 9. 13. den rechten Verstand dieses Gebotes angab und sprach: „Ihr könnet nicht GOtt sammt dem Mammon dienen“, und: „Machet euch Freunde mit dem ungerechten Mammon“, hörten die Phariseer dies alles, welche geizig waren, und verlachten ihn, und er sprach zu ihnen [V. 15.]: „Ihr seid es, die ihr euch selbst rechtfertiget vor den Menschen, aber GOtt kennet eure Herzen; denn was hoch ist unter den Menschen, das ist ein Greuel vor GOtt.“ Siehe, wie er die Herzen erfodert bei der Rechtfertigung, die aus seinen Geboten herkommt. Endlich folgt daselbst, wie das Gesetz geistlich ist, da er sagt [V. 16. f.]: „Das Gesetz und die Propheten weissagen bis auf Johannes; und von der Zeit an wird das Reich GOttes durch das Evangelium gepredigt, und jedermann bringet mit Gewalt hinein. Es ist aber leichter, daß Himmel und Erde vergehe, denn daß ein Tüttelchen vom Gesetz falle.“ Siehe, „Gewalt thun“ ist vonnöthen bei GOttes Geboten, denn nicht allein nicht die Schale der Worte, sondern nicht ein Tüttelchen wird vergehen etc.

Zweitens nach dem Geist und nach der Absicht des Gebietenden. Darnach verbietet er hier alles Begehren des Herzens in solcher Weise, daß es unmöglich ist, daß du es erfüllest, wenn du nicht geistlich arm bist, wie es heißt Matth. 5, 3.: „Selig sind, die da geistlich arm sind.“ Darum werden die Gläubigen Christi in den Psalmen gleichsam als mit ihren eigenen Namen „Arme“ genannt. Ps. 10, 2. [Vulg.]: „Weil der Gottlose Uebermuth treibt, muß der Arme leiden.“ Und wiederum [Ps. 9, 10.]: „Und der HErr ist des Armen Schutz“, und [Ps. 9, 19.]: „Denn er wird des Armen nicht so ganz vergessen, die Hoffnung der Elenden (pauperum) wird nicht verloren sein ewiglich.“ Darum ist der Geist dieses Buchstaben die Armuth im Geiste. Deshalb ist ganz offenbar, daß niemand ohne die Gnade Gottes dies Gebot erfüllt, da der Mensch von Natur begehrlich ist, woran die erste Sünde Schuld ist.

\subsection{Folgerungssatz}

Es folgt, daß die Erfüllung dieses Gesetzes dem Buchstaben nach nichts Anderes ist als eine Art Selbstbetrug, wie es heißt Ps. 5, 10. [Vulg.]: „Weil in ihrem Munde die Wahrheit nicht ist, so ist ihr Herz eitel.“ Im Hebräischen heißt es so: „Weil in ihrem Munde nicht das Rechte ist“, nämlich weil er in diesem Gebote nur den Schatten und den Buchstaben des Werkes lehrt, „so ist ihr Inwendiges Falschheit“ (insidiae), denn sie täuschen sich, da sie sich für gerecht halten, betrogen durch den Schein der Werke, so daß sie nicht das Uebel ihrer Begierde erkennen, deshalb auch nicht um Gnade seufzen, noch ihre Sünde bekennen, sondern in Sicherheit dahinleben (stant). Deshalb sagt Paulus Röm. 2, 21 wider sie: „Du predigst, man solle nicht stehlen, und du stiehlst“, nämlich durch Begehren, wenn auch nicht im Werke.

Die Juristen aber legen dies Gebot nicht so aus, daß es sich so weit erstrecke, sondern bleiben in den Grenzen des Diebstahls; aber vor GOtt wird alles das ein Thun genannt, womit der Wille zu schaffen hat, wenn auch das Werk nicht folgen sollte. Denn GOtt prüft Nieren und Herzen, Ps. 7, 10. Deshalb müssen wir über das ganze Gebiet (latitudinem) der Werke gehen, welche die Begierde in ihrem Gefolge hat. Aber ach, lieber GOtt, wie weit erstreckt sich dies! Denn wer hat nicht die vielen Klagen über den Geiz und seine verschiedenen Arten gehört?

Das erste Werk nun nach dem Wortlaut des Textes ist der Diebstahl selbst. Dieser zerfällt wieder in fünf Arten, nämlich einfacher Diebstahl, Kirchenraub, Unterschlagen von Geldern, Wegtreiben von Vieh, und Menschenraub, wovon die Juristen handeln. Denn dies Gebot handeln sie nach dem Buchstaben überaus reichlich und sorgfältig in allen seinen Arten, wie ihre eigene Rede bezeugt, daß es die juristische Facultät mit dem Broterwerb zu thun habe: Nimm Mein und Dein hinweg, dann ist kein Recht mehr vonnöthen; und in der That gibt dieses Begehren [des fremden Guts] diesen Leuten unendlich viel zu schaffen. Es ist nun der Diebstahl die Aneignung fremden Gutes wider den Willen und ohne das Wissen seines Herrn. Wenn dies nun geschieht mit den Gütern von Privatpersonen, so ist es einfacher Diebstahl; wenn es in heiligen Sachen geschieht, so ist es Kirchenraub (sacrilegium); wenn es mit öffentlichem Gut geschieht, ist es Unterschleif (peculatus); geschieht es durch Entwendung von Zugvieh und anderen Thieren, so ist es abigeatus; werden Menschen weggeführt, so ist es plagiatus. Wie bekannt ist, so bestrafen die Gesetze heutzutage
\quellseite{665}
diese fünf Stufen sehr hart, vor Alters aber wurde im Gesetze der Diebstahl nicht mit der Todesstrafe belegt, wie man 2 Mos. 22, 1. 4. liest, sondern für ein Schaf wurden vier wiedergegeben, und für einen Ochsen fünf, jedoch nur, wenn das gestohlene Thier verkauft oder getödtet worden war; wurde es aber noch unversehrt bei ihm gefunden, so wurde nur das Zwiefache wiedergegeben; wenn er dies aber nicht wiedergeben konnte, so wurde er verkauft. Dies Gesetz wäre sicherlich jetzt noch gut; denn es ist allzu strenge, um nicht zu sagen unrecht, daß man um Gutes willen das Leben nehme, und die Seele für das, was zum Lebensunterhalt dient. „Ist nicht der Leib besser als die Kleidung, und das Leben mehr, denn die Speise?“ [Matth. 6, 25.] Daher sagen auch die Sprichwörter Cap. 6, 30. ff., welche dem Gesetze Mosis folgen [Vulg.]: „Es ist nicht eine schwere Schuld, wenn jemand stiehlt“, und dies ist wahr, wenn es allein ein Diebstahl ist. Daher folgt, was für einen Diebstahl er meine, da er sagt: „Denn er stiehlt, seine Seele zu sättigen, weil ihn hungert. Und ob er begriffen wird, gibt er's siebenfältig wieder, und legt dar alles Gut in seinem Hause, und kann sich frei machen; wer aber ein Ehebrecher ist, der wird, weil sein Herz arm ist, sein Leben ins Verderben bringen.“ Aber der Ehebruch wird heutzutage ganz gelinde bestraft, was sicherlich ganz unbillig ist; auf den Diebstahl ist eine ungemein harte Strafe gesetzt; diese sollte viel mehr gelindert werden als die des Ehebruchs, besonders für den Diebstahl, der wegen leiblichen Mangels und aus Armuth geschieht.

Ferner wird der Kirchenraub noch greulicher gestraft als alle, nämlich mit dem Rade, während doch eine heilige Sache und eine weltliche dieselbe ist, was den Werth und den Gebrauch anbetrifft; aber weil die Ehrerbietung gegen das Heilige aus den Augen gesetzt ist, wird eine so scharfe Strafe angewandt. Aber auch hier ist eine ganz abergläubige Strenge bei den Herren, welche wider nichts so sehr donnern und fluchen als wider die, welche ihnen ihre erbärmlichen Sachen (recularum), ihre Einkünfte, ihr Besitzthum, nämlich Gold und Silber entwenden. Hier haben sie hundertmal so viel Augen als [der hundertäugige] Argus, und sind wachsamer als die Wachamisten. Aber daß unterdessen so viele tausend Seelen vernachlässigt werden, daß sie durch Schändung, Ehebruch und andere schändliche Unzucht die ganze Kirche verunreinigen, neben dem Pfuhl ihres unersättlichen Geizes, mit ihrer Simonie, und Ränken, und unmäßigen Betrügereien: das straft niemand, sondern man hält vielmehr diejenigen in Ehren, die solches thun. Oder bestrafen sie vielleicht die geringeren Dinge um deswillen so schwer, weil es ihnen nicht zusteht, so schwere Dinge zu bestrafen, sondern es Gott zukommt, sie aufs Allerschärfste zu strafen? Und ich glaube, daß es sicherlich wahr ist, daß GOtt sich in den überaus schweren Dingen eine um so härtere Strafe vorbehält, je härter sie dieselbe in geringeren Dingen in Anwendung bringen. Wie viel sicherer würden sie fahren, wenn sie der Diebe und der Kirchenräuber schonten, wie sie wollten, daß GOtt ihrer schonen möge. Aber wenn sie jener schonten, dann würden sie sicherlich auch durch Armuth gezwungen werden, ein besseres Leben zu führen. Darum schonen sie nicht, damit bei GOtt für sie auch kein Verschonen da sei. Doch laß sie fahren, sie sind blind. Wegen der anderen drei Arten des Diebstahls verweise ich auf die Juristen, welche mit dem bürgerlichen Rechte umgehen.

Das zweite Werk ist Straßenraub, welcher ein viel schwereres Verbrechen ist als Diebstahl, weil er nicht geschieht, da der Herr [des Gutes] nichts darum weiß, sondern in seiner Gegenwart, und mit offenbarer Gewalt fremdes Gut raubt. Doch sind diese vielmehr Räuber (raptores) als Mörder (latrones) zu nennen, weil latrocinium den Mord mit in sich begreift. Diese Sünde ist aber jetzt bei den Edelleuten eine Zierde und ein Ehrentitel, und wird an denen bisweilen gestraft, die sich nicht vertheidigen oder Widerstand leisten können. Die Kaufleute kennen diese Sünde genugsam, und wissen, wie schwer sie ist, nämlich diejenigen, welche schon einmal davon zu leiden hatten.

Das dritte Werk ist der Wucher, von dem der römische Censor Cato gesagt hat: Die kleinen (privati) Diebe müssen ihr Leben im Gefängniß und Fesseln hinbringen, die großen (publici) Diebe aber gehen in Gold und Purpur. Darum wird ein Wucherer ein Blutegel des Volks genannt, und er wird durch dies Gleichniß beschrieben: Wie ein Wurm in einem Apfel oder einer Nuß das ganze Innere zerstört, so verschlingt ein Wucherer das Vermögen einer Stadt
\quellseite{666}
in mancherlei seltsamer und heimlicher Weise. An diesem Laster leiden die Juden so sehr, daß sie auch meinen, sie dürften es thun, ja, es werde ihnen im Gesetze befohlen, wo es heißt 5 Mos. 23, 19. f.: „Du sollst an deinem Bruder nicht wuchern, weder mit Geld, noch mit Speise, noch mit allem, damit man wuchern kann, sondern an einem Fremden; deinem Bruder aber sollst du alles, was er bedarf, ohne Wucher leihen.“ Dies war damals den Juden um ihrer Härtigkeit willen nachgelassen, damit sie sich desto williger des Wucherns an den Juden enthielten, da sie ihn überaus frei an den Heiden ausüben konnten. Aber die Christen sind bedeutend schlechter als die Juden, weil sie auch an ihren Brüdern wuchern; und dieses Uebel ist so häufig, daß es im 15. Psalm [V. 5.] auch als ein seltenes Lob des Gerechten angeführt wird, daß er sein Geld nicht auf Wucher gibt. Wenn sich nun viele desselben enthalten, so thun sie dies aus Furcht, nicht von Herzen, denn der Geiz ist ein überaus tiefes Verderben. Das vierte Werk ist Betrug im Handel. Davon heißt es 1 Thess. 4, 6.: „Daß niemand seinen Bruder vortheile im Handel, denn der HErr ist der Rächer über das alles.“ Aber dies erstreckt sich so weit, daß es jetzt in den einzelnen Stücken nicht angezeigt werden kann, denn wie könnte ich aller einzelnen Hantierungen und Händel gedenken, wie man darin zu viel thut, in der rechten Mitte bleibt, oder zu wenig thut?

Deshalb soll man für alle insgemein die Regel der Liebe halten, welche der HErr Matth. 7, 12. gibt: „Alles, was ihr wollt, daß euch die Leute thun sollen, das thut ihr ihnen.“ Deshalb ist es nothwendig, daß ein jeglicher in allem, was er vornimmt, zuerst darauf sehe, wie er wollte, daß man ihm thäte, wenn er an des anderen Stelle wäre, und in solcher Weise auch handele, dann wird er sicher sein. Dies nennt man das Gesetz der Natur, ja, es ist das Gesetz der Liebe. Denn die Liebe ist es, welche das ansieht, was des anderen ist, und hat zuerst darauf Acht, was er wollte, daß man ihm thun solle, damit er dem Nächsten auch so thue. Denn um desswillen hat der HErr seiner Rede dies noch hinzugefügt: „Das ist das Gesetz und die Propheten“, denn das Gesetz und die Propheten werden allein durch die Liebe erfüllt. Darum, wenn [man sagt, daß] es ein Gesetz der Natur ist, so wird die gesunde und unverderbte Natur verstanden, welche dasselbe ist wie die Liebe. Wer sich aber an diese Regel nicht kehrt, der vollbringt auch leicht einen Betrug, ehe er es gewahr wird, weil in allen Geschäften die Gefahr nahe liegt, daß die Treue verletzt werde, wenn sie nicht mit aller Macht (vi aliqua) und mit allem Fleiß gehalten wird.

Aber wie viele Betrügereien gibt es doch heutzutage! Jedes Haus und jede Gasse ist derselben voll, und alle beklagen sich darüber. Denn auch die Bauern, die sonst einfältig sind, sind sehr geschickt im Betrügen. Und fürwahr, schon dies Eine Gebot zeigt genugsam an, was die Welt sei, daß sie nichts ist, als daß einer dem andern seine Güter raubt. Und, was noch mehr zu beklagen ist, der Betrug ist so sehr zur Gewohnheit geworden, daß die Leute sich nicht einmal ein Gewissen darüber machen, ja sogar meinen, etwas Gutes ausgerichtet zu haben, wenn einer den andern getäuscht und betrogen hat. Aber wehe der Welt!

So gehe du nun durch alle einzelnen Hantierungen in ihren Verrichtungen (officia artium) und zähle die Betrügereien, da täglich neue erdacht werden, so daß es etwas sehr Gefährliches ist, mit Menschen Verkehr zu haben: das Licht verkehren sie in Finsterniß und die Finsterniß in Licht. Ja, es ist dahin gekommen, daß alle Verkäufer mehr fordern, und die Käufer weniger bieten, als die Sache eigentlich werth ist, nämlich weil keiner dem andern aufs Wort glaubt, da es der Brauch ist, daß einer den andern betrügt.

Was aber gesagt ist von den Werken dieses Gebotes und von seinem geistlichen Verstande, daß nämlich die Armuth des Geistes seine wahre Absicht sei, damit die unersättliche Begierde getödtet werde, der verfluchte Hunger nach Gold, der Geiz, die Geldliebe und Geldgier, die Wurzel alles Uebels, wie der Apostel [1 Tim. 6, 10.] sagt, welche heutzutage unzählige Fallstricke gelegt hat und deren immer mehr macht, ist nicht zu verstehen von der Armuth in der That, sondern nach dem Geneigtheit dazu (affectu). Denn auch Bettler und Mönche können sehr reich sein in ihrem Begehren und innerlichen Geize, und vor GOtt die größten Diebe, indem sie die beneiden, die etwas haben, und deren Güter begehren. Und der Apostel nennt [Col. 3, 5.] dies Laster Abgötterei, wie auch der 78. Psalm [V. 58. Ps. 106, 36.] sagt, daß sie den Götzen
\quellseite{667}
(sculptilibus) Canaans dienten, das ist, den Pfennigen. Denn Canaan heißt ein Kaufmann und ein Händler; wie könnten wir nun glauben, daß durch dessen Götzen (sculptilia) etwas Anderes abgebildet sei als Gold und Silber, welches zusammengeschmolzen, gemünzt und geprägt ist?

Endlich, daß die Verehrer des Goldes auch die Knechte desselben sind, zeigt auch der 76. Psalm [2. 6. Vulg.], wo es heißt: „Und alle Männer des Reichthums fanden nichts in ihren Händen.“ Er sagt nicht: „Reichthum der Männer“, sondern „Männer des Reichthums“, damit du verstehest, daß die Männer [vom Reichthum] besessen werden (possessos viros), welche Reichthum besitzen (possidentes divitias); jene dienen, diesem wird gedient; darum sind jene Knechte, dieser der Herr.\footnote{Dieser Absatz fehlt in der alten Uebersetzung; vielleicht wegen der Schwierigkeit. Denn auch sonst sind schwierige Stellen weggelassen.}

Da dies Laster allen überaus wohl bekannt ist, will ich die Sache nun zum Beschluß kurz zusammenfassen. Dies Gebot wird erfüllt, wenn du dich nicht allein des Werkes enthältst, ja, es ist auch nicht genug, daß du mit dem Herzen nicht begehrst, sondern es [ist vonnöthen, daß du es] auch ganz gleichgültig erträgst, wenn man dir das Deine nimmt. Denn es nützt nichts, wenn du fremdes Gutes nicht begehrst, das deinige aber in gottloser Weise liebst. Gleichwie beim fünften Gebote jemand an sich selbst zum Todtschläger werden kann, wenn er auch keinen anderen Menschen tödtet, oder gegen niemanden zürnt, und im sechsten Gebote mit seiner eigenen Frau ein Ehebrecher sein, wie gesagt ist, so kann er auch in seinen eigenen Gütern ein Dieb und ein Geiziger sein, indem er das Seine allzusehr liebt. Doch, wer sollte nicht glauben, daß ein Mensch, der sich selbst tötet, nicht auch einen anderen tödten, oder ihm wenigstens das Leben nicht gönnen werde? und wer sein eigen Eheweib zur Unkeuschheit mißbraucht, der sollte von dem nicht voraussetzen, daß er auch eine andere, wenigstens in seiner bösen Lust, beflecken werde, da er ganz und gar von Unkeuschheit wüthet? Solcher Weise auch, wer wird glauben, daß derjenige, welcher das Seine allzusehr liebt, nicht auch fremdes Gut begehre? Deshalb soll ein Mensch in Bezug auf irgend welche Güter sich ganz gleichgültig verhalten (neuter esse): er soll sich nicht betrüben, wenn sie verloren gehen, er soll sich nicht freuen, wenn er sie erlangt, nach dem Worte des 62. Psalms [V. 11. Vulg.]: „Verlasset euch nicht auf Unrecht, und setzet euer Begehren nicht auf Raub; fällt euch Reichthum zu, so hänget das Herz nicht daran.“ Dies aber vermag die verderbte Natur nicht, sondern muß es bei Christo holen, der für uns arm geworden ist.

Ja, ein Mensch soll auch auf dieser Stufe nicht stehen bleiben, sondern auch zusehen, daß er viel mehr willig sei, von seinem Gut anderen mitzutheilen, als fremdes Gut zu begehren; denn das bringt die Liebe mit sich, welche allein das Gesetz erfüllt, daß er sich nicht allein völlig fern halte von den Gütern des Nächsten, sondern auch der Nächste werde dadurch, daß er sein Gut mittheilt, wenn er es nur irgend kann. Denn ohne Zweifel wollte jeder gern, daß man gegen ihn so handelte, nämlich, daß man ihm mittheilte, und niemand das Seine begehrte, sondern ihm alles aufrichtigen Herzens gönnte. Und es kann in diesem Gebote, welches für so verwickelte Händel gestellt ist, keine bessere Regel gegeben werden, als daß ein jeglicher in allem, was er mit dem Nächsten zu schaffen hat, sich vor allem dieses Wort Christi vor Augen stelle, welches die geistliche Auslegung aller Gebote ist [Matth. 7, 12.]: „Alles, was ihr wollt, daß euch die Leute thun sollen, das thut ihr ihnen“, und den Spruch des Tobias [Tob. 4, 16.]: „Was du nicht willst, daß man dir thue, das thue einem andern auch nicht.“ Denn ein jeglicher, der das ins Auge faßt, wird leicht ersehen können, was und wie er thun solle in irgend einem Falle, der sich bei diesem vielumfassenden Gebote zutragen mag. Denn wenn du dir Reichthum gönnst, aber nicht deinem Nächsten, dann hast du schon keine Liebe, sondern böses Begehren, und hast wider dies Gebot gesündigt. Denn du bist ein Dieb vor GOtt, weil du dem Nächsten das genommen hast, was du ihm schuldig bist, nämlich, daß du ihm sein Gut gönnest.

Hieraus ist offenbar, wie viele Leute einen Abscheu haben vor den Dieben und Diebstahl, die doch selbst viel größere Diebe sind, da ohne die Gnade niemand sein kann, dessen Herz nicht durch Geiz verunreinigt werde. Aber jene Diebe hängt der Kaiser, diese hingegen wird der Schöpfer aller Dinge verderben. Darum ist auch dieses Gebot nicht bloß ein Verbot, sondern ein
\quellseite{668}
überaus starkes Gebot, nämlich daß der Mensch mildthätig sein soll, willig zu geben und freigebig gegen den Nächsten; so kann er aber nicht beschaffen sein, wenn er im Geiste nicht ganz rein ist von Diebstahl. Darum ist dieses Verbot kräftiger, als wenn es in der Form eines Gebotes gestellt wäre, da jemand zwar freigebig sein könnte, nämlich wenn er anfängt, dies Gebot zu erfüllen (ut in statu incipientium), aber noch nicht ohne Begierde [nach fremdem Gut]; darum erfüllt er dieses Gebot noch nicht vollkommen. Ein vollkommener Erfüller des Gebotes muß freigebig sein aus ganz reinem Herzen.\footnote{Das Vorhergehende hat Luther am Tage der unschuldigen Kindlein (den 28. December) 1516 gepredigt. Es gehört dazu das Exordium, welches sich in unserer Ausgabe, Bd. XII, 1804, findet. Doch wissen wir nicht, an welcher Stelle Luther mit demselben eingesetzt hat, da Löscher es nicht angibt. (Weim. Ausg., Bd. I, 504.)}

\footnote{Löscher: „Das Folgende hat Luther am Fest der Beschneidung [1. Januar] im Jahre 1517 gepredigt.“ Das hiezu gehörende Exordium findet sich Walch, St. Louiser Ausgabe, Bd. X, 1276.} Man wirft auch die Frage auf: ob Spieler Diebe seien, und das von ihnen Gewonnene wieder zurückgeben müßten. Daß die Spieler durch Begehren und Habsucht sündigen (wenigstens die, welche des Gewinnes halben spielen), und deshalb vor GOtt Diebe sind und fremdes Gut begehren, ist ganz offenbar. Denn niemand spielt mit einem anderen deshalb, weil er ihm das Seine geben wollte (denn das könnte er thun, auch ohne zu spielen), auch nicht, damit er das Seine verliere, auch nicht, um des Andern Nutzen ebenso zu suchen, als wenn es sein eigener wäre, und deshalb ist das Spiel immer wider die Liebe und geht in böser Begierde daher, denn er sucht das Seine mit dem Schaden des anderen, oder wenigstens das, was des andern ist, nicht so wie das Seine. Aber dennoch wird er nicht zur Wiedererstattung angehalten, weil er vor Menschen nicht raubt, auch fremdes Gut nicht nimmt wider den Willen des Herrn, der vielmehr dazu seine Einwilligung gegeben hat durch das Uebereinkommen, daß sie spielen wollten (per pactum ludi). Denn beide setzen das Ihre mit Wissen und Willen in Gefahr, daher verliert der Verlierende mit Wissen und Willen. Denn es ist ein ganz gerechtes Gesetz, daß derjenige, welcher die Gefahr liebt, darin umkomme, wenn sie nicht etwa betrügerisch spielen. Doch die beste Strafe würde sein, wenn es geschehen könnte, daß beide verlieren müßten, denn dies wäre das richtige Recht des Spieles, weil jeder von beiden mit eigener Gefahr seiner Begierde gegen den Nächsten freien Lauf läßt; oder wenigstens, daß die Landeskasse (fiscus principalis) beiden ihren Gewinn abnähme, weil sie durch das Spielen gegen das Gebot [des Fürsten] und das geistliche Recht der Kirche handeln. Es ist aber offenbar, daß es nicht nöthig sei, solches wiederzugeben, weil keine andere Ursache dafür vorhanden wäre, als daß ein solcher Gewinn mit Sünde und wider das Gesetz erlangt wird. Aber wenn alle solche Dinge zurückgegeben werden sollten, die mit Sünden erlangt werden, dann müßten auch die Huren ihren Gewinn wieder herausgeben, ja, es müßte von den Kaufleuten und Händlern alles das wiedererstattet werden, was mit Lügen, falschen Beteurungen (wobei übrigens Betrug gegen den Nächsten und falsche Angabe [dolo] hinsichtlich der Sache, um die es sich handelt, ausgeschlossen sein müssen) und bösem Begehren verkauft oder gekauft wird. Dies wäre ganz ungereimt, weil in diesen Dingen dem Nächsten kein Unrecht geschehen ist, sondern GOtt beleidigt und das Gesetz übertreten worden ist. Daher ist er Wiedererstattung schuldig, zwar nicht dessen, was er an sich gebracht hat, sondern daß er seinen bösen Willen, den er daselbst seinen Lauf gelassen hat, fahren lasse; das geschieht durch Reue und Buße.

Dasselbe ist zu sagen von Priestern und Ordensleuten, welche entweder ihre Gebete nicht verrichten oder ihre Genugthuungen nicht vollbringen, oder mehr Pfründen als Eine besitzen; die sind nicht schuldig wiederzuerstatten, sondern daß sie Buße thun wegen des Mißbrauchs, denn der Pabst hat die Macht diese Güter zu verleihen; aber siehe du, der du dieselben erlangst oder dich hineindrängst, dich nur wohl vor. Vor Zeiten flohen die Leute vor dem priesterlichen Amte und man mußte sie hineinziehen, darum wurden sie von dem Volke geliebt; jetzt aber stürzen sie hinein, während das Volk sie gern davon abhalten möchte, darum werden sie gehaßt. Aber dieser Geiz wächst heutzutage in erschrecklicher Weise bei dem Einnehmen von Pfründen. Aber wenn Pfründen ohne Betrug erlangt werden, so ist keine Wiederherausgabe vonnöthen, wiewohl es nicht recht möglich scheint, daß sich jemand ohne Betrug selbst eindringen könne, besonders zu dem Amt der Seelsorge (curarum),
\quellseite{669}
denn es wird der Pabst und das Volk betrogen, wenn er nicht ein solcher Mann ist, als ein solches Amt erfordert. Wer aber möchte sich wohl unterstehen, das von sich selbst zu behaupten? Und doch schämt man sich nicht, ein solches Amt für sich zu erbitten. Endlich dient auch der Contract des Wiederkaufs jetzt dem Geize, zum großen Schaden der ganzen Welt, da Zinsen (census) nur um deswillen gekauft werden, damit des Gewinnes mehr werde und das Vermögen (substantia) sicher sei ohne Arbeit und Gefahr, nicht damit sie Gotte desto ungestörter dienen können oder der Noth abhelfen, während man doch sieht, daß des Pabsts Absicht dabei gewesen sei und noch sei, nicht, daß sie müßiggehen sollen; denn 1 Mos. 3, 19. sind Arbeit und Kummer allen Menschen aufgelegt: „Im Schweiß deines Angesichts sollst du dein Brod essen“, und Hiob 5, 7. heißt es [Vulg.]: „Der Mensch ist zur Arbeit geboren, wie der Vogel zum Fliegen.“ Darum wollen wir einige Ursachen anzeigen, wegen deren man mit Recht Zinsen kaufen kann. Nämlich erstens, wenn es alte Leute sind, Kinder und Kranke, denn diese haben Arbeit genug durch ihr [elend(es)] Leben; zweitens, die mit anderer Arbeit beschäftigt sind [durch welche sie nichts erwerben können],\footnote{Von uns hinzugesetzt.} nämlich die, welche GOtt dienen, als, Ordensleute, Priester und Prälaten, die mit Studiren und dem Worte GOttes zu thun haben; dann auch diejenigen, welche dem Gemeinwesen und fremder Noth dienen, als, Fürsten, Rathsherren, die Obersten in Städten und Dörfern. Sonst, wenn GOtt einem jeglichen wiedergeben wird nach seiner Arbeit, was werden die Müßiggänger empfangen, welche ihr Geld nur zum Gewinn mißbrauchen, den Gewinn nur zum Müßiggang, den Müßiggang nur zum Wohlleben, das Wohlleben nur zur Sünde? Die sind also nicht sicher, welche mit den erkauften Zinsen nur sich selbst dienen; nicht als ob Zinsen kaufen so gar böse wäre, sondern weil es Sünde ist, daß man darin dem Geize dient und nichts Anderes sucht als sicheren Reichthum.

Es folgt das achte Gebot.

\section{Das achte Gebot}

\begin{schriftzitat}
Du sollst nicht falsch Zeugniß reden wider deinen Nächsten.
\end{schriftzitat}

\footnote{Löscher. Das Folgende hat Luther am Sonntage nach dem Feste der Beschneidung [den 4. Januar] im Jahre 1517 gepredigt. Das Enchiridion hiezu ist nicht mehr vorhanden.} Dies ist das Laster der Zunge, welche der Apostel Jacobus eine Welt der Ungerechtigkeit nennt. Sie ist nach dem Herzen das allersündhafteste Glied, und kein Glied sündigt auf so vielerlei Weise, ausgenommen das Herz, welches in jeder Art sündigt. Denn die Sünde, welche die Zunge begeht mit Schwören, Geloben, Beten, Lehren, Mißbrauchen des Namens GOttes etc., ist in der ersten Tafel im zweiten Gebote verboten, das heißt, jedes Wort, welches geradezu gegen GOtt sündigt. Desgleichen ist sie verboten im dritten, vierten, fünften, sechsten, siebenten und jetzt im achten; durchaus in allen Geboten außer dem ersten, neunten und zehnten, welche nur die Seele betreffen. Denn wider das dritte thut man, wenn man nicht betet, nicht dankt, [GOttes Wort] nicht liest am Feiertage; wider das vierte, wenn man die Eltern hoffärtig anredet, ihnen unehrerbietig antwortet, sie schilt etc.; gegen das fünfte durch Fluchen, Uebelreden, Anschuldigen etc.; wider das sechste durch unkeusche Reden; wider das siebente durch Rath und Aufforderung zu stehlen, durch Entschuldigen und Vertheidigen des Diebstahls; wider das achte, wie wir jetzt sehen werden, durch Lügen, Schmeicheln etc. Wie es daher verboten ist, daß man jemandem Schaden thue an der eigenen Person, oder an den Personen, die ihm angehören, an seinen irdischen (corporalium) Gütern, so wird jetzt an den Gütern, welche den guten Ruf, das Lob, die Ehre, den Namen betreffen, alles Schaden thun verboten.

\subsection{Wider dies Gebot sündigen:}

Erstens und recht eigentlich die Gelehrten, besonders aber die Theologen und alle, die den
\quellseite{670}
Schülern falsche Lehren vortragen, und beide auf zwiefache Weise, nämlich entweder wissentlich oder unwissentlich, und darnach sündigen sie schwerer oder leichter. Zuerst nun will ich von den Theologen sagen, weil sie von allen am schwersten wider dies Gebot sündigen, wenn sie die heilige Schrift und das Wort GOttes verkehren und falsch auslegen oder nicht recht anwenden. Denn daher kommt es, daß etwas, das falsch geredet ist, die Seelen tödtet oder wenigstens die ungöttlichen Gemüthen martert.

Zu den ersten unter diesen gehören die Ketzer und diejenigen, welche sonderliche Meinungen aus ihrem eigenen Kopfe erdichten,\footnote{In allen Ausgaben: Monii singulares. Anstatt Monii schlägt Köcher Moni vor. Die Weimarische Ausgabe bemerkt: „Es ist wohl das griechische μονοί.“} welche Gottes Wort ihren Meinungen anpassen und in ihren Gehorsam gefangennehmen, während sie gerade entgegengesetzt handeln sollten, und zwar in den Dingen, die den Glauben und die Seligkeit antreffen. Dieses wird bei GOtt für eine so große Sünde gehalten, daß keine andere Sünde so häufig und so scharf in der Schrift gestraft wird durch alle Vorbilder und Prophetieen. Daher sagt der 4. Psalm, V. 3. [Vulg.]: „Ihr Menschenkinder, wie weit geht eure Herzenshärtigkeit! Wie habt ihr das Eitle so lieb und trachtet nach Lügen!“ und der 5. Psalm, V. 6. f.: „Du bist feind allen Uebelthätern; du bringest die Lügner um“, wo er von solchen Lehrern redet, wie aus dem Folgenden erhellt [V. 11. Vulg.]: „Laß sie fallen von ihren Gedanken“, das ist, von ihrem Vornehmen etc. Und im ersten Psalm heißt es: „Wohl dem, der nicht wandelt im Rath der Gottlosen, noch tritt auf den Weg der Sünder, noch sitzet, da die Spötter sitzen.“ Und überall redet die Schrift von den Falschen und von den falschen Zungen, und straft diese falschen Märtyrer, das ist, die falschen Zungen.

Deshalb sage ich: Sind sie denn etwa entschuldigt, da sie dies in guter Meinung und aus Eifer um GOtt thun? Denn so pflegen sich die Ketzer zu entschuldigen, wie unsere Nachbarn, die Picarden in Böhmen, thun. Ich antworte: Durchaus nicht. Vielmehr werden sie dadurch nur um so mehr angeklagt, weil dieser Eifer aus Hochmuth hervorgeht, daß sie sich nicht belehren lassen und nicht hören wollen, weil sie ihre eigenen Meister geworden sind. Daher heißt es Röm. 10, 2.: „Ich gebe ihnen das Zeugniß, daß sie eifern um GOtt, aber mit Unverstand“, und Gal. 5, 10.: „Wer euch irre macht, der wird sein Urtheil tragen“, und Gal. 1, 9.: „So jemand euch Evangelium prediget anders, der sei verflucht“ etc. Denn solche Leute haben einen guten Eifer, aber sie fürchten sich nicht, denn sie sind wegen solches ihres Eifers sicher, als ob sie nicht irrten, weil sie die gute Meinung haben, und schauen, daß alle Sünde im freien Willen stehe (est voluntarium), und wie die Meinung sei, so sei auch das Werk (intentio nominat opus). Aber wir, sagen sie, haben eine gute Meinung und einen gottseligen Willen; und so werden sie hartnäckig und sicher. Diesen muß man antworten, daß sie einen ganz bösen Willen und eine überaus giftige Meinung haben. Ich sage nicht von der, welche sie fühlen und geltend machen, sondern von der tiefern, welche sie stolz macht und taub, daß sie nichts Anderes hören wollen als ihre gute Meinung, und keinem anderen glauben als ihrem gottseligen Willen. Denn diesen ihren innerlichen Stolz wollen sie nicht sehen, und schmücken ihn mit der erdichteten und bei den Haaren herbeigezogenen (violentia) guten Meinung und dem gottseligen Willen. Denn die rechte gute Meinung und der gründlich gute Wille steht niemals in Sicherheit und ohne Furcht, denn er fürchtet sich, wie Hiob [Cap. 9, 28. Vulg.], wegen aller seiner Werke, ja, wegen seiner Worte und Gedanken, darum gründet er sich nicht auf seinen guten Willen und gottseelige Meinung, und ist immer mehr bereit, anderen zu weichen, als sich selbst zu glauben, einen wie überaus schönen Schein seine Meinung auch haben mag. Durch diese Furcht geschieht es, daß er nichts Neues auf die Bahn bringt, die Einigkeit bewahrt, Frieden hält und die Eintracht fördert.

Die zweiten sind die, welche zwar nicht Falsches lehren, so daß sie dadurch tödten, aber doch nur oberflächlich lehren, und weniger, als sie sollten. Darum sind sie Kreuziger der Gewissen und Peiniger der Seelen. Dazu gehören diejenigen, welche den Glauben an Christum zwar unangetastet lassen, aber die Gerechtigkeit der Werke lehren; selten den Glauben, häufiger die Werke. Von diesen scheint mir die Offenbarung [Cap. 9, 5. zu sagen, daß die Heuschrecken die Menschen nicht tödteten, sondern sie quälten fünf Monden lang, nämlich weil sie lehren, daß
\quellseite{671}
durch Werke für die Sünden genuggethan werde, und durch dieselben doch niemals Ruhe des Gewissens haben, welche der Glaube an Christum allein gibt.

Weil nun die heilige Schrift in Furcht und Demuth gehandelt sein will, und daß man in sie eindringe, mehr durch Studiren unter brünstigem Gebet als durch scharfen Verstand, so ist es deshalb unmöglich, daß diejenigen nicht sich selbst und anderen, die sie unterrichten, Schaden thun sollten, welche sich bloß auf ihren Verstand verlassen und mit ungemäßlichen Füßen, wie die Säue, über dieselbe herfallen, als ob es nur eine Art menschlicher Wissenschaft sei; so gar machen sie keinen Unterschied und verfahren mit ihr ohne alle Ehrerbietung. Daher kommt es, daß so viele sich so bald unterwinden, Lehrer zu sein, welche sofort, nachdem sie die Grammatik gelernt haben, ohne irgend etwas studirt zu haben, sich zu Lehrern der Theologie aufwerfen, und sprechen: Ach, die Bibel ist etwas ganz Leichtes. Besonders thun dies aber diejenigen, denen der Bauch geschwollen ist von den Träumen der Säue, das ist, der Philosophen. Daher klagt St. Hieronymus über diese Leute, welche die Schrift, die wider sie streitet, nach ihrem Belieben deuten, und schon ein Sprichwort daraus geworden ist: die Schrift habe eine wächserne Nase. Das haben diese abgeschmackten und ungereimten Träumer zuwege gebracht, welche mit dem buchstäblichen, bildlichen, sittlichen und höheren Verstande\footnote{Val. Tischreden Cap. 52, § 5. Walch, St. Louiser Ausg., Bd. XXII, 1341.} der Schrift ihr Spiel treiben, und scholastische Lehrer genannt werden, und zwar mit ihrem ganz eigentlichen und rechten Namen. Denn sie sind Scholastiker,\footnote{Σχολαστικός, ein Mensch, der nichts zu thun hat, herumgeht und träumt, oder auf alberne Einfälle geräth; Tropf, Einfaltspinsel, Narr.} das ist, lächerliche Leute und Spielkinder (lusores), ja, auch Leute, welche mit sich selbst und anderen ihr Gespött treiben (illusores), die nicht erkannt haben, weder was Buchstabe noch was Geist sei. Und es wäre das Behandeln der Schrift nach dieser vierfachen Auslegung noch leidlich, wenn es nur als ein erster Unterricht für die Anfänger angesehen würde. Nun aber setzen sie in diese Spielerei den Gipfel der Gelehrsamkeit, und kommen nie zu dem rechten Verstande, der ohne Reinheit des Herzens niemals erlangt werden kann. Aber auch die gottlosen Leute treiben dieses Biergespann. Siehe die Vorrede des heiligen Hieronymus zur Bibel über dieselbe Frevelhaftigkeit derer, welche die Schrift behandeln.

Zweitens ist von den Juristen zu reden, denn diese handeln wider dies Gebot nach dem Buchstaben, wenn sie in ihrem Labyrinthe, das ist, auf ihrem Hadermarkte (foro contentioso), nichts Anderes suchen, als zu gewinnen. Sie führen alles an, sie suchen nach allem, sie nehmen alles, was für sie dienen kann gegen den Widderpart, wo denn durch ein neues Wunder ein und derselbe Text für Sachen dienen muß, die einander gerade entgegengesetzt sind, und dieses Studium ist voll geworden von überaus vielen Gefahren. Aber sie sagen: Niemand ist schuldig, ein Bekenntniß über sich selbst abzulegen und sich selbst anzuklagen. Was sagt aber der HErr? „Du sollst nicht falsch Zeugniß reden gegen deinen Nächsten“, indem du das für dich anführst, was wider den ist, von dem du weißt, daß er eine gerechte Sache hat. Aber die Juristen haben ein großes Vorrecht, daß niemand wider sie reden oder sie richten darf; dies aber ist: daß niemand mit Recht das straft, was er nicht versteht. Aber der Verstand des Rechtes ist unendlich und unerschöpflich, und man hat noch keinen Juristen gefunden, der alles wüßte. Wenn man deshalb nicht eher über sie urtheilen dürfte, als bis man alle ihre Sachen verstanden hätte, so würde man niemals urtheilen können, denn sie werden immer sagen, man verstehe das nicht, was man straft.

Noch es muß eines gesagt werden, was auch die Rechte selbst enthalten, wiewohl man dessen fast gar nicht achtet, nämlich daß der Zweck, das Leben und die Kraft aller Gesetze darin besteht, daß Friede, Liebe und Eintracht hergestellt werde. Wenn jemand dies nicht im Auge hat, so gebraucht er nicht die Gesetze, sondern mißbraucht sie. Denn wer um deswillen eine Ausflucht sucht in einer Glosse wider den Text, um den Hader in die Länge zu ziehen und zu mehren: meinst du etwa, daß er das lebendige Gesetz habe? Vielmehr sucht er den Buchstaben des Todes und arbeitet wider den Frieden. Aber wer dies heutzutage nicht thut, da ist es zu verwundern, wenn er lange für einen Juristen gehalten wird, oder reich wird. Denn wenn sie den Frieden so sehr suchten als den Hader, so wären nicht so viele Bücher und so viele Arbeit vonnöthen. Darum siehe nicht bloß die Worte an, wie sie
\quellseite{672}
lauten, auch nicht wohin und wie du sie verdrehen und in der Schwebe halten (ponderare) könnest, sondern siehe darauf, was sie beabsichtigen, nämlich dem Streite ein Ende zu machen und den Zwiespalt zu tödten. Aber hier wäre es recht geeignet, das sechste Capitel des ersten Corintherbriefes recht anzuziehen, wo der Apostel sie tadelt [V. 3. 7.], daß sie um weltlicher, das ist, um zeitlicher Dinge willen, die zum Lebensunterhalt dienen, mit einander rechten und nicht viel lieber den Betrug über sich ergehen ließen. Desgleichen Christus befiehlt, daß man dem auch den Mantel geben soll, der den Rock nimmt etc. [Matth. 5, 40.] Aber es wird gesagt, daß dieses Räthe seien, nicht Gebote, die nur zur Bereitung des Herzens dienen sollen. Ich aber verstehe diese wunderliche Bereitung des Herzens nicht, die niemals zu Tage tritt und zum Werke kommt. Und St. Augustinus, welcher das Exempel Christi anführt, daß er, als er vor Hannas den Backenstreich empfangen habe, die andere Backe nicht auch dargereicht habe, wird, wie ich dahinthalte, nicht recht verstanden, und dadurch werden sie gar übel vertheidigt.

\footnote{Löscher: „Hier ist im Manuscript etwas von den Aerzten eingefügt.“ Dies hat Luther bei der Ueberarbeitung gestrichen.} Hieher mögen vielleicht auch die elenden Menschen gehören, die sich mit dem überaus falschen Vorgeben (nomen) der Philosophie des Aristoteles betrügen und abmühen, da sie den Aristoteles weder selbst verstehen, noch lehren,\footnote{Es scheint uns, daß statt doceantur in den Ausgaben doceant gelesen werden sollte. Danach haben wir übersetzt.} und unterdessen die an Christum gläubigen Seelen mit falschen und erdichteten Auslegungen einnehmen und belasten, ja vielmehr betrügen, da sie glauben, das sei Philosophie, was es doch nicht ist, und nur mit Narrenpossen zu thun haben. Und die Dichter und irgendwelche Fabelerzähler sind viel besser daran als sie, zumal da die Dichter, auch wenn sie Fabeln erzählen, wie St. Augustinus sagt, doch nicht glauben, daß sie wahr seien, sondern an der nichtigen Sache nur ihr Ergötzen haben; diese aber glauben, daß ihre Fabeln wahr seien. Denn wenn Aristoteles nach meinem Rathe gelesen würde, so könnte ihn ein Mensch von mittelmäßigem Verstande wohl in einem halben Jahre lernen; dann aber würde man ihn nicht, wie es bisher geschehen ist, mit demselben Glauben und derselben Andacht (religione) lesen als die heilige Schrift, oder gar noch mit größerer Ehrerbietung, sondern wie irgend einen anderen geringfügigen und [der Religion] fernliegenden Gegenstand (extranea materia), nur damit man es wüßte, nicht daß man es vertheidigte. Denn wahrlich, wir sind durch den Aristoteles und seine Fündlein verführt worden, und wenn kein anderer Grund da wäre [den Aristoteles bei Seite zu legen,\footnote{Von uns hinzugefügt.}] so wäre dies Eine Grund genug, daß in dieser heidnischen Bestie so viele Secten und Köpfe sind, die der Schlange (Hydrae) im lernäischen Sumpfe gleich ist; denn da sind die Thomisten, die Scotisten, die Albertisten, die Modernen; und Aristoteles ist vierköpfig geworden, und sein Reich in sich selbst zertheilt, und es ist zu verwundern, daß es nicht verwüstet wird; aber seine Verwüstung steht nahe bevor. Denn wie ist es möglich, daß die Wahrheit da sein sollte, wo so viele Meinungen verwirrt durch einander gehen? Denn da sie wider einander sind, müssen sie nothwendiger Weise auch falsch sein.

Aber wenn diese Lehrer und Meister sündigen und wider dies Gebot handeln, wo wollen denn jene Träumer und verruchten Menschen bleiben, die falsche Wunder predigen, lügenhafte Legenden, erdichtete Reliquien zeigen, ihre gotteslästerlichen Gesichte vorgeben? Hier rühmt einer mit unverschämter Stirn, er habe mit Christo und seiner Mutter geredet und gelacht; ich meine, er wird [mit ihr] auch getanzet haben. Ein anderer hat Heu, auf dem Christus gelegen hat; es ist aber in diesem Jahre gewachsen und in einer Scheuer aufgerafft. Ich habe gehört als eine Thatsache, daß ein gewisses Weib, auf des Teufels Anstiften und aus GOttes Verhängniß, damit sie einen so thörichten Anhänger eines solchen Aberglaubens zum Narren haben sollte, das Haar an ihrer Scham abschnitt, es ihm gab und ihn überredete, daß es das Haupthaar der heiligen Catharina wäre, das über das Meer gebracht worden wäre. Jener Mönch (cucullus) glaubte es und hielt es den Leuten als eine Reliquie hin zum Küssen und Verehren. Und siehe, was geschieht? Es fing das Haar von dem schandbaren Orte auch an, Wunder zu thun, und nicht mit Unrecht, damit die durch
\quellseite{673}
Lügen betrogen werden, die der Wahrheit nicht dienen [2 Thess. 2, 10. 11.]. Siehe, dies ist so unrein und schamlos, daß es vor züchtigen Leuten auch nicht einmal erzählt werden kann; dennoch geschieht es. Diese nun reden falsches Zeugniß, nicht gegen Menschen, sondern gegen GOtt und seine Heiligen, die schon in der Herrlichkeit sind (regnantes). Derartig ist auch das, daß kürzlich ein gewisser Mensch Heu gezeigt haben soll, welches er zufällig fand, und sagte, es sei das Heu, auf welchem Christus oder ich weiß nicht welcher Heilige geschlafen habe. Solche Fabelreden nun vervielfältigen die ungelehrten Priester überall, aber noch mehr die Mönche, welche als Reliquienkrämer umherziehen (stationarii),\footnote{Vgl. Tischreden Cap. 27, § 143. Walch, St. Louiser Ausg., Bd. XXII, 923.} und die Pfarrherren ergötzen sich gar sehr an den Fabeln der Legenden und des Ablasses und aller Dinge, die nur gesagt werden können, nur nicht an dem Evangelium, weil jene Dinge Gewinn bringen, während das Volt die Ehren von der Wahrheit wegwendet und sich zu den Fabeln kehrt. Aber wehe, wehe und abermal wehe über solche Fabelprediger! Hier wäre eine sehr große Reformation der Kirche vonnöthen, so daß es nicht gestattet werden sollte, irgend etwas noch so Göttliches und Gutes zu predigen, es sei denn die gewisse Wahrheit und dem rechten Glauben gemäß erfunden worden (authenticum et canonisatum). Denn unter solchem Schein der Göttlichkeit sind gleicherweise jene seltsamen Fabeln eingerissen, so daß die Prediger der Kirche die heidnischen Dichter fast übertreffen an Fabeln und Lügen. Ja, ich möchte mich wohl unterstehen zu sagen, es sei auch nicht ganz sicher, daß bewährte Legenden und Wunder so gepredigt werden, daß daraus der Brauch entstehe, das Evangelium gänzlich bei Seite zu setzen. Denn sie sind nicht für den hauptsächlichsten Theil der Predigt zu halten, sondern nur bei der Predigt des Evangelii mit einzustreuen, wo sie zur Sache dienen, aber nicht in einer ganzen Predigt zu handeln.

Zweitens wird jede Lüge verboten, das heißt, falsches Zeugniß wider den Nächsten. Die Lüge ist aber, wie St. Augustinus unterscheidet, dreifach: Die Scherzlüge, die Nothlüge (officiosum) und die schädliche Lüge.

Die Scherzlüge ist eigentlich nicht eine Lüge, weil man sofort wahrnimmt, daß es eine Lüge ist, und auch derjenige, welcher so lügt, ist beflissen, so zu reden, daß man sehen kann, er rede die Unwahrheit und scherze.

Dennoch ist es Sünde, denn es ist schandbare Rede oder Narrentheidinge und unnützes Geschwätz. Es kommt aber vor, daß etliche so einfältig sind, daß sie auch durch solche Dinge betrogen werden, besonders wenn dieser Scherz in heiligen Dingen (religionibus) geschieht, als, wenn gesagt wird: Hast du etwas verloren, so gib eine warme Semmel als Almosen, und streue geweihtes Salz darauf, — darauf, — nämlich auf die verlorene Sache, — so wirst du es finden. Das verstehen die Einfältigen so, daß das Salz auf die Semmel gestreut werden solle, und glauben, daß sie dann das Verlorene finden werden; dergleichen geschieht viel.

Die Nothlüge ist die, welche jemandem zu Dienste oder aus Gefälligkeit geschieht, für sich selbst oder zu des Nächsten Heil und Nutz, ohne irgend jemandem zu schaden, als, wenn die Keuschheit eines Weibes in Gefahr stände, und jemand sagte, sie wäre nicht da, wo sie ist, oder etwas Anderes der Art. Eine solche Lüge haben die Wehmütter der hebräischen Weiber in Egypten gesagt. Desgleichen Michal, Davids Weib, gegen Saul. Siehe beide Geschichten an ihren Orte [2 Mos. 1, 19. und 1 Sam. 19, 14.] nach. Auf die Frage, ob man sich ihrer bedienen möge, antwortet Augustinus: Keineswegs, sondern er sage die Wahrheit, soweit er kann, mit anderen Worten, wie er von einem gewissen Bischofe sagt in seinem Buche „von der Lüge“, der einem Richter antwortete, welcher ihn befragte wegen eines Angeklagten, der zu ihm geflohen war: Ich will weder lügen, noch ein Verräther sein. Doch ist dies keine Todsünde, wenn die Nothlüge aus einem vernünftigen Grunde geschieht und in plötzlicher Bestürzung.

Die schädliche Lüge ist die, welche Schaden thut. Denn es wird die Unwahrheit gesagt, um einem Anderen Schaden zu thun, und derjenige, welcher so lügt, betrügt den Nächsten in solcher Weise, wie er selbst nicht betrogen werden möchte. Darüber sagt der HErr durch Sacharja [Cap. 8, 16. Vulg.]: „Redet die Wahrheit, ein jeglicher mit seinem Nächsten“, und wiederum [V. 19.]: „Liebet Wahrheit und Frieden“ etc. Denn das ist, was mir gefällt. Dies Uebel reagiert heutzutage überaus weit, schädlich und gewaltig in Kaufhändeln, Hantierungen, Verträgen
\quellseite{674}
, Bündnissen, Briefen, Siegeln und in allem Verkehr der Menschen mit einander, wie Hosea [Cap. 4, 2.] sagt: „Lügen, Stehlen etc. hat überhand genommen“, weil die Menschen dieser Regel vergessen: Was du willst, daß man dir thue, das thue du auch einem anderen, und was du nicht willst, daß man dir thue, das thue einem anderen auch nicht, denn sie sind blind geworden durch Selbstsucht etc.

\footnote{Löscher: Das Folgende hat Luther am Tage der Bekehrung Pauli [25. Januar] im Jahre 1517 vorgetragen. Das Exordium hierzu findet sich Walch, St. Louiser Ausgabe, Bd. X, 976.} Siehe aber, was für ein Greuel das menschliche Herz ist. Wenn auch alle anderen Sünden den Menschen nicht verunreinigten, so ist doch dieses Laster am meisten wider die Natur des Menschen. Erstlich: Denn was begehrt ein jeder Mensch von Natur mehr als [die Wahrheit] zu wissen? Wovor hat er einen größeren Abscheu, als daß er betrogen werde? wie St. Augustinus sagt, daß die Seele noch nichts ein so großes Verlangen habe als nach der Wahrheit, und dennoch thut der Mensch dawider, und zwar gern. Zweitens: Wenngleich der Mensch mit vielen Thieren das Fressen oder die Schwelgerei, Zorn, Neid, Faulheit, Habsucht gemein hat, so übertrifft er sammt dem Teufel sie doch in diesem Laster, weil allein der Mensch ein Lügner ist, und der Teufel ein Vater der Lüge. Denn allein dem Menschen ist die Vernunft und die Sprache gegeben, deren Amt es ist, die Wahrheit zu sagen. Und deshalb glaube ich, daß unter allen Scheltwörtern keines bitterer ist, als wenn man sagt, jemand lüge ins Angesicht. Denn dies Scheltwort thut mehr, als wenn man sagt, daß jemand ein ausschweifender Mensch, ein Trunkenbold, ein Zorniger, ein Hoffärtiger sei, und ich glaube, daß die Natur dies um deswillen leichter empfindet, weil die Lüge für den Menschen das unnatürlichste Laster ist, ja, der größte Feind der menschlichen Gesellschaft.

Drittens wird wider dies Gebot durch Schmeichelei gesündigt, unter welcher der Ehrenhäsler, der zweizüngige Mensch und die falsche Zunge und alle Laster der Zunge begriffen sind. Anderes kannst du im fünften Gebote sehen. Von diesem Laster ist aber kein Mensch frei, es sei denn, er wolle sich bald todtschlagen lassen und den Haß aller Leute auf sich laden, nämlich wer, weil sein Herz recht steht, seinem Menschen die Wahrheit verschweiget, wie der 15. Psalm, V. 4., sagt: „Wer die Gottlosen nichts achtet“, das heißt, keinen bösen Menschen fürchtet, wie mächtig, reich, weise, gerecht er auch sein mag (denn vor ihm gilt kein Ansehen der Person und er kümmert sich auch um keinen Menschen), dagegen die Furchtsamen zu Ehren bringt, wie schwach, arm und thöricht sie auch immer sein mögen. Deshalb muß man der Welt entweder schmeicheln, oder den Haß der Welt leiden, was uns an Christo und allen Heiligen genugsam gezeigt worden ist. Wiewohl nun dies Laster überall herrscht, so hat es dennoch die größte, unbestrittene Herrschaft (triumphat) an den Höfen der Fürsten und im ganzen Lager der Bischöfe. Denn da sind die Leute, welche weiche Kleider tragen und gar zarte Ohren haben, und es ist doch gewißlich kein Ort in der Welt, von dem die Schmeichelei weiter entfernt sein sollte, und wo sie dennoch häufiger ist und größere Macht ausübt [als anderswo], da ja, wenn ein Fürst zum Bösen verführt wird, oder er vom Guten abgehalten wird, nothwendiger Weise alle Unterthanen Schaden leiden, und es kein schädlicheres Verderben für die Reiche und Länder gibt, als einen Schmeichler am Hofe. Denn was beklagen wir uns über Kriege oder Schwerter oder Waffen? Die Zunge des Schmeichlers geht über alle Schwerter; deshalb sollte man die Schmeichler so viel wie möglich von den Höfen wegtreiben und sie aufs schwerste bestrafen. Man liest, daß der Kaiser von Rom, Alexander, den Schmeichlern so feind gewesen sei, daß er, wenn er sah, daß jemand vor ihm in kriechender Weise (molliore gestu) schmeichelte, ihn nicht ungestraft ließ, oder ihn wenigstens verhaftete, wenn es ein Vornehmer war. Und als Diogenes gefragt wurde, welches von allen Thieren das schädlichste sei, antwortete er: Unter den wilden Thieren ein Tyrann, unter den zahmen ein Schmeichler.

Man findet in der Schrift wenig hervorragende Schmeichler und Exempel von ihnen, vielleicht, weil es ein so abscheuliches Laster ist, denn der erste und vornehmste vor anderen ist Doeg, der Edomiter, welcher, wie 1 Sam. 21, 7. und nachher [1 Sam. 22, 9. 10. 18. 22.] erzählt wird, Saul in allen Dingen schmeichelte und die Ursache war, daß alle Priester in Nobe erschlagen wurden, dann aber auch viele andere
\quellseite{675}
Leiden über David brachte. Danach folgen die Siphiter [1 Sam. 23, 19.], welche David verriethen, endlich Siba, der Knecht Mephiboseths, welcher auch den König David betrog [2 Sam. 16, 3.]. Aber man möge Acht darauf haben, was für ein Ende Doeg genommen hat sammt Saul [1 Sam. 31, 4. f.], dem er schmeichelte; aber auch sein Sohn, der unseliger Weise seinem Vater in der Schmeichelei nachfolgte, wurde sofort von David getödtet, 2 Sam. 1, 13. ff. Daher redet auch der ganze 52. Psalm Doeg’s Donnerworte wider diesen Doeg [V. 3.]: „Was trotzest du denn, daß du kannst Schaden thun“ etc. Es mögen die Schmeichler hören, was da folgt [V. 6. f.]: „Du redest gerne alles, was zum Verderben dient, mit falscher Zunge. Darum wird dich GOtt auch ganz und gar zerstören und zerschlagen und aus der Hütte reißen und aus dem Lande der Lebendigen ausrotten.“ Siehe, mit wie entsetzlichen Verwünschungen er wider den Schmeichler herausfährt! Denn sie lieben Worte des Verderbens, das heißt, Worte, durch welche sie ihre Nächsten stürzen mögen, sie zu Grunde richten, sie bei den Fürsten verdrängen (absorbent), damit sie nur allein herrschen und in Gunst stehen etc.

\footnote{Löscher: „Das Folgende hat Luther am 3. Sonntage nach Epiphania vorgetragen, im Jahre 1517.“ „Der 3. Sonntag nach Epiphania fiel 1517 auf den 25. Januar; aber schon das vorhergehende Stück gehört einer Predigt vom 25. Januar 1517 an: daher ist wohl ein Versehen bei Löscher anzunehmen; vielleicht hat im Manuscript ‚Dom. IV. post Epiph.‘ [1. Febr.] gestanden.“ (Weim. Ausg.) Das Exordium hierzu findet sich Walch, St. Louiser Ausg., Bd. XII, 1814.} Endlich wird durch falsches Zeugniß im öffentlichen Gerichte gesündigt, worüber auch oben geredet ist. Hier wird von der Gefahr gehandelt, in der die Juristen, Notarien, Rechtsbeistände, Advocaten und Richter stehen. Hier leiden selbst die überaus heiligen Rechte das Geschick der Rose, aus welcher die Biene Honig sammelt, die Spinne aber Gift, und sie sagen fein, um sich zu entschuldigen: Dem, der wachsam ist, kommen die Rechte zu Hülfe; und wiederum: Man muß seine Zuflucht nehmen zu dem Heilmittel des Rechts. Dies sind zwei sehr nützliche Regeln, erstlich, um die Vollstreckung der Gerechtigkeit hinauszuschieben, zweitens, um Hader zu säen, drittens, um den Geldbeutel der Rechtsanwälte zu füllen; nicht als ob jene Regeln schlecht oder falsch wären, sondern weil sie fast nirgends in gutem Brauche stehen. Dies wird offenbar, wenn du darauf Acht hast, was sie verstehen unter „Heilmittel des Rechts“ und „wohl Acht haben auf die Rechte“, nämlich dies, daß, wenn jemand auch nach dem Zeugniß seines eigenen Gewissens eine ganz schlechte oder wenigstens eine zweifelhafte Sache hat, er nicht auf den Rath des Geistes hört: Behalte das Gewisse und laß das Ungewisse fahren, vielmehr nur darauf bedacht ist, wie er aus der überaus schlechten Sache eine sehr gute machen könne, nämlich weil die Sache krank ist, so wendet er das Heilmittel des Rechts bei ihr an, und dies ist ihre Wachsamkeit. Denn „sie wachen über die Gottlosigkeit“, wie der Prophet [Jes. 29, 20. Vulg.] sagt, und so geschieht es, daß sie aus einem und demselben Texte des Rechts Gift und Honig schöpfen: der eine, welcher ungerechte Sachen vertheidigt, Gift, der andere, welcher der gerechten Sache Vorschub leistet, Honig.

Was soll ich aber hier sagen? Ich kann nur das Elend der Christen beklagen, welche all ihr Bestreben auf Hader, Prozesse und Gericht setzen, und keiner läßt sich dazu herbei, auf das Wort Christi Acht zu geben [Matth. 7, 12.]: „Alles, was ihr wollt, daß euch die Leute thun sollen, das thut ihr ihnen.“ Ich bitte dich, lieber Mensch, versuche, prüfe, frage dich selbst, wenn du etwa eine Sache hast wider deinen Nächsten. Wolltest du denn nicht, daß er dich in Frieden gehen ließe, nicht mit dir rechtete, daß er dir Glauben schenkte? Also das willst du, daß dir ein anderer thäte: warum thust du ihm denn nicht dasselbe und weichst ihm? Aber, sagst du, es ist ein Rath, nicht ein Gebot. Ich antworte: Wie kann es ein Rath sein und nicht ein Gebot, da Christus sofort hinzugefügt hat: „Denn das ist das Gesetz und die Propheten.“ Ist denn etwa ein Gesetz nur ein Rath? Bleibt denn also kein Gebot, kein Gesetz? Was folgt hieraus, als daß wir sowohl Christi als auch des Gesetzes der Natur vergessen und mitten in der Finsterniß tappen? Daher ist heutzutage das Studium des Rechts überaus gefährlich, da es nicht zum Frieden dient, nicht, den Hader beizulegen, sondern anzustiften und in die Länge zu ziehen. Aber, sprichst du, wenn man so handeln müßte, würde es bald dahin kommen, daß man uns mit Füßen träte und uns alles raubte, was wir haben. Hier sage ich, das ist es, was die
\quellseite{676}
Schrift will. Denn deshalb nennt sie uns Kindlein, Waisen und Witwen, Arme und Unterdrückte, weil nämlich wir, die wir das Evangelium bewahren, sehr viel Uebels leiden um seines Wortes willen, da wir uns nicht rächen, nicht Gleiches mit Gleichem vergelten, ja auch bitten für die, welche uns verfolgen und verleumden, auf daß wir Kinder seien unseres Vaters im Himmel, weil er auch gütig ist gegen die Undankbaren und Bösen [Matth. 5, 44. f.]. Und wir sollten uns freuen, daß wir des Namens „Kindlein“ werth geachtet werden, weil Christus ihr Beschützer ist. Wenn wir nun wiederfordern und haden, wie können wir dann Kindlein sein? Aber auch diese heilsamen Gebote haben sie uns in Rütte verwandelt, indem sie sagen, daß sie von der Bereitung des Herzens zu verzeihen seien, und beachten nicht, daß fort und fort Fälle kommen, diese Bereitung auch ins Werk zu bringen, und thun dies dennoch nicht. Und hier siehe GOttes überaus mächtige Weisheit und seine überaus weise Macht. Er hat gesagt [Luc. 21, 33.]: „Meine Worte werden nicht vergehen“ und [Matth. 5, 18.]: „Es wird nicht ein Tüttel vom Gesetz vergehen, bis daß es alles geschehe.“ Siehe, sage ich, wie wahr es sei [Ps. 33, 9.]: „So er spricht, so geschieht es.“

Diejenigen, welche dem, der ihnen den Mantel nimmt, nicht auch den Rock geben wollen, werden gezwungen, ihn\footnote{Wir haben mit der Baseler, der Wittenberger und der Jenaer Ausgabe illi angenommen, statt illis in den anderen Ausgaben.} wider ihren Willen auch das Hemd zu geben, weil sie den Rock nicht willig gaben, und so wird durch sie wider ihren Willen das Gebot erfüllt, welches mit ihrem Willen nicht erfüllt werden konnte. Denn es geschieht sehr häufig, wenn zwei über eine Sache streiten, daß sie beide durch einen Theil der streitigen Sache zufriedengestellt werden würden, indem sie einen gütlichen Vertrag mit einander machen, und jeder dem andern nachgäbe, während sie, wenn sie mit einander rechten, weit mehr darüber verzehren, und an die Advocaten doppelt so viel vergeuden, als sie dem Widerpart hätten nachlassen müssen. Wie geht es aber nun? Etliche müssen über das, was die Sache werth ist, auch noch das Doppelte darauf verwenden, und verlieren dennoch auch die Sache selbst. Hier ist es Zeit zu singen [Ps. 119, 137.]: „Herr, du bist gerecht, und dein Gericht ist recht.“ Gib einem unbarmherzigen, fremden Advocaten, der dich noch dazu verspottet, zwanzig Gulden für den Proceß, weil du deinem Bruder um Friedens willen nicht zehn Gulden hast nachlassen wollen. Denn wer ist hier anders zu tadeln als du ganz blinder Mensch, der du dein Geld verschleuderst an Unwürdige mit Sünden, welches du mit einem Verdienste hättest deinem Bruder lassen können, der dein Widersacher ist, den GOtt würdig geachtet hat, es zu haben? Wenn du nun auch gewinnen solltest, so hast du deinen Sieg dennoch mit Sünden, und wirst nackt erfunden werden in zwiefachem Tode.

Aber dies ist von Privatleuten gesagt, denn öffentliche Personen müssen wachen, daß niemand Unrecht geschehe, und, wenn sie dazu aufgefordert werden, zugefügtes Unrecht rächen. Dies zu fordern ist aber den Schwachen zugelassen, den Christen aber verboten.

Schließlich ist zu merken, daß, wie der Apostel Petrus sagt, niemand leide als ein Dieb, oder ein Uebelthäter [1 Petr. 4, 15.], das heißt, er soll nicht bekennen, daß ihm recht geschehe, wenn er unschuldig ist, sondern dem, der ihm etwas nimmt, seine Schuld anzeigen, und bei ihm den Stachel des Gewissens zurücklassen. Denn so haben die heiligen Märtyrer Attalus und seine Gefährten gethan. Da man die Beschuldigung erhob, daß alle Christen insgeheim Kinder fräßen, leugneten sie es beständiglich bis zum Tode. So sagte Sokrates, als er von dem Könige Joas getödtet wurde [2 Chron. 24, 22.]: „Der HErr wird’s sehen und suchen.“ Und St. Thomas von Canterbury wurde getödtet, nicht, weil er nicht wollte, daß man die Güter nähme, sondern weil er sagte, daß die, welche sie nahmen, kein Recht dazu hätten. Denn er ließ das Gut ohne Widerstand nehmen, und widersetzte sich auf keine Weise. Ja, auch St. Stephanus erklärte so stark, daß seine Verfolger nicht recht thäten, daß er auch für ihre Sünde betete. So schreibt St. Hieronymus von dem Weibe, auf welches [von dem Henker mit dem Schwerte] siebenmal geschlagen ward, daß sie den Tod und alles andere nicht fürchtete, sondern ihre Unschuld aufrecht erhalten wollte. Ja, um deswillen gab Christus dem Malchus eine harte Antwort, und nahm in Abrede, daß er übel geredet habe [Joh. 18, 23.], damit er
\quellseite{677}
nicht meinen möchte, er habe recht gehandelt. Deshalb antwortete er auch dem Caiaphas, er sei der Sohn GOttes, und widerlegte ihr falsches Zeugniß, und dennoch ließ er zu, daß man ihm das Leben nähme, ja, er wollte es. So sollen auch wir unser Gut fahren lassen, unsere gerechte Sache aber aufrecht erhalten.

Nun wollen wir es wiederholen und in eine Summa zusammenfassen: Dieses Gebot wird erfüllt, wenn wir nicht allein kein falsches Zeugniß reden, sondern auch nicht hören oder zulassen wollen, daß über jemanden geredet werde (wie diejenigen thun, welche entweder den Schmeichlern glauben, oder den Verleumdern beistimmen). Ja, derjenige ist noch dieses Gebotes schuldig, welcher nicht mit allem Fleiße darauf aus ist, seines Bruders Unrecht, Schmach und Schande zu entschuldigen, zuzudecken, oder wenigstens, wenn er nicht mehr thun kann, sie gering zu machen, denen, die überreden, anklagen, verleumden, zu widersprechen; denn ein jeglicher wollte, daß man so gegen ihn thäte. Und so wird er ein Kind des Geistes GOttes, welcher der Tröster genannt wird, weil er Gutes für uns redet, und unserem Geiste Zeugniß gibt wider alle Zeugnisse des Teufels und unseres eigenen Gewissens. Ferner soll man auch die belehren, welche GOtt nicht kennen und GOttes Willen nicht wissen, die Fremden ermahnen, den Verführern und Ketzern und andern Betrügern der Seelen Widerstand leisten, wie der Apostel sagt [Tit. 1, 9. f.]. Deshalb erfordert dies Gebot eine wahrhaftige, gerechte, heilsame, unschädliche und dem Bruder völlig zu Dienst stehende Zunge, und zwar so, daß man ihm von Grund des Herzens Gutes gönne. Dies ist die Liebe, die nicht Muthwillen treibt, nicht eifert [1 Cor. 13, 4.]. Wiewohl es daher nicht nothwendig ist, daß du dem Priester beichtest, wenn du deinen Bruder nicht entschuldiget hast, da man ihn angeklagt hat, geschändet mit falschem Zeugniß, unterdrückt durch Schmeichelei oder Verleumdung, und vielleicht damit zufrieden bist, daß du dies nicht gethan hast, so wirst du doch vor GOtt wegen eines so gottlosen Schweigens und bösen Nichtgebrauchs deiner Zunge, die dir gegeben worden ist für solche Noth deines Bruders und zu seinem Dienste, als ein Schuldiger verurtheilt werden. Aber wiederum, wenn du dies thust, wirst du dir alles auf den Hals laden und bald getödtet werden. Und dies ist es, was GOtt will, damit wir desto eher um seines Wortes willen aus der Welt hinausgestoßen und getödtet werden und die Krone der Gerechtigkeit desto eher empfangen. Dies ist der rechte und geistliche Verstand dieses Gebotes.

\section{Das neunte und zehnte Gebot}

\begin{quote}
Du sollst nicht begehren deines Nächsten Haus. Du sollst auch nicht begehren sein Weib, Knecht, Magd, Ochsen, Esel, noch alles, was sein ist.
\end{quote}

\footnote{Wahrscheinlich ist die Predigt am Tage der Reinigung Maria [2. Febr. 1517] das Exordium zu dem nachfolgenden Stück (Weim. Ausg.). Diese Einleitungspredigt findet sich Waldh. St. Louiser Ausg., Bd. XII, 1266.} Dies ist im sechsten und siebenten Gebote mitbegriffen, da nach der Auslegung Christi daselbst alles Verlangen und Begehren des Weibes und der Güter des Nächsten verboten wird. Denn „wer ein Weib ansiehet, ihrer zu begehren, der hat schon mit ihr die Ehe gebrochen in seinem Herzen“ [Matth. 5, 28.]. Ferner hat er den Geiz und das Begehren [der Güter] so gar verdammt, daß er auch befohlen hat, dem auch den Rock zu geben, der den Mantel genommen hat [Matth. 5, 40.].

Aber nach meinem Urtheil scheint mir durch diese beiden Gebote der Zunder und das unüberwindliche Begehren verboten zu werden, ich meine, die eigentliche Wurzel der bösen Gedanken; so daß es verstanden werden muß, daß im sechsten und siebenten Gebote das Zustimmen des Herzens und die Geberden der Glieder verboten sind, das Wort des Mundes und das Werk des bösen Leibes, hier aber auch selbst die ersten Regungen zusammen mit dem Zunder, welcher der Ursprung jener ist. Denn wir müssen so rein werden, ehe wir ins Himmelreich kommen können, daß auch keine bösen Regungen in uns seien, noch irgend ein Zunder, der sich zum Bösen hinneigt, sondern eine völlige Gesundheit des Leibes und der Seele, durchaus ohne alles Gebrechen, was freilich in diesem Leben nicht geschen
\quellseite{678}
wird, auch nicht in unserer Macht steht. Denn wer kann sich rühmen [Spr. 20, 9.], daß er ein reines Herz habe? Wer kann dieses tief innerliche Feuer auslöschen? Denn selbst der Apostel klagt über „das Gesetz in seinen Gliedern“ und über „das Gesetz der Sünde“ [Röm. 7, 23. f.]. Und deshalb wehren wir den Augen, Ohren und allen Sinnen inwendig und auswendig, damit die Sünde nicht in uns herrsche, aber der bösen Lust kann niemand wehren. Deshalb hat denn auch der barmherzige GOtt beschlossen, unsern Leib zu Asche zu machen und durch das große Feuer des jüngsten Tages (universalem) zu verbrennen, damit alles Gebrechen unseres Fleisches durch Feuer verzehrt werde: so sehr haßt GOtt diese Unreinigkeit. Und hieraus folgt, daß es allein diese beiden Gebote sind, die von keinem auch noch so heiligen Menschen auch nur einigermaßen erfüllt werden; alle anderen erfüllen sie, weil sie das Werk, das Wort, die Einwilligung mit aller Gewalt dämpfen, aber hier bleiben sie schuldig und Sünder, weil sie nichts von diesen Geboten erfüllen, weil sie angesteckt sind von der unüberwindlichen bösen Lust des Fleisches und dem Begehren der Güter: darum sind sie „allzumal Sünder und mangeln des Ruhms, den sie an GOtt haben sollten“ [Röm. 3, 23.]; darum beten alle: „Geheiliget werde dein Name, dein Wille geschehe, vergib uns unsere Schuld“, darum verzweifeln alle an ihren Verdiensten aus Furcht wegen ihrer Unreinigkeit und setzen ihr Vertrauen auf die Barmherzigkeit GOttes, damit in solcher Weise das fest bestehe, daß der HErr Gefallen hat an denen, die ihn fürchten, die auf seine Güte warten [Ps. 147, 11.]

Hierdurch wird auch dem Streit der scholastischen Lehrer ein Ende gemacht, ob diese beiden Gebote etwas Anderes gebieten als das sechste und siebente, desgleichen, ob es zwei verschiedene Gebote seien. Und sie mühen sich vergeblich ab, indem sie diese beiden Gebote von den Werken des Herzens verstehen, da sie doch nicht leugnen können, daß diese im sechsten und siebenten Gebote verboten seien, und daß dann hier dasselbe wiederholt sein würde, und diese beiden Gebote vergeblich gesetzt wären, da sie mit jenen einerlei sein würden. Sie können auch nicht sagen, daß dort das Werk, hier die Gedanken, oder dort die Hände, hier das Herz verboten werde, weil Christus selbst daneben lehrt, indem er sagt, daß jene beiden Gebote von beiden zu verstehen seien.

Es bleibt also nur übrig, daß wir dem Apostel Paulus folgen, der diese beiden Gebote in eins zusammenfaßt und sagt: „Ich wußte nicht von der Lust, daß sie Sünde sei, wo das Gesetz nicht gesagt hätte: Laß dich nicht gelüsten“, Röm. 7, 7. Und dort setzt er weitläufig auseinander, daß durch dies Gebot alle Heiligen betroffen werden, auch er selbst. Also wird hier das Uebel des Zunders und, daß ich so sage, die wesentliche und ursächliche Unreinigkeit in uns verboten. Wer sich daher dünken läßt, er habe alle vorhergehenden Gebote erfüllt, der sieht wenigstens hier, daß er unrein sei und der Reinigkeit Christi bedürfe, die für ihn geopfert und von GOtt, dem Vater der Barmherzigkeit, angenommen worden ist.

Aber das, was die Juden oder vielmehr die Säue hier verstehen, ist nicht werth, daß man es bei der Besprechung dieser Gebote erwähne. Denn sie sagen, daß im sechsten und siebenten Gebote nur das Werk verboten werde, und im neunten und zehnten ein äußerer Versuch und ein offenbarbares Anzeichen, daß man das Werk vorhabe, wie auch bei den Juristen der Versuch (conatus) bestraft wird, als, wenn jemand Vorbereitungen getroffen hat (inciperet), die Tochter jemandes zu entführen, oder einem andern nach dem Leben gestanden hat, wird er geköpft. Und so sind bei ihnen Gedanken und Worte frei, und voll sind sie von Tücke und Zorn und Bitterkeit und, wie Christus sagt [Matth. 23, 27. f.], „voll alles Unflaths“, weil sie das rein halten, was außen ist, aber ihr Inwendiges ist voller Geiz und aller Untugend. Darum ist die Synagoge das Weib, welches wegen seines schauderbaren Wesens den Scheidebrief erhalten hat und verstoßen worden ist.

Dies sind „die zehn Worte“, wie Moses sagt, in denen jedes heilsame Gebot völlig enthalten ist, und wiewohl sie auf eine geringere Anzahl hätten gebracht werden können, so hätte man sie doch auch in mehr Gebote theilen können. Doch es hat GOtte wohlgefallen, sie in der Zehnzahl zu geben, welche die Zahl eines abgeschlossenen Ganzen (universitatis) und einer vollkommenen Summe ist, so daß, gleichwie einst alles im Vorbild (figura) geschah, auch diese Zehnzahl der Gebote um deswillen gewählt worden ist, daß sie die Summe aller Gebote abbilden sollte.

\quellseite{679}
\footnote{Löscher: „Das Folgende hat Luther am Sonntage Septuagesimä [den 8. Februar] 1517 gepredigt.“ Das hierzu gehörende Exordium findet sich Walch, St. Louiser Ausgabe, Bd. XII, 1816.} Nachdem nun die zehn Gebote zu Ende gebracht sind, müssen die so vielen Arten von Sünden unter die Gebote eingereiht werden, weil es ja keine Sünde gibt, die nicht wider ein Gebot GOttes geht. Denn ich weiß nicht, ob es denen, die da beichten wollen, nütze sei, daß sie mit so vielen Unterschieden der Sünden ihr Gedächtniß beladen und den Priester ermüden, als da sind: 1. nach Thun und Unterlassen; 2. mit dem Herzen, mit dem Munde, mit dem Werk; 3. nach den fünf Sinnen; 4. nach den sechs Werken der Barmherzigkeit; 5. nach den sieben Sacramenten; 6. nach den sieben Todsünden; 7. nach den sieben Gaben [des Heiligen Geistes]; 8. nach den acht Seligpreisungen; 9. nach den neun fremden Sünden; 10. nach den zehn Geboten; 11. nach den zwölf Artikeln des Glaubens;\footnote{„Die zwölf Artikel des Glaubens“ sind das apostolische Glaubensbekenntniß. Vgl. Tischreden Cap. 54, § 13. Walch, St. Louiser Ausg., Bd. XXII, 1358.} 12. nach den zwölf Früchten des Geistes. Außer diesen sind noch die vier Haupttugenden und die drei theologischen Tugenden, desgleichen die stummen Sünden, die himmelschreienden Sünden und endlich die Sünden wider den Heiligen Geist. Ich bitte dich, was nützt doch solche Verwirrung und Zerstreuung des Gemüthes? Dies wird so beobachtet, daß man es auch für nothwendig hält, die Namen dieser verschiedenen Abtheilungen und Unterscheidungen aufzuzählen, natürlich nur, um die Zeit zu vergeuden, den Beichtiger zu betäuben und sich selbst zu verwirren, so daß man mehr Mühe darauf verwendet, diese Dinge im Gedächtniß zu haben, als an rechte Reue zu denken, sodann auch, andere zu hindern. Durch die Unwissenheit der Lehrer ist dies unordentliche Wesen (tumultus) in der Beichte entstanden, während die Beichte kurz und offen sein sollte, so daß beide [Beichtkind und Beichtiger] bald fertig werden könnten.

Nun wollen wir die Todsünden erstlich vor uns nehmen, deren Zahl man, wie es scheint, von Johannes Chrysostomus hergenommen hat. Denn er legt auch die sieben Völker der Cananiter 5 Mos. 7, 1. in sittlicher Weise (moraliter) als sieben Laster aus, indem er will, daß alle Sünden unter denselben begriffen werden, mehr auf Grund der Siebenzahl, welche ein Ganzes bedeutet, als daß die Namen solches in sich enthielten. Denn etliche machen acht, andere neun Todsünden, indem sie noch den Ungehorsam und die Eitelkehre hinzuthun; noch andere theilen die Mütter in Töchter, und die Gattungen in Arten, das Ganze in dazu gehörige (integrales) und untergeordnete (subjectivas) Theile: aus jenen haben sie unendlich viele, ja, ganze Meere von Sünden hergeleitet, sich aber unterdessen um den rechten Verstand der Gebote nicht bekümmert.

Darum gibt es nicht bloß sieben Todsünden, sondern dies ist eine willkürliche Theilung, die keinen nothwendigen Grund noch Bewährung [aus der Schrift] hat, sondern es können sowohl weniger als auch mehr sein, je nachdem in der Schrift eine verschiedene figürliche Rede (tropologia) vorkommt. Daher hat St. Augustinus sie in zwei Theile getheilt, nämlich ein angezündetes und ein unterdrücktes Feuer, das heißt, die Liebe, die übel entzündet, und die Furcht, die in böser Weise demüthigt.

Die erste Todsünde nun ist die Hoffahrt, welche zwiefach ist: die äußere, welche in leiblichen Gütern stattfindet, und die innere, welche in geistlichen Gütern sich äußert, und deshalb immer wider das erste Gebot ist. Das ist offenbar, denn keinen anderen Gott haben, das heißt, auf nichts vertrauen, kein Gefallen haben an irgend etwas, sich über nichts freuen, sich an nichts ergötzen, keines Dinges genießen als allein in GOtt. Aber die Hoffahrt setzt ihr Vertrauen bald auf Reichthum, bald auf Stärke, bald auf Schönheit und Kleidung, bald auf Macht und Ehre, bald auf edeles Geschlecht; desgleichen innerlich auf Weisheit, Kunst, Verstand, Gerechtigkeit, Tugend, Heiligkeit, hat ihr Gefallen daran, und ehrt diese erschaffenen Güter in solcher Weise (eo actu), wie sie allein Gotte gebührt. Denn die Hoffahrt bringt immer zwei Dinge mit sich: erstlich, daß man sich selbst gefalle, zweitens, daß man andere hochmüthig verachte. Deshalb schließet sie die eitle Ehre inwendig in sich; wenn diese nach außen herausbricht, ist es um so schlimmer. Darum ehrt der Hoffärtige, der an sich etwas Gutes sieht, darin nicht GOtt, dankt ihm auch nicht dafür, indem er GOtt preist, sondern steht auf sich selbst, hat sein Ergötzen an sich selbst, und ist eitel in seinen Gedanken. Dann hat er aber zugleich Acht auf den, der solche Gaben nicht hat; wenn er den nun mit sich vergleicht, so folgt mit Nothmendigkeit
\quellseite{680}
, daß er ihn verachtet und spricht: er sei etwas, jener sei nichts, und es wird das Wort Röm. 1, 22. erfüllt: „Da sie sprachen, sie seien weise, sind sie zu Narren geworden“; indem sie so reden (nämlich inwendig, zu sich selbst), sie seien gerecht, schön, reich etc., sind sie ungerecht, häßlich, arm geworden. Offenb. 3, 17.: „Du sprichst: Ich bin reich, und habe gar satt, und weißt nicht, daß du bist arm und jämmerlich und bloß.“

So hat also die Hoffahrt kein eigentliches für sich selbst bestehendes Wesen (substantiam), es sei denn, daß sich der Hoffärtige mit einem Geringeren vergleichen kann, und gleichsam von oben nach unten herabsieht. Dagegen ist das Wesen der Demuth nur da, wo der Demüthige sich mit einem Besseren vergleicht, und gleichsam nach oben zu den höchsten Dingen hinaufsieht. Darum hat die Hoffahrt, dieses geile und feiste Kalb, vier Füße, nämlich:

Unkenntniß seines Bösen. Kenntniß seines Guten Unkenntniß des fremden Guten. Kenntniß des fremden Bösen Dies ist die rechte Unkenntniß Gottes seiner selbst.

Das Gegentheil dieses Vierfüßlers hat vier Füße der Demuth, nämlich:

Unkenntniß seines Guten. Kenntniß seines Bösen Unkenntniß des fremden Bösen. Kenntniß des fremden Guten Dies ist die rechte Kenntniß Gottes seiner selbst.

Darum ist es nicht vonnöthen, daß du beichtest, du seiest geneigt zur Hoffahrt, denn wir sind immer alle hoffärtig, und kein Mensch ist ohne alle Hoffahrt, sondern nur dann [mußt du beichten], wenn du ihrer bösen Lust gehorcht hast, und eingewilligt durch Wort, Werk oder Herz. Das Uebrige muß man Gotte klagen in heimlichem Seufzen und verborgener Beichte im Kämmerlein, daß er durch seine Gnade dies Uebel, die alte Wurzel, ausrotte. Es ist gesagt worden von der Hoffahrt, wie sie unter dem ersten Gebote verboten worden sei, denn die Hoffahrt ist der Anfang der Sünde, oder, was dasselbe ist: der Anfang der Sünde ist, von GOtt abfallen, nämlich, daß man sich selbst ein Abgott wird, sich selbst wohlgefällt, an sich selbst ein größeres Ergötzen hat als an dem HErrn, und so in sich selbst einen anderen Gott ehrt. Dies ist die allerschwerste und erste Sünde, wiewohl keiner da ist, der von Adam herkommt, so lange wir leben, der nicht in irgend einem Stücke diese Abgötterei in sich treibe, und um desswillen immer Buße thun und seufzen müsse darüber, daß er sich selbst gefalle.

\footnote{Löscher: „Das Folgende hat Luther als Predigt gehalten am Sonntage Septuagesimä“ [den 15. Februar 1517]. Das Exordium hierzu findet sich Walch, St. Louiser Ausg., Bd. XII, 1822.} Jetzt muß vom Geize geredet werden, welcher ein Gelüsten ist und, wie der Apostel

[1 Tim. 6, 9. f.] sagt, gleicherweise „eine Wurzel alles Uebels“, „denn, die da reich werden wollen, die fallen in die Stricke des Teufels und viele thörichte und schädliche Lüste“. Dieses Laster wird in zwei, ja in drei Geboten verboten, nämlich im siebenten: „Du sollst nicht stehlen“ und im letzten: „Du sollst nicht begehren deines Nächsten Gut“ und im ersten. Denn da in dem siebenten Gebote (wie gesagt ist) alles Anziehen (tractatio) fremden Gutes verboten wird, und ein jegliches Gebot (wie auch gesagt ist), so geistlich ist, daß es nicht allein das Werk und die Frucht und die Blätter und Zweige, das ist, die Werke, Worte und Geberden verbietet, sondern auch die Wurzel und den Saft, das heißt, die böse Lust, aus welcher solches herkommt, so ist offenbar, daß der Geiz unter dem Diebstahl verboten werde, da bei einem Diebe und einem Geizigen dieselbe Wurzel da ist, nämlich das Gelüsten. Desgleichen, da in dem letzten Gebote nach dem Buchstaben das Begehren des Gutes des Nächsten verboten wird, so ist offenbar, daß auch die Quelle und das Haupt dieses Begehrens im Geiste verboten wird; aber dies ist der Geiz. Darum verbietet das allerliebste Gebot recht eigentlich das letzte und allerlebendigste Laster, welches bei den alten Leuten wieder jung wird und bei den Sterbenden lebt. Aber der Apostel Paulus, welcher Eph. 5, 5.
\quellseite{681}
in seinem Sinne höher einhergeht, sagt, daß ein Geiziger ein Diener der Götzen sei, das heißt, ein Verehrer derselben, und so stellt er dies unter das erste Gebot. Dasselbe thut auch Baruch, Cap. 3, 17.: „Die Silber und Gold sammeln, darauf die Menschen ihr Vertrauen setzen.“ Und der 78. Psalm, V. 58. [Ps. 106, 38.], sagt im Geiste: „Und sie dienten den Götzen Canaan's.“ Canaan heißt ein Händler; was könnten deren Götzen wohl passender bedeuten, als Geld und Münzen, goldene und silberne Gefäße, Häuser und andere Besitztümer? Ja, deshalb wird uns im Gesetze vornehmlich verboten, daß wir nicht goldene und silberne Götter machen sollen, und der erste Götzendienst des Volkes ist mit dem Golde geschehen, unter dem goldenen Kalbe [2 Mos. 32, 4.], aber in gleicher Weise auch der größte Götzendienst zur Zeit Jerobeams, unter den goldenen Kälbern [1 Kön. 12, 28.]. Daher hat der Apostel richtig unter Bild einen Götzen verstanden, indem er den Geizigen einen Verehrer der Götzen nennt, und im Lateinischen klingt aurum fast wie auarum; auri auro, auari auaro sind bloß in Einem Buchstaben verschieden, so daß aurum [Gold] von auare [heftig begehren] seinen Namen erhalten zu haben scheint. Da nun der Apostel den Geizigen einen Götzendiener nennt, so legt er dadurch das erste Gebot genugsam aus, daß das Herz des Menschen allein auf Gott vertrauen solle. Denn auf dieselbe Meinung sagt er, daß auch die Fresser und die Unkeuschen Verehrer des Bauches seien, und sagt [Phil. 3, 19.]: „Welchen der Bauch ihr Gott ist.“ Aber in diesem Sinne sind alle Gebote in dem ersten Gebote enthalten als in ihrem Haupte. Denn jede Sünde ist eine Verachtung GOttes, und eben dadurch wider das erste Gebot, und es kann auch kein Gebot verletzt werden, ohne daß zugleich auch das erste verletzt werde, welches alle in sich enthält. Denn GOttes Wille wird in allen Geboten übertreten und unserem Willen nachgesetzt [, wenn man wider dieselben thut], und so wird GOtt nicht verehrt, sondern vielmehr ein anderer [Gott] wird anstatt seiner geehrt. Doch der Apostel hat recht eigentlich angegeben, daß der Geiz wider dasselbe sei, weil er sich ergötzt an den Bildnissen des Goldes und Silbers, sehr ähnlich der rechten Abgötterei. Alle anderen Sünden aber haben solche Götzenbilder und Bildnisse nicht.

Deshalb müssen wir wiederum damit schließen, daß niemand ganz völlig rein ist von Geiz, wie Jesaias und Jeremias sagen [Jes. 56, 11. Jer. 6, 13.]: „Sie geizen allesamt, Klein und Groß.“ Aber selig ist der, welcher um deswillen dieser Neigung seiner verderbten Natur weniger folgt. Alle müssen deshalb vor GOtt seufzen, aber nur das Werk, es sei innerlich oder äußerlich, vor dem Priester beichten. Und wer sagen würde, daß er dies Begehren nicht habe, der würde sich gewißlich in großem Hochmuth die höchste Vollkommenheit der Reinheit anmaßen. Darum sind diejenigen selig, welche Leid tragen [Matth. 5, 4.], denn sie haben Grund (materiam) in sich selbst, über sich Leid zu tragen, wie der HErr zu den Weibern sagte, die ihm nachfolgten [Luc. 23, 28.].

\footnote{Bücher: „Das Folgende hat Luther vorgetragen am Tage Matthiae [den 24. Februar] im Jahre 1517.“ Das Exordium hiezu findet sich Walch, St. Louiser Ausg., Bd. XII, 1762, und noch einmal, in verbesserter Uebersetzung, Bd. XIX, 754.} Die dritte Todsünde ist die Unkeuschheit, deren Arten unter dem sechsten Gebote: „Du sollst nicht ehebrechen“ genugsam angezeigt sind, und niemand zweifelt daran, daß sie dahin gehöre; darum mag es genügen, daß dies angezeigt ist.

Die vierte ist Fressen und Saufen, die Schwelgerei, ja, die Anreizung und der Freier und Diener der Unkeuschheit, wie auch der Heide sagt: Sine Cerere et Baccho friget Venus [Ohne Speise und Trank gedeiht die Unkeuschheit nicht]. Dieses Laster setzt der Apostel Phil. 3, 19. unter das erste Gebot, indem er sagt: „Welchen der Bauch ihr Gott ist.“ Dasselbe thut er Röm. 16, 18.: „Denn solche dienen nicht Christo, sondern ihrem Bauch.“ Es kann aber einfach unter das sechste Gebot gesetzt werden, denn wo die Unkeuschheit verboten wird, da wird sicherlich auch aller Anlaß und Reizung zur Unkeuschheit verboten. Aber unter allen Reizungen zur Unkeuschheit ist das Größte und Stärkste Fressen und Saufen. Die Ursache davon ist diese, daß dies etwas Innerliches ist; alles Andere, als Anschauen, Unterredung, Hören, Betasten, ist äußerlich. Fressen und Saufen aber macht die Adern schwellen und bewegt den ganzen Leib zur Unkeuschheit. Denn auch der Heide Plinius sagt, daß auf die Trunkenheit eine entsetzliche Unkeuschheit folge. Und St. Hieronymus sagt
\quellseite{682}
in seiner Schrift an den Eustochius und anderswo, daß nichts so zur Unkeuschheit bewege als unverdaute Speise. Darum gibt sich ein gefräßiger Mensch (gulosus) williglich der Unkeuschheit hin. So sind auch Müßiggang und Faulheit, schläfriges Wesen, weiches Lager und alle anderen Dinge verboten, die da Unkeuschheit erwecken, was ohne Zweifel daraus [aus dem sechsten Gebote] erkannt werden kann. Und wenn sich ein Mensch, der in Essen und Trinken zu viel gethan hat (gulosus), nicht sofort mit viel Arbeit beschäftigt, so wird er von unüberwindlicher Unkeuschheit und Kitzel eingenommen werden. Daher sagt St. Hieronymus: Ein Bauch der glüht von Weine verschäumt bald in Unkeuschheit. Siehe, wie die Unkeuschheit gewissermaßen gleichsam der natürliche Schaum der Völlerei ist, gleichwie von Natur beim Feuer der Rauch aufgeht, und der Schaum aufsteigt bei eben eingegossenem Weine, und das siedende Wasser Blasen aufwirft. Wer kann dies verhindern, es sei denn mit der größten Gewalt? Wer daher Keuschheit gebietet, der gebietet ohne Zweifel auch die Mittel, die zur Keuschheit dienen, als Mäßigkeit, Wachen, Arbeit, Beten, Lesen, Nachdenken, Studiren, dem Nächsten dienen, Kälte, Hitze, Armut etc.

Hier ist jedoch dies zu bemerken: Viele Leute machen sich über das Erbrechen ein größeres Gewissen, als über das Vollsaufen und Trunkenheit, gleichwie viele die nächtliche Befleckung für etwas Größeres achten als die Ursache davon, die bösen Gedanken und die Völlerei. Aber sicherlich ist das Erbrechen nichts Böses, noch auch jemals eine Sünde, da es sowohl wider Willen als auch unter Schmerzen statt hat. Deshalb sollen die Trunkenen das Erbrechen nicht allein nicht fliehen, sondern sogar suchen, wie Sirach Cap. 31, 25. sagt [Vulg.]: „Wenn du gezwungen worden bist viel zu trinken,\footnote{In der Vulgata: in edendo.} so stehe auf und wirf es aus, und es wird dich erfrischen, und du wirst bei deinem Leibe keine Krankheit zuziehen.“ Diesen Rath des weisen Mannes soll man annehmen, und nicht ohne Grund. Denn er sagt: „Es wird dich erfrischen“, das heißt, es wird dich erleichtern. Denn es dient diese Entledigung dazu, daß die Adern nicht so aufgebläht werden und Kitzel des Fleisches entstehe, da das Wenige, was zurückbleibt, leichter ver-

daut wird. Denn es ist besser durch Erbrechen der Unkeuschheit zu entgehen, als durch Völlerei in Unkeuschheit zu gerathen.

Zweitens, weil Krankheit des Leibes vermieden wird, als, Verletzung des Kopfes und aller Sinne und Kräfte. Wer aber sollte nicht glauben, daß es eine größere Sünde sei, daß der ganze Leib verletzt werde durch Vollsein und Behalten des Genossenen, als durch Erbrechen von dieser Gefahr befreit zu werden? Denn das Leben verkürzen ist eine größere Sünde, als sich erbrechen, ja, sich erbrechen ist keine Sünde, aber das ist böse, daß du dich so gefüllt und überladen hast, daß du dich erbrechen mußt. Und wollte GOtt, daß ein Gesetz gegeben würde, daß alle Trunkenen und vollen Leute gezwungen würden zu erbrechen! Ich hoffe, daß sie endlich die Trunkenheit verlassen würden, da sie wüßten, daß sie das bald wieder ausspien müßten, was sie tranken. Wiewohl sie auch bald das Wasser lassen und zu Stuhl gehen müssen, so kehren sie sich doch nicht daran. Daß aber die Schrift bisweilen das Speien straft, als [2 Petr. 2, 22.]: „Der Hund frißt wieder, was er gespieen hat“, und Jes. 28, 8.: „Alle Tische sind voll Speiens“, wird verstanden von dem Auspeien des guten und heilsamen Wortes, welches er hätte behalten sollen.

Drittens. Wenn jemand Gift oder etwas Anderes, was schädlich ist, getrunken hat, und es ohne Sünde ausspeiet, warum sollte der mit Sünde speien, der Wein oder einen anderen Trank so getrunken hat, daß es einen dem Gifte ähnlichen Schaden ausübt? Warum muß man sich ein Gewissen machen über das Hineinschütten, nicht über das Von sichgeben, da man dort vielem Bösen dient wider GOttes Gebot, hier aber vielem Guten, wider kein Gebot, ja, für GOttes Gebot.

Die Arten und Unterschiede der Völlerei übergehe ich, weil sie nicht immer Todsünden sind, da ja auch Fressen und Saufen nicht immer eine Todsünde ist, ja selten, es sei denn zu einer Gewohnheit geworden.

Die fünfte Todsünde ist der Zorn, über den bei dem fünften Gebote genug gesagt worden ist.

Die sechste ist der Haß, der vom Zorn nicht hätte [als eine besondere Sünde] unterschieden werden sollen, da er nichts ist als ein eingewurzelter Zorn, wie St. Augustinus in seiner Regel sagt: Laß ja nicht den Zorn zum Haß werden,
